% 来自2026-10-21统一复核稿；题号、图号和解答保留。
\clearpage
\begingroup
\fancyhead[R]{\small 模块五：磁场、电磁感应与机电耦合}
\clubpenalty=10000\widowpenalty=10000\displaywidowpenalty=10000
\chapter{模块五：磁场、电磁感应与机电耦合}\label{m5:ux6a21ux5757ux4e94ux78c1ux573aux7535ux78c1ux611fux5e94ux4e0eux673aux7535ux8026ux5408}

本模块讨论洛伦兹力、典型磁场、法拉第定律、自感互感与磁场储能。选讲部分加入磁矢势、移动边界、磁共能和磁扩散。

{\def\LTcaptype{none} % do not increment counter
\begin{longtable}[]{@{}
  >{\raggedright\arraybackslash}p{(\linewidth - 4\tabcolsep) * \real{0.3333}}
  >{\raggedright\arraybackslash}p{(\linewidth - 4\tabcolsep) * \real{0.3333}}
  >{\raggedright\arraybackslash}p{(\linewidth - 4\tabcolsep) * \real{0.3333}}@{}}
\toprule\noalign{}
\begin{minipage}[b]{\linewidth}\raggedright
层级
\end{minipage} & \begin{minipage}[b]{\linewidth}\raggedright
内容
\end{minipage} & \begin{minipage}[b]{\linewidth}\raggedright
阅读定位
\end{minipage} \\
\midrule\noalign{}
\endhead
\bottomrule\noalign{}
\endlastfoot
基础必讲 &
带电粒子运动、毕奥---萨伐尔定律、安培环路定理、磁介质基础、电磁感应、自感互感、线性无源元件
& 构成磁场与电磁感应题的主干 \\
拔高选讲与决赛拓展 &
磁矢势与规范、非轴对称感生电场、磁能与磁共能、电磁---机械耦合、磁扩散与趋肤效应、转动约束、正则动量
& 处理空间分布、强耦合和一般化动力学模型 \\
\end{longtable}
}

\par\noindent\begin{minipage}{\linewidth}
\begin{center}
\includegraphics[width=\linewidth,height=.40\textheight,keepaspectratio]{\detokenize{assets/module-05/theory/tikz-v11-module5-method-map.pdf}}
\par\smallskip{\small 图示：模块五连接磁场来源、电磁感应、储能和机械运动}
\end{center}
\end{minipage}\par

\section{第一单元：带电粒子、稳恒磁场与磁介质}\label{m5:ux7b2cux4e00ux5355ux5143ux5e26ux7535ux7c92ux5b50ux7a33ux6052ux78c1ux573aux4e0eux78c1ux4ecbux8d28}

专题四十一至四十五讨论带电粒子运动、磁场来源、磁介质和磁矢势。

\section{专题四十一：洛伦兹力与带电粒子运动（基础必讲与拔高选讲）}\label{m5:ux4e13ux9898ux56dbux5341ux4e00ux6d1bux4f26ux5179ux529bux4e0eux5e26ux7535ux7c92ux5b50ux8fd0ux52a8ux57faux7840ux5fc5ux8bb2ux4e0eux62d4ux9ad8ux9009ux8bb2}

\begin{quote}
阅读准备：两个一阶方程可以组成转动 匀强磁场中的 \(v_x,v_y\)
相互耦合。把其中一式再求导并代入另一式，就得到简谐型二阶方程；带符号角频率决定转向，绝对值决定周期。初始速度给出两个分量的初值。
\end{quote}

\par\addvspace{.4\baselineskip}
\noindent\begin{minipage}{\linewidth}

\textbf{例 M5-001：速度分量组成回旋}

设平面速度满足
\(\dot v_x=\Omega v_y\text{、}\dot v_y=-\Omega v_x\)，其中带符号常量
\(\Omega\ne0\)，初值为
\(v_x(0)=v_0>0\text{、}v_y(0)=0\)。消去一个分量得到二阶方程，证明速率不变，并判断
\(\Omega>0\) 时速度在 \(v_xv_y\) 平面内的转向。

\end{minipage}\par

\par\noindent\begin{minipage}{\linewidth}
\textbf{解答}

\begin{center}
\includegraphics[width=\linewidth,height=.40\textheight,keepaspectratio]{\detokenize{assets/module-05/examples/M5-001.pdf}}
\par\smallskip{\small 图 M5-001：速度分量组成回旋}
\end{center}
\end{minipage}\par

图 M5-001 在速度平面上表示过程；红箭头是 \(\Omega>0\)
时的转向。由两分量方程消元可得简谐方程，蓝圆则表示速率守恒。

\(\dot v_x=\Omega v_y\)、\(\dot v_y=-\Omega v_x\) 推出
\(\ddot v_x+\Omega^2v_x=0\)，同时
\(d(v_x^2+v_y^2)/dt=0\)，故速率不变而方向旋转。

由初值进一步得到

\[
v_x=v_0\cos\Omega t,\qquad v_y=-v_0\sin\Omega t.
\]

\(\Omega>0\) 时速度矢量先向 \(-v_y\)
方向偏转，在速度平面内顺时针旋转；\(\Omega<0\) 时反向。

以下讨论经典点电荷在给定场中的运动，包括匀强场方程、交叉场、选速与荷质比测量、回旋加速、霍尔效应和缓变磁场。除非特别说明，粒子均视为经典点粒子，忽略辐射损失、碰撞和粒子自身对外场的反作用。

\subsection{一、洛伦兹力、功率与适用条件（基础必讲）}\label{m5:ux4e00ux6d1bux4f26ux5179ux529bux529fux7387ux4e0eux9002ux7528ux6761ux4ef6ux57faux7840ux5fc5ux8bb2}

电场可以通过静止试探电荷所受力定义；磁力则与电荷速度有关。电荷量为 \(q\)
的粒子在给定电磁场中的动力学方程是

\[
\boxed{
\frac{d\mathbf p}{dt}
=q\left(\mathbf E+\mathbf v\times\mathbf B\right)
}.
\]

非相对论条件 \(v\ll c\)
下，\(\mathbf p=m\mathbf v\)；相对论条件下则必须使用
\(\mathbf p=\gamma m\mathbf v\)。将动力学方程点乘速度，得到

\[
\frac{dK}{dt}
=\frac{d\mathbf p}{dt}\cdot\mathbf v
=q\mathbf E\cdot\mathbf v,
\]

因为

\[
q(\mathbf v\times\mathbf B)\cdot\mathbf v=0.
\]

因此，\textbf{磁场可以改变动量方向，却不能直接改变粒子的动能}。若速率发生变化，能量必来自电场、碰撞或其他非磁力；``磁场变强所以粒子被加速''不是充分理由。

若电场是与时间无关的静电场，\(\mathbf E=-\nabla\varphi\)，则

\[
\frac{d}{dt}(K+q\varphi)
=q\mathbf E\cdot\mathbf v
+q\nabla\varphi\cdot\mathbf v
=0.
\]

故粒子的机械能满足

\[
\boxed{K+q\varphi=\text{常量}}.
\]

\par\addvspace{.4\baselineskip}
\noindent\begin{minipage}{\linewidth}

\textbf{例 M5-002：静电势降与动能增量}

质量为 \(m\)、电荷量为 \(q>0\) 的经典粒子由 \(A\) 运动至
\(B\)。外电场为静电场，\(U=\varphi_A-\varphi_B>0\)；同时可有任意静磁场。忽略碰撞、辐射和其他非电磁力，求这段运动的动能增量，并说明结果是否依赖粒子的实际弯曲路径。

\end{minipage}\par

\par\noindent\begin{minipage}{\linewidth}
\textbf{解答}

\begin{center}
\includegraphics[width=\linewidth,height=.40\textheight,keepaspectratio]{\detokenize{assets/module-05/examples/M5-002.pdf}}
\par\smallskip{\small 图 M5-002：静电势降与动能增量}
\end{center}
\end{minipage}\par

图 M5-002 的 \(A\text{、}B\)
是电势降的两端；蓝色弯曲路径不影响静电力的功，静磁场只改变路径方向。

正电荷从高电势处经过电势降 \(U\) 到达低电势处，其动能增加
\(qU\)；同时存在的静磁场可以改变到达路径和方向，却不改变这一能量增量。

\par\addvspace{.4\baselineskip}
\noindent\begin{minipage}{\linewidth}

\textbf{例 M5-003：电场力的瞬时功率}

经典带电粒子在某时刻 \(t_0\)
的速度非零，且与给定非零电场垂直。求此时电场力的瞬时功率；若后来速度方向发生改变，能否据此断言全过程电场不做功？用受力与轨迹说明原因。

\end{minipage}\par

\par\noindent\begin{minipage}{\linewidth}
\textbf{解答}

\begin{center}
\includegraphics[width=\linewidth,height=.40\textheight,keepaspectratio]{\detokenize{assets/module-05/examples/M5-003.pdf}}
\par\smallskip{\small 图 M5-003：电场力的瞬时功率}
\end{center}
\end{minipage}\par

图 M5-003
中初始速度向上、电场力向右，直角标记对应零瞬时功率；蓝轨迹随后向右弯曲，不能据初始直角推出全过程不做功。

若某瞬间 \(\mathbf E\perp\mathbf v\)，则即使电场力不为零，瞬时功率仍为
\(q\mathbf E\cdot\mathbf v=0\)；但随后的速度方向改变后，电场通常会开始做功。因此``瞬时不做功''不等于``全过程动能不变''。

\par\noindent\begin{minipage}{\linewidth}
\begin{center}
\includegraphics[width=\linewidth,height=.40\textheight,keepaspectratio]{\detokenize{assets/module-05/original/tikz-m5-01-lorentz-force-direction.png}}
\par\smallskip{\small 图示：洛伦兹力方向由速度与磁场张成的平面及电荷符号共同确定，且磁场力不改变动能}
\end{center}
\end{minipage}\par

判定磁场力方向必须分两步：先用右手定则求
\(\mathbf v\times\mathbf B\)，再乘 \(q\)
的符号。正电荷同向，负电荷反向。结果必须同时垂直于 \(\mathbf v\) 和
\(\mathbf B\)。

\subsection{二、匀强磁场中的运动方程（基础必讲）}\label{m5:ux4e8cux5300ux5f3aux78c1ux573aux4e2dux7684ux8fd0ux52a8ux65b9ux7a0bux57faux7840ux5fc5ux8bb2}

设

\[
\mathbf B=B\mathbf e_z,
\qquad B>0,
\]

并暂取 \(\mathbf E=\mathbf 0\)。非相对论运动方程按分量写成

\[
\dot v_x=\Omega v_y,
\qquad
\dot v_y=-\Omega v_x,
\qquad
\dot v_z=0,
\qquad
\Omega=\frac{qB}{m}.
\]

这里的 \(\Omega\) 是\textbf{带符号回旋角频率}；其绝对值

\[
\omega_c=|\Omega|=\frac{|q|B}{m}
\]

决定周期，而其符号决定绕行方向。由初始速度 \((v_{x0},v_{y0},v_{z0})\) 得

\[
\begin{aligned}
v_x(t)&=v_{x0}\cos\Omega t+v_{y0}\sin\Omega t,\\
v_y(t)&=v_{y0}\cos\Omega t-v_{x0}\sin\Omega t,\\
v_z(t)&=v_{z0}.
\end{aligned}
\]

积分后可将横向轨迹写成以固定点为圆心的圆：

\[
\begin{aligned}
x_c&=x_0+\frac{v_{y0}}{\Omega},\\
y_c&=y_0-\frac{v_{x0}}{\Omega},
\end{aligned}
\qquad
(x-x_c)^2+(y-y_c)^2=R^2,
\]

其中

\[
\boxed{
R=\frac{mv_{\perp}}{|q|B},
\qquad
T=\frac{2\pi m}{|q|B}
}.
\]

半径由 \(|q|\) 决定，圆心在哪一侧则由 \(q\)
的符号决定。直接使用``向心力指向圆心''是检查叉乘方向最可靠的办法。

\par\noindent\begin{minipage}{\linewidth}
\begin{center}
\includegraphics[width=\linewidth,height=.40\textheight,keepaspectratio]{\detokenize{assets/module-05/theory/tikz-m5-topic41-v9-from38-v1-01-charge-sign-orbit.pdf}}
\par\smallskip{\small 图示：正负电荷的轨道与圆心}
\end{center}
\end{minipage}\par

\par\addvspace{.4\baselineskip}
\noindent\begin{minipage}{\linewidth}

\textbf{例 M5-004：电荷符号与圆心位置}

质量均为 \(m\)、电荷量分别为 \(+|q|\) 和 \(-|q|\)
的两粒子，分别从各自原点以 \(v\mathbf e_x\) 射入匀强磁场
\(B\mathbf e_z\)，其中
\(v\text{、}B>0\)。忽略电场及其他作用，求初始受力方向、圆心坐标、半径和绕行方向，采用非相对论模型。

\end{minipage}\par

\par\noindent\begin{minipage}{\linewidth}
\textbf{解答}

\begin{center}
\includegraphics[width=\linewidth,height=.40\textheight,keepaspectratio]{\detokenize{assets/module-05/examples/M5-004.pdf}}
\par\smallskip{\small 图 M5-004：电荷符号与圆心位置}
\end{center}
\end{minipage}\par

图 M5-004
的两圆以各自初始点为局部原点；正电荷的圆心在下、负电荷的圆心在上。初始向心力方向决定圆心所在一侧。

粒子从原点以 \(\mathbf v_0=v\mathbf e_x\) 射入
\(\mathbf B=B\mathbf e_z\)。若 \(q>0\)，初始磁力为
\(-qvB\mathbf e_y\)，圆心为 \((0,-mv/(qB))\)；若
\(q<0\)，圆心改在上方。两者半径相同，绕行方向相反。

\subsubsection{1. 螺旋运动}\label{m5:ux87baux65cbux8fd0ux52a8}

将速度分解为

\[
\mathbf v=\mathbf v_{\perp}+v_{\parallel}\mathbf e_z.
\]

\(v_{\parallel}\) 保持不变，\(\mathbf v_{\perp}\)
作匀速圆周运动。两者叠加成沿磁场方向推进的螺旋线，螺距为

\[
\boxed{
h=v_{\parallel} T
=\frac{2\pi m v_{\parallel}}{|q|B}
}.
\]

\par\noindent\begin{minipage}{\linewidth}
\begin{center}
\includegraphics[width=\linewidth,height=.40\textheight,keepaspectratio]{\detokenize{assets/module-05/theory/tikz-m5-topic41-v9-from38-v1-02-helix-projections.pdf}}
\par\smallskip{\small 图示：螺旋运动的两个投影}
\end{center}
\end{minipage}\par

读空间图时应先看两个投影：沿磁场观察得到圆；侧面观察得到轴向匀速叠加周期振动。这样可以直接读出半径、周期和螺距，无须从透视图猜测前后关系。

\subsubsection{2.
相对论修正（拔高选讲）}\label{m5:ux76f8ux5bf9ux8bbaux4feeux6b63ux62d4ux9ad8ux9009ux8bb2}

当粒子速度不可忽略于光速时，磁场仍不做功，因而 \(\gamma\)
保持不变；但半径和回旋频率应改写为

\[
\boxed{
R=\frac{p_{\perp}}{|q|B}
=\frac{\gamma mv_{\perp}}{|q|B},
\qquad
\omega_c=\frac{|q|B}{\gamma m}
}.
\]

因此高能粒子的回旋频率会降低。非相对论公式 \(R=mv_{\perp}/(|q|B)\) 只在
\(\gamma\simeq1\) 时成立。

\subsection{三、匀强交叉电磁场与导引中心漂移（基础必讲与拔高选讲）}\label{m5:ux4e09ux5300ux5f3aux4ea4ux53c9ux7535ux78c1ux573aux4e0eux5bfcux5f15ux4e2dux5fc3ux6f02ux79fbux57faux7840ux5fc5ux8bb2ux4e0eux62d4ux9ad8ux9009ux8bb2}

先考虑

\[
\mathbf E\perp\mathbf B,
\qquad
E<cB,
\]

且两场均匀、恒定。定义

\[
\boxed{
\mathbf v_E=\frac{\mathbf E\times\mathbf B}{B^2}
}.
\]

利用三重积

\[
\mathbf v_E\times\mathbf B
=\frac{(\mathbf E\times\mathbf B)\times\mathbf B}{B^2}
=-\mathbf E,
\]

再令 \(\mathbf u=\mathbf v-\mathbf v_E\)，运动方程化为

\[
m\frac{d\mathbf u}{dt}=q\mathbf u\times\mathbf B.
\]

所以 \(\mathbf u\) 作回旋运动，而整个回旋中心以 \(\mathbf v_E\)
匀速漂移：

\[
\boxed{
\mathbf v_{\text{漂}}=\frac{\mathbf E\times\mathbf B}{B^2}
}.
\]

这一漂移与 \(m\)、\(|q|\)
以及电荷正负均无关；正负电荷只是在漂移中心周围反向回旋。

\par\noindent\begin{minipage}{\linewidth}
\begin{center}
\includegraphics[width=\linewidth,height=.40\textheight,keepaspectratio]{\detokenize{assets/module-05/theory/tikz-m5-topic41-v9-from38-v1-03-exb-drift.pdf}}
\par\smallskip{\small 图示：交叉场中的回旋与漂移}
\end{center}
\end{minipage}\par

\par\addvspace{.4\baselineskip}
\noindent\begin{minipage}{\linewidth}

\textbf{例 M5-005：交叉场的漂移方向}

空间有恒定匀强交叉场
\(\mathbf E=E\mathbf e_x\text{、}\mathbf B=B\mathbf e_z\)，其中
\(E\text{、}B>0\) 且
\(E/B\ll c\)。对非相对论带电粒子，求导引中心的平均漂移方向与速度，判断它是否依赖电荷正负；若另加平行电场，说明轴向运动如何变化。

\end{minipage}\par

\par\noindent\begin{minipage}{\linewidth}
\textbf{解答}

\begin{center}
\includegraphics[width=\linewidth,height=.40\textheight,keepaspectratio]{\detokenize{assets/module-05/examples/M5-005.pdf}}
\par\smallskip{\small 图 M5-005：交叉场的漂移方向}
\end{center}
\end{minipage}\par

图 M5-005
用从静止释放的轨迹示意漂移；蓝线每一回旋周期向下推进，绿色箭头就是
\(\mathbf E\times\mathbf B/B^2\)，不依赖电荷符号。

先找一个横向匀速特殊解 \(\mathbf v_D\)，使
\(\mathbf E_{\perp}+\mathbf v_D\times\mathbf B=0\)。由叉乘恒等式，可验
\(\mathbf v_D=\mathbf E_{\perp}\times\mathbf B/B^2\)。令
\(\mathbf w=\mathbf v_{\perp}-\mathbf v_D\)，运动方程变为
\(m\dot{\mathbf w}=q\mathbf w\times\mathbf B\)，即绕匀速解叠加回旋。图中轨迹仍起伏，但一周期的回旋位移平均为零，留下的平均速度才是漂移速度；它不随
\(q\) 的符号改变。

若 \(\mathbf E=E\mathbf e_x\)、\(\mathbf B=B\mathbf e_z\)，则

\[
\mathbf v_E
=\frac{E\mathbf e_x\times B\mathbf e_z}{B^2}
=-\frac EB\mathbf e_y.
\]

无论粒子带正电还是负电，其平均漂移都沿 \(-y\) 方向。

若还存在 \(\mathbf E_{\parallel}\)，横向运动仍包含
\(\mathbf E\times\mathbf B\) 漂移，而平行方向满足

\[
m\frac{dv_{\parallel}}{dt}=qE_{\parallel}.
\]

此时电场会改变动能和螺距。

\subsection{四、速度选择器、质谱仪与回旋加速器（基础必讲与拔高选讲）}\label{m5:ux56dbux901fux5ea6ux9009ux62e9ux5668ux8d28ux8c31ux4eeaux4e0eux56deux65cbux52a0ux901fux5668ux57faux7840ux5fc5ux8bb2ux4e0eux62d4ux9ad8ux9009ux8bb2}

\subsubsection{1.
速度选择器与磁偏转}\label{m5:ux901fux5ea6ux9009ux62e9ux5668ux4e0eux78c1ux504fux8f6c}

取正离子束沿 \(+x\) 方向运动，电场沿 \(-y\)，磁场 \(\mathbf B_1\)
垂直纸面向里。电场力向下，磁场力向上。只有满足

\[
qE=qvB_1
\]

的粒子不偏转，因此

\[
\boxed{v=\frac{E}{B_1}}.
\]

进入只有磁场 \(B_2\) 的分析区后，轨道半径满足

\[
R=\frac{mv}{|q|B_2}.
\]

于是

\[
\boxed{
\frac{m}{|q|}
=\frac{B_2R}{v}
=\frac{B_1B_2R}{E}
}.
\]

\par\noindent\begin{minipage}{\linewidth}
\begin{center}
\includegraphics[width=\linewidth,height=.40\textheight,keepaspectratio]{\detokenize{assets/module-05/theory/tikz-m5-topic41-v9-from38-v1-04-mass-spectrometer.pdf}}
\par\smallskip{\small 图示：选速后的磁偏转与探测位置}
\end{center}
\end{minipage}\par

\par\addvspace{.4\baselineskip}
\noindent\begin{minipage}{\linewidth}

\textbf{例 M5-006：离子半径比与质荷比}

两种同号离子通过同一速度选择器后速度相同，并进入同一个匀强分析磁场。测得
\(R_2=2R_1\)。离子质量、电荷量分别为
\(m_i\text{、}q_i\)，电荷态尚未知。求两种质荷比 \(m_i/|q_i|\)
之比，并说明能否单独确定质量比。

\end{minipage}\par

\par\noindent\begin{minipage}{\linewidth}
\textbf{解答}

\begin{center}
\includegraphics[width=\linewidth,height=.40\textheight,keepaspectratio]{\detokenize{assets/module-05/examples/M5-006.pdf}}
\par\smallskip{\small 图 M5-006：离子半径比与质荷比}
\end{center}
\end{minipage}\par

图 M5-006 中两条半圆的半径对应
\(m_i/|q_i|\)，探测点的差异不能在未知电荷态时直接解释为质量差。

两种同号离子经过同一选速器后速度相同。在同一分析磁场中，若第二种离子的轨道半径是第一种的两倍，则

\[
\frac{m_2/|q_2|}{m_1/|q_1|}=\frac{R_2}{R_1}=2.
\]

这只能确定质荷比之比；若电荷态未知，不能只凭半径断定质量比。

\subsubsection{2.
回旋加速器（拔高选讲）}\label{m5:ux56deux65cbux52a0ux901fux5668ux62d4ux9ad8ux9009ux8bb2}

非相对论粒子在均匀磁场中的回旋频率

\[
f_c=\frac{|q|B}{2\pi m}
\]

与半径和速率无关。两个空心电极内部近似没有加速电场，粒子只在缝隙处被交变电场加速；每跨越一次电势差为
\(U\) 的缝隙，动能增加 \(|q|U\)。

\par\noindent\begin{minipage}{\linewidth}
\begin{center}
\includegraphics[width=\linewidth,height=.40\textheight,keepaspectratio]{\detokenize{assets/module-05/theory/tikz-m5-topic41-v9-from38-v1-05-cyclotron-safe-v7.pdf}}
\par\smallskip{\small 图示：回旋加速器的半圆轨道与窄缝}
\end{center}
\end{minipage}\par

若设备最大轨道半径为 \(R_{\max}\)，则非相对论最大动能为

\[
\boxed{
K_{\max}
=\frac{q^2B^2R_{\max}^2}{2m}
}.
\]

当 \(\gamma\) 明显偏离 \(1\) 时，回旋频率降为
\(|q|B/(2\pi\gamma m)\)，固定频率电场与粒子逐渐失谐；这给经典回旋加速器设定了能量上限。

\subsection{五、电流受力、磁偶极矩与外场作用（基础必讲）}\label{m5:ux4e94ux7535ux6d41ux53d7ux529bux78c1ux5076ux6781ux77e9ux4e0eux5916ux573aux4f5cux7528ux57faux7840ux5fc5ux8bb2}

体电流密度在磁场中的力密度为

\[
\mathbf f=\mathbf J\times\mathbf B.
\]

对横截面积为 \(A\)、长度为 \(d\mathbf l\)
的细导线微元，取电流密度在截面上近似均匀，则

\[
d\mathbf F
=(\mathbf J\times\mathbf B)A\,dl
=(JA)\,d\mathbf l\times\mathbf B.
\]

利用 \(I=JA\)，得到

\[
\boxed{
d\mathbf F=I\,d\mathbf l\times\mathbf B
}.
\]

闭合平面电流环的磁偶极矩定义为

\[
\boxed{
\mathbf m=IS\mathbf n
},
\]

其中 \(\mathbf n\) 按电流方向的右手定则确定；若线圈有 \(N\) 匝，则
\(\mathbf m=NIS\mathbf n\)。在匀强外磁场中，闭合环的合力为零，但力矩一般不为零：

\[
\boxed{
\boldsymbol\tau=\mathbf m\times\mathbf B,
\qquad
U=-\mathbf m\cdot\mathbf B
}.
\]

以矩形线圈为例，若其面积为 \(S=ab\)，线圈法向与磁场夹角为
\(\theta\)，两条与磁场垂直的边各受力 \(F=IaB\)，力臂的有效间距为
\(b\sin\theta\)，故

\[
\tau=F b\sin\theta
=IabB\sin\theta
=mB\sin\theta.
\]

这正是矢量式 \(\boldsymbol\tau=\mathbf m\times\mathbf B\) 的大小形式。

\par\noindent\begin{minipage}{\linewidth}
\begin{center}
\includegraphics[width=\linewidth,height=.40\textheight,keepaspectratio]{\detokenize{assets/module-05/theory/tikz-m5-topic41-v9-from38-v1-07-loop-torque.pdf}}
\par\smallskip{\small 图示：均匀外场中的电流环：合力与力矩}
\end{center}
\end{minipage}\par

若 \(\theta\) 是 \(\mathbf m\) 与 \(\mathbf B\) 的夹角，则

\[
\tau=mB\sin\theta,
\qquad
U=-mB\cos\theta.
\]

\(\mathbf m\) 与 \(\mathbf B\) 同向是稳定平衡，反向是不稳定平衡。

\par\addvspace{.4\baselineskip}
\noindent\begin{minipage}{\linewidth}

\textbf{例 M5-007：磁偶极子转动所做的功}

固定大小磁偶极矩为 \(m\) 的小电流环置于恒定匀强外磁场 \(B\)
中。保持电流和环面积不变，使磁矩与外场的夹角从 \(90^{\circ}\) 缓慢转至
\(0\)。求磁场力矩所做的功，并明确力矩公式使用外场还是电流环自身场。

\end{minipage}\par

\par\noindent\begin{minipage}{\linewidth}
\textbf{解答}

\begin{center}
\includegraphics[width=\linewidth,height=.40\textheight,keepaspectratio]{\detokenize{assets/module-05/examples/M5-007.pdf}}
\par\smallskip{\small 图 M5-007：磁偶极子转动所做的功}
\end{center}
\end{minipage}\par

图 M5-007 中绿色磁矩随电流环侧视红线转动，\(\theta\)
从外场方向量起。比较转动前后的势能即可求磁场力矩所做的功。

电流环从 \(\theta=90^{\circ}\) 缓慢转到 \(\theta=0\)，磁场所做的功为

\[
W_B=U_i-U_f=0-(-mB)=mB.
\]

\par\noindent\begin{minipage}{\linewidth}
\begin{center}
\includegraphics[width=\linewidth,height=.40\textheight,keepaspectratio]{\detokenize{assets/module-05/process/M5-007a.pdf}}
\par\smallskip{\small 图 M5-007a：小电流环自身场与外场力矩的区别}
\end{center}
\end{minipage}\par

图 M5-007a
的蓝色闭合曲线表示小电流环的自身远场；前图中施加力矩的平行磁场则是外场。两者作用不同，不能把自身磁感线当成环的受力轨迹。

小电流环自身产生的远场只由总磁偶极矩决定：

图中的闭合磁感线表示电流环的\textbf{自身场}，而力矩公式中的
\(\mathbf B\)
是\textbf{外场}，两者不能混为一张``受力轨迹图''。在缓慢变化的外场中，小磁偶极子的平动势能仍为
\(U=-\mathbf m\cdot\mathbf B\)；若 \(\mathbf m\) 近似固定，则合力可写为

\[
\mathbf F=\nabla(\mathbf m\cdot\mathbf B).
\]

\subsection{六、霍尔效应的大小与符号（基础必讲）}\label{m5:ux516dux970dux5c14ux6548ux5e94ux7684ux5927ux5c0fux4e0eux7b26ux53f7ux57faux7840ux5fc5ux8bb2}

设一种载流子的数密度为 \(n\)、电荷量为带符号的 \(q\)、平均漂移速度为
\(\mathbf v_d\)。电流密度为

\[
\mathbf J=nq\mathbf v_d.
\]

横向稳态建立后，霍尔电场满足

\[
q\left(\mathbf E_H+\mathbf v_d\times\mathbf B\right)=0.
\]

因此

\[
\boxed{
\mathbf E_H
=-\mathbf v_d\times\mathbf B
=\frac1{nq}\mathbf B\times\mathbf J
},
\]

霍尔系数为

\[
\boxed{R_H=\frac1{nq}}.
\]

它保留 \(q\) 的符号：单一正载流子对应 \(R_H>0\)，电子对应 \(R_H<0\)。

\par\noindent\begin{minipage}{\linewidth}
\begin{center}
\includegraphics[width=\linewidth,height=.40\textheight,keepaspectratio]{\detokenize{assets/module-05/theory/tikz-m5-topic41-v9-from38-v1-06-hall-sign-margin12-v11.pdf}}
\par\smallskip{\small 图示：常规电流和磁场相同时，正载流子与电子具有相反的漂移速度、积累电荷和霍尔电场方向}
\end{center}
\end{minipage}\par

取样品电流沿 \(+x\)，磁场沿 \(+z\)，宽度沿 \(y\)，厚度为 \(t\)。定义

\[
U_H\equiv\varphi(y=-w/2)-\varphi(y=+w/2)=E_{Hy}w.
\]

由 \(j_x=I/(wt)\) 可得带符号公式

\[
\boxed{
U_H=R_H\frac{IB}{t}
=\frac{IB}{nqt}
}.
\]

若只求大小，则

\[
|U_H|=\frac{IB}{n|q|t}.
\]

换用相反的电压端次序、磁场方向或电流方向，霍尔电压符号都会反转。因此判号前必须先写明端子定义。

\par\addvspace{.4\baselineskip}
\noindent\begin{minipage}{\linewidth}

\textbf{例 M5-008：霍尔电压随磁场反向}

单一载流子薄片保持纵向电流不变。将垂直薄片的磁场从 \(+B\) 反为
\(-B\)，接触引线位置和电压参考方向不变。判断霍尔场、霍尔电压以及理想线性纵向电阻压降的奇偶性，并写出用两次电压测量提取霍尔信号的方法。

\end{minipage}\par

\par\noindent\begin{minipage}{\linewidth}
\textbf{解答}

\begin{center}
\includegraphics[width=\linewidth,height=.40\textheight,keepaspectratio]{\detokenize{assets/module-05/examples/M5-008.pdf}}
\par\smallskip{\small 图 M5-008：霍尔电压随磁场反向}
\end{center}
\end{minipage}\par

图 M5-008
以正载流子示意上下侧的霍尔极性；保持纵向电流后反转磁场，上下电荷与霍尔电压一起反号，可据此提取奇对称信号。

反转磁场时保持电流与接触条件不变。在寄生偏置 \(U_{\rm off}\)
不反转的模型下，\(U(+B)=U_{\rm off}+U_H\)、\(U(-B)=U_{\rm off}-U_H\)。因此图中的提取式为
\[
U_H=\frac{U(+B)-U(-B)}2.
\] 半和给出不反转的偏置，半差才保留随磁场反转的霍尔信号。

同一薄片中把磁场反向而电流不变，则 \(\mathbf E_H\) 和 \(U_H\)
同时反向；普通纵向电阻压降在理想线性近似下不随
\(\mathbf B\to-\mathbf B\)
反号。实验上可利用这一奇偶性分离霍尔信号与接触不对称造成的背景电压。

\par\noindent\begin{minipage}{\linewidth}
\begin{center}
\includegraphics[width=\linewidth,height=.40\textheight,keepaspectratio]{\detokenize{assets/module-05/theory/tikz-m5-topic41-v9-from38-v8-01.pdf}}
\par\smallskip{\small 图示：洛伦兹力判向与轨道分析}
\end{center}
\end{minipage}\par

\subsection{七、缓变非均匀磁场、绝热磁矩与磁镜（拔高选讲）}\label{m5:ux4e03ux7f13ux53d8ux975eux5747ux5300ux78c1ux573aux7eddux70edux78c1ux77e9ux4e0eux78c1ux955cux62d4ux9ad8ux9009ux8bb2}

若磁场在一个回旋半径内变化很小，即

\[
\frac{R_L}{L_B}\ll1,
\qquad
L_B\sim\frac{B}{|\nabla B|},
\]

并且在一个回旋周期内变化缓慢，则横向回旋运动具有绝热不变量

\[
\boxed{
\mu_{\rm ad}=\frac{mv_{\perp}^2}{2B}
}.
\]

这里 \(\mu_{\rm ad}\)
是正的标量绝热磁矩，不应与任意方向固定的外加磁偶极矩混淆。静态磁场不做功，粒子总动能保持不变，于是

\[
\frac12mv_{\parallel}^2+\mu_{\rm ad}B=K.
\]

沿磁场方向的平均运动满足

\[
m\frac{dv_{\parallel}}{dt}
=-\mu_{\rm ad}\nabla_{\parallel} B.
\]

粒子进入强场区时，\(v_{\perp}\) 增大而 \(v_{\parallel}\) 减小。当
\(v_{\parallel}=0\) 时到达镜点并反射。

\par\noindent\begin{minipage}{\linewidth}
\begin{center}
\includegraphics[width=\linewidth,height=.40\textheight,keepaspectratio]{\detokenize{assets/module-05/theory/tikz-m5-topic41-v9-from38-v1-08-magnetic-mirror.pdf}}
\par\smallskip{\small 图示：磁镜反射与绝热不变量}
\end{center}
\end{minipage}\par

设粒子在 \(B_0\) 处的速度与磁场夹角为投掷角 \(\alpha_0\)，则

\[
v_{\perp0}=v\sin\alpha_0.
\]

在镜点 \(B_m\)
处，\(v_{\parallel}=0\)、\(v_{\perp}=v\)，由绝热磁矩守恒得

\[
\boxed{
B_m=\frac{B_0}{\sin^2\alpha_0}
}.
\]

\par\addvspace{.4\baselineskip}
\noindent\begin{minipage}{\linewidth}

\textbf{例 M5-009：磁镜反射的临界角}

无电场的缓变磁镜装置在中心的磁场为 \(B_0\)，两端最大磁场均为
\(4B_0\)。粒子从中心出发，投掷角 \(\alpha_0\in[0,\pi]\) 是速度与
\(+\mathbf B\)
的夹角。假定总动能和绝热磁矩守恒，求两端都不能逃逸的投掷角范围；另写出朝
\(+\mathbf B\) 运动时的临界角。

\end{minipage}\par

\par\noindent\begin{minipage}{\linewidth}
\textbf{解答}

\begin{center}
\includegraphics[width=\linewidth,height=.40\textheight,keepaspectratio]{\detokenize{assets/module-05/examples/M5-009.pdf}}
\par\smallskip{\small 图 M5-009：磁镜反射的临界角}
\end{center}
\end{minipage}\par

图 M5-009 的 \(B_m\)
是平行速度降为零的镜点。要求它不超出装置最大场给出临界角；若把投掷角限定在
\(0\) 至 \(90^{\circ}\)，只需写 \(\alpha_0\ge30^{\circ}\)。

这里使用无电场、磁场沿轨道缓慢变化的近似：总动能守恒，磁矩
\(\mu_m=mv_{\perp}^2/(2B)\) 近似不变。于是 \[
v_{\parallel}^2=v_0^2\left(1-\frac{B}{B_0}\sin^2\alpha_0\right).
\] 镜点满足 \(v_{\parallel}=0\)，所以
\(B_m=B_0/\sin^2\alpha_0\)。将图中这一必要场强与装置能提供的
\(B_{\max}\) 比较，才得到下面的反射角条件。

若磁镜装置中最大磁场为 \(B_{\max}=4B_0\)，则能够被反射的粒子须满足

\[
\sin^2\alpha_0\ge\frac{B_0}{B_{\max}}=\frac14,
\]

即
\(\alpha_0\ge30^{\circ}\)。投掷角太小的粒子沿磁力线穿过强场区，不会反射。

若采用题设的完整投掷角区间
\([0,\pi]\)，反射条件还要求不落入另一端的损失锥，所以完整范围为

\[
30^{\circ}\le\alpha_0\le150^{\circ}.
\]

若只考虑朝 \(+\mathbf B\) 出发的锐角分支，条件可简写为
\(\alpha_0\ge30^{\circ}\)。

\subsubsection{梯度与曲率漂移（决赛拓展）}\label{m5:ux68afux5ea6ux4e0eux66f2ux7387ux6f02ux79fbux51b3ux8d5bux62d3ux5c55}

缓变磁场还会使导引中心产生与电荷符号有关的漂移。记
\(\mathbf b=\mathbf B/B\)、\(\boldsymbol\kappa=(\mathbf b\cdot\nabla)\mathbf b\)，则

\[
\boxed{
\mathbf v_{\nabla B}
=\frac{mv_{\perp}^2}{2qB^3}\mathbf B\times\nabla B
},
\]

\[
\boxed{
\mathbf v_{\text{曲率}}
=\frac{mv_{\parallel}^2}{qB}\mathbf b\times\boldsymbol\kappa
}.
\]

两式都与 \(1/q\)
成正比，因此正、负电荷沿相反方向漂移；这与不依赖电荷符号的
\(\mathbf E\times\mathbf B\) 漂移形成鲜明对比。

\subsection{八、练习与解答}\label{m5:ux516bux7ec3ux4e60ux4e0eux89e3ux7b54}

\par\addvspace{.4\baselineskip}
\noindent\begin{minipage}{\linewidth}

\subsubsection{作业一：圆轨道的受力、半径与圆心（基础）}\label{m5:ux4f5cux4e1aux4e00ux5706ux8f68ux9053ux7684ux53d7ux529bux534aux5f84ux4e0eux5706ux5fc3ux57faux7840}

\textbf{题号：M5-010}

\textbf{层级：} 基础　　\textbf{方法：} 洛伦兹力、圆周运动

设粒子质量为 \(m\)，采用非相对论模型，忽略电场、重力、碰撞与辐射。

正电荷 \(q\) 在原点以速度 \(\mathbf v=v\mathbf e_x\) 运动，匀强磁场为
\(\mathbf B=B\mathbf e_z\)，其中
\(v,B>0\)。求初始受力、轨道半径与圆心位置。

\end{minipage}\par

\par\noindent\begin{minipage}{\linewidth}
\textbf{解答}

\begin{center}
\includegraphics[width=\linewidth,height=.40\textheight,keepaspectratio]{\detokenize{assets/module-05/examples/M5-010.pdf}}
\par\smallskip{\small 图 M5-010：圆轨道的受力、半径与圆心}
\end{center}
\end{minipage}\par

图 M5-010 中 \(O\)
为入射点，向右的速度与向纸外的磁场叉乘后向下；红色初始受力指向
\(C\)，半径由这个向心力确定。

\[
\mathbf F=q\mathbf v\times\mathbf B
=qvB(\mathbf e_x\times\mathbf e_z)
=-qvB\mathbf e_y.
\]

粒子初始向 \(-y\) 偏转，且磁力始终垂直速度，速率不变。由 \(qvB=mv^2/R\)
得

\[
R=\frac{mv}{qB}.
\]

初始向心加速度指向 \(-\mathbf e_y\)，故圆心位于

\[
\boxed{\mathbf r_c=-R\mathbf e_y}.
\]

这与一般公式 \(x_c=x_0+v_{y0}/\Omega\)、\(y_c=y_0-v_{x0}/\Omega\) 一致。

\par\addvspace{.4\baselineskip}
\noindent\begin{minipage}{\linewidth}

\subsubsection{作业二：交叉场中的粒子轨迹（复赛）}\label{m5:ux4f5cux4e1aux4e8cux4ea4ux53c9ux573aux4e2dux7684ux7c92ux5b50ux8f68ux8ff9ux590dux8d5b}

\textbf{题号：M5-011}

\textbf{层级：} 复赛　　\textbf{方法：} 洛伦兹力

粒子质量为
\(m\)，\(E\text{、}B\text{、}q>0\)，\(E/B\ll c\)，采用经典非相对论模型，忽略重力与碰撞。

正电荷 \(q\) 在 \(t=0\) 时位于原点并静止。空间存在

\[
\mathbf E=E\mathbf e_x,
\qquad
\mathbf B=B\mathbf e_z.
\]

求速度和位置随时间的关系，并指出平均漂移方向。

\end{minipage}\par

\par\noindent\begin{minipage}{\linewidth}
\textbf{解答}

\begin{center}
\includegraphics[width=\linewidth,height=.40\textheight,keepaspectratio]{\detokenize{assets/module-05/examples/M5-011.pdf}}
\par\smallskip{\small 图 M5-011：交叉场中的粒子轨迹}
\end{center}
\end{minipage}\par

图 M5-011
的摆线从原点静止出发；在每个尖点速度重新为零。绿色下移方向来自对一个回旋周期取平均，不能把整条轨迹当作圆。

带符号回旋频率为 \(\Omega=qB/m>0\)。分量方程为

\[
\dot v_x=\frac{qE}{m}+\Omega v_y,
\qquad
\dot v_y=-\Omega v_x.
\]

取漂移速度

\[
\mathbf v_E=-\frac EB\mathbf e_y,
\]

并令 \(\mathbf u=\mathbf v-\mathbf v_E\)。初始条件为

\[
u_x(0)=0,
\qquad
u_y(0)=\frac EB.
\]

故

\[
u_x=\frac EB\sin\Omega t,
\qquad
u_y=\frac EB\cos\Omega t.
\]

于是

\[
\boxed{
v_x=\frac EB\sin\Omega t,
\qquad
v_y=\frac EB(\cos\Omega t-1)
}.
\]

\par\noindent\begin{minipage}{\linewidth}
\begin{center}
\includegraphics[width=\linewidth,height=.40\textheight,keepaspectratio]{\detokenize{assets/module-05/process/M5-011a.pdf}}
\par\smallskip{\small 图 M5-011a：交叉场速度分量随回旋相位变化}
\end{center}
\end{minipage}\par

图 M5-011a
的蓝、红曲线分别按标出的不同纵标归一化；红线的非正性说明摆线沿 \(-y\)
逐周期推进。

再积分并使用 \(x(0)=y(0)=0\)：

\[
\boxed{
x(t)=\frac{E}{B\Omega}(1-\cos\Omega t),
\qquad
y(t)=\frac{E}{B\Omega}\sin\Omega t-\frac EBt
}.
\]

轨迹是沿 \(-y\) 方向推进的摆线。对一个周期取平均，振荡项消失：

\[
\boxed{
\langle\mathbf v\rangle
=-\frac EB\mathbf e_y
=\frac{\mathbf E\times\mathbf B}{B^2}
}.
\]

\par\addvspace{.4\baselineskip}
\noindent\begin{minipage}{\linewidth}

\subsubsection{作业三：同位素质谱分离（复赛）}\label{m5:ux4f5cux4e1aux4e09ux540cux4f4dux7d20ux8d28ux8c31ux5206ux79bbux590dux8d5b}

\textbf{题号：M5-012}

\textbf{层级：} 复赛　　\textbf{方法：} 洛伦兹力

设 \(E\text{、}B_1\text{、}B_2>0\)
为场强大小，电场与选速器磁场方向按不偏转条件配置，所有离子非相对论；进入分析区时入口位置及方向完全相同。

两种单电离正离子质量分别为 \(m_1\)、\(m_2\)。它们经过电场强度为
\(E\)、磁场为 \(B_1\) 的速度选择器，再进入磁场为 \(B_2\)
的分析区，均沿半圆轨道打到同一直线探测板上。求两条落点之间的距离。

\end{minipage}\par

\par\noindent\begin{minipage}{\linewidth}
\textbf{解答}

\begin{center}
\includegraphics[width=\linewidth,height=.40\textheight,keepaspectratio]{\detokenize{assets/module-05/examples/M5-012.pdf}}
\par\smallskip{\small 图 M5-012：同位素质谱分离}
\end{center}
\end{minipage}\par

图 M5-012
的探测板与入口位于同一直线上，每条半圆从入口到落点跨过一条直径。因此
\(P_1P_2\) 是半径差的两倍。

选速器给出共同速度

\[
v=\frac{E}{B_1}.
\]

单电离离子有 \(|q|=e\)，分析区半径为

\[
R_i=\frac{m_iv}{eB_2}
=\frac{m_iE}{eB_1B_2}.
\]

半圆轨道的入口到落点距离为直径 \(2R_i\)，故两落点间距为

\[
\boxed{
\Delta x
=2|R_2-R_1|
=\frac{2E|m_2-m_1|}{eB_1B_2}
}.
\]

若两种离子的电荷态不同，应将 \(m_i/e\) 改为
\(m_i/|q_i|\)，不能直接把半径比解释为质量比。

\par\addvspace{.4\baselineskip}
\noindent\begin{minipage}{\linewidth}

\subsubsection{作业四：回旋加速器的圈数与最大动能（复赛）}\label{m5:ux4f5cux4e1aux56dbux56deux65cbux52a0ux901fux5668ux7684ux5708ux6570ux4e0eux6700ux5927ux52a8ux80fdux590dux8d5b}

\textbf{题号：M5-013}

\textbf{层级：} 复赛　　\textbf{方法：} 洛伦兹力

假定过窄缝时总能获得固定能量
\(|q|U\)，缝内偏转可忽略，电压相位与回旋运动同步。以``首次达到或超过半径
\(R\)''定义跨缝次数；若恰好达到，可令连续公式为整数。

非相对论回旋加速器的磁场为 \(B\)，两电极缝隙间的电势差幅值为
\(U\)，粒子电荷量绝对值为 \(|q|\)、质量为
\(m\)。忽略初动能。求粒子达到轨道半径 \(R\)
前跨越缝隙的次数以及此时动能。

\end{minipage}\par

\par\noindent\begin{minipage}{\linewidth}
\textbf{解答}

\begin{center}
\includegraphics[width=\linewidth,height=.40\textheight,keepaspectratio]{\detokenize{assets/module-05/examples/M5-013.pdf}}
\par\smallskip{\small 图 M5-013：回旋加速器的跨缝次数与动能}
\end{center}
\end{minipage}\par

图 M5-013 的窄缝承担加速，两个电极内部只发生磁偏转。每次过缝增加
\(|q|U\)，以此求首次达到指定半径所需的整数次数；橙色各段是半径逐次增大的半圆，红色短段是过缝加速。

设跨过窄缝的次数为整数 \(N\ge0\)。每次同步过缝增能 \(|q|U\)，所以

\[
K_N=N|q|U,\qquad
R_N=\frac{\sqrt{2mN|q|U}}{|q|B}.
\]

欲首次达到或超过半径 \(R\)，要求 \(R_N\ge R\)。由此

\[
\boxed{N_{\min}=\left\lceil\frac{|q|B^2R^2}{2mU}\right\rceil},
\qquad
\boxed{K_{\rm first}=N_{\min}|q|U}.
\]

若指定半径恰好属于某次过缝后的允许轨道，则连续值本身就是整数，且

\[
K(R)=\frac{q^2B^2R^2}{2m}.
\]

\par\noindent\begin{minipage}{\linewidth}
\begin{center}
\includegraphics[width=\linewidth,height=.40\textheight,keepaspectratio]{\detokenize{assets/module-05/process/M5-013a.pdf}}
\par\smallskip{\small 图 M5-013a：每次过缝的动能是离散增加的}
\end{center}
\end{minipage}\par

图 M5-013a 的红点只取整数
\(N\)，虚线用于读出连续估算。若要求首次达到目标半径，必须取不小于连续估算的最小整数。

若装置要求轨道绝不超过边缘半径，则应改为向下取整求允许的最大次数；两种问法不可混用。以上只适用于非相对论回旋同步模型。

\par\addvspace{.4\baselineskip}
\noindent\begin{minipage}{\linewidth}

\subsubsection{作业五：磁镜的损失锥（决赛）}\label{m5:ux4f5cux4e1aux4e94ux78c1ux955cux7684ux635fux5931ux9525ux51b3ux8d5b}

\textbf{题号：M5-014}

\textbf{层级：} 决赛　　\textbf{方法：} 洛伦兹力

假定场在一个回旋尺度上缓慢变化，忽略电场、碰撞与辐射，粒子总动能及绝热磁矩守恒；投掷角与
\(+\mathbf B\) 的夹角取 \(0\le\alpha_0\le\pi\)。

粒子在磁场最弱处 \(B_0\) 的速度大小为 \(v\)，投掷角为
\(\alpha_0\)。沿磁力线方向的最大磁场为
\(B_{\max}\)。求能被磁镜反射的投掷角范围，并解释所谓``损失锥''。

\end{minipage}\par

\par\noindent\begin{minipage}{\linewidth}
\textbf{解答}

\begin{center}
\includegraphics[width=\linewidth,height=.40\textheight,keepaspectratio]{\detokenize{assets/module-05/examples/M5-014.pdf}}
\par\smallskip{\small 图 M5-014：磁镜的损失锥}
\end{center}
\end{minipage}\par

图 M5-014 的粒子先向强场区运动，再在 \(v_{\parallel}=0\)
处反射。最大场不足以形成镜点的轨道逃逸，速度空间中对应两端的损失锥。

反射要求理论镜点不超过装置最大磁场：

\[
B_m=\frac{B_0}{\sin^2\alpha_0}
\le B_{\max}.
\]

\par\noindent\begin{minipage}{\linewidth}
\begin{center}
\includegraphics[width=\linewidth,height=.40\textheight,keepaspectratio]{\detokenize{assets/module-05/process/M5-014a.pdf}}
\par\smallskip{\small 图 M5-014a：磁镜：平行速度在镜点降为零}
\end{center}
\end{minipage}\par

图 M5-014a 表示 \(v_{\parallel}^2\) 随磁场线性下降；横坐标归一化后应在
\(B/B_m=1\) 达零，最终图的终点就是镜点。

因此

\[
\boxed{
\sin^2\alpha_0\ge\frac{B_0}{B_{\max}}
}.
\]

定义临界角

\[
\alpha_c=\arcsin\sqrt{\frac{B_0}{B_{\max}}}.
\]

沿 \(+\mathbf B\) 方向运动的粒子若
\(0\le\alpha_0<\alpha_c\)，平行速度过大，在到达 \(B_{\max}\)
前不能降为零，会从端部逃逸；速度空间中围绕磁场方向的这一区域称为损失锥。另一端同样存在关于
\(-\mathbf B\) 方向的损失锥。

\subsection{九、使用条件}\label{m5:ux4e5dux4f7fux7528ux6761ux4ef6}

\begin{enumerate}
\def\labelenumi{\arabic{enumi}.}
\tightlist
\item
  \(R=mv_{\perp}/(|q|B)\)、\(\omega_c=|q|B/m\)
  是非相对论公式；高能粒子必须改用 \(R=p_{\perp}/(|q|B)\) 和
  \(\omega_c=|q|B/(\gamma m)\)。
\item
  ``磁场不做功''只说明磁力本身不改动能，不代表含电场、碰撞或时变约束的系统动能一定不变。
\item
  \(\mathbf E\times\mathbf B\)
  漂移公式要求场在一个回旋尺度上近似均匀；若
  \(E\ge cB\)，不能在经典低速图像中简单理解为速度 \(E/B\)
  的导引中心漂移。
\item
  霍尔系数 \(R_H=1/(nq)\)
  只适用于单一载流子模型；多载流子材料中霍尔系数还取决于各类载流子的密度和迁移率。
\item
  磁镜公式依赖绝热条件
  \(R_L/L_B\ll1\)；在磁场突变、接近零场或强散射区域，绝热磁矩可能不守恒。
\item
  电流环力矩公式中的 \(\mathbf B\)
  是外磁场；电流环自身的磁场不能直接代入计算自身所受总力矩。
\end{enumerate}

\section{专题四十二：毕奥-萨伐尔定律与典型磁场（基础必讲与拔高选讲）}\label{m5:ux4e13ux9898ux56dbux5341ux4e8cux6bd5ux5965-ux8428ux4f10ux5c14ux5b9aux5f8bux4e0eux5178ux578bux78c1ux573aux57faux7840ux5fc5ux8bb2ux4e0eux62d4ux9ad8ux9009ux8bb2}

毕奥---萨伐尔定律用于由稳恒电流分布计算空间磁场。积分前先写清源点到场点的位矢，用叉乘判断磁场微元方向，再利用对称性筛选需要积分的分量。

\subsection{一、毕奥-萨伐尔定律的对象、方向与适用条件}\label{m5:ux4e00ux6bd5ux5965-ux8428ux4f10ux5c14ux5b9aux5f8bux7684ux5bf9ux8c61ux65b9ux5411ux4e0eux9002ux7528ux6761ux4ef6}

\subsubsection{1. 线电流形式}\label{m5:ux7ebfux7535ux6d41ux5f62ux5f0f}

设源点为 \(\mathbf r'\)，场点为 \(\mathbf r\)，从源点指向场点的位矢为

\[
\mathbf R=\mathbf r-\mathbf r',
\qquad
R=|\mathbf R|.
\]

稳恒细导线中，电流元 \(I\,d\mathbf l'\) 在场点产生

\[
\boxed{
d\mathbf B(\mathbf r)
=\frac{\mu_0}{4\pi}
\frac{I\,d\mathbf l'\times\mathbf R}{R^3}
}.
\]

其中 \(d\mathbf l'\) 沿约定的正电流方向。大小为

\[
dB
=\frac{\mu_0}{4\pi}
\frac{I\,dl'\sin\theta}{R^2},
\]

\(\theta\) 是 \(d\mathbf l'\) 与 \(\mathbf R\) 的夹角，方向由
\(d\mathbf l'\times\mathbf R\) 的右手定则确定。

\par\noindent\begin{minipage}{\linewidth}
\begin{center}
\includegraphics[width=\linewidth,height=.40\textheight,keepaspectratio]{\detokenize{assets/module-05/theory/tikz-m5-topic42-v9-from39-v1-01-current-element-direction.pdf}}
\par\smallskip{\small 图示：分别沿叉乘方向和从侧面观察电流元、源点到场点位矢与磁场微元，明确磁场垂直于二者张成的平面}
\end{center}
\end{minipage}\par

图中的圆点表示矢量指向观察者，叉号表示矢量背离观察者。只背``右手螺旋''容易在倾斜空间图中判错；可靠做法是先把
\(d\mathbf l'\) 与 \(\mathbf R\) 放入同一平面，再确定该平面的法向。

\par\addvspace{.4\baselineskip}
\noindent\begin{minipage}{\linewidth}

\textbf{例 M5-015：径向引线的磁场贡献}

一段稳恒直导线通电流 \(I\)，场点 \(P\)
位于该线段所在直线上且不在导线材料内。只计算这一段导线对 \(P\)
的贡献，用毕奥---萨伐尔定律判断磁场是否为零，并说明结论能否推广为空间任意场点。

\end{minipage}\par

\par\noindent\begin{minipage}{\linewidth}
\textbf{解答}

\begin{center}
\includegraphics[width=\linewidth,height=.40\textheight,keepaspectratio]{\detokenize{assets/module-05/examples/M5-015.pdf}}
\par\smallskip{\small 图 M5-015：径向引线的磁场贡献}
\end{center}
\end{minipage}\par

图 M5-015 中场点 \(P\)
位于直导线的延长线上，绿色电流元与红色位移矢量共线，叉乘为零；这只说明该段导线在
\(P\) 的贡献为零。

若某段直导线与场点位于同一直线上，则对该段每个电流元都有
\(d\mathbf l'\parallel\mathbf R\) 或
\(d\mathbf l'\parallel(-\mathbf R)\)，因而

\[
d\mathbf l'\times\mathbf R=\mathbf0.
\]

该段导线对这个场点没有磁场贡献。结论只针对这个场点，并不表示该导线在空间处处不产生磁场。

\subsubsection{2.
面电流与体电流形式}\label{m5:ux9762ux7535ux6d41ux4e0eux4f53ux7535ux6d41ux5f62ux5f0f}

若电流分布不能近似成细导线，分别用面电流密度 \(\mathbf K\) 和体电流密度
\(\mathbf J\)：

\[
\boxed{
\mathbf B(\mathbf r)
=\frac{\mu_0}{4\pi}
\int_S
\frac{\mathbf K(\mathbf r')\times(\mathbf r-\mathbf r')}
{|\mathbf r-\mathbf r'|^3}\,dS'
}
\]

以及

\[
\boxed{
\mathbf B(\mathbf r)
=\frac{\mu_0}{4\pi}
\int_V
\frac{\mathbf J(\mathbf r')\times(\mathbf r-\mathbf r')}
{|\mathbf r-\mathbf r'|^3}\,dV'
}.
\]

三种形式本质相同，只是源的理想化层次不同。理想细导线模型在导线本身处给出发散磁场；若题目要求导线内部或表面附近的有限场，必须保留真实横截面和电流密度分布。

\subsubsection{3.
``电流元''不是孤立的物理源}\label{m5:ux7535ux6d41ux5143ux4e0dux662fux5b64ux7acbux7684ux7269ux7406ux6e90}

稳恒电流满足电荷不随时间积累，体电流区域内有

\[
\nabla\cdot\mathbf J=0.
\]

因此真实稳恒电流必须形成闭合回路，或从研究区域边界流入再流出。\(I\,d\mathbf l'\)
是积分中的局部片段，不能把一小段电流从回路中孤立出来而忽略回流路径；只有当其余部分由对称性抵消、贡献为零或远小于所求量时，才可不写。

\subsection{二、有限长直导线与无限长极限}\label{m5:ux4e8cux6709ux9650ux957fux76f4ux5bfcux7ebfux4e0eux65e0ux9650ux957fux6781ux9650}

\subsubsection{1.
有符号端点公式}\label{m5:ux6709ux7b26ux53f7ux7aefux70b9ux516cux5f0f}

取直导线沿 \(z\) 轴，场点 \(P\) 到导线的垂足为 \(O\)，垂直距离为
\(a\)。以 \(O\) 为 \(z=0\)，导线两端坐标分别为 \(z_-\)、\(z_+\)，其中
\(z_-<z_+\)。电流沿 \(+z\) 方向时，

\[
d\mathbf l'=dz\,\mathbf e_z,
\qquad
\mathbf R=a\mathbf e_x-z\mathbf e_z.
\]

所以

\[
d\mathbf l'\times\mathbf R
=a\,dz\,\mathbf e_y,
\]

所有磁场微元方向相同，只需积分大小：

\[
dB
=\frac{\mu_0I}{4\pi}
\frac{a\,dz}{(a^2+z^2)^{3/2}}.
\]

利用

\[
\frac{d}{dz}
\left(
\frac{z}{\sqrt{a^2+z^2}}
\right)
=\frac{a^2}{(a^2+z^2)^{3/2}},
\]

得到最不容易误判正负号的端点式

\[
\boxed{
B
=\frac{\mu_0I}{4\pi a}
\left[
\frac{z}{\sqrt{a^2+z^2}}
\right]_{z_-}^{z_+}
}.
\]

方向沿以导线为轴的圆周切向，由右手定则决定。

\par\noindent\begin{minipage}{\linewidth}
\begin{center}
\includegraphics[width=\linewidth,height=.40\textheight,keepaspectratio]{\detokenize{assets/module-05/theory/tikz-m5-topic42-v9-from39-v1-02-finite-wire.pdf}}
\par\smallskip{\small 图示：有限长直导线的端点张角、有符号端点式以及半无限长和无限长极限}
\end{center}
\end{minipage}\par

若垂足 \(O\) 位于两端之间，并把从 \(PO\) 到两条端点连线的锐角记作
\(\alpha_+\)、\(\alpha_-\)，则

\[
\boxed{
B=\frac{\mu_0I}{4\pi a}
(\sin\alpha_++\sin\alpha_-)
}.
\]

若垂足位于导线延长线上而不在两端之间，直接套``两个正弦相加''容易出错；此时应使用前面的有符号端点式，或者给角度本身规定正负号。

\par\addvspace{.4\baselineskip}
\noindent\begin{minipage}{\linewidth}

\textbf{例 M5-016：半无限长导线的磁场}

一条通电流 \(I\) 的细直导线从 \(z=0\) 沿 \(+z\) 半无限延伸，电流指向
\(+z\)。场点位于 \((a,0,0)\)，其中
\(a>0\)。在真空中求该段导线的磁场大小和方向，并与同轴无限长导线比较。回流支路足够远，仅讨论给定段的贡献。

\end{minipage}\par

\par\noindent\begin{minipage}{\linewidth}
\textbf{解答}

\begin{center}
\includegraphics[width=\linewidth,height=.40\textheight,keepaspectratio]{\detokenize{assets/module-05/examples/M5-016.pdf}}
\par\smallskip{\small 图 M5-016：半无限长导线的磁场}
\end{center}
\end{minipage}\par

图 M5-016 的垂足正是半无限导线的端点；场点到端点的距离为
\(a\)。端点因子为 \(0\) 和 \(1\)，故结果是无限长导线的一半。

导线沿 \(+z\) 通电，观测点在正 \(x\) 轴。每个电流元的
\(d\boldsymbol\ell\times(\mathbf r-\mathbf r')\) 均沿 \(+y\)；在图的
\(xz\) 纸面中，正 \(y\) 方向为入纸面。因此带方向的结果是 \[
\boxed{\mathbf B(P)=\frac{\mu_0I}{4\pi a}\mathbf e_y}.
\] 大小为相同距离处无限长导线场的一半；方向须由电流元叉乘确定。

若一端正好位于垂足，另一端伸向无穷远，则
\(z_-=0\)、\(z_+\to+\infty\)，故

\[
B=\frac{\mu_0I}{4\pi a}.
\]

这恰好是无限长直导线场强的一半。

\subsubsection{2.
无限长直导线}\label{m5:ux65e0ux9650ux957fux76f4ux5bfcux7ebf}

令 \(z_-\to-\infty\)、\(z_+\to+\infty\)，便有

\[
B
=\frac{\mu_0I}{4\pi a}[1-(-1)]
=\boxed{\frac{\mu_0I}{2\pi a}}.
\]

磁场绕导线形成同心圆。这里的 \(1/a\)
衰减来自无限延伸的源；不能把它误当成所有局域电流分布的远场规律。

\subsection{三、圆弧、径向引线与组合电流}\label{m5:ux4e09ux5706ux5f27ux5f84ux5411ux5f15ux7ebfux4e0eux7ec4ux5408ux7535ux6d41}

半径为 \(R\)、圆心角为 \(\varphi\) 的圆弧通有电流 \(I\)。在圆心
\(O\)，所有电流元满足

\[
d\mathbf l'\perp\mathbf R,
\qquad
R=\text{常量},
\]

并且各 \(d\mathbf B\) 方向相同。由于弧长 \(s=R\varphi\)，其中
\(\varphi\) 必须用弧度，

\[
B_O
=\frac{\mu_0I}{4\pi R^2}
\int_{\text{圆弧}}dl'
=\boxed{
\frac{\mu_0I\varphi}{4\pi R}
}.
\]

方向垂直圆弧所在平面，由右手定则确定。

\par\noindent\begin{minipage}{\linewidth}
\begin{center}
\includegraphics[width=\linewidth,height=.40\textheight,keepaspectratio]{\detokenize{assets/module-05/theory/tikz-m5-topic42-v9-from39-v1-03-arc-radial.pdf}}
\par\smallskip{\small 图示：圆弧中心磁场与径向导线零贡献分别用平面图解释}
\end{center}
\end{minipage}\par

连接圆弧的径向直导线在圆心处满足
\(d\mathbf l'\parallel\mathbf R\)，因此贡献为零。这一性质常把看似复杂的扇形回路化成只计算圆弧。

\par\addvspace{.4\baselineskip}
\noindent\begin{minipage}{\linewidth}

\textbf{例 M5-017：半圆弧与径向引线}

半径为 \(R\)
的半圆弧位于纸面内，两端各接沿共同直径方向的径向直引线，稳恒电流沿半圆逆时针流动。求共同圆心
\(O\)
处的磁场，说明径向段为何无贡献；再讨论两条同心圆弧闭合成扇环时各段贡献如何叠加。

\end{minipage}\par

\par\noindent\begin{minipage}{\linewidth}
\textbf{解答}

\begin{center}
\includegraphics[width=\linewidth,height=.40\textheight,keepaspectratio]{\detokenize{assets/module-05/examples/M5-017.pdf}}
\par\smallskip{\small 图 M5-017：半圆弧与径向引线}
\end{center}
\end{minipage}\par

图 M5-017 中圆弧电流逆时针，故 \(O\) 处磁场向纸外。径向引线的电流元与到
\(O\) 的位移共线，只需积分半圆弧。

取 \(\varphi=\pi\)，两条径向线段均无贡献，故

\[
B_O=\frac{\mu_0I}{4R}.
\]

若电流沿内、外两条同心圆弧流动，并由径向线段连接，则两条圆弧中的电流绕向相反。设内、外半径分别为
\(R_1<R_2\)，圆心角均为 \(\varphi\)，则圆心场强大小为

\[
\boxed{
B_O
=\frac{\mu_0I\varphi}{4\pi}
\left(
\frac1{R_1}-\frac1{R_2}
\right)
}.
\]

方向由半径较小的圆弧决定。

\subsection{四、圆形电流环的轴线场}\label{m5:ux56dbux5706ux5f62ux7535ux6d41ux73afux7684ux8f74ux7ebfux573a}

半径为 \(R\) 的圆环通电流 \(I\)，场点 \(P\) 位于轴线上，距圆心
\(x\)。圆环上每个电流元到 \(P\) 的距离均为

\[
r=\sqrt{R^2+x^2}.
\]

而且 \(d\mathbf l'\) 与 \(\mathbf R\) 垂直，所以

\[
dB=\frac{\mu_0I}{4\pi}\frac{dl'}{r^2}.
\]

不能只写``由对称性''便直接给出答案。应把关于轴线对称的两个电流元配对：其垂直轴线的磁场分量大小相等、方向相反，轴向分量方向相同。

\par\noindent\begin{minipage}{\linewidth}
\begin{center}
\includegraphics[width=\linewidth,height=.40\textheight,keepaspectratio]{\detokenize{assets/module-05/theory/tikz-m5-topic42-v9-from39-v1-04-loop-double-view-margin12-v11.pdf}}
\par\smallskip{\small 图示：圆电流环用沿轴线观察和轴截面两幅平面图说明电流方向、横向分量抵消与轴向分量叠加，并给出轴线场曲线}
\end{center}
\end{minipage}\par

轴向投影因子为

\[
\cos\alpha=\frac{R}{r},
\]

故

\[
dB_x
=dB\cos\alpha
=\frac{\mu_0I}{4\pi}
\frac{R\,dl'}{(R^2+x^2)^{3/2}}.
\]

沿全环积分 \(\oint dl'=2\pi R\)，得到

\[
\boxed{
B_x(x)
=\frac{\mu_0IR^2}{2(R^2+x^2)^{3/2}}
}.
\]

多匝紧密线圈若每匝半径和位置相同，只需乘以匝数 \(N\)。圆心处

\[
\boxed{
B_0=\frac{\mu_0NI}{2R}
}.
\]

\par\addvspace{.4\baselineskip}
\noindent\begin{minipage}{\linewidth}

\textbf{例 M5-018：圆电流轴线上z=R处的磁场}

半径为 \(R\)、匝数为 \(N\) 的紧密圆线圈位于真空中，通稳恒电流
\(I\)。沿右手法向记轴坐标为 \(z\)，圆心为 \(z=0\)。求轴线上 \(z=R\)
处的场强及其与圆心场强之比。

\end{minipage}\par

\par\noindent\begin{minipage}{\linewidth}
\textbf{解答}

\begin{center}
\includegraphics[width=\linewidth,height=.40\textheight,keepaspectratio]{\detokenize{assets/module-05/examples/M5-018.pdf}}
\par\smallskip{\small 图 M5-018：圆电流轴线上z=R处的磁场}
\end{center}
\end{minipage}\par

图 M5-018 的 \(P\) 在圆线圈轴线上；令轴向距离等于线圈半径 \(R\)
后，几何因子是 \((R^2+R^2)^{3/2}\)。

代入上式，

\[
B(R)
=\frac{\mu_0NI}{2^{5/2}R}
=\frac{B_0}{2^{3/2}}.
\]

\subsubsection{1.
圆心附近的展开}\label{m5:ux5706ux5fc3ux9644ux8fd1ux7684ux5c55ux5f00}

把轴线场写成

\[
\frac{B_x}{B_0}
=\left(1+\frac{x^2}{R^2}\right)^{-3/2}.
\]

当 \(|x|\ll R\) 时，二项式展开给出

\[
\boxed{
\frac{B_x}{B_0}
=1-\frac32\frac{x^2}{R^2}
+\frac{15}{8}\frac{x^4}{R^4}
+O\!\left(\frac{x^6}{R^6}\right)
}.
\]

轴线场关于圆心是偶函数，因此没有奇次项；但二次项不为零，所以单个圆环只能在很小范围内近似产生匀强场。

\subsubsection{2.
远轴场与磁偶极矩}\label{m5:ux8fdcux8f74ux573aux4e0eux78c1ux5076ux6781ux77e9}

当 \(|x|\gg R\) 时，

\[
B_x
=\frac{\mu_0NIR^2}{2|x|^3}
\left[
1-\frac{3R^2}{2x^2}
+O\!\left(\frac{R^4}{x^4}\right)
\right].
\]

定义线圈磁偶极矩大小

\[
m=NI\pi R^2,
\]

则轴线远场的首项为

\[
\boxed{
B_x\simeq\frac{\mu_0m}{2\pi|x|^3}
}.
\]

这里得到的只是磁偶极子轴线场；完整空间方向将在专题四十五中由磁矢势统一推导。

\subsection{五、有限螺线管的轴线场}\label{m5:ux4e94ux6709ux9650ux87baux7ebfux7ba1ux7684ux8f74ux7ebfux573a}

设螺线管半径为 \(R\)，单位长度匝数为 \(n\)，电流为 \(I\)，两端坐标为
\(z_-\)、\(z_+\)，长度 \(L=z_+-z_-\)。取厚度为 \(dz'\) 的薄片，其中包含
\(n\,dz'\) 匝，可视为一个匝数为 \(n\,dz'\) 的圆环。

由圆环轴线场，薄片在轴线上坐标 \(z\) 处产生

\[
dB_z
=\frac{\mu_0(n\,dz')IR^2}
{2[R^2+(z-z')^2]^{3/2}}.
\]

\par\noindent\begin{minipage}{\linewidth}
\begin{center}
\includegraphics[width=\linewidth,height=.40\textheight,keepaspectratio]{\detokenize{assets/module-05/theory/tikz-m5-topic42-v9-from39-v1-05-finite-solenoid.pdf}}
\par\smallskip{\small 图示：有限螺线管在轴截面中分割成圆环薄片，再沿轴线叠加并检验中心、长管内部与端面中心三个极限}
\end{center}
\end{minipage}\par

积分得

\[
B_z(z)
=\frac{\mu_0nI}{2}
\int_{z_-}^{z_+}
\frac{R^2\,dz'}{[R^2+(z-z')^2]^{3/2}},
\]

即

\[
\boxed{
B_z(z)
=\frac{\mu_0nI}{2}
\left[
\frac{z-z_-}{\sqrt{R^2+(z-z_-)^2}}
-\frac{z-z_+}{\sqrt{R^2+(z-z_+)^2}}
\right]
}.
\]

若把两个端点到场点的连线与轴线夹角分别记为
\(\theta_-\)、\(\theta_+\)，并按图形给余弦规定方向，上式也常写成

\[
B_z=\frac{\mu_0nI}{2}
(\cos\theta_-+\cos\theta_+).
\]

端点式的符号更不容易因角度定义不同而混乱。

若螺线管关于原点对称，\(z_{\pm}=\pm L/2\)，中心磁场为

\[
\boxed{
B(0)
=\frac{\mu_0nIL}
{2\sqrt{R^2+(L/2)^2}}
}.
\]

\par\addvspace{.4\baselineskip}
\noindent\begin{minipage}{\linewidth}

\textbf{例 M5-019：长度等于直径的螺线管}

半径为 \(R\)、总长为 \(L=2R\) 的有限均匀螺线管，每单位长度有 \(n\)
匝，通稳恒电流
\(I\)，内部和外部均为真空。忽略导线截面，求轴线中心磁场，并与 \(L\gg R\)
时的中心及端面场比较。

\end{minipage}\par

\par\noindent\begin{minipage}{\linewidth}
\textbf{解答}

\begin{center}
\includegraphics[width=\linewidth,height=.40\textheight,keepaspectratio]{\detokenize{assets/module-05/examples/M5-019.pdf}}
\par\smallskip{\small 图 M5-019：长度等于直径的螺线管}
\end{center}
\end{minipage}\par

图 M5-019 的中心位于绕组两端之间；题中 \(L=2R\)
时，两端张角均不能忽略，不能直接用无限长螺线管的结果。

此时

\[
B(0)=\frac{\mu_0nI}{\sqrt2}.
\]

当 \(L\gg R\) 时，中心场趋于
\(\mu_0nI\)；若场点位于足够长螺线管的端面中心，则只有一侧线圈贡献，场趋于
\(\mu_0nI/2\)。这些极限既能检验积分，也说明理想无限长螺线管结论不应直接套到有限管口附近。

\subsection{六、亥姆霍兹线圈（拔高选讲）}\label{m5:ux516dux4ea5ux59c6ux970dux5179ux7ebfux5708ux62d4ux9ad8ux9009ux8bb2}

两只完全相同的同轴圆线圈，每只半径为 \(R\)、匝数为 \(N\)、电流为
\(I\)，圆心位于 \(z=\pm d/2\)，电流绕向相同。轴线上总场为

\[
B(z)
=\frac{\mu_0NIR^2}{2}
\left\{
\frac1{[R^2+(z-d/2)^2]^{3/2}}
+\frac1{[R^2+(z+d/2)^2]^{3/2}}
\right\}.
\]

\par\noindent\begin{minipage}{\linewidth}
\begin{center}
\includegraphics[width=\linewidth,height=.40\textheight,keepaspectratio]{\detokenize{assets/module-05/theory/tikz-m5-topic42-v9-from39-v1-06-helmholtz.pdf}}
\par\smallskip{\small 图示：两只同轴线圈的轴截面、电流方向和中心附近轴线场平台，说明间距等于半径时场最平坦}
\end{center}
\end{minipage}\par

由于 \(B(z)\) 是偶函数，自动有

\[
B'(0)=0.
\]

再求二阶导数：

\[
B''(0)
=\frac{3\mu_0NIR^2(d^2-R^2)}
{(R^2+d^2/4)^{7/2}}.
\]

要让中心附近二次不均匀项也消失，应取

\[
\boxed{d=R}.
\]

这就是亥姆霍兹线圈条件。此时

\[
\boxed{
B(0)
=\left(\frac45\right)^{3/2}
\frac{\mu_0NI}{R}
},
\]

并且由偶对称性和 \(B''(0)=0\)，中心附近满足

\[
B(z)=B(0)+O(z^4).
\]

\par\addvspace{.4\baselineskip}
\noindent\begin{minipage}{\linewidth}

\textbf{例 M5-020：反向亥姆霍兹线圈的场}

两只半径、匝数均相同的同轴圆线圈位于
\(z=\pm d/2\)，通等大反向电流，电流符号以同一个 \(+z\)
轴向场为准。求中心场并判断中心附近是近似匀强场还是梯度场；与同向亥姆霍兹线圈作比较。

\end{minipage}\par

\par\noindent\begin{minipage}{\linewidth}
\textbf{解答}

\begin{center}
\includegraphics[width=\linewidth,height=.40\textheight,keepaspectratio]{\detokenize{assets/module-05/examples/M5-020.pdf}}
\par\smallskip{\small 图 M5-020：反向亥姆霍兹线圈的场}
\end{center}
\end{minipage}\par

图 M5-020
两线圈电流代数符号相反，中心处的轴向场箭头抵消；离开中心后两者大小不等，留下的是梯度场。

若两线圈的电流反向，轴线场关于中心成为奇函数，中心处 \(B(0)=0\)，但
\(B'(0)\) 一般不为零，得到的是近似线性梯度场而不是近似匀强场。

\subsection{七、正多边形线圈（拔高选讲）}\label{m5:ux4e03ux6b63ux591aux8fb9ux5f62ux7ebfux5708ux62d4ux9ad8ux9009ux8bb2}

正 \(n\) 边形线圈的外接圆半径为 \(R\)，中心到每条边的垂直距离为

\[
\rho=R\cos\frac{\pi}{n}.
\]

对其中一条边，中心到两个端点的张角关于垂线对称，均为
\(\pi/n\)。由有限直导线公式，一条边产生

\[
B_1
=\frac{\mu_0I}{2\pi\rho}
\sin\frac{\pi}{n}.
\]

各边磁场方向相同，故

\[
\boxed{
B_n
=\frac{\mu_0In}{2\pi R}
\tan\frac{\pi}{n}
}.
\]

\par\noindent\begin{minipage}{\linewidth}
\begin{center}
\includegraphics[width=\linewidth,height=.40\textheight,keepaspectratio]{\detokenize{assets/module-05/theory/tikz-m5-topic42-v9-from39-v1-07-polygon-limit-safe-v7.pdf}}
\par\smallskip{\small 图示：以正六边形说明单边距离和中心磁场，再比较三角形、正方形、正六边形与正十二边形趋近圆环的过程}
\end{center}
\end{minipage}\par

\par\addvspace{.4\baselineskip}
\noindent\begin{minipage}{\linewidth}

\textbf{例 M5-021：正方形线圈的中心磁场}

真空中边长为 \(a\) 的单匝正方形线圈通逆时针稳恒电流 \(I\)，圆心为
\(O\)。求 \(O\)
处磁场，注明外接圆半径与边长的关系；并用固定外接圆半径的正多边形极限检验圆环结果。

\end{minipage}\par

\par\noindent\begin{minipage}{\linewidth}
\textbf{解答}

\begin{center}
\includegraphics[width=\linewidth,height=.40\textheight,keepaspectratio]{\detokenize{assets/module-05/examples/M5-021.pdf}}
\par\smallskip{\small 图 M5-021：正方形线圈的中心磁场}
\end{center}
\end{minipage}\par

图 M5-021 的四条边在圆心产生同向磁场，\(R\)
为外接圆半径。先求单边场再乘四，并以圆环极限检验。

正方形外接圆半径 \(R=a/\sqrt2\)，所以

\[
B_4
=\frac{2\mu_0I}{\pi R}
=\frac{2\sqrt2\mu_0I}{\pi a}.
\]

当 \(n\to\infty\) 时，利用 \(\tan(\pi/n)\sim\pi/n\)，得到

\[
\lim_{n\to\infty}B_n
=\frac{\mu_0I}{2R},
\]

正好回到半径为 \(R\) 的圆电流环中心场。

\subsection{八、毕奥---萨伐尔定律的使用条件}\label{m5:ux516bux6bd5ux5965ux8428ux4f10ux5c14ux5b9aux5f8bux7684ux4f7fux7528ux6761ux4ef6}

\subsubsection{1.
标准计算路线}\label{m5:ux6807ux51c6ux8ba1ux7b97ux8defux7ebf}

\begin{enumerate}
\def\labelenumi{\arabic{enumi}.}
\tightlist
\item
  明确电流是线分布、面分布还是体分布，并写出源点 \(\mathbf r'\)、场点
  \(\mathbf r\) 与 \(\mathbf R=\mathbf r-\mathbf r'\)。
\item
  先用 \(d\mathbf l'\times\mathbf R\)
  判定每个磁场微元的空间方向，再利用镜面对称、旋转对称或配对对称判断哪些分量抵消。
\item
  只对能够同向叠加的标量分量积分；几何量和积分上下限必须随源点位置一起变化。
\item
  用量纲、方向、特殊位置和远近极限检查结果。闭合局域电流回路的远场首项通常按
  \(1/r^3\) 衰减，而无限长直导线因源不局域而按 \(1/r\) 衰减。
\end{enumerate}

\subsubsection{2.
何时换用其他方法}\label{m5:ux4f55ux65f6ux6362ux7528ux5176ux4ed6ux65b9ux6cd5}

毕奥-萨伐尔定律适合``电流分布已知且几何积分可做''的问题。若场具有足够高的平移或旋转对称性，专题四十三的安培环路定理通常更快；若需要处理多个回路、远场多极展开或边界问题，专题四十五的磁矢势更系统。

\subsubsection{3.
不可忽略的边界}\label{m5:ux4e0dux53efux5ffdux7565ux7684ux8fb9ux754c}

\begin{enumerate}
\def\labelenumi{\arabic{enumi}.}
\tightlist
\item
  本专题公式按严格意义适用于稳恒电流。缓慢变化电流可在准静态条件下近似使用；变化快到传播延迟不可忽略时，需要使用含推迟时间的电磁场公式。
\item
  理想线电流在导线轴线上给出奇异场，不能用于求真实导线内部场。导线内部问题必须给出半径和
  \(\mathbf J\) 的分布。
\item
  有限导线公式必须包含完整闭合回路的其余部分。只有其余部分在所求点严格为零、由对称性抵消或有明确小参数时，才能忽略。
\item
  无限长直导线公式只适用于场点到端点的距离远大于横向距离的中段区域；靠近端点必须回到有限导线公式。
\item
  圆环近轴展开要求 \(|x|/R\ll1\)，远轴展开要求
  \(R/|x|\ll1\)。若小参数不小，保留首项可能比原公式更不准确。
\item
  把有限螺线管看成连续圆环要求线圈绕制足够密；若匝间距不可忽略，应逐匝叠加。
\end{enumerate}

\subsection{练习与解答}\label{m5:ux7ec3ux4e60ux4e0eux89e3ux7b54}

\par\addvspace{.4\baselineskip}
\noindent\begin{minipage}{\linewidth}

\subsubsection{作业一：圆形线圈中心场}\label{m5:ux4f5cux4e1aux4e00ux5706ux5f62ux7ebfux5708ux4e2dux5fc3ux573a}

\textbf{题号：M5-022}

\textbf{层级：} 基础　　\textbf{方法：} 毕奥---萨伐尔

线圈为单匝细线回路，空间取真空，导线截面和回流附加段的磁场影响忽略。

半径为 \(R\) 的圆形线圈通有电流 \(I\)，取线圈法向按右手定则为
\(+\mathbf n\)。求圆心处磁感应强度。

\end{minipage}\par

\par\noindent\begin{minipage}{\linewidth}
\textbf{解答}

\begin{center}
\includegraphics[width=\linewidth,height=.40\textheight,keepaspectratio]{\detokenize{assets/module-05/examples/M5-022.pdf}}
\par\smallskip{\small 图 M5-022：圆电流的中心场}
\end{center}
\end{minipage}\par

图 M5-022
的闭合圆周在本题表示电流路径，取逆时针电流时圆心场沿纸外正法向。各电流元到
\(O\) 的距离均为 \(R\)，贡献同向相加。

圆心到每个电流元的距离均为 \(R\)，且
\(d\mathbf l'\perp\mathbf R\)。毕奥-萨伐尔定律给出

\[
dB=\frac{\mu_0}{4\pi}
\frac{I\,dl'}{R^2}.
\]

各 \(d\mathbf B\) 均沿 \(+\mathbf n\)，因此可直接积分大小：

\[
\mathbf B(0)
=\frac{\mu_0I}{4\pi R^2}
\oint dl'\,\mathbf n
=\frac{\mu_0I}{4\pi R^2}(2\pi R)\mathbf n.
\]

故

\[
\boxed{
\mathbf B(0)=\frac{\mu_0I}{2R}\mathbf n
}.
\]

本题可直接把大小相加，是因为距离、夹角和方向都相同；离开圆心后，至少方向不再相同，必须先投影再积分。

\par\addvspace{.4\baselineskip}
\noindent\begin{minipage}{\linewidth}

\subsubsection{作业二：有限直导线及其极限}\label{m5:ux4f5cux4e1aux4e8cux6709ux9650ux76f4ux5bfcux7ebfux53caux5176ux6781ux9650}

\textbf{题号：M5-023}

\textbf{层级：} 基础　　\textbf{方法：} 毕奥---萨伐尔

导线沿 \(z\) 轴从 \(-L\) 到 \(+L\)，电流沿 \(+z\)，场点为
\((a,0,0)\)，空间取真空；仅计算该直线段的贡献。

长度为 \(2L\) 的直导线通有电流
\(I\)，场点位于导线中点的垂线上，到导线距离为 \(a\)。求磁场，并检验
\(L\ll a\) 与 \(L\gg a\) 两个极限。

\end{minipage}\par

\par\noindent\begin{minipage}{\linewidth}
\textbf{解答}

\begin{center}
\includegraphics[width=\linewidth,height=.40\textheight,keepaspectratio]{\detokenize{assets/module-05/examples/M5-023.pdf}}
\par\smallskip{\small 图 M5-023：有限直导线及其极限}
\end{center}
\end{minipage}\par

图 M5-023 中上下端点分别为 \(z=\pm L\)，场点 \(P\)
位于中垂线上。两端张角相同，所以积分的两个端点项相加。

取中点为 \(z=0\)，两端为 \(z_-= -L\)、\(z_+=L\)。代入有符号端点式：

\[
B
=\frac{\mu_0I}{4\pi a}
\left[
\frac{z}{\sqrt{a^2+z^2}}
\right]_{-L}^{L}.
\]

于是

\[
\boxed{
B
=\frac{\mu_0IL}
{2\pi a\sqrt{a^2+L^2}}
}.
\]

方向垂直于导线和场点构成的平面，由右手定则确定。

当 \(L\ll a\) 时，

\[
B\simeq\frac{\mu_0I(2L)}{4\pi a^2}.
\]

这与把整段导线近似看成距离为 \(a\)、总长度为 \(2L\) 的短电流段一致。当
\(L\gg a\) 时，

\[
B\to\frac{\mu_0I}{2\pi a},
\]

回到无限长直导线结果。

\par\addvspace{.4\baselineskip}
\noindent\begin{minipage}{\linewidth}

\subsubsection{作业三：同心圆弧组成的扇环}\label{m5:ux4f5cux4e1aux4e09ux540cux5fc3ux5706ux5f27ux7ec4ux6210ux7684ux6247ux73af}

\textbf{题号：M5-024}

\textbf{层级：} 基础　　\textbf{方法：} 毕奥---萨伐尔

取
\(0<\varphi<2\pi\)，真空中忽略导线截面；为明确方向，规定外弧电流沿逆时针、内弧沿顺时针。

两条同心圆弧的半径分别为 \(R_1<R_2\)，圆心角均为
\(\varphi\)，两端由径向直导线连接成闭合回路，通电流
\(I\)。求共同圆心处磁场。

\end{minipage}\par

\par\noindent\begin{minipage}{\linewidth}
\textbf{解答}

\begin{center}
\includegraphics[width=\linewidth,height=.40\textheight,keepaspectratio]{\detokenize{assets/module-05/examples/M5-024.pdf}}
\par\smallskip{\small 图 M5-024：同心圆弧组成的扇环}
\end{center}
\end{minipage}\par

图 M5-024 的外弧逆时针、内弧顺时针；两径向连接段在 \(O\)
无贡献。内弧半径较小，所以净场沿内弧对应的纸内方向。

两条径向直导线的 \(d\mathbf l'\) 都与指向圆心的 \(\mathbf R\) 共线，所以

\[
d\mathbf l'\times\mathbf R=\mathbf0,
\]

没有贡献。闭合回路中内、外圆弧电流绕向相反，磁场相减。内圆弧贡献较大，故总场方向由内圆弧决定，大小为

\[
B_O
=\frac{\mu_0I\varphi}{4\pi R_1}
-\frac{\mu_0I\varphi}{4\pi R_2}.
\]

因此

\[
\boxed{
B_O
=\frac{\mu_0I\varphi}{4\pi}
\left(
\frac1{R_1}-\frac1{R_2}
\right)
}.
\]

若
\(R_1\to R_2\)，两圆弧几乎重合且电流反向，总场趋零；这是必要的极限检验。

\par\addvspace{.4\baselineskip}
\noindent\begin{minipage}{\linewidth}

\subsubsection{作业四：有限螺线管中心场}\label{m5:ux4f5cux4e1aux56dbux6709ux9650ux87baux7ebfux7ba1ux4e2dux5fc3ux573a}

\textbf{题号：M5-025}

\textbf{层级：} 基础　　\textbf{方法：} 毕奥---萨伐尔

空间为真空，绕组以连续均匀电流面近似，中心设为轴坐标原点。

半径为 \(R\)、长度为 \(L\)、单位长度匝数为 \(n\) 的有限螺线管通电流
\(I\)。由圆环轴线场出发求中心磁场，并给出 \(L\gg R\) 的极限。

\end{minipage}\par

\par\noindent\begin{minipage}{\linewidth}
\textbf{解答}

\begin{center}
\includegraphics[width=\linewidth,height=.40\textheight,keepaspectratio]{\detokenize{assets/module-05/examples/M5-025.pdf}}
\par\smallskip{\small 图 M5-025：有限螺线管中心场}
\end{center}
\end{minipage}\par

图 M5-025 把绕组看作沿轴分布的圆环，中心为积分原点；每一薄片含 \(n\,dz\)
匝，积分区间为 \([-L/2,L/2]\)。

把轴向坐标为 \(z'\)、厚度为 \(dz'\) 的薄片看成 \(n\,dz'\)
匝圆环。中心位于 \(z=0\)，故

\[
dB
=\frac{\mu_0(n\,dz')IR^2}
{2(R^2+z'^2)^{3/2}}.
\]

\par\noindent\begin{minipage}{\linewidth}
\begin{center}
\includegraphics[width=\linewidth,height=.40\textheight,keepaspectratio]{\detokenize{assets/module-05/process/M5-025a.pdf}}
\par\smallskip{\small 图 M5-025a：把连续绕组分成圆电流薄片}
\end{center}
\end{minipage}\par

图 M5-025a
中每一条蓝色细环代表一个薄绕组片；各片在中心的轴向场同向，积分范围对中心对称。

积分范围为 \(-L/2\le z'\le L/2\)：

\[
B(0)
=\frac{\mu_0nI}{2}
\int_{-L/2}^{L/2}
\frac{R^2\,dz'}{(R^2+z'^2)^{3/2}}.
\]

利用

\[
\int\frac{R^2\,dz'}{(R^2+z'^2)^{3/2}}
=\frac{z'}{\sqrt{R^2+z'^2}},
\]

得到

\[
\boxed{
B(0)
=\frac{\mu_0nIL}
{2\sqrt{R^2+(L/2)^2}}
}.
\]

当 \(L\gg R\) 时，分母趋于 \(L/2\)，因此

\[
\boxed{B(0)\to\mu_0nI}.
\]

\par\addvspace{.4\baselineskip}
\noindent\begin{minipage}{\linewidth}

\subsubsection{作业五：亥姆霍兹条件}\label{m5:ux4f5cux4e1aux4e94ux4ea5ux59c6ux970dux5179ux6761ux4ef6}

\textbf{题号：M5-026}

\textbf{层级：} 基础　　\textbf{方法：} 毕奥---萨伐尔

两线圈均为细而紧密的圆绕组，空间取真空，场只在共同轴线上讨论；取中心为
\(z=0\)。

两只半径均为 \(R\)、匝数均为 \(N\) 的同轴圆线圈，圆心间距为
\(d\)，通同向电流 \(I\)。证明使中心附近磁场二次不均匀项消失的条件为
\(d=R\)，并求中心场。

\end{minipage}\par

\par\noindent\begin{minipage}{\linewidth}
\textbf{解答}

\begin{center}
\includegraphics[width=\linewidth,height=.40\textheight,keepaspectratio]{\detokenize{assets/module-05/examples/M5-026.pdf}}
\par\smallskip{\small 图 M5-026：亥姆霍兹条件}
\end{center}
\end{minipage}\par

图 M5-026
两绕组的轴向场同向，中心一阶导数由反射对称性自动消失。令二阶项也为零才得到亥姆霍兹间距
\(d=R\)。

两线圈圆心位于 \(z=\pm d/2\)，轴线场为

\[
B(z)
=\frac{\mu_0NIR^2}{2}
\left\{
[R^2+(z-d/2)^2]^{-3/2}
+[R^2+(z+d/2)^2]^{-3/2}
\right\}.
\]

该函数为偶函数，所以 \(B'(0)=0\)。对

\[
f(z)=[R^2+(z-a)^2]^{-3/2}
\]

求两次导数，并取 \(a=d/2\)、\(z=0\)，可得两线圈总场的二阶导数

\[
B''(0)
=\frac{3\mu_0NIR^2(d^2-R^2)}
{(R^2+d^2/4)^{7/2}}.
\]

\par\noindent\begin{minipage}{\linewidth}
\begin{center}
\includegraphics[width=\linewidth,height=.40\textheight,keepaspectratio]{\detokenize{assets/module-05/process/M5-026a.pdf}}
\par\smallskip{\small 图 M5-026a：中心二次项消失后的均匀区}
\end{center}
\end{minipage}\par

图 M5-026a
的中心平台来自两线圈轴线场相加；相对于中心的两侧等高，一阶与二阶项都消失。

要求 \(B''(0)=0\)，且 \(d>0\)，因此

\[
\boxed{d=R}.
\]

代入 \(z=0\)、\(d=R\)：

\[
B(0)
=\frac{\mu_0NIR^2}
{(R^2+R^2/4)^{3/2}}
=\boxed{
\left(\frac45\right)^{3/2}
\frac{\mu_0NI}{R}
}.
\]

\par\addvspace{.4\baselineskip}
\noindent\begin{minipage}{\linewidth}

\subsubsection{作业六：正多边形线圈与圆环极限}\label{m5:ux4f5cux4e1aux516dux6b63ux591aux8fb9ux5f62ux7ebfux5708ux4e0eux5706ux73afux6781ux9650}

\textbf{题号：M5-027}

\textbf{层级：} 基础　　\textbf{方法：} 毕奥---萨伐尔

\(n\ge3\)，回路为单匝细线且空间为真空；取沿电流的右手法向为磁场正方向。

正 \(n\) 边形线圈的外接圆半径为 \(R\)，通电流 \(I\)。求中心磁场，并证明
\(n\to\infty\) 时结果趋于圆电流环中心场。

\end{minipage}\par

\par\noindent\begin{minipage}{\linewidth}
\textbf{解答}

\begin{center}
\includegraphics[width=\linewidth,height=.40\textheight,keepaspectratio]{\detokenize{assets/module-05/examples/M5-027.pdf}}
\par\smallskip{\small 图 M5-027：正多边形线圈与圆环极限}
\end{center}
\end{minipage}\par

图 M5-027 用六边形示意一般正多边形；积分距离是圆心到边的垂距
\(\rho\)，不能把它误当外接圆半径 \(R\)。

中心到一条边的距离为

\[
\rho=R\cos\frac{\pi}{n}.
\]

该边两端相对于垂线的张角均为 \(\pi/n\)，所以一边产生

\[
B_1
=\frac{\mu_0I}{4\pi\rho}
\left(
2\sin\frac{\pi}{n}
\right)
=\frac{\mu_0I}{2\pi\rho}
\sin\frac{\pi}{n}.
\]

由旋转对称性，\(n\) 条边在中心产生的磁场方向相同，因此

\[
B_n=nB_1
=\boxed{
\frac{\mu_0In}{2\pi R}
\tan\frac{\pi}{n}
}.
\]

当 \(n\to\infty\) 时，\(\tan(\pi/n)\sim\pi/n\)，故

\[
\lim_{n\to\infty}B_n
=\frac{\mu_0I}{2R},
\]

这就是半径为 \(R\) 的圆形线圈中心场。

\section{专题四十三：安培环路定理、典型应用与磁边界条件（基础必讲与拔高选讲）}\label{m5:ux4e13ux9898ux56dbux5341ux4e09ux5b89ux57f9ux73afux8defux5b9aux7406ux5178ux578bux5e94ux7528ux4e0eux78c1ux8fb9ux754cux6761ux4ef6ux57faux7840ux5fc5ux8bb2ux4e0eux62d4ux9ad8ux9009ux8bb2}

安培环路定理给出磁场沿闭合曲线的\textbf{环量}。它对稳恒电流总成立；只有对称性足以确定路径上的切向场时，环量方程才能直接给出局部场强。界面处的场量跳变还要另由薄回路或边界条件求出。

\subsection{一、安培环路定理、定向与场源}\label{m5:ux4e00ux5b89ux57f9ux73afux8defux5b9aux7406ux5b9aux5411ux4e0eux573aux6e90}

\subsubsection{\texorpdfstring{1. 真空中的 \(\mathbf B\)
形式}{1. 真空中的 \textbackslash mathbf B 形式}}\label{m5:ux771fux7a7aux4e2dux7684-mathbf-b-ux5f62ux5f0f}

对稳恒电流，

\[
\boxed{\oint_C\mathbf B\cdot d\mathbf l
=\mu_0I_{\mathrm{enc}}
=\mu_0\int_S\mathbf J\cdot d\mathbf S}\text{。}
\]

其中 \(S\) 是以 \(C\) 为边界的任一跨越曲面。曲线 \(C\) 的正向与曲面法向
\(\mathbf n\) 必须按右手定则配对：四指沿 \(C\) 的正向弯曲时，拇指指向
\(\mathbf n\)。穿过 \(S\) 且与 \(\mathbf n\) 同向的电流取正，反向取负。

用 Stokes 定理

\[
\oint_C\mathbf B\cdot d\mathbf l
=\int_S(\nabla\times\mathbf B)\cdot d\mathbf S,
\]

可得局部形式

\[
\boxed{\nabla\times\mathbf B=\mu_0\mathbf J}\text{。}
\]

稳恒时
\(\nabla\cdot\mathbf J=0\)，因此同一闭合曲线所张的不同曲面给出相同的净穿越电流。若电场随时间变化，不能再单独使用上式；在真空中应改用麦克斯韦---安培定律

\[
\oint_C\mathbf B\cdot d\mathbf l
=\mu_0I_{\mathrm{enc}}
+\mu_0\varepsilon_0\frac{d\Phi_E}{dt}\text{。}
\]

\subsubsection{\texorpdfstring{2. 磁介质中的 \(\mathbf H\)
形式}{2. 磁介质中的 \textbackslash mathbf H 形式}}\label{m5:ux78c1ux4ecbux8d28ux4e2dux7684-mathbf-h-ux5f62ux5f0f}

宏观磁介质问题中，用 \(\mathbf H\) 只计自由电流更便利：

\[
\boxed{\oint_C\mathbf H\cdot d\mathbf l=I_{f,\mathrm{enc}}},
\qquad
\boxed{\nabla\times\mathbf H=\mathbf J_f}\text{。}
\]

\(\mathbf B\)
形式中的电流则应包括自由电流和束缚电流。在均匀线性各向同性介质中才可使用
\(\mathbf B=\mu\mathbf H\)；若 \(\mu\) 随位置变化，不能把一个统一的
\(\mu\) 从全环路上随意提出。

\par\addvspace{.4\baselineskip}
\noindent\begin{minipage}{\linewidth}

\textbf{例 M5-028：环路定向与包围电流}

有向闭曲线 \(C\)
按逆时针取向，配对曲面正法向指向纸外。穿过曲面的两条电流分别为纸外
\(I_1>0\)、纸内
\(I_2>0\)。写出真空静磁安培环路定理的代数右端，并讨论同时反转环路和法向时等式怎样变化。

\end{minipage}\par

\par\noindent\begin{minipage}{\linewidth}
\textbf{解答}

\begin{center}
\includegraphics[width=\linewidth,height=.40\textheight,keepaspectratio]{\detokenize{assets/module-05/examples/M5-028.pdf}}
\par\smallskip{\small 图 M5-028：环路定向与包围电流}
\end{center}
\end{minipage}\par

图 M5-028
中圆点电流沿正法向、叉号电流逆正法向，环量右端必须作代数相减；反转环路及法向后，两边同时反号。

若正向法线指向纸外，曲面内有纸外电流 \(I_1\) 和纸内电流 \(I_2\)，则
\(I_{\mathrm{enc}}=I_1-I_2\)。同时反转 \(C\) 与 \(\mathbf n\)
时，等式两边同时变号，物理结论不变。

\par\noindent\begin{minipage}{\linewidth}
\begin{center}
\includegraphics[width=\linewidth,height=.40\textheight,keepaspectratio]{\detokenize{assets/module-05/theory/tikz-m5-topic43-v9-from40-v1-01-orientation-symmetry-safe-v7.pdf}}
\par\smallskip{\small 图示：安培环量与包围电流}
\end{center}
\end{minipage}\par

\subsection{二、什么时候能用环路定理直接求场}\label{m5:ux4e8cux4ec0ux4e48ux65f6ux5019ux80fdux7528ux73afux8defux5b9aux7406ux76f4ux63a5ux6c42ux573a}

选取安培环路前应先做对称性分析：

\begin{enumerate}
\def\labelenumi{\arabic{enumi}.}
\tightlist
\item
  由旋转、平移或镜像对称确定场的方向。
\item
  选择场强模长为常量的路径段。
\item
  使 \(\mathbf B\parallel d\mathbf l\) 的路径段产生 \(B\,l\)，使
  \(\mathbf B\perp d\mathbf l\) 的路径段贡献为零。
\item
  分区计算环路真正包围的电流，而不是机械地代入全部电流。
\end{enumerate}

任意形状的环路仍然满足安培环路定理，但若 \(B\)
沿程变化且方向不定，一个环量等式不足以求出每一点的 \(B\)。

\par\addvspace{.4\baselineskip}
\noindent\begin{minipage}{\linewidth}

\textbf{例 M5-029：无限长直导线的环向场}

无限长细直导线沿 \(z\) 轴通电流
\(I>0\)，空间为真空。先用右手定则或毕奥---萨伐尔的叉乘方向确定磁场沿环向，再用平移与旋转对称性说明大小只依赖半径
\(r\)，再选同轴半径 \(r>0\) 的有向圆环求场强，说明为何可把 \(B\)
提出积分。

\end{minipage}\par

\par\noindent\begin{minipage}{\linewidth}
\textbf{解答}

\begin{center}
\includegraphics[width=\linewidth,height=.40\textheight,keepaspectratio]{\detokenize{assets/module-05/examples/M5-029.pdf}}
\par\smallskip{\small 图 M5-029：无限长直导线的环向场}
\end{center}
\end{minipage}\par

图 M5-029
的逆时针圆环与纸外电流组成右手取向。圆柱对称性保证场处处切向且等大，才可把
\(B\) 提出环量积分。

圆柱对称性保证 \(\mathbf B=B(r)\mathbf e_{\phi}\)。在半径为 \(r\)
的同轴圆上，

\[
B(r)(2\pi r)=\mu_0I,
\qquad
\boxed{B(r)=\frac{\mu_0I}{2\pi r}}\text{。}
\]

这与毕奥---萨伐尔积分的结果一致，但简化的根本原因是圆柱对称性，而不是圆形环路本身。

\subsection{三、一般柱对称电流与实心导线}\label{m5:ux4e09ux4e00ux822cux67f1ux5bf9ux79f0ux7535ux6d41ux4e0eux5b9eux5fc3ux5bfcux7ebf}

设无限长柱对称电流为

\[
\mathbf J=J_z(r)\mathbf e_z\text{。}
\]

对称性给出 \(\mathbf B=B_{\varphi}(r)\mathbf e_{\phi}\)。半径 \(r\)
的同轴圆包围的电流为

\[
I_{\mathrm{enc}}(r)=2\pi\int_0^rJ_z(s)s\,ds,
\]

因而

\[
\boxed{B_{\varphi}(r)
=\frac{\mu_0I_{\mathrm{enc}}(r)}{2\pi r}
=\frac{\mu_0}{r}\int_0^rJ_z(s)s\,ds}\text{。}
\]

若电流只分布在 \(r<R\)，则 \(r\ge R\) 时

\[
B_{\varphi}(r)=\frac{\mu_0I}{2\pi r}\text{。}
\]

对有限的 \(J_z(0)\)，轴线附近
\(B_{\varphi}\sim\mu_0J_z(0)r/2\)，磁场不会在轴线发散。反过来，若已知柱对称磁场，还可反求电流密度：

\[
\boxed{J_z(r)=\frac{1}{\mu_0r}\frac{d}{dr}\!\left[rB_{\varphi}(r)\right]}\text{。}
\]

\par\addvspace{.4\baselineskip}
\noindent\begin{minipage}{\linewidth}

\textbf{例 M5-030：均匀实心导线的分区场}

半径 \(R\) 的无限长实心圆柱导线沿 \(+z\) 均匀通总电流 \(I\)，磁导率取
\(\mu_0\)。分别求 \(0\le r<R\) 和 \(r\ge R\)
的环向磁场，并判断场强及其径向导数在材料表面是否连续。

\end{minipage}\par

\par\noindent\begin{minipage}{\linewidth}
\textbf{解答}

\begin{center}
\includegraphics[width=\linewidth,height=.40\textheight,keepaspectratio]{\detokenize{assets/module-05/examples/M5-030.pdf}}
\par\smallskip{\small 图 M5-030：均匀实心导线的分区场}
\end{center}
\end{minipage}\par

图 M5-030 的材料圆边界为
\(R\)；安培环半径移入导体后，只包围按面积比例分配的电流，故内外表达式不同。

对半径 \(R\) 、总电流 \(I\) 的均匀实心导线，

\[
B_{\varphi}(r)=
\begin{cases}
\dfrac{\mu_0Ir}{2\pi R^2},&0\le r<R,\\[6pt]
\dfrac{\mu_0I}{2\pi r},&r\ge R.
\end{cases}
\]

\(B\) 在 \(r=R\) 连续，但斜率改变；磁场在导线表面取最大值
\(\mu_0I/(2\pi R)\)。

\subsection{四、同轴电缆：回流如何封闭磁场}\label{m5:ux56dbux540cux8f74ux7535ux7f06ux56deux6d41ux5982ux4f55ux5c01ux95edux78c1ux573a}

设内导体半径为 \(a\)，均匀输送电流 \(I\)；外导体占据
\(b<r<c\)，并均匀返回电流 \(-I\)。对四个区域分别计算包围电流：

\[
I_{\mathrm{enc}}(r)=
\begin{cases}
I\dfrac{r^2}{a^2},&0\le r<a,\\[6pt]
I,&a\le r<b,\\[6pt]
I\left[1-\dfrac{r^2-b^2}{c^2-b^2}\right],&b\le r<c,\\[8pt]
0,&r\ge c.
\end{cases}
\]

于是

\[
\boxed{
B_{\varphi}(r)=\frac{\mu_0I}{2\pi}
\begin{cases}
\dfrac{r}{a^2},&0\le r<a,\\[6pt]
\dfrac1r,&a\le r<b,\\[6pt]
\dfrac{c^2-r^2}{r(c^2-b^2)},&b\le r<c,\\[8pt]
0,&r\ge c.
\end{cases}}
\]

因为这里把电流设为有限体电流，\(B(r)\) 在 \(a,b,c\)
处都连续。若把电流理想化为无厚度面电流，切向 \(B\) 则可发生有限跃变。

\par\addvspace{.4\baselineskip}
\noindent\begin{minipage}{\linewidth}

\textbf{例 M5-031：同轴回流使外场抵消}

理想无限长同轴电缆内导体输送电流 \(+I\)，外导体完整返回
\(-I\)，外表面半径为 \(c\)，电流分布轴对称。求 \(r>c\)
处磁场，并指出``净包围电流为零''与对称性在推理中各起什么作用。

\end{minipage}\par

\par\noindent\begin{minipage}{\linewidth}
\textbf{解答}

\begin{center}
\includegraphics[width=\linewidth,height=.40\textheight,keepaspectratio]{\detokenize{assets/module-05/examples/M5-031.pdf}}
\par\smallskip{\small 图 M5-031：同轴回流使外场抵消}
\end{center}
\end{minipage}\par

图 M5-031
中外侧安培环同时包围内导体出流与外导体回流；在圆柱对称条件下，净电流为零才进一步给出外场为零。

本题还必须使用理想无限长同轴结构的对称性：磁场沿环向，且同一同心圆周上等大，环量才能写成
\(2\pi rB_{\phi}\)。在包围两股等大反向电流的外部， \[
2\pi rB_{\phi}=\mu_0(I-I)=0,
\]
才得到局部场为零。一般几何中，净包围电流为零只保证环量为零，不能直接断言每一点磁场都为零。

当 \(r>c\) 时，环路包围的净电流为
\(I-I=0\)，故理想无限长同轴电缆的外部磁场严格为零。这是回流路径使磁场主要封闭在结构内部的典型例子。

\par\noindent\begin{minipage}{\linewidth}
\begin{center}
\includegraphics[width=\linewidth,height=.40\textheight,keepaspectratio]{\detokenize{assets/module-05/theory/tikz-m5-topic43-v9-from40-v1-03-coaxial-cable-margin12-v11.pdf}}
\par\smallskip{\small 同轴电缆的分区包围电流与磁场曲线}
\end{center}
\end{minipage}\par

\subsection{五、无限电流面与双电流面}\label{m5:ux4e94ux65e0ux9650ux7535ux6d41ux9762ux4e0eux53ccux7535ux6d41ux9762}

设无限平面上的面电流密度为 \(\mathbf K\)，单位为
\(\mathrm{A/m}\)。令单位法向 \(\mathbf n\)
由下侧指向上侧。对称性与窄矩形环路给出

\[
\boxed{\mathbf B_+=\frac{\mu_0}{2}\mathbf K\times\mathbf n},
\qquad
\boxed{\mathbf B_-=-\frac{\mu_0}{2}\mathbf K\times\mathbf n}\text{。}
\]

单面两侧的磁场大小都为 \(\mu_0K/2\)，方向相反。跨越电流面时，

\[
\mathbf n\times(\mathbf B_+-\mathbf B_-)=\mu_0\mathbf K\text{。}
\]

\par\addvspace{.4\baselineskip}
\noindent\begin{minipage}{\linewidth}

\textbf{例 M5-032：双电流面的叠加}

两张无限平行电流面位于 \(y=\pm d/2\)。上面面电流
\(K\mathbf e_z\)，下面为 \(-K\mathbf e_z\)，其中
\(K>0\)，介质是真空。求两面之间及外侧的磁场，并说明若改成同向电流，内部与外部的叠加关系怎样交换。

\end{minipage}\par

\par\noindent\begin{minipage}{\linewidth}
\textbf{解答}

\begin{center}
\includegraphics[width=\linewidth,height=.40\textheight,keepaspectratio]{\detokenize{assets/module-05/examples/M5-032.pdf}}
\par\smallskip{\small 图 M5-032：双电流面的叠加}
\end{center}
\end{minipage}\par

图 M5-032
的上层电流向纸外、下层向纸内；右手定则给出两层在中间的场都向右，外侧则相互抵消。

两个平行无限电流面的面电流大小均为
\(K\)、方向相反时，两面之间的磁场同向叠加为
\(\mu_0K\)，外部则相互抵消。若两面电流同向，结果恰好相反：内部抵消，外部叠加。

\par\noindent\begin{minipage}{\linewidth}
\begin{center}
\includegraphics[width=\linewidth,height=.40\textheight,keepaspectratio]{\detokenize{assets/module-05/theory/tikz-m5-topic43-v9-from40-v1-04-current-sheets.pdf}}
\par\smallskip{\small 单个无限电流面与反向双电流面的磁场叠加}
\end{center}
\end{minipage}\par

\subsection{六、理想长螺线管：环路定理先给内外场之差}\label{m5:ux516dux7406ux60f3ux957fux87baux7ebfux7ba1ux73afux8defux5b9aux7406ux5148ux7ed9ux5185ux5916ux573aux4e4bux5dee}

单位长度匝数为 \(n\) 的长螺线管通电流 \(I\)，紧密绕组可看成面电流
\(K_f=nI\)。取一条跨越管壁的矩形安培环路，其内外两条长边长度均为
\(l\)。两条短边与磁场垂直，不贡献环量，因而

\[
(H_{\mathrm{in}}-H_{\mathrm{out}})l=(nl)I,
\]

即

\[
\boxed{H_{\mathrm{in}}-H_{\mathrm{out}}=nI}\text{。}
\]

这一步\textbf{只得到内外场之差}。对无限长螺线管，再加上``没有外加均匀场，且距轴线足够远处的外部场趋零''这一边界条件，才有

\[
H_{\mathrm{out}}=0,
\qquad
\boxed{H_{\mathrm{in}}=nI}\text{。}
\]

真空中

\[
\boxed{B_{\mathrm{in}}=\mu_0nI}\text{。}
\]

若内部完全填充均匀线性磁介质，在理想无限长近似下

\[
\boxed{H_{\mathrm{in}}=nI,
\qquad B_{\mathrm{in}}=\mu nI}\text{。}
\]

\par\addvspace{.4\baselineskip}
\noindent\begin{minipage}{\linewidth}

\textbf{例 M5-033：同轴双层螺线管的叠加}

在同一线性磁介质中，同轴紧密绕制两层理想长螺线管，单位长度匝数为
\(n_1\text{、}n_2\)。共同轴向取 \(+z\)，各电流 \(I_1\text{、}I_2\)
按产生正轴向场的绕向取代数值。求共同内部区域的 \(H\) 和 \(B\)，磁导率为
\(\mu\)。

\end{minipage}\par

\par\noindent\begin{minipage}{\linewidth}
\textbf{解答}

\begin{center}
\includegraphics[width=\linewidth,height=.40\textheight,keepaspectratio]{\detokenize{assets/module-05/examples/M5-033.pdf}}
\par\smallskip{\small 图 M5-033：同轴双层螺线管的叠加}
\end{center}
\end{minipage}\par

图 M5-033
中蓝、橙两层绕组共享轴线；先按各自绕向确定轴向场的符号，再把两个 \(nI\)
项相加。

若同轴紧密绕有两层长螺线管，单位长度匝数分别为
\(n_1,n_2\)，电流按给定轴向磁场的右手定则取代数值 \(I_1,I_2\)，则

\[
H_{\mathrm{in}}=n_1I_1+n_2I_2\text{。}
\]

线性叠加时必须先判断每层绕组的方向，再代数相加。

有限长螺线管只在远离端部的中段近似满足上述结果；端部附近的场线弯曲且外部场不为零，应使用专题四十二的有限螺线管公式或直接积分。

\subsection{七、理想环形螺线管：必须分区讨论}\label{m5:ux4e03ux7406ux60f3ux73afux5f62ux87baux7ebfux7ba1ux5fc5ux987bux5206ux533aux8ba8ux8bba}

设理想紧密绕制的环形螺线管共有 \(N\) 匝，电流为
\(I\)，磁芯内、外半径分别为 \(R_1,R_2\)。旋转对称性表明
\(\mathbf H=H_{\varphi}(r)\mathbf e_{\phi}\)。但不同半径的同心圆链住的净电流不同：

\[
H_{\varphi}(r)=
\begin{cases}
0,&r<R_1,\\[4pt]
\dfrac{NI}{2\pi r},&R_1<r<R_2,\\[8pt]
0,&r>R_2,
\end{cases}
\]

其中内孔与外部的结果采用了理想连续、紧密绕组的对称近似。在均匀线性磁芯中

\[
\boxed{B_{\varphi}(r)=\frac{\mu NI}{2\pi r}},
\qquad R_1<r<R_2\text{。}
\]

若磁导率沿半径变化，自由电流决定的仍是

\[
H_{\varphi}(r)=\frac{NI}{2\pi r},
\]

而

\[
B_{\varphi}(r)=\mu(r)H_{\varphi}(r)\text{。}
\]

只有当磁芯径向厚度远小于平均半径 \(R\)时，才可用

\[
B\approx\frac{\mu NI}{2\pi R}
\]

代替严格的 \(1/r\) 分布。

\par\addvspace{.4\baselineskip}
\noindent\begin{minipage}{\linewidth}

\textbf{例 M5-034：厚环形磁芯的内外场比}

线性各向同性厚环形磁芯占据 \(R_1<r<R_2\)，环向均匀绕有 \(N\)
匝线圈，通电流 \(I\)。忽略漏磁及端部效应，求内侧 \(R_1^+\) 与外侧
\(R_2^-\) 的磁场比值，并判断能否将厚磁芯中的场视为常量。

\end{minipage}\par

\par\noindent\begin{minipage}{\linewidth}
\textbf{解答}

\begin{center}
\includegraphics[width=\linewidth,height=.40\textheight,keepaspectratio]{\detokenize{assets/module-05/examples/M5-034.pdf}}
\par\smallskip{\small 图 M5-034：厚环形磁芯的内外场比}
\end{center}
\end{minipage}\par

图 M5-034 的磁芯位于
\(R_1<r<R_2\)，每条同轴环路包围同样的安匝数，但周长随 \(r\)
增大，场强因此按 \(1/r\) 减小。

均匀磁芯内侧与外侧的磁场比值为

\[
\frac{B(R_1^+)}{B(R_2^-)}=\frac{R_2}{R_1}\text{。}
\]

因而对厚环形磁芯，``磁芯内磁场处处相同''不成立。

\subsection{\texorpdfstring{八、磁边界条件：法向看 \(\mathbf B\)，切向看
\(\mathbf H\)}{八、磁边界条件：法向看 \textbackslash mathbf B，切向看 \textbackslash mathbf H}}\label{m5:ux516bux78c1ux8fb9ux754cux6761ux4ef6ux6cd5ux5411ux770b-mathbf-bux5207ux5411ux770b-mathbf-h}

令单位法向 \(\mathbf n\) 从介质 1 指向介质 2。

\subsubsection{\texorpdfstring{1. \(\mathbf B\)
的法向分量}{1. \textbackslash mathbf B 的法向分量}}\label{m5:mathbf-b-ux7684ux6cd5ux5411ux5206ux91cf}

由
\(\nabla\cdot\mathbf B=0\)，在界面上取高度趋零的薄高斯柱，侧面通量趋零，故

\[
\boxed{\mathbf n\cdot(\mathbf B_2-\mathbf B_1)=0},
\qquad
\boxed{B_{2n}=B_{1n}}\text{。}
\]

这一结论与界面是否有面电流无关；它表明不存在孤立磁荷，磁通量不会在界面上凭空中断。

\subsubsection{\texorpdfstring{2. \(\mathbf H\)
的切向分量}{2. \textbackslash mathbf H 的切向分量}}\label{m5:mathbf-h-ux7684ux5207ux5411ux5206ux91cf}

由 \(\nabla\times\mathbf H=\mathbf J_f\)，取高度趋零的跨界窄矩形回路，得

\[
\boxed{\mathbf n\times(\mathbf H_2-\mathbf H_1)=\mathbf K_f}\text{。}
\]

若界面上无自由面电流，则

\[
\boxed{H_{2t}=H_{1t}}\text{。}
\]

对两侧均为线性各向同性介质的情形，

\[
\boxed{\frac{B_{2t}}{\mu_2}=\frac{B_{1t}}{\mu_1}}\text{。}
\]

注意：没有自由面电流时连续的是 \(H_t\)，而不一定是
\(B_t\)。磁化产生的束缚面电流已经包含在介质的磁化响应中。若直接用
\(\mathbf B\) 和总面电流 \(\mathbf K_{\mathrm{tot}}\)，则磁静态边界式为

\[
\mathbf n\times(\mathbf B_2-\mathbf B_1)=\mu_0\mathbf K_{\mathrm{tot}}\text{。}
\]

\par\addvspace{.4\baselineskip}
\noindent\begin{minipage}{\linewidth}

\textbf{例 M5-035：磁边界的法向与切向}

两线性各向同性磁介质的磁导率为
\(\mu_1\text{、}\mu_2\)，静止平面界面无自由面电流。分别讨论磁场完全垂直界面、完全平行界面两种极限：哪些
\(B\) 分量连续，哪些 \(H\) 分量连续，另一量怎样变化？

\end{minipage}\par

\par\noindent\begin{minipage}{\linewidth}
\textbf{解答}

\begin{center}
\includegraphics[width=\linewidth,height=.40\textheight,keepaspectratio]{\detokenize{assets/module-05/examples/M5-035.pdf}}
\par\smallskip{\small 图 M5-035：磁边界的法向与切向}
\end{center}
\end{minipage}\par

图 M5-035 的虚线为界面法线，实线为界面；法向比较 \(B\)，切向比较
\(H\)，再用两侧各自的磁导率转换。

无自由面电流时，若磁场垂直入射界面，则 \(B_n\) 相等而 \(H_n=B_n/\mu\)
一般不等；若磁场完全沿界面切向，则 \(H_t\) 相等而 \(B_t\)
按磁导率成比例变化。

\par\noindent\begin{minipage}{\linewidth}
\begin{center}
\includegraphics[width=\linewidth,height=.40\textheight,keepaspectratio]{\detokenize{assets/module-05/theory/tikz-m5-topic43-v9-from40-v1-07-boundary-conditions.pdf}}
\par\smallskip{\small 图示：磁介质界面的法向与切向条件}
\end{center}
\end{minipage}\par

\subsection{九、磁力线在不同磁导率介质间的折转}\label{m5:ux4e5dux78c1ux529bux7ebfux5728ux4e0dux540cux78c1ux5bfcux7387ux4ecbux8d28ux95f4ux7684ux6298ux8f6c}

设界面上无自由面电流，两侧为线性各向同性介质。令 \(\alpha_1,\alpha_2\)
分别为 \(\mathbf B_1,\mathbf B_2\) 与界面法线的夹角。由

\[
B_{2n}=B_{1n},
\qquad
\frac{B_{2t}}{\mu_2}=\frac{B_{1t}}{\mu_1},
\]

得

\[
\boxed{\frac{\tan\alpha_2}{\tan\alpha_1}
=\frac{\mu_2}{\mu_1}}\text{。}
\]

因此，当 \(\mu_2>\mu_1\) 时，按``相对法线的夹角''衡量，介质 2
中的磁力线更偏向界面切向。这一说法依赖于角度的定义；若角度从界面量起，比值将取倒数。

\par\addvspace{.4\baselineskip}
\noindent\begin{minipage}{\linewidth}

\textbf{例 M5-036：磁力线折转的角度}

两种正磁导率线性各向同性介质满足
\(\mu_2=4\mu_1\)，界面无自由面电流。介质1中的磁场与由1指向2的法线夹角为
\(\alpha_1=30^{\circ}\)。由边界条件求介质2中的角度
\(\alpha_2\)，两个角均从法线量起。

\end{minipage}\par

\par\noindent\begin{minipage}{\linewidth}
\textbf{解答}

\begin{center}
\includegraphics[width=\linewidth,height=.40\textheight,keepaspectratio]{\detokenize{assets/module-05/examples/M5-036.pdf}}
\par\smallskip{\small 图 M5-036：磁力线折转的角度}
\end{center}
\end{minipage}\par

图 M5-036 中两个角都从同一法线量起；进入磁导率较大的介质后，切向 \(B\)
增强而法向 \(B\) 连续，所以 \(\alpha_2\) 增大。

若 \(\mu_2=4\mu_1\)、\(\alpha_1=30^{\circ}\)，则

\[
\tan\alpha_2=4\tan30^{\circ}=\frac4{\sqrt3},
\]

故 \(\alpha_2\approx66.6^{\circ}\)。

\subsection{十、使用条件与常见错误}\label{m5:ux5341ux4f7fux7528ux6761ux4ef6ux4e0eux5e38ux89c1ux9519ux8bef}

\begin{enumerate}
\def\labelenumi{\arabic{enumi}.}
\tightlist
\item
  \textbf{定理成立不等于能直接求场。}
  安培环路定理对稳恒电流总成立，但只有高对称场才能把 \(B\) 或 \(H\)
  提出积分号。
\item
  \textbf{必须区分 \(\mathbf B\) 与 \(\mathbf H\) 的场源。}
  \(\mathbf H\) 的环量只计自由电流；\(\mathbf B\) 形式中应计入总电流。
\item
  \textbf{必须先分区，再算包围电流。}
  实心导线、同轴电缆和环形螺线管都不能用一个未分区公式覆盖所有 \(r\)。
\item
  \textbf{安培环路先给的可能只是差值。} 长螺线管中先得
  \(H_{\mathrm{in}}-H_{\mathrm{out}}=nI\)，还需要无限长极限的边界条件才能令
  \(H_{\mathrm{out}}=0\)。
\item
  \textbf{面电流可使切向场跳变。}
  体电流分布在有限厚度内时，磁场通常连续；零厚度面电流是场量有限跳变的理想化来源。
\item
  \textbf{时变情形不能遗漏位移电流。}
  充电电容器等问题必须使用麦克斯韦---安培定律。
\end{enumerate}

\subsection{练习与解答}\label{m5:ux7ec3ux4e60ux4e0eux89e3ux7b54-1}

\par\addvspace{.4\baselineskip}
\noindent\begin{minipage}{\linewidth}

\subsubsection{作业一：均匀载流实心圆柱}\label{m5:ux4f5cux4e1aux4e00ux5747ux5300ux8f7dux6d41ux5b9eux5fc3ux5706ux67f1}

\textbf{题号：M5-037}

\textbf{层级：} 基础　　\textbf{方法：} 安培环路、分区讨论

磁导率取
\(\mu_0\)，电流密度无环向分量，忽略端部效应；用正向电流的右手环向定义
\(+\mathbf e_{\phi}\)。

半径为 \(R\) 的无限长实心圆柱导线均匀通以总电流
\(I\)。求柱内外磁场的大小，并说明为何本题可用安培环路定理直接解出
\(B\)。

\end{minipage}\par

\par\noindent\begin{minipage}{\linewidth}
\textbf{解答}

\begin{center}
\includegraphics[width=\linewidth,height=.40\textheight,keepaspectratio]{\detokenize{assets/module-05/examples/M5-037.pdf}}
\par\smallskip{\small 图 M5-037：均匀载流实心圆柱}
\end{center}
\end{minipage}\par

图 M5-037 中 \(R\)
是载流区半径，红色同轴环路可以分别取在柱内和柱外。由旋转和平移对称性先确定
\(B\) 沿环向，再求环量。

取与导线同轴、半径为 \(r\)
的圆形环路。圆柱对称性使磁场沿环路切向且环上大小相同，故

\[
B(2\pi r)=\mu_0I_{\mathrm{enc}}(r)\text{。}
\]

当 \(r<R\) 时，\(J=I/(\pi R^2)\)，因而

\[
I_{\mathrm{enc}}(r)=J\pi r^2=I\frac{r^2}{R^2},
\]

\[
\boxed{B(r)=\frac{\mu_0Ir}{2\pi R^2}},
\qquad r<R\text{。}
\]

当 \(r\ge R\) 时，\(I_{\mathrm{enc}}=I\)，故

\[
\boxed{B(r)=\frac{\mu_0I}{2\pi r}},
\qquad r\ge R\text{。}
\]

两式在 \(r=R\)
相等。安培环路定理对任意稳恒电流都成立，但本题之所以能直接求出
\(B\)，是因为圆柱对称性保证 \(\mathbf B\parallel d\mathbf l\)
且同一圆周上 \(B\) 为常量。

\par\addvspace{.4\baselineskip}
\noindent\begin{minipage}{\linewidth}

\subsubsection{作业二：同轴电缆的四个区域}\label{m5:ux4f5cux4e1aux4e8cux540cux8f74ux7535ux7f06ux7684ux56dbux4e2aux533aux57df}

\textbf{题号：M5-038}

\textbf{层级：} 基础　　\textbf{方法：} 安培环路

设 \(0<a<b<c\)，内外导体均沿轴向均匀载流，介质磁导率取
\(\mu_0\)，电缆无限长。

内导体半径为 \(a\)，均匀输送电流 \(I\)；外导体占据
\(b<r<c\)，均匀返回电流 \(-I\)。求全空间的磁场。

\end{minipage}\par

\par\noindent\begin{minipage}{\linewidth}
\textbf{解答}

\begin{center}
\includegraphics[width=\linewidth,height=.40\textheight,keepaspectratio]{\detokenize{assets/module-05/examples/M5-038.pdf}}
\par\smallskip{\small 图 M5-038：同轴电缆的四个区域}
\end{center}
\end{minipage}\par

图 M5-038 的 \(a\text{、}b\text{、}c\)
划分四区。外导体内只需扣除环路已经包围的回流部分；到 \(r>c\)
时才恰好扣完总电流。

圆柱对称性给出 \(\mathbf B=B_{\varphi}(r)\mathbf e_{\phi}\)，且

\[
B_{\varphi}(r)=\frac{\mu_0I_{\mathrm{enc}}(r)}{2\pi r}\text{。}
\]

内导体中 \(J_1=I/(\pi a^2)\)；外导体中 \(J_2=-I/[\pi(c^2-b^2)]\)。因此

\[
I_{\mathrm{enc}}(r)=
\begin{cases}
I r^2/a^2,&0\le r<a,\\
I,&a\le r<b,\\
I-I(r^2-b^2)/(c^2-b^2),&b\le r<c,\\
0,&r\ge c.
\end{cases}
\]

\par\noindent\begin{minipage}{\linewidth}
\begin{center}
\includegraphics[width=\linewidth,height=.40\textheight,keepaspectratio]{\detokenize{assets/module-05/process/M5-038a.pdf}}
\par\smallskip{\small 图 M5-038a：同轴电缆四区的连续磁场曲线}
\end{center}
\end{minipage}\par

图 M5-038a 用 \(b=2a\text{、}c=3a\)
的几何比例示意四区形状；折点分别对应导体界面，场在所有界面连续。

代入得

\[
\boxed{
B_{\varphi}(r)=\frac{\mu_0I}{2\pi}
\begin{cases}
r/a^2,&0\le r<a,\\
1/r,&a\le r<b,\\
(c^2-r^2)/[r(c^2-b^2)],&b\le r<c,\\
0,&r\ge c.
\end{cases}}
\]

四段在 \(a,b,c\) 处两两连续。

\par\addvspace{.4\baselineskip}
\noindent\begin{minipage}{\linewidth}

\subsubsection{作业三：反向平行电流面}\label{m5:ux4f5cux4e1aux4e09ux53cdux5411ux5e73ux884cux7535ux6d41ux9762}

\textbf{题号：M5-039}

\textbf{层级：} 基础　　\textbf{方法：} 安培环路

空间为真空，\(d\text{、}K>0\)，采用标准右手直角坐标系。

两无限平行电流面分别位于 \(y=d/2\) 和 \(y=-d/2\)。上层面电流为
\(K\mathbf e_z\)，下层为 \(-K\mathbf e_z\)。求三个区域内的磁场。

\end{minipage}\par

\par\noindent\begin{minipage}{\linewidth}
\textbf{解答}

\begin{center}
\includegraphics[width=\linewidth,height=.40\textheight,keepaspectratio]{\detokenize{assets/module-05/examples/M5-039.pdf}}
\par\smallskip{\small 图 M5-039：反向平行电流面}
\end{center}
\end{minipage}\par

图 M5-039 中上层为 \(+z\) 电流、下层为 \(-z\) 电流，图的向右方向取
\(+x\)。分别叠加两电流面的场，三个区域的符号便确定。

对单个电流面，若取 \(\mathbf n=\mathbf e_y\) 由下指向上，则

\[
\mathbf B_+=\frac{\mu_0}{2}\mathbf K\times\mathbf n,
\qquad
\mathbf B_-=-\frac{\mu_0}{2}\mathbf K\times\mathbf n\text{。}
\]

在 \(y>d/2\) 区域，上层产生 \(-(\mu_0K/2)\mathbf e_x\)，下层产生
\(+(\mu_0K/2)\mathbf e_x\)，二者抵消。在 \(-d/2<y<d/2\) 区域，两层都产生
\(+(\mu_0K/2)\mathbf e_x\)。在 \(y<-d/2\) 区域，二者再次抵消。故

\[
\boxed{
\mathbf B=
\begin{cases}
\mathbf0,&y>d/2,\\
\mu_0K\mathbf e_x,&-d/2<y<d/2,\\
\mathbf0,&y<-d/2.
\end{cases}}
\]

\par\addvspace{.4\baselineskip}
\noindent\begin{minipage}{\linewidth}

\subsubsection{作业四：厚环形磁芯中的磁通量}\label{m5:ux4f5cux4e1aux56dbux539aux73afux5f62ux78c1ux82afux4e2dux7684ux78c1ux901aux91cf}

\textbf{题号：M5-040}

\textbf{层级：} 基础　　\textbf{方法：} 安培环路、磁通积分

磁导率为给定常量
\(\mu\)，\(0<R_1<R_2\)，忽略漏磁与边缘场；径向截面各处轴向高度均为
\(h\)。

均匀线性环形磁芯的内、外半径为 \(R_1,R_2\)，轴向高度为 \(h\)，绕有 \(N\)
匝线圈，通电流 \(I\)。求磁芯中的 \(B(r)\) 及穿过每一匝线圈的磁通量。

\end{minipage}\par

\par\noindent\begin{minipage}{\linewidth}
\textbf{解答}

\begin{center}
\includegraphics[width=\linewidth,height=.40\textheight,keepaspectratio]{\detokenize{assets/module-05/examples/M5-040.pdf}}
\par\smallskip{\small 图 M5-040：厚环形磁芯中的磁通量}
\end{center}
\end{minipage}\par

图 M5-040 的一匝磁通穿过磁芯的径向截面，面积元是 \(h\,dr\)；积分必须从
\(R_1\) 到 \(R_2\)，不能直接用某个平均场乘面积。

对 \(R_1<r<R_2\) 的同心圆环路，

\[
H_{\varphi}(r)(2\pi r)=NI,
\]

所以

\[
\boxed{B_{\varphi}(r)=\frac{\mu NI}{2\pi r}}\text{。}
\]

磁场沿环向，穿过磁芯径向---轴向截面。截取宽 \(dr\)、高 \(h\) 的面元
\(dS=h\,dr\)，得每匝所链磁通量

\[
\Phi=\int_{R_1}^{R_2}B_{\varphi}(r)h\,dr
=\frac{\mu NIh}{2\pi}\int_{R_1}^{R_2}\frac{dr}{r}\text{。}
\]

因而

\[
\boxed{\Phi=\frac{\mu NIh}{2\pi}\ln\frac{R_2}{R_1}}\text{。}
\]

这一结果也说明，厚磁芯不能用一个平均 \(B\) 粗略代替全截面积分。

\par\addvspace{.4\baselineskip}
\noindent\begin{minipage}{\linewidth}

\subsubsection{作业五：磁力线折转}\label{m5:ux4f5cux4e1aux4e94ux78c1ux529bux7ebfux6298ux8f6c}

\textbf{题号：M5-041}

\textbf{层级：} 基础　　\textbf{方法：} 安培环路

两个角均从界面法线量起，磁导率为正常数；忽略界面附近的自由体电流层。

两种线性各向同性磁介质的磁导率满足
\(\mu_2=4\mu_1\)，界面上无自由面电流。已知介质 1
中的磁场与界面法线夹角为 \(30^{\circ}\)，求介质 2 中的夹角。

\end{minipage}\par

\par\noindent\begin{minipage}{\linewidth}
\textbf{解答}

\begin{center}
\includegraphics[width=\linewidth,height=.40\textheight,keepaspectratio]{\detokenize{assets/module-05/examples/M5-041.pdf}}
\par\smallskip{\small 图 M5-041：磁力线折转}
\end{center}
\end{minipage}\par

图 M5-041 以虚线标出法线，防止把题给的 \(30^{\circ}\)
误认为与界面的夹角。法向与切向边界条件联立后给出折转角。

无自由面电流时，

\[
B_{2n}=B_{1n},
\qquad
\frac{B_{2t}}{\mu_2}=\frac{B_{1t}}{\mu_1}\text{。}
\]

因此

\[
\frac{\tan\alpha_2}{\tan\alpha_1}=\frac{\mu_2}{\mu_1}=4\text{。}
\]

\par\noindent\begin{minipage}{\linewidth}
\begin{center}
\includegraphics[width=\linewidth,height=.40\textheight,keepaspectratio]{\detokenize{assets/module-05/process/M5-041a.pdf}}
\par\smallskip{\small 图 M5-041a：磁边界两侧的分量对应}
\end{center}
\end{minipage}\par

图 M5-041a
的两个角都相对于虚线法向定义；把矢量分成界面法向和切向后，磁导率只进入切向
\(B\) 的比例。

代入 \(\alpha_1=30^{\circ}\)，

\[
\tan\alpha_2=\frac4{\sqrt3},
\]

故

\[
\boxed{\alpha_2=\arctan\frac4{\sqrt3}\approx66.6^{\circ}}\text{。}
\]

\par\addvspace{.4\baselineskip}
\noindent\begin{minipage}{\linewidth}

\subsubsection{作业六：由磁场反求电流密度}\label{m5:ux4f5cux4e1aux516dux7531ux78c1ux573aux53cdux6c42ux7535ux6d41ux5bc6ux5ea6}

\textbf{题号：M5-042}

\textbf{层级：} 基础　　\textbf{方法：} 安培环路

\(B_0\text{、}R>0\)，导线磁导率取
\(\mu_0\)，给定场仅有环向分量；表面电流按允许分布型跳变一并检验。

某无限长柱对称稳恒电流产生的磁场为

\[
B_{\varphi}(r)=
\begin{cases}
B_0(r/R)^3,&0\le r<R,\\[4pt]
B_0R/r,&r\ge R.
\end{cases}
\]

求体电流密度 \(J_z(r)\)、是否存在表面电流，以及总电流 \(I\)。

\end{minipage}\par

\par\noindent\begin{minipage}{\linewidth}
\textbf{解答}

\begin{center}
\includegraphics[width=\linewidth,height=.40\textheight,keepaspectratio]{\detokenize{assets/module-05/examples/M5-042.pdf}}
\par\smallskip{\small 图 M5-042：由磁场反求电流密度}
\end{center}
\end{minipage}\par

图 M5-042 的导体圆周处 \(B_{\phi}\) 连续，故无面电流跳变；在内部对
\(rB_{\phi}\) 求导，再对截面积积分检验总电流。

柱对称下

\[
J_z(r)=\frac1{\mu_0r}\frac{d}{dr}\left[rB_{\varphi}(r)\right]\text{。}
\]

对 \(r<R\)，

\[
rB_{\varphi}(r)=\frac{B_0r^4}{R^3},
\]

故

\[
\boxed{J_z(r)=\frac{4B_0r^2}{\mu_0R^3}},
\qquad r<R\text{。}
\]

对 \(r>R\)，\(rB_{\varphi}=B_0R\) 为常量，因而

\[
\boxed{J_z(r)=0},
\qquad r>R\text{。}
\]

内外两式在 \(r=R\) 都给出
\(B_{\varphi}=B_0\)，切向磁场无跃变，因此不存在额外的理想表面电流。由外部磁场

\[
B_0\frac Rr=\frac{\mu_0I}{2\pi r},
\]

得

\[
\boxed{I=\frac{2\pi RB_0}{\mu_0}}\text{。}
\]

也可直接积分体电流密度验算：

\[
2\pi\int_0^R J_z(r)r\,dr
=\frac{8\pi B_0}{\mu_0R^3}\int_0^Rr^3\,dr
=\frac{2\pi RB_0}{\mu_0}\text{。}
\]

\section{专题四十四：磁介质、磁化、束缚电流与退磁场（基础必讲与拔高选讲）}\label{m5:ux4e13ux9898ux56dbux5341ux56dbux78c1ux4ecbux8d28ux78c1ux5316ux675fux7f1aux7535ux6d41ux4e0eux9000ux78c1ux573aux57faux7840ux5fc5ux8bb2ux4e0eux62d4ux9ad8ux9009ux8bb2}

宏观磁场同时由自由电流和材料内部的磁偶极矩决定，因此不能把真空磁场简单乘以一个常数。由磁偶极矩的体密度可以推导束缚电流和磁标势两种表述；有限样品还会因表面磁化分布产生退磁场。

\subsection{一、从磁偶极矩到磁化强度}\label{m5:ux4e00ux4eceux78c1ux5076ux6781ux77e9ux5230ux78c1ux5316ux5f3aux5ea6}

\subsubsection{1.
电流回路的磁偶极矩}\label{m5:ux7535ux6d41ux56deux8defux7684ux78c1ux5076ux6781ux77e9}

面积为 \(S\)、电流为 \(I\) 的平面小回路，其磁偶极矩定义为

\[
\boxed{\mathbf m=IS\hat{\mathbf n}},
\]

其中 \(\hat{\mathbf n}\)
与电流方向按右手定则配对。一般局域稳恒电流分布的磁偶极矩为

\[
\boxed{\mathbf m=\frac12\int\mathbf r\times\mathbf J(\mathbf r)\,dV}.
\]

\subsubsection{2. 磁化强度}\label{m5:ux78c1ux5316ux5f3aux5ea6}

把足够小但仍含有大量微观磁偶极子的体积元记为 \(\Delta V\)，定义

\[
\boxed{\mathbf M(\mathbf r)
=\lim_{\Delta V\to0}\frac{\sum_{\Delta V}\mathbf m_i}{\Delta V}}.
\]

\(\mathbf M\) 是单位体积的磁偶极矩，单位为 \(\mathrm{A/m}\)。体积元
\(dV\) 的磁偶极矩为

\[
d\mathbf m=\mathbf M\,dV.
\]

磁化强度描述材料响应，既不是磁感应强度 \(\mathbf B\)，也不是磁场强度
\(\mathbf H\)。

\par\noindent\begin{minipage}{\linewidth}
\begin{center}
\includegraphics[width=\linewidth,height=.40\textheight,keepaspectratio]{\detokenize{assets/module-05/theory/tikz-m5-topic44-v9-from41-v1-01-magnetization.pdf}}
\par\smallskip{\small 图示：微观磁偶极矩的平均形成宏观磁化强度}
\end{center}
\end{minipage}\par

\par\addvspace{.4\baselineskip}
\noindent\begin{minipage}{\linewidth}

\textbf{例 M5-043：均匀磁化样品的总磁矩}

体积为 \(V\) 的有限样品具有均匀给定的磁化强度
\(\mathbf M\)。根据磁化强度的定义求样品的总磁偶极矩，并说明形状是否会改变这一积分关系，以及形状还可能通过什么量影响实际磁化。

\end{minipage}\par

\par\noindent\begin{minipage}{\linewidth}
\textbf{解答}

\begin{center}
\includegraphics[width=\linewidth,height=.40\textheight,keepaspectratio]{\detokenize{assets/module-05/examples/M5-043.pdf}}
\par\smallskip{\small 图 M5-043：均匀磁化样品的总磁矩}
\end{center}
\end{minipage}\par

图 M5-043 用圆柱样品示意均匀磁化；体内各小磁矩平行叠加为
\(V\mathbf M\)，结果并不限于圆柱。

若体积为 \(V\) 的小样品内磁化强度均匀，则总磁偶极矩为

\[
\boxed{\mathbf m_{\mathrm{\text{总}}}=\mathbf M V}.
\]

这个结论与样品形状无关；形状影响的是样品内部的退磁场，而不是总偶极矩的定义。

\subsection{二、磁化的等效束缚电流}\label{m5:ux4e8cux78c1ux5316ux7684ux7b49ux6548ux675fux7f1aux7535ux6d41}

\subsubsection{1.
体束缚电流与面束缚电流}\label{m5:ux4f53ux675fux7f1aux7535ux6d41ux4e0eux9762ux675fux7f1aux7535ux6d41}

磁偶极子可以看成微小闭合电流回路。相邻体积元内部的电流彼此抵消后，磁化分布对外产生的场可等效为

\[
\boxed{\mathbf J_b=\nabla\times\mathbf M},
\]

以及磁介质外表面上的

\[
\boxed{\mathbf K_b=\mathbf M\times\mathbf n},
\]

其中 \(\mathbf n\) 是磁介质的外法向，\(\mathbf J_b\) 的单位为
\(\mathrm{A/m^2}\)，\(\mathbf K_b\) 的单位为 \(\mathrm{A/m}\)。

这两式也可由磁矢势严格得到。磁化体的磁矢势为

\[
\mathbf A(\mathbf r)
=\frac{\mu_0}{4\pi}
\int_V
\frac{\mathbf M(\mathbf r')\times(\mathbf r-\mathbf r')}
{|\mathbf r-\mathbf r'|^3}\,dV'.
\]

对源坐标积分分部后，

\[
\boxed{
\mathbf A(\mathbf r)
=\frac{\mu_0}{4\pi}
\int_V\frac{\mathbf J_b(\mathbf r')}{|\mathbf r-\mathbf r'|}\,dV'
+\frac{\mu_0}{4\pi}
\oint_{\partial V}\frac{\mathbf K_b(\mathbf r')}{|\mathbf r-\mathbf r'|}\,da'
}.
\]

因此，\(\mathbf J_b\) 与 \(\mathbf K_b\)
不是额外假设，而是磁偶极分布的电流表示。

\subsubsection{2.
均匀磁化时为何只剩表面电流}\label{m5:ux5747ux5300ux78c1ux5316ux65f6ux4e3aux4f55ux53eaux5269ux8868ux9762ux7535ux6d41}

若 \(\mathbf M\) 在材料内部为常矢量，

\[
\nabla\times\mathbf M=\mathbf0,
\]

故体束缚电流为零；但一般仍有

\[
\mathbf K_b=\mathbf M\times\mathbf n\ne\mathbf0.
\]

\par\noindent\begin{minipage}{\linewidth}
\begin{center}
\includegraphics[width=\linewidth,height=.40\textheight,keepaspectratio]{\detokenize{assets/module-05/theory/tikz-m5-topic44-v9-from41-v1-02-bound-current.pdf}}
\par\smallskip{\small 图示：均匀磁化的束缚电流}
\end{center}
\end{minipage}\par

\par\addvspace{.4\baselineskip}
\noindent\begin{minipage}{\linewidth}

\textbf{例 M5-044：轴向磁化圆柱的束缚电流}

半径为 \(a\)、长度为 \(\ell\) 的有限圆柱具有均匀轴向磁化
\(\mathbf M=M\mathbf e_z\)。以材料外法向为
\(\mathbf n\)，求体束缚电流和侧面、两个端面的束缚面电流，并判断其等效绕流方向。

\end{minipage}\par

\par\noindent\begin{minipage}{\linewidth}
\textbf{解答}

\begin{center}
\includegraphics[width=\linewidth,height=.40\textheight,keepaspectratio]{\detokenize{assets/module-05/examples/M5-044.pdf}}
\par\smallskip{\small 图 M5-044：轴向磁化圆柱的束缚电流}
\end{center}
\end{minipage}\par

图 M5-044
中侧面法向与轴向磁化垂直，叉乘给出环向束缚电流；端面法向平行或反平行磁化，面电流为零。

设圆柱内

\[
\mathbf M=M\mathbf e_z.
\]

侧面外法向为 \(\mathbf e_r\)，所以

\[
\mathbf K_b=M\mathbf e_z\times\mathbf e_r
=M\mathbf e_{\phi}.
\]

两个端面上 \(\mathbf M\parallel\mathbf n\)，故
\(\mathbf K_b=0\)。圆柱侧面的环向束缚电流与面电流密度为 \(K=M\)
的螺线管完全等效。

\subsection{\texorpdfstring{三、\(\mathbf B\)、\(\mathbf H\)
与自由电流}{三、\textbackslash mathbf B、\textbackslash mathbf H 与自由电流}}\label{m5:ux4e09mathbf-bmathbf-h-ux4e0eux81eaux7531ux7535ux6d41}

总电流分为自由电流与束缚电流：

\[
\mathbf J_{\mathrm{\text{总}}}=\mathbf J_f+\mathbf J_b.
\]

磁静态安培定律为

\[
\nabla\times\mathbf B
=\mu_0(\mathbf J_f+\mathbf J_b)
=\mu_0\mathbf J_f+\mu_0\nabla\times\mathbf M.
\]

定义

\[
\boxed{\mathbf H=\frac{\mathbf B}{\mu_0}-\mathbf M},
\]

便得到

\[
\boxed{\nabla\times\mathbf H=\mathbf J_f},
\qquad
\boxed{\oint_C\mathbf H\cdot d\mathbf l=I_{f,\mathrm{enc}}}.
\]

因此

\[
\boxed{\mathbf B=\mu_0(\mathbf H+\mathbf M)}.
\]

\(\mathbf H\) 的旋度只由自由电流决定，但这不表示 \(\mathbf H\)
只由自由电流单独决定；材料形状、磁化和边界条件仍会通过无旋部分影响
\(\mathbf H\)。

\par\addvspace{.4\baselineskip}
\noindent\begin{minipage}{\linewidth}

\textbf{例 M5-045：线性磁芯中的H、M与B}

理想无限长螺线管单位长度匝数为 \(n\)，通自由电流 \(I\)，内部填满磁导率
\(\mu=\mu_0(1+\chi_m)\) 的线性介质。忽略外部背景场，先由自由电流求内部
\(H\)，再求 \(M\) 和 \(B\)，说明束缚电流是否应重复加进 \(H\) 的环路式。

\end{minipage}\par

\par\noindent\begin{minipage}{\linewidth}
\textbf{解答}

\begin{center}
\includegraphics[width=\linewidth,height=.40\textheight,keepaspectratio]{\detokenize{assets/module-05/examples/M5-045.pdf}}
\par\smallskip{\small 图 M5-045：线性磁芯中的H、M与B}
\end{center}
\end{minipage}\par

图 M5-045 的绕组自由电流先决定轴向 \(H\)；磁化对 \(B\)
的额外贡献通过本构关系进入，不能重复计入安匝数。

理想无限长、均匀绕制的螺线管单位长度匝数为 \(n\)，自由电流为 \(I\)。由
\(\mathbf H\) 的环路定理，

\[
H_{\mathrm{in}}=nI.
\]

若内部填充线性磁介质，则

\[
B_{\mathrm{in}}=\mu H_{\mathrm{in}}=\mu nI.
\]

这里自由电流先确定 \(H\)，本构关系再把 \(H\) 转换为 \(B\)。

\subsection{四、磁标势与等效``磁荷''}\label{m5:ux56dbux78c1ux6807ux52bfux4e0eux7b49ux6548ux78c1ux8377}

在不含自由电流的单连通区域，

\[
\nabla\times\mathbf H=\mathbf0,
\]

因而可以定义磁标势

\[
\boxed{\mathbf H=-\nabla\varphi_m}.
\]

由

\[
\nabla\cdot\mathbf B
=\mu_0\nabla\cdot(\mathbf H+\mathbf M)=0
\]

得到

\[
\nabla\cdot\mathbf H=-\nabla\cdot\mathbf M.
\]

定义等效磁荷密度

\[
\boxed{\rho_m=-\nabla\cdot\mathbf M},
\qquad
\boxed{\sigma_m=\mathbf M\cdot\mathbf n},
\]

则磁标势可写成

\[
\boxed{
\varphi_m(\mathbf r)
=\frac1{4\pi}\int_V\frac{\rho_m(\mathbf r')}{|\mathbf r-\mathbf r'|}\,dV'
+\frac1{4\pi}\oint_{\partial V}\frac{\sigma_m(\mathbf r')}{|\mathbf r-\mathbf r'|}\,da'
}.
\]

这里的``磁荷''只是求解 \(\mathbf H\)
的数学等效源，不表示自然界存在孤立磁单极。束缚电流表示直接求
\(\mathbf B\)，磁荷表示常用于无自由电流区域求
\(\mathbf H\)；两种方法必须给出相同的物理磁场。

\par\noindent\begin{minipage}{\linewidth}
\begin{center}
\includegraphics[width=\linewidth,height=.40\textheight,keepaspectratio]{\detokenize{assets/module-05/theory/tikz-m5-topic44-v9-from41-v1-03-two-representations.pdf}}
\par\smallskip{\small 图示：束缚电流与等效磁荷}
\end{center}
\end{minipage}\par

\par\addvspace{.4\baselineskip}
\noindent\begin{minipage}{\linewidth}

\textbf{例 M5-046：磁化圆柱的等效磁荷}

均匀轴向磁化圆柱的磁化强度为 \(M\mathbf e_z\)。采用等效磁荷法，按
\(\rho_m=-\nabla\cdot\mathbf M\text{、}\sigma_m=\mathbf M\cdot\mathbf n\)
求体内和上下端面的等效磁荷，并说明它们是否是真实磁单极子。

\end{minipage}\par

\par\noindent\begin{minipage}{\linewidth}
\textbf{解答}

\begin{center}
\includegraphics[width=\linewidth,height=.40\textheight,keepaspectratio]{\detokenize{assets/module-05/examples/M5-046.pdf}}
\par\smallskip{\small 图 M5-046：磁化圆柱的等效磁荷}
\end{center}
\end{minipage}\par

图 M5-046 上下端面的等效磁荷符号由 \(\mathbf M\cdot\mathbf n\)
决定。它们是计算磁场的等效源，并不表示存在真实磁单极子。

均匀轴向磁化圆柱内部 \(\nabla\cdot\mathbf M=0\)，故
\(\rho_m=0\)；上端面有 \(\sigma_m=+M\)，下端面有
\(\sigma_m=-M\)。这与前面的侧面环流表示看似不同，却描述同一个磁场。

\subsection{五、线性磁介质的使用条件}\label{m5:ux4e94ux7ebfux6027ux78c1ux4ecbux8d28ux7684ux4f7fux7528ux6761ux4ef6}

在线性、各向同性、局域且无明显磁滞的介质中，

\[
\boxed{\mathbf M=\chi_m\mathbf H},
\]

于是

\[
\boxed{\mathbf B=\mu_0(1+\chi_m)\mathbf H
=\mu\mathbf H},
\]

其中

\[
\boxed{\mu_r=1+\chi_m},
\qquad
\boxed{\mu=\mu_0\mu_r}.
\]

\begin{itemize}
\tightlist
\item
  抗磁介质通常有 \(\chi_m<0\)；
\item
  顺磁介质通常有 \(\chi_m>0\) 且数值较小；
\item
  铁磁材料常具有强非线性、磁滞和饱和，不能在任意场强下用常数 \(\mu\)
  描述。
\end{itemize}

若介质各向异性，应写成

\[
\mathbf B=\boldsymbol\mu\,\mathbf H,
\]

此时 \(\mathbf B\) 与 \(\mathbf H\) 不一定平行。

\par\addvspace{.4\baselineskip}
\noindent\begin{minipage}{\linewidth}

\textbf{例 M5-047：弱磁介质的磁化强度}

一个线性各向同性弱顺磁样品内的实际宏观磁场强度为
\(H=4.0\times10^3\ \mathrm{A/m}\)，磁化率
\(\chi_m=2.0\times10^{-3}\)。求磁化强度 \(M\)、磁感应强度 \(B\) 及相对于
\(\mu_0H\) 的增幅；此处 \(H\) 已是退磁修正后的内场。

\end{minipage}\par

\par\noindent\begin{minipage}{\linewidth}
\textbf{解答}

\begin{center}
\includegraphics[width=\linewidth,height=.40\textheight,keepaspectratio]{\detokenize{assets/module-05/examples/M5-047.pdf}}
\par\smallskip{\small 图 M5-047：弱磁介质的磁化强度}
\end{center}
\end{minipage}\par

图 M5-047 的绿色磁化与实际内场 \(H\) 同向；本题磁化量远小于 \(H\)，将
\(H+M\) 代入 \(B\) 时只得到千分之二的相对增量。

某线性介质 \(\chi_m=2.0\times10^{-3}\)，在
\(H=4.0\times10^3\ \mathrm{A/m}\) 的均匀场中，

\[
M=\chi_mH=8.0\ \mathrm{A/m},
\]

\[
B=\mu_0(H+M)\approx1.002\,\mu_0H.
\]

因此弱顺磁介质对 \(B\)
的增幅很小，不能把``顺磁''误解为任意大的磁场增强。

\subsection{六、磁介质边界条件}\label{m5:ux516dux78c1ux4ecbux8d28ux8fb9ux754cux6761ux4ef6}

令单位法向 \(\mathbf n\) 从介质 1 指向介质
2。专题四十三已经从薄柱和窄回路得到

\[
\boxed{\mathbf n\cdot(\mathbf B_2-\mathbf B_1)=0},
\]

\[
\boxed{\mathbf n\times(\mathbf H_2-\mathbf H_1)=\mathbf K_f}.
\]

若界面上无自由面电流，

\[
B_{2n}=B_{1n},
\qquad
H_{2t}=H_{1t}.
\]

束缚面电流已经通过磁化强度计入。若改用 \(\mathbf B\)
的跳变式，则应使用总面电流。按本节法向约定，界面上的束缚面电流为

\[
\boxed{
\mathbf K_b
=\mathbf n\times(\mathbf M_2-\mathbf M_1)
=(\mathbf M_1-\mathbf M_2)\times\mathbf n
}.
\]

因而

\[
\boxed{
\mathbf n\times(\mathbf B_2-\mathbf B_1)
=\mu_0(\mathbf K_f+\mathbf K_b)
}.
\]

不能在 \(\mathbf H\) 的边界条件右端再次加入
\(\mathbf K_b\)，否则会重复计数。

\subsection{七、退磁场：样品形状为何重要}\label{m5:ux4e03ux9000ux78c1ux573aux6837ux54c1ux5f62ux72b6ux4e3aux4f55ux91cdux8981}

有限磁化体表面的等效磁荷或束缚电流会反过来产生一个通常与 \(\mathbf M\)
反向的内部场，称为退磁场。对均匀磁化椭球，内部退磁场也是均匀的：

\[
\boxed{\mathbf H_d=-\mathbf N\mathbf M}.
\]

沿椭球三个主轴，

\[
H_{d,i}=-N_iM_i,
\qquad
\boxed{N_x+N_y+N_z=1}.
\]

常见极限为：

\[
\begin{array}{c|c}
\text{形状与磁化方向}&N\\ \hline
\text{球体}&1/3\\
\text{很长圆柱，沿轴磁化}&\approx0\\
\text{很薄平板，垂直板面磁化}&\approx1
\end{array}
\]

\par\noindent\begin{minipage}{\linewidth}
\begin{center}
\includegraphics[width=\linewidth,height=.40\textheight,keepaspectratio]{\detokenize{assets/module-05/theory/tikz-m5-topic44-v9-from41-v1-04-demagnetizing-field.pdf}}
\par\smallskip{\small 图示：不同形状的退磁场}
\end{center}
\end{minipage}\par

若外加磁场为 \(\mathbf H_0\)，线性介质内

\[
\mathbf H=\mathbf H_0+\mathbf H_d.
\]

沿某一主轴磁化时，

\[
H=H_0-NM,
\qquad M=\chi_mH.
\]

消去 \(H\) 得

\[
\boxed{
M=\frac{\chi_m}{1+N\chi_m}H_0
},
\]

\[
\boxed{
H=\frac{H_0}{1+N\chi_m}
}.
\]

这说明同一种材料做成长棒、球或薄片，内部磁化程度可以不同。\(\mu\)
是材料本构参数，\(N\) 是几何参数，二者不能混为一谈。

\par\addvspace{.4\baselineskip}
\noindent\begin{minipage}{\linewidth}

\textbf{例 M5-048：磁介质球的退磁场}

均匀线性各向同性磁介质球置于真空外加匀强场 \(\mathbf H_0\)
中，材料磁化率为 \(\chi_m\)。已知球的退磁因子 \(N=1/3\)，利用
\(\mathbf H_{\rm in}=\mathbf H_0-N\mathbf M\)
求内场与磁化，说明不能直接用外场代替本构关系中的内场。

\end{minipage}\par

\par\noindent\begin{minipage}{\linewidth}
\textbf{解答}

\begin{center}
\includegraphics[width=\linewidth,height=.40\textheight,keepaspectratio]{\detokenize{assets/module-05/examples/M5-048.pdf}}
\par\smallskip{\small 图 M5-048：磁介质球的退磁场}
\end{center}
\end{minipage}\par

图 M5-048 的红色退磁场反向于绿色磁化，蓝色 \(H_0\)
是没有样品时的外加场。材料本构关系应使用实际内场 \(H_{\rm in}\)。

球体 \(N=1/3\)，故

\[
\boxed{
H_{\mathrm{in}}=\frac{3}{3+\chi_m}H_0
},
\]

\[
\boxed{
M=\frac{3\chi_m}{3+\chi_m}H_0
}.
\]

即使材料内部满足 \(\mathbf M=\chi_m\mathbf H\)，也不能直接写成
\(M=\chi_mH_0\)，因为 \(H_0\) 是没有样品时的外加场，\(H_{\mathrm{in}}\)
才是材料实际感受到的宏观场。

\subsection{八、带气隙磁芯与磁路近似}\label{m5:ux516bux5e26ux6c14ux9699ux78c1ux82afux4e0eux78c1ux8defux8fd1ux4f3c}

对截面积近似恒定、漏磁和边缘效应可忽略的磁芯---气隙串联系统，磁通连续使磁芯和气隙中的
\(B\) 近似相同。若磁芯平均长度为 \(\ell_c\)、相对磁导率为
\(\mu_r\)，气隙长度为 \(g\)，线圈匝数为 \(N\)，则

\[
NI=H_c\ell_c+H_gg
=\frac{B\ell_c}{\mu_0\mu_r}
+\frac{Bg}{\mu_0}.
\]

所以

\[
\boxed{
B=\frac{\mu_0NI}{g+\ell_c/\mu_r}
}.
\]

磁通量和电感为

\[
\boxed{\Phi=BS},
\qquad
\boxed{
L=\frac{N\Phi}{I}
=\frac{\mu_0N^2S}{g+\ell_c/\mu_r}
}.
\]

\par\noindent\begin{minipage}{\linewidth}
\begin{center}
\includegraphics[width=\linewidth,height=.40\textheight,keepaspectratio]{\detokenize{assets/module-05/theory/tikz-m5-topic44-v9-from41-v1-05-air-gap-core.pdf}}
\par\smallskip{\small 图示：单气隙磁路与吸力}
\end{center}
\end{minipage}\par

气隙虽短，但因磁导率低，可能承担主要安匝降和主要磁场能。在线性近似下，

\[
\frac{U_g}{U_c}
=\frac{[B^2/(2\mu_0)]Sg}
{[B^2/(2\mu_0\mu_r)]S\ell_c}
=\boxed{\frac{\mu_rg}{\ell_c}}.
\]

\par\addvspace{.4\baselineskip}
\noindent\begin{minipage}{\linewidth}

\textbf{例 M5-049：气隙与磁芯的储能分配}

同截面磁路由相对磁导率 \(\mu_r=2000\)、长度 \(\ell_c\)
的线性磁芯与长度为 \(g\)
的空气隙串联构成，\(g/\ell_c=10^{-3}\)。忽略漏磁和边缘效应，两段通过相同磁通。求气隙与磁芯储能之比及气隙占总磁能的比例。

\end{minipage}\par

\par\noindent\begin{minipage}{\linewidth}
\textbf{解答}

\begin{center}
\includegraphics[width=\linewidth,height=.40\textheight,keepaspectratio]{\detokenize{assets/module-05/examples/M5-049.pdf}}
\par\smallskip{\small 图 M5-049：气隙与磁芯的储能分配}
\end{center}
\end{minipage}\par

图 M5-049 的磁芯和气隙串联，通过同一磁通；能量密度与 \(1/\mu\)
成正比，因此很窄的气隙也可储存大部分磁能。

若 \(\mu_r=2000\)、\(g/\ell_c=10^{-3}\)，则
\(U_g/U_c=2\)。气隙只占磁路长度的千分之一，却储存约三分之二的总磁能。

\subsection{九、非线性、磁滞与能量}\label{m5:ux4e5dux975eux7ebfux6027ux78c1ux6edeux4e0eux80fdux91cf}

对无磁滞但可能非线性的可逆磁化过程，

\[
\boxed{
u_B(B)=\int_0^B H(B')\,dB'
},
\]

磁共能密度为

\[
\boxed{
u_B'(H)=\int_0^H B(H')\,dH'
}.
\]

线性介质中二者均为

\[
u_B=u_B'=\frac12BH.
\]

若存在磁滞，系统经历一个闭合 \(B\)---\(H\) 回线后，每周期单位体积耗散为

\[
\boxed{
w_{\mathrm{hys}}=\oint H\,dB
},
\]

即磁滞回线所围的有向面积。此时不能把全部输入功都解释成可恢复磁场能。

\subsection{十、使用条件与常见错误}\label{m5:ux5341ux4f7fux7528ux6761ux4ef6ux4e0eux5e38ux89c1ux9519ux8bef-1}

\begin{enumerate}
\def\labelenumi{\arabic{enumi}.}
\tightlist
\item
  \(\mathbf B=\mu\mathbf H\)
  只适用于相应的线性本构区间；铁磁材料不能默认 \(\mu\) 为常数。
\item
  \(\mathbf H\) 的环量只计自由电流，但 \(\mathbf H\)
  仍受磁化和边界条件影响。
\item
  \(\mathbf J_b=\nabla\times\mathbf M\)、\(\mathbf K_b=\mathbf M\times\mathbf n\)
  与
  \(\rho_m=-\nabla\cdot\mathbf M\)、\(\sigma_m=\mathbf M\cdot\mathbf n\)
  是两种等价表示，不能把两套等效源同时叠加。
\item
  磁标势只可在无自由电流的单连通区域直接定义；环绕自由电流的区域不能使用全局单值磁标势。
\item
  退磁因子是几何量。只有椭球体在均匀外场中能够严格保持内部均匀磁化；其他形状通常只能近似使用平均退磁因子。
\item
  带气隙磁路公式忽略漏磁和边缘效应；气隙与截面尺度可比时必须修正有效面积。
\item
  磁荷表示不意味着存在真实磁单极。
\end{enumerate}

\subsection{练习与解答}\label{m5:ux7ec3ux4e60ux4e0eux89e3ux7b54-2}

\par\addvspace{.4\baselineskip}
\noindent\begin{minipage}{\linewidth}

\subsubsection{作业一：均匀磁化圆柱的等效源}\label{m5:ux4f5cux4e1aux4e00ux5747ux5300ux78c1ux5316ux5706ux67f1ux7684ux7b49ux6548ux6e90}

\textbf{题号：M5-050}

\textbf{层级：} 基础　　\textbf{方法：} 介质边界、等效变换

磁化在圆柱内部严格均匀、外部为零，\(M_0\text{、}a\text{、}\ell>0\)；磁荷是等效表示，不是真实自由磁荷。

半径为 \(a\)、长度为 \(\ell\)
的圆柱体均匀磁化，\(\mathbf M=M_0\mathbf e_z\)。求体束缚电流、各表面的束缚面电流、等效磁荷及圆柱总磁偶极矩。

\end{minipage}\par

\par\noindent\begin{minipage}{\linewidth}
\textbf{解答}

\begin{center}
\includegraphics[width=\linewidth,height=.40\textheight,keepaspectratio]{\detokenize{assets/module-05/examples/M5-050.pdf}}
\par\smallskip{\small 图 M5-050：均匀磁化圆柱的等效源}
\end{center}
\end{minipage}\par

图 M5-050 把侧面环向电流和上下等效磁荷放在同一几何上，分别使用
\(\mathbf M\times\mathbf n\) 与
\(\mathbf M\cdot\mathbf n\)；总磁矩再由体积求得。

因为 \(\mathbf M\) 在圆柱内部为常矢量，

\[
\boxed{\mathbf J_b=\nabla\times\mathbf M=\mathbf0}.
\]

侧面外法向为 \(\mathbf n=\mathbf e_r\)，故

\[
\boxed{
\mathbf K_b
=\mathbf M\times\mathbf n
=M_0\mathbf e_{\phi}
}.
\]

上、下端面的法向分别为 \(\pm\mathbf e_z\)，均有

\[
\boxed{\mathbf K_b=\mathbf0}.
\]

等效体磁荷为

\[
\boxed{\rho_m=-\nabla\cdot\mathbf M=0}.
\]

上端面

\[
\boxed{\sigma_m=+M_0},
\]

下端面

\[
\boxed{\sigma_m=-M_0},
\]

侧面 \(\sigma_m=0\)。圆柱体积为 \(\pi a^2\ell\)，因此总磁偶极矩

\[
\boxed{
\mathbf m=\mathbf M V
=M_0\pi a^2\ell\,\mathbf e_z
}.
\]

侧面束缚电流表示和端面磁荷表示描述的是同一个磁场，不能把它们当作两组独立物理源重复相加。

\par\addvspace{.4\baselineskip}
\noindent\begin{minipage}{\linewidth}

\subsubsection{作业二：线性磁介质球置于均匀外场}\label{m5:ux4f5cux4e1aux4e8cux7ebfux6027ux78c1ux4ecbux8d28ux7403ux7f6eux4e8eux5747ux5300ux5916ux573a}

\textbf{题号：M5-051}

\textbf{层级：} 基础　　\textbf{方法：} 介质边界

假设可采用连续线性本构，球内无自由电流、无磁滞；外加场是移去样品后真空中给定的场。

半径为 \(a\)、磁化率为 \(\chi_m\)
的线性各向同性磁介质球置于真空均匀外加场 \(\mathbf H_0=H_0\mathbf e_z\)
中。求球内的 \(\mathbf H\)、\(\mathbf M\)、\(\mathbf B\)
以及球的总磁偶极矩。

\end{minipage}\par

\par\noindent\begin{minipage}{\linewidth}
\textbf{解答}

\begin{center}
\includegraphics[width=\linewidth,height=.40\textheight,keepaspectratio]{\detokenize{assets/module-05/examples/M5-051.pdf}}
\par\smallskip{\small 图 M5-051：线性磁介质球置于均匀外场}
\end{center}
\end{minipage}\par

图 M5-051 中均匀内场由外场与退磁场相加而成。求出绿色磁化后，再乘球体积
\(4\pi a^3/3\) 得到总偶极矩。

球体沿任意方向的退磁因子均为

\[
N=\frac13.
\]

球内退磁场为

\[
\mathbf H_d=-\frac13\mathbf M.
\]

因此

\[
\mathbf H=\mathbf H_0-\frac13\mathbf M.
\]

\par\noindent\begin{minipage}{\linewidth}
\begin{center}
\includegraphics[width=\linewidth,height=.40\textheight,keepaspectratio]{\detokenize{assets/module-05/process/M5-051a.pdf}}
\par\smallskip{\small 图 M5-051a：外加场与退磁场的自洽叠加}
\end{center}
\end{minipage}\par

图 M5-051a
中红色退磁场反向于磁化，必须与外加场先合成实际内场，再代回材料本构关系。

再用线性本构关系

\[
\mathbf M=\chi_m\mathbf H,
\]

得到

\[
\mathbf H
=\mathbf H_0-\frac{\chi_m}{3}\mathbf H.
\]

所以

\[
\boxed{
\mathbf H_{\mathrm{in}}
=\frac{3}{3+\chi_m}\mathbf H_0
}.
\]

磁化强度为

\[
\boxed{
\mathbf M
=\frac{3\chi_m}{3+\chi_m}\mathbf H_0
}.
\]

球内磁感应强度为

\[
\mathbf B_{\mathrm{in}}
=\mu_0(\mathbf H_{\mathrm{in}}+\mathbf M),
\]

即

\[
\boxed{
\mathbf B_{\mathrm{in}}
=\frac{3\mu_0(1+\chi_m)}{3+\chi_m}\mathbf H_0
}.
\]

球体积为 \(4\pi a^3/3\)，故总磁偶极矩

\[
\boxed{
\mathbf m
=\frac{4\pi a^3}{3}\mathbf M
=\frac{4\pi a^3\chi_m}{3+\chi_m}\mathbf H_0
}.
\]

球外附加场等于位于球心、偶极矩为上述 \(\mathbf m\) 的磁偶极场。

\par\addvspace{.4\baselineskip}
\noindent\begin{minipage}{\linewidth}

\subsubsection{作业三：带气隙环形磁芯}\label{m5:ux4f5cux4e1aux4e09ux5e26ux6c14ux9699ux73afux5f62ux78c1ux82af}

\textbf{题号：M5-052}

\textbf{层级：} 基础　　\textbf{方法：} 介质边界、安培环路

材料线性且参数不随电流变化；\(\ell_c\) 仅计铁芯段长度，\(g\)
仅计一个空气隙，二者同截面。

截面积为 \(S\) 的环形磁芯平均磁路长度为 \(\ell_c\)，相对磁导率为
\(\mu_r\)；磁芯上开有长度为 \(g\) 的窄气隙，并绕有 \(N\) 匝线圈，通电流
\(I\)。忽略漏磁、边缘效应和磁芯非线性，求磁通密度、磁通量、电感以及气隙与磁芯的储能比。

\end{minipage}\par

\par\noindent\begin{minipage}{\linewidth}
\textbf{解答}

\begin{center}
\includegraphics[width=\linewidth,height=.40\textheight,keepaspectratio]{\detokenize{assets/module-05/examples/M5-052.pdf}}
\par\smallskip{\small 图 M5-052：带气隙环形磁芯}
\end{center}
\end{minipage}\par

图 M5-052
的铁芯段和空气隙沿同一闭合磁路串联，安培环路包含两段磁势降；同一 \(B\)
和截面 \(S\) 使磁通连续。

在忽略漏磁且截面积相同的近似下，磁芯和气隙中的磁通密度均为
\(B\)。安培环路定理给出

\[
NI=H_c\ell_c+H_gg.
\]

由

\[
H_c=\frac{B}{\mu_0\mu_r},
\qquad
H_g=\frac{B}{\mu_0},
\]

得到

\[
NI
=\frac{B}{\mu_0}
\left(\frac{\ell_c}{\mu_r}+g\right).
\]

因此

\[
\boxed{
B=\frac{\mu_0NI}{g+\ell_c/\mu_r}
}.
\]

磁通量为

\[
\boxed{
\Phi=BS
=\frac{\mu_0NIS}{g+\ell_c/\mu_r}
}.
\]

磁链为 \(N\Phi=LI\)，故

\[
\boxed{
L=\frac{\mu_0N^2S}{g+\ell_c/\mu_r}
}.
\]

气隙和磁芯中的能量分别为

\[
U_g=\frac{B^2}{2\mu_0}Sg,
\]

\[
U_c=\frac{B^2}{2\mu_0\mu_r}S\ell_c.
\]

所以

\[
\boxed{
\frac{U_g}{U_c}=\frac{\mu_rg}{\ell_c}
}.
\]

当 \(\mu_rg\gg\ell_c\) 时，系统的主要安匝降和主要磁场能都位于气隙中。

\section{专题四十五：磁矢势、规范与磁偶极子（拔高选讲）}\label{m5:ux4e13ux9898ux56dbux5341ux4e94ux78c1ux77e2ux52bfux89c4ux8303ux4e0eux78c1ux5076ux6781ux5b50ux62d4ux9ad8ux9009ux8bb2}

磁矢势不是为了给磁场换一种记号。它把``磁通量''转化为``边界上的环量''，把稳恒电流问题转化为泊松方程，并把局域电流分布的远场结构自然组织成磁偶极展开。使用磁矢势时必须同时说明定义域、规范和边界条件，否则同一个公式可能对应完全不同的问题。

\subsection{三种势函数的适用区域}\label{m5:ux4e09ux79cdux52bfux51fdux6570ux7684ux9002ux7528ux533aux57df}

静电势、磁标势和磁矢势的形式相似，但它们来自不同的微分约束。

{\def\LTcaptype{none} % do not increment counter
\begin{longtable}[]{@{}
  >{\raggedright\arraybackslash}p{(\linewidth - 8\tabcolsep) * \real{0.2000}}
  >{\raggedright\arraybackslash}p{(\linewidth - 8\tabcolsep) * \real{0.2000}}
  >{\raggedright\arraybackslash}p{(\linewidth - 8\tabcolsep) * \real{0.2000}}
  >{\raggedright\arraybackslash}p{(\linewidth - 8\tabcolsep) * \real{0.2000}}
  >{\raggedright\arraybackslash}p{(\linewidth - 8\tabcolsep) * \real{0.2000}}@{}}
\toprule\noalign{}
\begin{minipage}[b]{\linewidth}\raggedright
势函数
\end{minipage} & \begin{minipage}[b]{\linewidth}\raggedright
场的表示
\end{minipage} & \begin{minipage}[b]{\linewidth}\raggedright
可用区域
\end{minipage} & \begin{minipage}[b]{\linewidth}\raggedright
源与边界
\end{minipage} & \begin{minipage}[b]{\linewidth}\raggedright
额外限制
\end{minipage} \\
\midrule\noalign{}
\endhead
\bottomrule\noalign{}
\endlastfoot
静电势 \(\varphi\) & \(\mathbf E=-\nabla\varphi\) & 静电场的无旋区域 &
\(\nabla^2\varphi=-\rho/\varepsilon_0\)，匹配电势与法向场 &
多连通区域仍需检查闭路环量 \\
磁标势 \(\psi_m\) & \(\mathbf H=-\nabla\psi_m\) & \(\mathbf J_f=0\)
的局部区域 & 由磁化源或磁边界条件确定 &
绕自由电流的闭路使全局单值标势失效 \\
磁矢势 \(\mathbf A\) & \(\mathbf B=\nabla\times\mathbf A\) &
包含普通电流源的磁场区域 & 库仑规范下
\(\nabla^2\mathbf A=-\mu_0\mathbf J\) & 存在
\(\mathbf A\to\mathbf A+\nabla\chi\) 的规范自由度 \\
\end{longtable}
}

选择势函数前应先问``哪一个旋度为零''以及区域是否含孔。磁感应强度的无散性质对应磁矢势；磁标势则来自无自由电流区域中
\(\nabla\times\mathbf H=0\)，两者不能互换。

\par\noindent\begin{minipage}{\linewidth}
\begin{center}
\includegraphics[width=\linewidth,height=.40\textheight,keepaspectratio]{\detokenize{assets/module-05/theory/tikz-v11-topic45-potential-comparison.pdf}}
\par\smallskip{\small 图示：静电势磁标势和磁矢势具有不同定义域与拓扑条件}
\end{center}
\end{minipage}\par

\subsection{一、磁矢势的存在性与磁通量}\label{m5:ux4e00ux78c1ux77e2ux52bfux7684ux5b58ux5728ux6027ux4e0eux78c1ux901aux91cf}

\subsubsection{1.
局部存在与整体存在}\label{m5:ux5c40ux90e8ux5b58ux5728ux4e0eux6574ux4f53ux5b58ux5728}

麦克斯韦方程给出

\[
\nabla\cdot\mathbf B=0.
\]

因此，在每个不含奇点的充分小邻域内，都可以引入磁矢势 \(\mathbf A\)，使

\[
\boxed{\mathbf B=\nabla\times\mathbf A}.
\]

若定义域是星形区域或更一般的可缩区域，\(\nabla\cdot\mathbf B=0\)
还保证存在覆盖整个区域的光滑磁矢势。对带孔区域，仅有散度为零并不足以直接断言存在单值、光滑的整体矢势；还要检查每个不能在定义域内缩成一点的闭曲面是否满足

\[
\oiint_{\Sigma}\mathbf B\cdot d\mathbf S=0.
\]

这里的本质是：散度为零给出局部条件，定义域的拓扑决定局部矢势能否拼成整体矢势。``单连通''只排除了不可缩闭曲线，在三维空间中并不自动排除不可缩闭曲面，因而不能代替完整的存在性判断。

普通电磁学中的物理磁场没有磁单极子，在包含全部普通电流源的全空间中不会出现非零闭合磁通。假设存在非零闭合磁通的数学反例，只用于说明拓扑障碍，不代表经典电磁学中存在磁荷。

\par\noindent\begin{minipage}{\linewidth}
\begin{center}
\includegraphics[width=\linewidth,height=.40\textheight,keepaspectratio]{\detokenize{assets/module-05/theory/tikz-m5-topic45-v9-from42-v8-01.pdf}}
\par\smallskip{\small 图示：磁矢势的旋度与边界环量}
\end{center}
\end{minipage}\par

\subsubsection{2.
磁通量与环量}\label{m5:ux78c1ux901aux91cfux4e0eux73afux91cf}

由斯托克斯定理，若曲面 \(S\) 的定向法向量与边界 \(C=\partial S\)
的绕行方向满足右手规则，则

\[
\boxed{
\Phi_B(S)=\iint_S\mathbf B\cdot d\mathbf S
=\oint_{\partial S}\mathbf A\cdot d\mathbf l
}.
\]

若两个曲面 \(S_1,S_2\)
具有同一条边界，并且二者围成的区域完全位于磁场的光滑定义域内，则

\[
\iint_{S_1}\mathbf B\cdot d\mathbf S
-\iint_{S_2}\mathbf B\cdot d\mathbf S
=\oiint_{S_1\cup(-S_2)}\mathbf B\cdot d\mathbf S=0.
\]

所以边界相同的曲面具有相同磁通量。由 \(\mathbf B=\nabla\times\mathbf A\)
还可看出

\[
[\mathbf A]=[\mathbf B]\,[\text{长度}]=\mathrm{T\cdot m}.
\]

\par\addvspace{.4\baselineskip}
\noindent\begin{minipage}{\linewidth}

\textbf{例 M5-053：均匀磁场中的圆周}

匀强静磁场 \(\mathbf B=B_0\mathbf e_z\) 中，半径为 \(a\) 的圆周位于
\(xy\) 平面并取逆时针正向。设 \(\mathbf B=\nabla\times\mathbf A\)，求
\(\oint_C\mathbf A\cdot d\mathbf l\)，注明曲面法向及结果是否依赖规范。

\end{minipage}\par

\par\noindent\begin{minipage}{\linewidth}
\textbf{解答}

\begin{center}
\includegraphics[width=\linewidth,height=.40\textheight,keepaspectratio]{\detokenize{assets/module-05/examples/M5-053.pdf}}
\par\smallskip{\small 图 M5-053：均匀磁场中的圆周}
\end{center}
\end{minipage}\par

图 M5-053
的逆时针有向圆周与纸外正法向配对；应用斯托克斯定理后，矢势环量等于圆内磁通，而非任意一点的矢势值。

若 \(\mathbf B=B_0\mathbf e_z\)，圆周 \(C\) 位于 \(xy\) 平面、半径为
\(a\)，并取逆时针正向，则

\[
\oint_C\mathbf A\cdot d\mathbf l
=\iint_S B_0\,dS
=B_0\pi a^2.
\]

这一结论不依赖于选用了哪一种规范，只依赖于圆周所围的磁通量。

\subsection{二、规范自由度}\label{m5:ux4e8cux89c4ux8303ux81eaux7531ux5ea6}

\subsubsection{1.
静磁场中的规范变换}\label{m5:ux9759ux78c1ux573aux4e2dux7684ux89c4ux8303ux53d8ux6362}

若 \(\chi(\mathbf r)\) 是二阶连续可微的单值标量函数，则

\[
\mathbf A'=\mathbf A+\nabla\chi
\]

给出

\[
\nabla\times\mathbf A'
=\nabla\times\mathbf A+\nabla\times\nabla\chi
=\mathbf B.
\]

因此磁矢势不是唯一的。反过来，在单连通区域内，若
\(\nabla\times\mathbf A'=\nabla\times\mathbf A\)，则
\(\mathbf A'-\mathbf A\) 的旋度为零，可以写成某个单值标量函数的梯度。

对闭合回路，单值规范函数还保证

\[
\oint_C\mathbf A'\cdot d\mathbf l
-\oint_C\mathbf A\cdot d\mathbf l
=\oint_C\nabla\chi\cdot d\mathbf l=0.
\]

所以闭合回路上的矢势环量与磁通量都不随这类规范变换而改变。

\subsubsection{2.
库仑规范及其剩余自由度}\label{m5:ux5e93ux4ed1ux89c4ux8303ux53caux5176ux5269ux4f59ux81eaux7531ux5ea6}

静磁问题中常选

\[
\boxed{\nabla\cdot\mathbf A=0},
\]

这称为库仑规范。若从任意 \(\mathbf A\) 出发，希望通过
\(\mathbf A'=\mathbf A+\nabla\chi\) 达到库仑规范，\(\chi\) 应满足

\[
\nabla^2\chi=-\nabla\cdot\mathbf A.
\]

但库仑规范本身仍不能唯一确定 \(\mathbf A\)。若 \(\mathbf A\) 与
\(\mathbf A'\) 都满足库仑规范，则

\[
\nabla^2\chi=0.
\]

只有再给出边界条件、无穷远行为或周期条件，才能消除相应的调和函数自由度。

\par\addvspace{.4\baselineskip}
\noindent\begin{minipage}{\linewidth}

\textbf{例 M5-054：同一均匀磁场的两种库仑规范}

对均匀静磁场 \(\mathbf B=B_0\mathbf e_z\)，考虑
\(\mathbf A_1=\tfrac12\mathbf B\times\mathbf r\) 与
\(\mathbf A_2=B_0x\mathbf e_y\)。验证两者都满足库仑规范并产生同一磁场，求使
\(\mathbf A_2=\mathbf A_1+\nabla\chi\) 的一个规范函数。

\end{minipage}\par

\par\noindent\begin{minipage}{\linewidth}
\textbf{解答}

\begin{center}
\includegraphics[width=\linewidth,height=.40\textheight,keepaspectratio]{\detokenize{assets/module-05/examples/M5-054.pdf}}
\par\smallskip{\small 图 M5-054：同一均匀磁场的两种库仑规范}
\end{center}
\end{minipage}\par

图 M5-054
的蓝、绿矢量取自两个不同规范，局部方向不同，但对闭合曲线求环量及旋度后都给出同一个磁场。

对 \(\mathbf B=B_0\mathbf e_z\)，可取对称规范

\[
\mathbf A_{\mathrm{s}}
=\frac12\mathbf B\times\mathbf r
=\frac{B_0}{2}(-y\mathbf e_x+x\mathbf e_y),
\]

也可取

\[
\mathbf A_{\mathrm{L}}=B_0x\mathbf e_y.
\]

二者均满足 \(\nabla\cdot\mathbf A=0\) 与
\(\nabla\times\mathbf A=B_0\mathbf e_z\)，并且

\[
\mathbf A_{\mathrm{L}}-\mathbf A_{\mathrm{s}}
=\nabla\left(\frac12B_0xy\right).
\]

\subsubsection{3.
时变场中的完整规范变换}\label{m5:ux65f6ux53d8ux573aux4e2dux7684ux5b8cux6574ux89c4ux8303ux53d8ux6362}

时变电磁场中

\[
\mathbf B=\nabla\times\mathbf A,
\qquad
\mathbf E=-\nabla\varphi-\frac{\partial\mathbf A}{\partial t}.
\]

要使 \(\mathbf E\) 与 \(\mathbf B\) 同时不变，规范变换必须写成

\[
\boxed{
\mathbf A'=\mathbf A+\nabla\chi,
\qquad
\varphi'=\varphi-\frac{\partial\chi}{\partial t}
}.
\]

只改变 \(\mathbf A\) 而不相应改变
\(\varphi\)，一般会改变电场。处理辐射问题时常采用洛伦茨规范

\[
\nabla\cdot\mathbf A+\frac{1}{c^2}\frac{\partial\varphi}{\partial t}=0,
\]

此时两种势分别满足非齐次波动方程。静磁极限下
\(\partial/\partial t=0\)，它退化为库仑规范。

\subsection{三、库仑规范下的矢势方程}\label{m5:ux4e09ux5e93ux4ed1ux89c4ux8303ux4e0bux7684ux77e2ux52bfux65b9ux7a0b}

以下讨论真空中的稳恒电流。安培定律为

\[
\nabla\times\mathbf B=\mu_0\mathbf J,
\qquad
\nabla\cdot\mathbf J=0.
\]

代入 \(\mathbf B=\nabla\times\mathbf A\)：

\[
\nabla\times(\nabla\times\mathbf A)=\mu_0\mathbf J.
\]

利用恒等式

\[
\nabla\times(\nabla\times\mathbf A)
=\nabla(\nabla\cdot\mathbf A)-\nabla^2\mathbf A,
\]

并采用库仑规范，可得矢量泊松方程

\[
\boxed{\nabla^2\mathbf A=-\mu_0\mathbf J}.
\]

它表示三个笛卡尔分量分别满足标量泊松方程。对上式取散度又得到

\[
0=-\mu_0\nabla\cdot\mathbf J,
\]

说明该方程与稳恒电流的连续性条件相容。

泊松方程只给出局部微分关系，并不单独确定解。若 \(\mathbf J\)
局域分布于有限区域，并规定 \(\mathbf A\to\mathbf0\) 当
\(r\to\infty\)，才得到常用的自由空间积分解；若存在导体边界、周期边界或无限延伸电流，还必须改用与边界条件相符的格林函数，或补加满足
\(\nabla^2\mathbf A_{\mathrm{h}}=0\) 的齐次解。

\par\noindent\begin{minipage}{\linewidth}
\begin{center}
\includegraphics[width=\linewidth,height=.40\textheight,keepaspectratio]{\detokenize{assets/module-05/original/tikz-m5-topic45-v9-from42-v8-03.png}}
\par\smallskip{\small 图示：矢量泊松方程必须与电流连续性及边界条件一起使用}
\end{center}
\end{minipage}\par

\par\addvspace{.4\baselineskip}
\noindent\begin{minipage}{\linewidth}

\textbf{例 M5-055：无电流区域}

在不含电流的局部单连通区域内，静磁矢势采用库仑规范。由
\(\nabla^2\mathbf A=-\mu_0\mathbf J\)
判断矢势满足什么方程；这能否推出矢势或磁场必为零？给出均匀磁场矢势作为反例，并指出边界条件的作用。

\end{minipage}\par

\par\noindent\begin{minipage}{\linewidth}
\textbf{解答}

\begin{center}
\includegraphics[width=\linewidth,height=.40\textheight,keepaspectratio]{\detokenize{assets/module-05/examples/M5-055.pdf}}
\par\smallskip{\small 图 M5-055：无电流区域}
\end{center}
\end{minipage}\par

图 M5-055
的均匀磁场矢势可在无电流的局部区域内非零，且分量满足拉普拉斯方程。边界与外部电流决定具体解。

在 \(\mathbf J=0\) 的区域内，库仑规范下有

\[
\nabla^2\mathbf A=0.
\]

这并不意味着 \(\mathbf A=0\)，也不意味着
\(\mathbf B=0\)。例如均匀磁场的对称规范满足
\(\nabla^2\mathbf A_{\mathrm{s}}=0\)，但其旋度仍为 \(B_0\mathbf e_z\)。

\subsection{四、自由空间积分解与毕奥-萨伐尔定律}\label{m5:ux56dbux81eaux7531ux7a7aux95f4ux79efux5206ux89e3ux4e0eux6bd5ux5965-ux8428ux4f10ux5c14ux5b9aux5f8b}

\subsubsection{1. 积分解}\label{m5:ux79efux5206ux89e3}

利用

\[
\nabla^2\frac{1}{|\mathbf r-\mathbf r'|}
=-4\pi\delta^{(3)}(\mathbf r-\mathbf r'),
\]

局域稳恒电流在自由空间中的库仑规范解为

\[
\boxed{
\mathbf A(\mathbf r)
=\frac{\mu_0}{4\pi}
\int\frac{\mathbf J(\mathbf r')}{|\mathbf r-\mathbf r'|}\,dV'
}.
\]

记

\[
\mathbf R=\mathbf r-\mathbf r',
\qquad R=|\mathbf R|.
\]

对局域稳恒电流，\(\nabla'\cdot\mathbf J=0\)，且无穷远表面项为零。于是

\[
\begin{aligned}
\nabla\cdot\mathbf A
&=\frac{\mu_0}{4\pi}
\int\mathbf J(\mathbf r')\cdot\nabla\frac1R\,dV'\\
&=-\frac{\mu_0}{4\pi}
\int\mathbf J(\mathbf r')\cdot\nabla'\frac1R\,dV'=0,
\end{aligned}
\]

所以该积分解自动满足库仑规范。

对闭合细导线回路，\(\mathbf J\,dV'=I\,d\mathbf l'\)，因此

\[
\boxed{
\mathbf A(\mathbf r)
=\frac{\mu_0I}{4\pi}
\oint\frac{d\mathbf l'}{|\mathbf r-\mathbf r'|}
}.
\]

这里必须是闭合稳恒电流或含有明确回流路径的完整电路。把一段孤立开口导线直接当作稳恒电流源，会在端点破坏
\(\nabla\cdot\mathbf J=0\)。

\subsubsection{2.
从矢势恢复磁场}\label{m5:ux4eceux77e2ux52bfux6062ux590dux78c1ux573a}

对积分解取旋度：

\[
\begin{aligned}
\mathbf B(\mathbf r)
&=\nabla\times\mathbf A\\
&=\frac{\mu_0}{4\pi}
\int\nabla\left(\frac1R\right)\times\mathbf J(\mathbf r')\,dV'\\
&=\frac{\mu_0}{4\pi}
\int\frac{\mathbf J(\mathbf r')\times\mathbf R}{R^3}\,dV'.
\end{aligned}
\]

这正是体电流形式的毕奥-萨伐尔定律。

\par\noindent\begin{minipage}{\linewidth}
\begin{center}
\includegraphics[width=\linewidth,height=.40\textheight,keepaspectratio]{\detokenize{assets/module-05/theory/tikz-m5-topic45-v9-from42-v8-04.pdf}}
\par\smallskip{\small 图示：对矢势取旋度恢复毕奥---萨伐尔公式}
\end{center}
\end{minipage}\par

\par\addvspace{.4\baselineskip}
\noindent\begin{minipage}{\linewidth}

\textbf{例 M5-056：量纲检验}

自由空间中的局域稳恒电流密度为
\(\mathbf J(\mathbf r^{\prime})\)。已知库仑规范积分式
\(\mathbf A=\frac{\mu_0}{4\pi}\int\mathbf J(\mathbf r^{\prime})/|\mathbf r-\mathbf r^{\prime}|\,dV^{\prime}\)。用SI单位检验
\(\mathbf A\) 的量纲，并验证取一次旋度后的单位与磁感应强度一致。

\end{minipage}\par

\par\noindent\begin{minipage}{\linewidth}
\textbf{解答}

\begin{center}
\includegraphics[width=\linewidth,height=.40\textheight,keepaspectratio]{\detokenize{assets/module-05/examples/M5-056.pdf}}
\par\smallskip{\small 图 M5-056：量纲检验}
\end{center}
\end{minipage}\par

图 M5-056 中积分变量属于源点 \(Q\)，分母 \(R\) 是源点到场点 \(P\)
的距离。检查 \(\mu_0J\,dV/R\) 的单位，再由一次空间微分得到磁场单位。

由积分解

\[
[\mu_0J\,dV/R]
=\frac{\mathrm{T\cdot m}}{\mathrm A}
\frac{\mathrm A}{\mathrm{m^2}}
\frac{\mathrm{m^3}}{\mathrm m}
=\mathrm{T\cdot m},
\]

与 \([\mathbf A]=\mathrm{T\cdot m}\)
一致；再取一次空间导数便得到特斯拉。

\subsection{五、两个典型模型}\label{m5:ux4e94ux4e24ux4e2aux5178ux578bux6a21ux578b}

\subsubsection{1.
无限长直导线}\label{m5:ux65e0ux9650ux957fux76f4ux5bfcux7ebf-1}

设无限长直导线沿 \(z\) 轴，电流为 \(I\)。圆柱对称性允许取

\[
\mathbf A=A_z(r)\mathbf e_z.
\]

圆柱坐标中

\[
B_{\varphi}=-\frac{dA_z}{dr}.
\]

由 \(B_{\varphi}=\mu_0I/(2\pi r)\) 得

\[
\boxed{
A_z(r)=-\frac{\mu_0I}{2\pi}\ln\frac{r}{r_0}
}.
\]

\(r_0\) 是人为选择的参考长度。无限长电流不是局域源，\(A_z\)
无法在无穷远趋于零，因此不能直接套用局域电流的自由空间积分解。改变
\(r_0\) 只给 \(A_z\) 加上常量，不改变磁场。

\par\noindent\begin{minipage}{\linewidth}
\begin{center}
\includegraphics[width=\linewidth,height=.40\textheight,keepaspectratio]{\detokenize{assets/module-05/theory/tikz-m5-topic45-v9-from42-v8-05.pdf}}
\par\smallskip{\small 图示：直导线的轴向磁矢势}
\end{center}
\end{minipage}\par

\par\addvspace{.4\baselineskip}
\noindent\begin{minipage}{\linewidth}

\textbf{例 M5-057：给定轴向规范中的矢势差}

真空无限长直导线沿 \(+z\) 通电流 \(I\)，外部轴向矢势取
\(A_z(r)=-\mu_0I\ln(r/r_0)/(2\pi)\)。求半径 \(r_1\text{、}r_2>0\)
两点间的矢势差，说明任意参考半径 \(r_0\) 对差值和磁场有什么影响。

\end{minipage}\par

\par\noindent\begin{minipage}{\linewidth}
\textbf{解答}

\begin{center}
\includegraphics[width=\linewidth,height=.40\textheight,keepaspectratio]{\detokenize{assets/module-05/examples/M5-057.pdf}}
\par\smallskip{\small 图 M5-057：给定轴向规范中的矢势差}
\end{center}
\end{minipage}\par

图 M5-057 中矢势沿导线轴线，环向磁场由其径向导数给出。参考半径 \(r_0\)
的选择改变零点，两个半径的矢势差不受影响。

任取 \(r_1,r_2\)，有

\[
A_z(r_2)-A_z(r_1)
=-\frac{\mu_0I}{2\pi}\ln\frac{r_2}{r_1},
\]

其中参考长度 \(r_0\) 已消去。

上述差值只在题设轴向规范族中不依赖参考半径
\(r_0\)。一般规范变换仍可能改变两点的矢势差；例如 \(\chi=cxz\) 会使
\(A_z\) 增加 \(cx\)，同时 \(A_x\) 增加 \(cz\)，磁场保持不变。

\subsubsection{2.
无限长理想螺线管}\label{m5:ux65e0ux9650ux957fux7406ux60f3ux87baux7ebfux7ba1}

设螺线管半径为 \(R\)，内部磁场为
\(B_0\mathbf e_z\)，外部磁场为零。由轴对称性取

\[
\mathbf A=A_{\varphi}(r)\mathbf e_{\phi}.
\]

对半径为 \(r\) 的同心圆周使用斯托克斯定理：

\[
2\pi rA_{\varphi}(r)=\Phi_B(r).
\]

当 \(r<R\) 时，\(\Phi_B=B_0\pi r^2\)；当 \(r>R\) 时，圆周包围的总磁通为
\(B_0\pi R^2\)。故

\[
\boxed{
A_{\varphi}(r)=
\begin{cases}
\dfrac12B_0r, & r<R,\\[6pt]
\dfrac{B_0R^2}{2r}, & r>R.
\end{cases}
}
\]

在 \(r=R\) 处 \(A_{\varphi}\)
连续，而径向导数发生变化，对应管壁上的面电流。管外虽然
\(\mathbf B=0\)，但

\[
\oint_C\mathbf A\cdot d\mathbf l
=B_0\pi R^2\ne0.
\]

因此在环绕螺线管的整个带孔外部区域内，不能用一个单值 \(\chi\) 把
\(\mathbf A\)
整体规范为零；但在不环绕螺线管的局部小区域内可以这样做。这是局部微分结论与整体拓扑结论的典型区别。

\par\noindent\begin{minipage}{\linewidth}
\begin{center}
\includegraphics[width=\linewidth,height=.40\textheight,keepaspectratio]{\detokenize{assets/module-05/original/tikz-m5-topic45-v9-from42-v8-06.png}}
\par\smallskip{\small 图示：理想螺线管内外的磁场、磁矢势和环量}
\end{center}
\end{minipage}\par

\par\addvspace{.4\baselineskip}
\noindent\begin{minipage}{\linewidth}

\textbf{例 M5-058：外部环量不随圆周半径改变}

理想无限长螺线管半径为 \(R\)，内部场
\(B_0\mathbf e_z\)、外部场为零。已知管外 \(A_{\phi}=B_0R^2/(2r)\)。对
\(r>R\) 的同轴圆周求矢势环量，并解释为何管外局部磁场为零而环量仍非零。

\end{minipage}\par

\par\noindent\begin{minipage}{\linewidth}
\textbf{解答}

\begin{center}
\includegraphics[width=\linewidth,height=.40\textheight,keepaspectratio]{\detokenize{assets/module-05/examples/M5-058.pdf}}
\par\smallskip{\small 图 M5-058：外部环量不随圆周半径改变}
\end{center}
\end{minipage}\par

图 M5-058 的外环围住蓝色螺线管截面；扩大外环时 \(A_{\phi}\) 按 \(1/r\)
减小，周长按 \(r\) 增大，乘积保持不变。

对任意 \(r>R\)，

\[
\oint_C\mathbf A\cdot d\mathbf l
=2\pi r\frac{B_0R^2}{2r}
=B_0\pi R^2.
\]

半径增大时 \(A_{\varphi}\propto1/r\)，但圆周长度正比于
\(r\)，二者乘积保持为被包围的磁通量。

\subsection{六、局域电流的磁偶极展开}\label{m5:ux516dux5c40ux57dfux7535ux6d41ux7684ux78c1ux5076ux6781ux5c55ux5f00}

\subsubsection{1. 磁偶极矩}\label{m5:ux78c1ux5076ux6781ux77e9}

对局域稳恒体电流，定义磁偶极矩

\[
\boxed{
\mathbf m=\frac12\int\mathbf r'\times\mathbf J(\mathbf r')\,dV'
}.
\]

对平面细导线闭合回路，利用面积矢量恒等式

\[
\frac12\oint\mathbf r'\times d\mathbf l'=\mathbf S,
\]

可得

\[
\boxed{\mathbf m=I\mathbf S=IS\hat{\mathbf n}}.
\]

\(\hat{\mathbf n}\) 由电流方向按右手规则确定。

\subsubsection{2. 远场展开}\label{m5:ux8fdcux573aux5c55ux5f00}

设电流分布的尺度为 \(a\)，场点满足 \(r\gg a\)。展开核函数：

\[
\frac1{|\mathbf r-\mathbf r'|}
=\frac1r+\frac{\mathbf r\cdot\mathbf r'}{r^3}
+O\!\left(\frac{a^2}{r^3}\right).
\]

对闭合细导线回路，代入矢势积分式：

\[
\mathbf A
=\frac{\mu_0I}{4\pi}
\left[
\frac1r\oint d\mathbf l'
+\frac1{r^3}\oint(\mathbf r\cdot\mathbf r')d\mathbf l'
+\cdots
\right].
\]

第一项因 \(\oint d\mathbf l'=0\) 而消失。再用

\[
\oint(\mathbf r\cdot\mathbf r')d\mathbf l'
=\mathbf S\times\mathbf r,
\]

得到磁偶极项

\[
\boxed{
\mathbf A_{\mathrm{dip}}(\mathbf r)
=\frac{\mu_0}{4\pi}
\frac{\mathbf m\times\mathbf r}{r^3}
=\frac{\mu_0}{4\pi}
\frac{\mathbf m\times\hat{\mathbf r}}{r^2}
}.
\]

对一般局域稳恒体电流，\(\nabla\cdot\mathbf J=0\) 与局域性保证
\(\int\mathbf J\,dV'=0\)，其最低非零远场项同样由上述 \(\mathbf m\)
决定。磁偶极矢势按 \(1/r^2\) 衰减，磁场再多一次空间导数，按 \(1/r^3\)
衰减。

\par\noindent\begin{minipage}{\linewidth}
\begin{center}
\includegraphics[width=\linewidth,height=.40\textheight,keepaspectratio]{\detokenize{assets/module-05/theory/tikz-m5-topic45-v9-from42-v8-07.pdf}}
\par\smallskip{\small 图示：局域闭合电流的磁偶极远场}
\end{center}
\end{minipage}\par

\par\addvspace{.4\baselineskip}
\noindent\begin{minipage}{\linewidth}

\textbf{例 M5-059：圆形小电流环}

真空中半径为 \(a\) 的单匝小圆电流环通逆时针电流 \(I\)，法向为 \(+z\)。在
\(r\gg a\)
的远场求最低非零磁多极矩及其矢势主项，说明近导线处是否仍可使用该近似。

\end{minipage}\par

\par\noindent\begin{minipage}{\linewidth}
\textbf{解答}

\begin{center}
\includegraphics[width=\linewidth,height=.40\textheight,keepaspectratio]{\detokenize{assets/module-05/examples/M5-059.pdf}}
\par\smallskip{\small 图 M5-059：圆形小电流环}
\end{center}
\end{minipage}\par

图 M5-059
的绿色轴向磁矩由小圆环右手定则确定，远场只保留这一最低非零多极矩；近导线区域不能使用偶极公式。

半径为 \(a\)、电流为 \(I\) 的圆环有

\[
\mathbf m=I\pi a^2\mathbf e_z.
\]

在球坐标中，其远场矢势只有方位分量：

\[
A_{\varphi}
=\frac{\mu_0m\sin\theta}{4\pi r^2}
=\frac{\mu_0Ia^2\sin\theta}{4r^2}.
\]

\subsection{七、磁偶极子的磁场}\label{m5:ux4e03ux78c1ux5076ux6781ux5b50ux7684ux78c1ux573a}

\subsubsection{1.
由矢势计算旋度}\label{m5:ux7531ux77e2ux52bfux8ba1ux7b97ux65cbux5ea6}

在 \(r\ne0\) 处，令 \(\mathbf F=\mathbf r/r^3\)。因 \(\mathbf m\)
为常矢量，

\[
\nabla\times(\mathbf m\times\mathbf F)
=\mathbf m(\nabla\cdot\mathbf F)
-(\mathbf m\cdot\nabla)\mathbf F.
\]

又有

\[
\nabla\cdot\frac{\mathbf r}{r^3}=0,
\qquad
(\mathbf m\cdot\nabla)\frac{\mathbf r}{r^3}
=\frac{\mathbf m}{r^3}
-\frac{3\mathbf r(\mathbf m\cdot\mathbf r)}{r^5}.
\]

于是

\[
\boxed{
\mathbf B_{\mathrm{dip}}(\mathbf r)
=\frac{\mu_0}{4\pi r^3}
\left[
3(\mathbf m\cdot\hat{\mathbf r})\hat{\mathbf r}-\mathbf m
\right],
\qquad r\ne0
}.
\]

若 \(\mathbf m=m\mathbf e_z\)，则球坐标分量为

\[
\boxed{
B_r=\frac{\mu_0}{4\pi}\frac{2m\cos\theta}{r^3},
\qquad
B_{\theta}=\frac{\mu_0}{4\pi}\frac{m\sin\theta}{r^3},
\qquad
B_{\varphi}=0
}.
\]

磁感线满足

\[
\frac{dr}{r\,d\theta}=\frac{B_r}{B_{\theta}}=2\cot\theta,
\]

积分得

\[
\boxed{r=C\sin^2\theta}.
\]

\par\noindent\begin{minipage}{\linewidth}
\begin{center}
\includegraphics[width=\linewidth,height=.40\textheight,keepaspectratio]{\detokenize{assets/module-05/original/tikz-m5-topic45-v9-from42-v8-08.png}}
\par\smallskip{\small 图示：磁偶极子的矢势方向、磁场分量和磁感线}
\end{center}
\end{minipage}\par

\par\addvspace{.4\baselineskip}
\noindent\begin{minipage}{\linewidth}

\textbf{例 M5-060：轴线与赤道面的比较}

一个沿 \(+z\) 的固定磁偶极矩 \(\mathbf m=m\mathbf e_z\)
位于原点。真空远场中，在相同距离 \(r\)
的正轴线上和赤道面上比较磁场大小及方向，忽略偶极近似的原点接触项。

\end{minipage}\par

\par\noindent\begin{minipage}{\linewidth}
\textbf{解答}

\begin{center}
\includegraphics[width=\linewidth,height=.40\textheight,keepaspectratio]{\detokenize{assets/module-05/examples/M5-060.pdf}}
\par\smallskip{\small 图 M5-060：轴线与赤道面的比较}
\end{center}
\end{minipage}\par

图 M5-060 的轴线上红箭头与磁矩同向，赤道处红箭头反向；比较的是相同 \(r\)
处的大小，不能只比方向或不同距离的场。

在偶极轴线上

\[
\mathbf B_{\mathrm{\text{轴}}}
=\frac{\mu_0}{4\pi}\frac{2\mathbf m}{r^3},
\]

方向与 \(\mathbf m\) 相同；在赤道面上

\[
\mathbf B_{\mathrm{\text{赤道}}}
=-\frac{\mu_0}{4\pi}\frac{\mathbf m}{r^3},
\]

方向与 \(\mathbf m\) 相反。同一距离处，轴线上磁场大小是赤道面上的两倍。

\subsubsection{2.
适用范围与原点项}\label{m5:ux9002ux7528ux8303ux56f4ux4e0eux539fux70b9ux9879}

上述偶极场是有限尺寸电流分布在 \(r\gg a\)
区域内的首项，不能带回电流环内部或 \(r=0\)
使用。若把偶极子理想化为严格的点源，并要求磁场在分布意义下仍满足麦克斯韦方程，则完整表达式还包含接触项

\[
\mathbf B_{\mathrm{\text{点偶极}}}
=\frac{\mu_0}{4\pi}
\frac{3(\mathbf m\cdot\hat{\mathbf r})\hat{\mathbf r}-\mathbf m}{r^3}
+\frac{2\mu_0}{3}\mathbf m\,\delta^{(3)}(\mathbf r),
\]

其中第一项在原点处须按分布理解。普通有限线圈的场没有真正的 \(\delta\)
奇点；竞赛常规远场题只使用 \(r\ne0\) 的偶极公式。

磁偶极场与电偶极场在源外具有相同的角分布：

\[
\mathbf E_{\mathrm{dip}}(\mathbf r)
=\frac{1}{4\pi\varepsilon_0r^3}
\left[3(\mathbf p\cdot\hat{\mathbf r})\hat{\mathbf r}-\mathbf p\right].
\]

但两者的源结构不同：电偶极矩可由正负电荷分离形成，而经典磁偶极矩来自闭合电流，不能解释为两个孤立磁荷。

\subsection{八、磁偶极子在外磁场中的力学效应}\label{m5:ux516bux78c1ux5076ux6781ux5b50ux5728ux5916ux78c1ux573aux4e2dux7684ux529bux5b66ux6548ux5e94}

若外磁场在偶极子尺度内变化缓慢，则磁偶极矩受到的力矩为

\[
\boxed{\boldsymbol\tau=\mathbf m\times\mathbf B_{\mathrm{ext}}},
\]

相应的取向势能为

\[
\boxed{U=-\mathbf m\cdot\mathbf B_{\mathrm{ext}}}.
\]

对大小和方向固定的局域磁偶极矩，其平动力为

\[
\boxed{
\mathbf F=\nabla(\mathbf m\cdot\mathbf B_{\mathrm{ext}})
}.
\]

在偶极子所在的无自由电流区域，\(\nabla\times\mathbf B_{\mathrm{ext}}=0\)，还可写成

\[
\mathbf F=(\mathbf m\cdot\nabla)\mathbf B_{\mathrm{ext}}.
\]

\par\addvspace{.4\baselineskip}
\noindent\begin{minipage}{\linewidth}

\textbf{例 M5-061：均匀场只能转向，不能平移}

固定大小磁偶极矩 \(\mathbf m\)
处于空间均匀的静外磁场中，磁矩与外场不平行。求合力和力矩，判断它是否会出现由该场导致的质心平移加速度，并用势能说明其转动趋势。

\end{minipage}\par

\par\noindent\begin{minipage}{\linewidth}
\textbf{解答}

\begin{center}
\includegraphics[width=\linewidth,height=.40\textheight,keepaspectratio]{\detokenize{assets/module-05/examples/M5-061.pdf}}
\par\smallskip{\small 图 M5-061：均匀场只能转向，不能平移}
\end{center}
\end{minipage}\par

图 M5-061
的蓝色外场处处相同，环上合力相消但力矩可非零；绿色磁矩会趋向蓝色场方向。

在均匀外场中 \(\nabla\mathbf B_{\mathrm{ext}}=0\)，故
\(\mathbf F=0\)；若 \(\mathbf m\) 与 \(\mathbf B_{\mathrm{ext}}\)
不平行，则仍有力矩使 \(\mathbf m\) 趋向与磁场同向，势能由高值向最低值
\(-mB\) 变化。

\subsection{九、使用条件与常见错误}\label{m5:ux4e5dux4f7fux7528ux6761ux4ef6ux4e0eux5e38ux89c1ux9519ux8bef}

\begin{enumerate}
\def\labelenumi{\arabic{enumi}.}
\tightlist
\item
  \(\mathbf A\) 依赖规范，\(\mathbf B\)
  才是经典力学中直接进入洛伦兹力的场；但单值规范变换不改变闭合回路上的
  \(\oint\mathbf A\cdot d\mathbf l\)。
\item
  \(\nabla\cdot\mathbf B=0\)
  首先保证局部矢势存在。讨论整体矢势时必须检查定义域及穿过不可缩闭曲面的磁通，不能只写``区域单连通''。
\item
  库仑规范不是完整边界条件。\(\nabla\cdot\mathbf A=0\) 后仍可能叠加满足
  \(\nabla^2\chi=0\) 的规范函数。
\item
  自由空间积分解要求局域稳恒电流及适当的无穷远条件。无限长直导线、理想螺线管和周期电流必须单独处理边界行为。
\item
  稳恒细导线电流必须闭合或包含回流路径。孤立有限线段不满足稳恒连续性方程。
\item
  \(\mathbf A\) 的方向不是磁感线方向。磁场由 \(\mathbf A\)
  的空间旋度决定，而不是由 \(\mathbf A\) 的大小直接决定。
\item
  磁偶极近似要求
  \(r\gg a\)；越靠近源，高阶多极项越不能忽略，更不能把远场式用于原点。
\item
  时变问题中规范变换必须同时改变 \(\varphi\) 与
  \(\mathbf A\)，否则一般不能保持电场不变。
\end{enumerate}

\subsection{练习与解答}\label{m5:ux7ec3ux4e60ux4e0eux89e3ux7b54-3}

\par\addvspace{.4\baselineskip}
\noindent\begin{minipage}{\linewidth}

\subsubsection{作业一：规范变换与库仑规范}\label{m5:ux4f5cux4e1aux4e00ux89c4ux8303ux53d8ux6362ux4e0eux5e93ux4ed1ux89c4ux8303}

\textbf{题号：M5-062}

\textbf{层级：} 决赛　　\textbf{方法：} 势与规范

所有场与规范函数在所研究区域内光滑；若有边界约束，规范变换还须保持这些约束。

设某一区域内 \(\mathbf B=\nabla\times\mathbf A\)。令
\(\mathbf A'=\mathbf A+\nabla\chi\)，其中 \(\chi\)
是二阶连续可微的标量函数。证明两种矢势给出同一磁场，并说明这一变换不表示磁场发生改变。若题目同时给定库仑规范，写出
\(\chi\) 还应满足的条件。

\end{minipage}\par

\par\noindent\begin{minipage}{\linewidth}
\textbf{解答}

\begin{center}
\includegraphics[width=\linewidth,height=.40\textheight,keepaspectratio]{\detokenize{assets/module-05/examples/M5-062.pdf}}
\par\smallskip{\small 图 M5-062：规范变换与库仑规范}
\end{center}
\end{minipage}\par

图 M5-062
的局部矢势可以因规范不同而改变；真正用于求磁场的是旋度。若还要求相同规范约束，规范函数须另满足调和条件及边界条件。

\[
\nabla\times\mathbf A'
=\nabla\times\mathbf A+\nabla\times(\nabla\chi)
=\nabla\times\mathbf A
=\mathbf B.
\]

第二步使用混合偏导可交换时成立的恒等式
\(\nabla\times(\nabla\chi)=\mathbf0\)。因此两种矢势只是同一磁场的不同表示，磁场没有发生改变。

若 \(\mathbf A\) 与 \(\mathbf A'\) 都满足库仑规范，则

\[
0=\nabla\cdot\mathbf A'
=\nabla\cdot\mathbf A+\nabla^2\chi
=\nabla^2\chi.
\]

故 \(\chi\) 必须是定义域内的调和函数，并同时满足题目给定的边界条件。

\par\addvspace{.4\baselineskip}
\noindent\begin{minipage}{\linewidth}

\subsubsection{作业二：均匀磁场的两种矢势}\label{m5:ux4f5cux4e1aux4e8cux5747ux5300ux78c1ux573aux7684ux4e24ux79cdux77e2ux52bf}

\textbf{题号：M5-063}

\textbf{层级：} 决赛　　\textbf{方法：} 势与规范

\(B_0\) 为常量，空间采用标准右手直角坐标。

均匀磁场为 \(\mathbf B=B_0\mathbf e_z\)。验证

\[
\mathbf A_1=\frac12\mathbf B\times\mathbf r,
\qquad
\mathbf A_2=B_0x\mathbf e_y
\]

都满足库仑规范并给出同一磁场，求连接二者的规范函数 \(\chi\)。

\end{minipage}\par

\par\noindent\begin{minipage}{\linewidth}
\textbf{解答}

\begin{center}
\includegraphics[width=\linewidth,height=.40\textheight,keepaspectratio]{\detokenize{assets/module-05/examples/M5-063.pdf}}
\par\smallskip{\small 图 M5-063：均匀磁场的两种矢势}
\end{center}
\end{minipage}\par

图 M5-063 的蓝色对称规范绕中心转动，绿色非对称规范沿 \(y\)
方向；逐项验证散度和旋度后，用二者之差求规范函数。

写成笛卡尔分量：

\[
\mathbf A_1=-\frac{B_0y}{2}\mathbf e_x
+\frac{B_0x}{2}\mathbf e_y,
\qquad
\mathbf A_2=B_0x\mathbf e_y.
\]

二者的散度分别为

\[
\nabla\cdot\mathbf A_1=0,
\qquad
\nabla\cdot\mathbf A_2=0.
\]

其旋度的 \(z\) 分量为

\[
(\nabla\times\mathbf A_1)_z
=\frac{\partial}{\partial x}\frac{B_0x}{2}
-\frac{\partial}{\partial y}\left(-\frac{B_0y}{2}\right)=B_0,
\]

\[
(\nabla\times\mathbf A_2)_z
=\frac{\partial(B_0x)}{\partial x}=B_0,
\]

其余分量均为零。又

\[
\mathbf A_2-\mathbf A_1
=\frac{B_0y}{2}\mathbf e_x
+\frac{B_0x}{2}\mathbf e_y
=\nabla\left(\frac{B_0xy}{2}\right).
\]

因此可取

\[
\boxed{\chi=\frac12B_0xy+C}.
\]

\par\addvspace{.4\baselineskip}
\noindent\begin{minipage}{\linewidth}

\subsubsection{作业三：无限长直导线的磁矢势}\label{m5:ux4f5cux4e1aux4e09ux65e0ux9650ux957fux76f4ux5bfcux7ebfux7684ux78c1ux77e2ux52bf}

\textbf{题号：M5-064}

\textbf{层级：} 复赛　　\textbf{方法：} 势与规范、安培环路

\(r\text{、}r_0>0\)，所研究区域在细导线之外，磁导率为 \(\mu_0\)；\(I>0\)
表示沿 \(+z\)。

无限长直导线沿 \(z\) 轴，电流为 \(I\)。在导线外取
\(\mathbf A=A_z(r)\mathbf e_z\)，并规定 \(A_z(r_0)=0\)。求
\(\mathbf A\)，再由旋度验证磁场方向和大小。

\end{minipage}\par

\par\noindent\begin{minipage}{\linewidth}
\textbf{解答}

\begin{center}
\includegraphics[width=\linewidth,height=.40\textheight,keepaspectratio]{\detokenize{assets/module-05/examples/M5-064.pdf}}
\par\smallskip{\small 图 M5-064：无限长直导线的磁矢势}
\end{center}
\end{minipage}\par

图 M5-064 中电流沿纸外
\(+z\)，磁场沿逆时针环向；为恢复这一方向，圆柱坐标旋度中的 \(-dA_z/dr\)
负号不能省略。

圆柱坐标中，当 \(\mathbf A=A_z(r)\mathbf e_z\) 时

\[
(\nabla\times\mathbf A)_{\varphi}=-\frac{dA_z}{dr}.
\]

由安培定律

\[
B_{\varphi}=\frac{\mu_0I}{2\pi r},
\]

故

\[
\frac{dA_z}{dr}=-\frac{\mu_0I}{2\pi r}.
\]

从 \(r_0\) 积分到 \(r\)：

\[
A_z(r)-A_z(r_0)
=-\frac{\mu_0I}{2\pi}\ln\frac r{r_0}.
\]

利用 \(A_z(r_0)=0\)，得到

\[
\boxed{
\mathbf A(r)
=-\frac{\mu_0I}{2\pi}\ln\frac r{r_0}\,\mathbf e_z
}.
\]

于是

\[
\nabla\times\mathbf A
=-\frac{dA_z}{dr}\mathbf e_{\phi}
=\frac{\mu_0I}{2\pi r}\mathbf e_{\phi},
\]

方向符合右手螺旋规则。

\par\addvspace{.4\baselineskip}
\noindent\begin{minipage}{\linewidth}

\subsubsection{作业四：理想螺线管内外的矢势}\label{m5:ux4f5cux4e1aux56dbux7406ux60f3ux87baux7ebfux7ba1ux5185ux5916ux7684ux77e2ux52bf}

\textbf{题号：M5-065}

\textbf{层级：} 决赛　　\textbf{方法：} 势与规范、安培环路

不叠加奇异轴线矢势，选取中心正则、轴对称的环向矢势；\(B_0>0\)。

无限长理想螺线管半径为 \(R\)，内部磁场为
\(B_0\mathbf e_z\)，外部磁场为零。利用磁通量与矢势环量的关系求
\(A_{\varphi}(r)\)，并验证管外虽有
\(\mathbf B=0\)，但环绕螺线管的矢势环量不为零。

\end{minipage}\par

\par\noindent\begin{minipage}{\linewidth}
\textbf{解答}

\begin{center}
\includegraphics[width=\linewidth,height=.40\textheight,keepaspectratio]{\detokenize{assets/module-05/examples/M5-065.pdf}}
\par\smallskip{\small 图 M5-065：理想螺线管内外的矢势}
\end{center}
\end{minipage}\par

图 M5-065 的圆周分别取在蓝色磁区内外；内区所围面积随 \(r^2\)
变，外区包围磁通固定，给出两段不同的矢势。

由轴对称性取 \(\mathbf A=A_{\varphi}(r)\mathbf e_{\phi}\)。对半径为
\(r\) 的圆周，

\[
2\pi rA_{\varphi}(r)=\Phi_B(r).
\]

当 \(r<R\) 时，

\[
\Phi_B(r)=B_0\pi r^2,
\qquad
A_{\varphi}(r)=\frac12B_0r.
\]

当 \(r>R\) 时，

\[
\Phi_B(r)=B_0\pi R^2,
\qquad
A_{\varphi}(r)=\frac{B_0R^2}{2r}.
\]

\par\noindent\begin{minipage}{\linewidth}
\begin{center}
\includegraphics[width=\linewidth,height=.40\textheight,keepaspectratio]{\detokenize{assets/module-05/process/M5-065a.pdf}}
\par\smallskip{\small 图 M5-065a：螺线管环向矢势的两区变化}
\end{center}
\end{minipage}\par

图 M5-065a 的矢势在管壁连续，外区虽然仍有矢势，但 \(rA_{\phi}\)
为常量，所以外区旋度为零。

因此

\[
\boxed{
\mathbf A=
\begin{cases}
\dfrac12B_0r\,\mathbf e_{\phi}, & r<R,\\[6pt]
\dfrac{B_0R^2}{2r}\,\mathbf e_{\phi}, & r>R.
\end{cases}
}
\]

圆柱坐标中

\[
B_z=\frac1r\frac{d}{dr}(rA_{\varphi}).
\]

代入内外两式，分别得到 \(B_z=B_0\) 与 \(B_z=0\)。但对任意 \(r>R\)，

\[
\oint_C\mathbf A\cdot d\mathbf l
=2\pi rA_{\varphi}(r)
=B_0\pi R^2\ne0.
\]

\par\addvspace{.4\baselineskip}
\noindent\begin{minipage}{\linewidth}

\subsubsection{作业五：从矢势积分解推出毕奥-萨伐尔定律}\label{m5:ux4f5cux4e1aux4e94ux4eceux77e2ux52bfux79efux5206ux89e3ux63a8ux51faux6bd5ux5965-ux8428ux4f10ux5c14ux5b9aux5f8b}

\textbf{题号：M5-066}

\textbf{层级：} 决赛　　\textbf{方法：} 势与规范、毕奥---萨伐尔

电流连续、局域且稳恒，\(\nabla^{\prime}\cdot\mathbf J=0\)，并在无穷远足够快衰减，允许交换积分与空间微分。

设 \(\mathbf J(\mathbf r')\) 是局域稳恒电流。证明

\[
\mathbf A(\mathbf r)=\frac{\mu_0}{4\pi}
\int\frac{\mathbf J(\mathbf r')}{|\mathbf r-\mathbf r'|}\,dV'
\]

满足库仑规范，并由它推出毕奥-萨伐尔定律。

\end{minipage}\par

\par\noindent\begin{minipage}{\linewidth}
\textbf{解答}

\begin{center}
\includegraphics[width=\linewidth,height=.40\textheight,keepaspectratio]{\detokenize{assets/module-05/examples/M5-066.pdf}}
\par\smallskip{\small 图 M5-066：从矢势积分解推出毕奥-萨伐尔定律}
\end{center}
\end{minipage}\par

图 M5-066 中 \(Q\) 是被积分的源点，\(P\)
是取导数的场点。先用稳恒性和局域性验证库仑规范，再对 \(1/R\)
取旋度恢复叉乘。

令 \(R=|\mathbf r-\mathbf r'|\)。因

\[
\nabla\frac1R=-\nabla'\frac1R,
\]

故

\[
\begin{aligned}
\nabla\cdot\mathbf A
&=\frac{\mu_0}{4\pi}
\int\mathbf J\cdot\nabla\frac1R\,dV'\\
&=-\frac{\mu_0}{4\pi}
\int\left[
\nabla'\cdot\left(\frac{\mathbf J}{R}\right)
-\frac{\nabla'\cdot\mathbf J}{R}
\right]dV'.
\end{aligned}
\]

局域性使第一项化成的无穷远表面积分为零，稳恒条件使
\(\nabla'\cdot\mathbf J=0\)，所以 \(\nabla\cdot\mathbf A=0\)。

再取旋度：

\[
\begin{aligned}
\mathbf B
&=\nabla\times\mathbf A\\
&=\frac{\mu_0}{4\pi}
\int\nabla\left(\frac1R\right)\times\mathbf J\,dV'\\
&=\frac{\mu_0}{4\pi}
\int\frac{\mathbf J(\mathbf r')\times(\mathbf r-\mathbf r')}{|\mathbf r-\mathbf r'|^3}\,dV'.
\end{aligned}
\]

这就是毕奥-萨伐尔定律。

\par\addvspace{.4\baselineskip}
\noindent\begin{minipage}{\linewidth}

\subsubsection{作业六：圆形电流环的磁偶极远场}\label{m5:ux4f5cux4e1aux516dux5706ux5f62ux7535ux6d41ux73afux7684ux78c1ux5076ux6781ux8fdcux573a}

\textbf{题号：M5-067}

\textbf{层级：} 决赛　　\textbf{方法：} 势与规范、多极展开

逆时针是从 \(+z\) 方向正视线圈时的绕向；\(r\text{、}\theta\) 为以 \(+z\)
为极轴的球坐标。

半径为 \(a\) 的圆形线圈位于 \(xy\) 平面，通有逆时针电流 \(I\)。在
\(r\gg a\)
的区域求磁偶极矩、矢势和磁场的球坐标分量，并比较同一距离处轴线与赤道面上的磁场大小。

\end{minipage}\par

\par\noindent\begin{minipage}{\linewidth}
\textbf{解答}

\begin{center}
\includegraphics[width=\linewidth,height=.40\textheight,keepaspectratio]{\detokenize{assets/module-05/examples/M5-067.pdf}}
\par\smallskip{\small 图 M5-067：圆形电流环的磁偶极远场}
\end{center}
\end{minipage}\par

图 M5-067 中绿色磁矩沿
\(+z\)；轴线与赤道面的红色磁场方向用于检验球坐标分量的正负号和二比一的大小关系。

由右手规则，

\[
\mathbf m=I\pi a^2\mathbf e_z.
\]

偶极矢势为

\[
\mathbf A
=\frac{\mu_0}{4\pi}
\frac{\mathbf m\times\hat{\mathbf r}}{r^2}.
\]

由于 \(\mathbf e_z\times\hat{\mathbf r}=\sin\theta\,\mathbf e_{\phi}\)，

\[
\boxed{
A_{\varphi}=\frac{\mu_0Ia^2}{4r^2}\sin\theta
}.
\]

取旋度或直接使用偶极场公式，得

\[
\boxed{
B_r=\frac{\mu_0Ia^2}{2r^3}\cos\theta,
\qquad
B_{\theta}=\frac{\mu_0Ia^2}{4r^3}\sin\theta,
\qquad
B_{\varphi}=0
}.
\]

轴线上 \(\theta=0\)：

\[
B_{\mathrm{\text{轴}}}=\frac{\mu_0Ia^2}{2r^3}.
\]

赤道面上 \(\theta=\pi/2\)：

\[
B_{\mathrm{\text{赤道}}}=\frac{\mu_0Ia^2}{4r^3}.
\]

所以

\[
\boxed{B_{\mathrm{\text{轴}}}:B_{\mathrm{\text{赤道}}}=2:1}.
\]

轴线处磁场沿 \(+\mathbf e_z\)，赤道面处磁场沿 \(-\mathbf e_z\)。

\par\addvspace{.4\baselineskip}
\noindent\begin{minipage}{\linewidth}

\subsubsection{作业七：满足麦克斯韦约束的非均匀外场}\label{m5:ux4f5cux4e1aux4e03ux6ee1ux8db3ux9ea6ux514bux65afux97e6ux7ea6ux675fux7684ux975eux5747ux5300ux5916ux573a}

\textbf{题号：M5-068}

\textbf{层级：} 决赛　　\textbf{方法：} 势与规范、麦克斯韦方程

\(B_0\text{、}G\text{、}m\)
为常量，磁矩固定而非感应磁矩；只计算外场对它的作用，偶极近似有效。

原点附近的外磁场为

\[
\mathbf B_{\mathrm{ext}}
=B_0\mathbf e_z
+G\left(z\mathbf e_z-\frac{x}{2}\mathbf e_x-\frac{y}{2}\mathbf e_y\right).
\]

一个固定磁偶极矩 \(\mathbf m=m\mathbf e_z\)
位于原点。验证该外场在原点附近满足
\(\nabla\cdot\mathbf B_{\mathrm{ext}}=0\) 与
\(\nabla\times\mathbf B_{\mathrm{ext}}=0\)，求偶极子所受力和力矩。

\end{minipage}\par

\par\noindent\begin{minipage}{\linewidth}
\textbf{解答}

\begin{center}
\includegraphics[width=\linewidth,height=.40\textheight,keepaspectratio]{\detokenize{assets/module-05/examples/M5-068.pdf}}
\par\smallskip{\small 图 M5-068：无源梯度场与固定磁偶极子}
\end{center}
\end{minipage}\par

图 M5-068 在本题只把绿色箭头视为原点处固定偶极矩；给定外场的轴向梯度
\(G\) 导致轴向力，原点处磁矩与外场平行使力矩为零。

散度为

\[
\nabla\cdot\mathbf B_{\mathrm{ext}}
=-\frac G2-\frac G2+G=0.
\]

各分量只依赖于相应坐标，且没有交叉偏导，因此

\[
\nabla\times\mathbf B_{\mathrm{ext}}=\mathbf0.
\]

偶极子势能为

\[
U=-\mathbf m\cdot\mathbf B_{\mathrm{ext}}
=-m(B_0+Gz).
\]

所以

\[
\boxed{
\mathbf F=-\nabla U=mG\mathbf e_z
}.
\]

在原点处 \(\mathbf B_{\mathrm{ext}}=B_0\mathbf e_z\) 与 \(\mathbf m\)
平行，故

\[
\boxed{\boldsymbol\tau=\mathbf m\times\mathbf B_{\mathrm{ext}}=\mathbf0}.
\]

\section{第二单元：电磁感应、储能与机电耦合}\label{m5:ux7b2cux4e8cux5355ux5143ux7535ux78c1ux611fux5e94ux50a8ux80fdux4e0eux673aux7535ux8026ux5408}

专题四十六至五十四讨论磁通变化、自感互感、磁能和机械运动之间的关系。

\section{专题四十六：电磁感应、动生电动势与方向判定（基础必讲与拔高选讲）}\label{m5:ux4e13ux9898ux56dbux5341ux516dux7535ux78c1ux611fux5e94ux52a8ux751fux7535ux52a8ux52bfux4e0eux65b9ux5411ux5224ux5b9aux57faux7840ux5fc5ux8bb2ux4e0eux62d4ux9ad8ux9009ux8bb2}

\begin{quote}
阅读准备：移动边界会扫过新的面积 对固定曲面求 \(d\Phi_B/dt\) 只需考虑
\(\partial\mathbf B/\partial t\)；回路移动或变形时，还要计入边界在时间
\(dt\)
内扫过的窄条面积。窄条的有向面积元由边界速度与线元共同决定，这正是动生项出现的几何来源。
\end{quote}

\par\addvspace{.4\baselineskip}
\noindent\begin{minipage}{\linewidth}

\textbf{例 M5-069：滑动导体扫过的面积}

长度为 \(\ell\) 的直金属杆沿两平行导轨以恒速 \(v\)
垂直自身移动，磁场垂直导轨平面且大小为 \(B\)。忽略边缘场，求时间 \(dt\)
内扫过的面积、磁通变化率大小及动生电动势大小。

\end{minipage}\par

\par\noindent\begin{minipage}{\linewidth}
\textbf{解答}

\begin{center}
\includegraphics[width=\linewidth,height=.40\textheight,keepaspectratio]{\detokenize{assets/module-05/examples/M5-069.pdf}}
\par\smallskip{\small 图 M5-069：滑动导体扫过的面积}
\end{center}
\end{minipage}\par

图 M5-069 的滑杆向右运动，导轨间距为 \(\ell\)。时间 \(dt\)
内扫过的窄条面积是 \(\ell v\,dt\)，这就是动生电动势的几何来源。

长度为 \(\ell\) 的杆以速度 \(v\)
垂直自身运动，在匀强垂直磁场中每秒扫过面积
\(\ell v\)，因而磁通变化率的大小为 \(B\ell v\)。

电磁感应题需要区分空间中的感生电场、运动导体中的磁力项和闭合回路的总电动势。先确定参考系、回路、各段速度和有向磁通，再从洛伦兹力定义计算。

\subsection{一、电动势的定义}\label{m5:ux4e00ux7535ux52a8ux52bfux7684ux5b9aux4e49}

\subsubsection{1.
闭合回路上的单位电荷做功}\label{m5:ux95edux5408ux56deux8defux4e0aux7684ux5355ux4f4dux7535ux8377ux505aux529f}

在实验室参考系中，设细导线闭合回路为 \(C(t)\)，导线元的局部速度为
\(\mathbf u(\mathbf r,t)\)。沿选定正向的总电动势定义为

\[
\boxed{
\mathcal E
=\oint_{C(t)}
(\mathbf E+\mathbf u\times\mathbf B)\cdot d\mathbf l
}.
\]

这里的 \(\mathbf E\)
包含静电场和由变化磁场产生的感生电场。载流子的总速度还含有沿导线的漂移速度，但漂移速度平行于
\(d\mathbf l\)，其磁力与 \(d\mathbf l\)
垂直，对上式的线积分没有贡献。因此在细导线近似下，动生项只需使用导线元速度
\(\mathbf u\)。

电动势的单位是伏特，但它不是``回路某两点之间的静电势差''。当
\(\nabla\times\mathbf E\ne0\) 时，不能在整个区域内用单值静电势表示
\(\mathbf E\)，闭合回路上的电动势可以非零。

\subsubsection{2.
两种常见极限}\label{m5:ux4e24ux79cdux5e38ux89c1ux6781ux9650}

若回路静止，\(\mathbf u=\mathbf0\)，则

\[
\mathcal E=\oint_C\mathbf E\cdot d\mathbf l.
\]

若磁场在实验室系中不随时间变化，而导体运动，则

\[
\mathcal E=\oint_C(\mathbf u\times\mathbf B)\cdot d\mathbf l.
\]

在一般问题中，\(\mathbf E\) 与 \(\mathbf u\times\mathbf B\)
可以同时存在，不能把``感生电动势''和``动生电动势''当作互相排斥的两条定律。

\par\noindent\begin{minipage}{\linewidth}
\begin{center}
\includegraphics[width=\linewidth,height=.40\textheight,keepaspectratio]{\detokenize{assets/module-05/original/tikz-m5-topic46-v9-from43-v8-01.png}}
\par\smallskip{\small 图示：静止回路中的涡旋电场与静磁场中的运动导体都服从洛伦兹力定义}
\end{center}
\end{minipage}\par

\par\addvspace{.4\baselineskip}
\noindent\begin{minipage}{\linewidth}

\textbf{例 M5-070：只做存在性判断}

比较两个情形：固定圆环位于随时间增强的轴向匀强磁场中；开路直金属杆在稳恒匀强磁场中平移且
\(\mathbf v\times\mathbf B\)
沿杆。分别判断电动势能否存在，说明各自由电场项还是运动项产生。

\end{minipage}\par

\par\noindent\begin{minipage}{\linewidth}
\textbf{解答}

\begin{center}
\includegraphics[width=\linewidth,height=.40\textheight,keepaspectratio]{\detokenize{assets/module-05/examples/M5-070.pdf}}
\par\smallskip{\small 图 M5-070：只做存在性判断}
\end{center}
\end{minipage}\par

图 M5-070
的圆周固定不动，增强的轴向磁场仍产生顺时针电场；此情形不需要导线有整体运动。第二个情形的开路杆单独画在图
M5-070a 中。

\begin{center}
\includegraphics[width=\linewidth,height=.40\textheight,keepaspectratio]{\detokenize{assets/module-05/process/M5-070a.pdf}}\par\smallskip{\small 图 M5-070a：开路杆中的电荷分离与端电压}
\end{center}

在图 M5-070a 中，杆沿 \(+x\) 运动，磁场沿 \(+z\)，故正电荷所受磁力沿
\(-y\)，积累在 \(A\) 端。开路稳态下 \(I=0\)，杆内电场沿 \(+y\)，满足
\(\mathbf E+\mathbf v\times\mathbf B=0\)；这并不消去运动项的电动势，而是使它与静电力平衡。两端电压为
\(V_A-V_B=B\ell v\)。

静止圆环放在随时间增强的轴向磁场中，即使导线内电荷没有随导线整体运动，也有
\(\oint\mathbf E\cdot d\mathbf l\ne0\)。相反，一根金属杆在稳恒磁场中平移时，实验室中的
\(\partial\mathbf B/\partial t=0\)，但杆内仍可由
\(\mathbf u\times\mathbf B\) 产生端电压。

\subsection{二、固定回路的法拉第定律}\label{m5:ux4e8cux56faux5b9aux56deux8defux7684ux6cd5ux62c9ux7b2cux5b9aux5f8b}

麦克斯韦方程

\[
\nabla\times\mathbf E=-\frac{\partial\mathbf B}{\partial t}
\]

对一个固定曲面 \(S\) 使用斯托克斯定理，得到

\[
\boxed{
\oint_C\mathbf E\cdot d\mathbf l
=-\int_S\frac{\partial\mathbf B}{\partial t}\cdot d\mathbf S
}.
\]

因为此时边界与曲面都不运动，才可以写成

\[
\boxed{
\mathcal E=-\frac{d\Phi_B}{dt}},
\qquad
\Phi_B=\int_S\mathbf B\cdot d\mathbf S.
\]

固定回路中的电动势来自空间本身的非保守电场，并不要求导线实际存在。即使把导线拿走，变化磁场周围的涡旋电场仍然存在。

\par\addvspace{.4\baselineskip}
\noindent\begin{minipage}{\linewidth}

\textbf{例 M5-071：均匀变化磁场中的固定圆环}

半径为 \(r\)
的固定圆形导线环全部位于轴向均匀磁场中，\(\mathbf B=B(t)\mathbf e_z\) 且
\(\dot B=\beta>0\)。以 \(+z\)
为正法向、逆时针为正环向，求磁通量与有符号感应电动势。

\end{minipage}\par

\par\noindent\begin{minipage}{\linewidth}
\textbf{解答}

\begin{center}
\includegraphics[width=\linewidth,height=.40\textheight,keepaspectratio]{\detokenize{assets/module-05/examples/M5-071.pdf}}
\par\smallskip{\small 图 M5-071：均匀变化磁场中的固定圆环}
\end{center}
\end{minipage}\par

图 M5-071
的磁场向纸外增强；若取逆时针为正向，实际顺时针电场就对应负电动势。

半径为 \(r\) 的圆环所在区域内，磁场垂直环面且满足
\(dB/dt=\beta\)。取正法向与 \(\mathbf B\) 同向，则

\[
\mathcal E=-\frac{d}{dt}(B\pi r^2)=-\pi r^2\beta.
\]

负号表示实际感生方向与所选回路正向相反。

\subsection{三、运动回路的磁通输运公式}\label{m5:ux4e09ux8fd0ux52a8ux56deux8defux7684ux78c1ux901aux8f93ux8fd0ux516cux5f0f}

\subsubsection{1.
移动曲面的全导数}\label{m5:ux79fbux52a8ux66f2ux9762ux7684ux5168ux5bfcux6570}

令 \(S(t)\) 随边界 \(C(t)\) 运动，边界元速度为 \(\mathbf u\)。在时间
\(dt\)
内，边界扫过一条狭窄曲面。把磁通变化分成``固定位置上的磁场变化''和``边界扫过面积''两部分，可得

\[
\boxed{
\frac{d}{dt}\int_{S(t)}\mathbf B\cdot d\mathbf S
=\int_{S(t)}\frac{\partial\mathbf B}{\partial t}\cdot d\mathbf S
-\oint_{C(t)}(\mathbf u\times\mathbf B)\cdot d\mathbf l
}.
\]

符号可由局部扫掠面积检查。对按右手定则定向的边界元
\(d\mathbf l\)，新增曲面的面积矢量为

\[
d\mathbf S_{\mathrm{\text{扫}}}=\mathbf u\,dt\times d\mathbf l,
\]

所以

\[
\mathbf B\cdot d\mathbf S_{\mathrm{\text{扫}}}
=-(\mathbf u\times\mathbf B)\cdot d\mathbf l\,dt.
\]

这里的 \(d/dt\) 是对移动曲面的全导数，一般不能替换成
\(\partial/\partial t\)。

\subsubsection{2.
与麦克斯韦方程合并}\label{m5:ux4e0eux9ea6ux514bux65afux97e6ux65b9ux7a0bux5408ux5e76}

由固定空间曲线上的法拉第定律

\[
\oint_{C(t)}\mathbf E\cdot d\mathbf l
=-\int_{S(t)}\frac{\partial\mathbf B}{\partial t}\cdot d\mathbf S,
\]

再与磁通输运公式相加，得到

\[
\boxed{
\oint_{C(t)}
(\mathbf E+\mathbf u\times\mathbf B)\cdot d\mathbf l
=-\frac{d}{dt}\int_{S(t)}\mathbf B\cdot d\mathbf S
}.
\]

\par\noindent\begin{minipage}{\linewidth}
\begin{center}
\includegraphics[width=\linewidth,height=.40\textheight,keepaspectratio]{\detokenize{assets/module-05/theory/tikz-m5-topic46-v9-from43-v8-02.pdf}}
\par\smallskip{\small 图示：移动曲面的磁通全导数}
\end{center}
\end{minipage}\par

最后一个等号的使用条件比洛伦兹力积分更严格：所选几何边界的局部速度必须与每一段导体的实际速度一致。普通移动或变形线框、滑杆与固定导轨组成的分段回路都可据此处理；有滑动电刷的转盘等装置中，几何回路未必随导体材料一起运动，此时直接积分
\(\mathbf E+\mathbf u\times\mathbf B\) 更可靠。

\par\addvspace{.4\baselineskip}
\noindent\begin{minipage}{\linewidth}

\textbf{例 M5-072：半径变化且磁场也变化的圆环}

一个平面圆环半径为 \(a(t)>0\)，处于垂直环面的空间均匀磁场
\(B(t)\mathbf e_z\) 中。正法向取
\(+z\)，边界与其配对。求完整电动势，并区分固定边界的感生项和边界径向运动的动生项。

\end{minipage}\par

\par\noindent\begin{minipage}{\linewidth}
\textbf{解答}

\begin{center}
\includegraphics[width=\linewidth,height=.40\textheight,keepaspectratio]{\detokenize{assets/module-05/examples/M5-072.pdf}}
\par\smallskip{\small 图 M5-072：半径变化且磁场也变化的圆环}
\end{center}
\end{minipage}\par

图 M5-072 的绿色箭头表示圆环边界向外运动，红色箭头示意
\(\dot B>0\text{、}\dot a>0\) 时的感应方向。对 \(\pi a^2B\)
求全导数，得到两个来源。

圆环半径为 \(a(t)\)，垂直于空间均匀磁场
\(B(t)\)。若取磁场方向为正法向，则

\[
\Phi_B=\pi a^2B,
\]

\[
\boxed{
\mathcal E
=-\pi a^2\frac{dB}{dt}
-2\pi aB\frac{da}{dt}
}.
\]

第一项是感生贡献，第二项是边界运动贡献。只保留其中任意一项都会漏掉一般情形。

\subsection{四、方向、符号与楞次定律}\label{m5:ux56dbux65b9ux5411ux7b26ux53f7ux4e0eux695eux6b21ux5b9aux5f8b}

方向题统一按以下顺序处理：

\begin{enumerate}
\def\labelenumi{\arabic{enumi}.}
\tightlist
\item
  任意规定闭合回路的正向。
\item
  用右手定则确定与该正向配对的正法向 \(\mathbf n\)。
\item
  写有向磁通
\end{enumerate}

\[
\Phi_B=\int_S\mathbf B\cdot\mathbf n\,dS.
\]

\begin{enumerate}
\def\labelenumi{\arabic{enumi}.}
\setcounter{enumi}{3}
\tightlist
\item
  计算 \(\mathcal E=-d\Phi_B/dt\)。
\item
  若 \(\mathcal E>0\)，实际电动势沿规定正向；若
  \(\mathcal E<0\)，实际方向与规定正向相反。
\end{enumerate}

\par\noindent\begin{minipage}{\linewidth}
\begin{center}
\includegraphics[width=\linewidth,height=.40\textheight,keepaspectratio]{\detokenize{assets/module-05/theory/tikz-m5-topic46-v9-from43-v8-03.pdf}}
\par\smallskip{\small 图示：法拉第负号与一致的取向约定}
\end{center}
\end{minipage}\par

楞次定律不是额外添加的一条``反向手势规则''。它说明感应效应总是抵抗引起它的磁通变化，正是法拉第定律负号的物理内容。若同时翻转回路正向与面积法向，\(\Phi_B\)
和 \(\mathcal E\) 的符号都会改变，但实际电流方向不变。

\par\addvspace{.4\baselineskip}
\noindent\begin{minipage}{\linewidth}

\textbf{例 M5-073：向纸外的磁场增强}

固定闭合线圈的正环向为逆时针，正法向指向纸外。向纸外的磁场随时间增强。判断
\(d\Phi_B/dt\)、电动势的符号以及电阻回路中的实际电流方向，并说明其感应磁场怎样抵抗磁通变化。

\end{minipage}\par

\par\noindent\begin{minipage}{\linewidth}
\textbf{解答}

\begin{center}
\includegraphics[width=\linewidth,height=.40\textheight,keepaspectratio]{\detokenize{assets/module-05/examples/M5-073.pdf}}
\par\smallskip{\small 图 M5-073：向纸外的磁场增强}
\end{center}
\end{minipage}\par

图 M5-073
的蓝色圆点表示纸外磁场，红色顺时针箭头对应负环量，正好抵抗纸外磁通的增加。

规定逆时针为正向，则正法向指向纸外，\(d\Phi_B/dt>0\)，所以
\(\mathcal E<0\)。实际感应方向为顺时针，其产生的磁场指向纸内，抵抗向纸外磁通量的增加。

\subsection{五、变化磁场产生的涡旋电场}\label{m5:ux4e94ux53d8ux5316ux78c1ux573aux4ea7ux751fux7684ux6da1ux65cbux7535ux573a}

设只有半径 \(a\) 的圆柱区域内存在均匀轴向磁场
\(B(t)\mathbf e_z\)，圆柱外磁场为零。由轴对称性，感生电场只能取环向形式

\[
\mathbf E=E_{\varphi}(r)\mathbf e_{\phi}.
\]

对半径为 \(r\) 的同心圆使用法拉第定律。

当 \(r<a\) 时，

\[
2\pi rE_{\varphi}
=-\pi r^2\frac{dB}{dt},
\]

故

\[
E_{\varphi}(r)=-\frac r2\frac{dB}{dt}.
\]

当 \(r>a\) 时，圆周包围的变化磁通只有 \(\pi a^2B\)，所以

\[
2\pi rE_{\varphi}
=-\pi a^2\frac{dB}{dt},
\]

从而

\[
\boxed{
E_{\varphi}(r)=
\begin{cases}
-\dfrac r2\dfrac{dB}{dt}, & r<a,\\[6pt]
-\dfrac{a^2}{2r}\dfrac{dB}{dt}, & r>a.
\end{cases}
}.
\]

\par\noindent\begin{minipage}{\linewidth}
\begin{center}
\includegraphics[width=\linewidth,height=.40\textheight,keepaspectratio]{\detokenize{assets/module-05/theory/tikz-m5-topic46-v9-from43-v8-04.pdf}}
\par\smallskip{\small 图示：有限圆形磁区的感生电场}
\end{center}
\end{minipage}\par

\(E_{\varphi}\) 在 \(r=a\) 处连续：内区按 \(r\) 增大，外区按 \(1/r\)
衰减。这一结果还说明，某点即使满足
\(\partial\mathbf B/\partial t=0\)，也可能存在非零感生电场；局部旋度方程必须结合整个环路所包围的变化磁通理解。

\par\addvspace{.4\baselineskip}
\noindent\begin{minipage}{\linewidth}

\textbf{例 M5-074：圆柱外的感生电场}

真空中半径 \(a\) 的无限长圆柱磁区内 \(\dot B_z=\beta>0\)，区外
\(B=0\)。无其他电荷和外加静电场，系统轴对称。求外部 \(r=2a\)
处感生电场的大小和方向，说明为何局部无磁场不等于感生电场为零。

\end{minipage}\par

\par\noindent\begin{minipage}{\linewidth}
\textbf{解答}

\begin{center}
\includegraphics[width=\linewidth,height=.40\textheight,keepaspectratio]{\detokenize{assets/module-05/examples/M5-074.pdf}}
\par\smallskip{\small 图 M5-074：圆柱外的感生电场}
\end{center}
\end{minipage}\par

图 M5-074 的磁区半径为 \(a\)；本题场点在红色虚圆上的 \(r=2a\)
处，所围磁通仍是 \(\pi a^2B\)，环量除以 \(2\pi r\) 得场强。

若 \(dB/dt=\beta>0\)，在 \(r=2a\) 处

\[
E_{\varphi}(2a)=-\frac a4\beta.
\]

负号表示相对于逆时针 \(\mathbf e_{\phi}\)，实际电场沿顺时针。

\subsection{六、直线运动导体与滑杆模型}\label{m5:ux516dux76f4ux7ebfux8fd0ux52a8ux5bfcux4f53ux4e0eux6ed1ux6746ux6a21ux578b}

\subsubsection{1.
开路金属杆的电荷分离}\label{m5:ux5f00ux8defux91d1ux5c5eux6746ux7684ux7535ux8377ux5206ux79bb}

长度为 \(\ell\) 的直杆以速度 \(\mathbf v\)
在匀强磁场中平移，并满足杆方向与 \(\mathbf v\times\mathbf B\)
平行。正电荷先受到磁力
\(q\mathbf v\times\mathbf B\)，在杆两端积累异号电荷。开路稳态时，杆内静电场满足

\[
q(\mathbf E+\mathbf v\times\mathbf B)=\mathbf0,
\]

即

\[
\boxed{\mathbf E=-\mathbf v\times\mathbf B}.
\]

若从低电势端到高电势端的杆方向与 \(\mathbf v\times\mathbf B\) 同向，则

\[
\boxed{V_{\mathrm{\text{高}}}-V_{\mathrm{\text{低}}}=B\ell v}.
\]

\par\noindent\begin{minipage}{\linewidth}
\begin{center}
\includegraphics[width=\linewidth,height=.40\textheight,keepaspectratio]{\detokenize{assets/module-05/original/tikz-m5-topic46-v9-from43-v8-05.png}}
\par\smallskip{\small 图示：开路运动金属杆中电荷分离产生的静电场最终平衡磁力}
\end{center}
\end{minipage}\par

开路杆上存在端电压，但没有持续电流。接入负载后，电荷不再处于单纯静电平衡，必须使用闭合回路电动势和电路方程。

\subsubsection{2.
闭合滑杆回路}\label{m5:ux95edux5408ux6ed1ux6746ux56deux8def}

设杆、导轨与电阻构成闭合回路。忽略杆和导轨电阻以及回路自感，匀强磁场垂直导轨平面，杆长、速度和磁场两两垂直。动生电动势大小为

\[
\boxed{
\mathcal E
=\int_{\mathrm{\text{杆}}}(\mathbf v\times\mathbf B)\cdot d\mathbf l
=B\ell v
}.
\]

若总电阻为 \(R\)，则

\[
I=\frac{B\ell v}{R}.
\]

载流杆受到的安培力大小为

\[
F_B=I\ell B
=\frac{B^2\ell^2v}{R},
\]

方向与杆的运动方向相反。外力维持匀速时

\[
P_{\mathrm{\text{外}}}=F_Bv
=\frac{B^2\ell^2v^2}{R},
\]

而焦耳热功率为

\[
P_J=I^2R
=\frac{B^2\ell^2v^2}{R}.
\]

因此

\[
\boxed{P_{\mathrm{\text{外}}}=P_J}.
\]

\par\noindent\begin{minipage}{\linewidth}
\begin{center}
\includegraphics[width=\linewidth,height=.40\textheight,keepaspectratio]{\detokenize{assets/module-05/theory/tikz-m5-topic46-v9-from43-v8-06.pdf}}
\par\smallskip{\small 图示：闭合滑杆的动生电动势与能量}
\end{center}
\end{minipage}\par

磁力对单个带电粒子的瞬时功率
\(q\mathbf v_q\cdot(\mathbf v_q\times\mathbf B)\)
为零，但这不意味着动生电动势不能传递能量。外力通过运动导体和电磁场向电路输送能量，导体内的电场再对漂移电荷做功，最终形成焦耳热。

若不能忽略回路自感 \(L\)，电流不再瞬时等于 \(B\ell v/R\)，而应满足

\[
L\frac{dI}{dt}+RI=B\ell v.
\]

这将与下一专题的暂态过程衔接。

\subsubsection{3.
整个闭合线框平移为何没有电动势}\label{m5:ux6574ux4e2aux95edux5408ux7ebfux6846ux5e73ux79fbux4e3aux4f55ux6ca1ux6709ux7535ux52a8ux52bf}

若刚性闭合线框以同一速度 \(\mathbf v\) 在均匀稳恒磁场中整体平移，则

\[
\mathcal E
=\oint_C(\mathbf v\times\mathbf B)\cdot d\mathbf l
=(\mathbf v\times\mathbf B)\cdot\oint_Cd\mathbf l=0.
\]

\par\noindent\begin{minipage}{\linewidth}
\begin{center}
\includegraphics[width=\linewidth,height=.40\textheight,keepaspectratio]{\detokenize{assets/module-05/original/tikz-m5-topic46-v9-from43-v8-07.png}}
\par\smallskip{\small 图示：刚性闭合线框在均匀静磁场中平移时各段动生电动势严格抵消}
\end{center}
\end{minipage}\par

零电动势并不表示每个电荷都不受磁力，而是各段动生项在闭合积分中抵消。若磁场不均匀、线框转动或发生形变，\(\mathbf v\times\mathbf B\)
不能作为常矢量提出积分号，结果可以非零。

\subsection{七、转动导体与转动线圈}\label{m5:ux4e03ux8f6cux52a8ux5bfcux4f53ux4e0eux8f6cux52a8ux7ebfux5708}

\subsubsection{1.
绕端点转动的金属杆}\label{m5:ux7ed5ux7aefux70b9ux8f6cux52a8ux7684ux91d1ux5c5eux6746}

长度为 \(a\) 的直杆在垂直杆所在平面的匀强磁场中，以角速度
\(\boldsymbol\omega\) 绕一端 \(O\) 转动。若 \(\boldsymbol\omega\) 与
\(\mathbf B\) 同向，则距转轴 \(r\) 处导体元的速度为

\[
\mathbf u=\boldsymbol\omega\times\mathbf r,
\]

且 \(\mathbf u\times\mathbf B\) 沿杆由轴心指向外端。微元动生电动势为

\[
d\mathcal E=B\omega r\,dr.
\]

积分得到

\[
\boxed{
V_{\mathrm{\text{边}}}-V_O
=\int_0^aB\omega r\,dr
=\frac12B\omega a^2
}.
\]

\par\noindent\begin{minipage}{\linewidth}
\begin{center}
\includegraphics[width=\linewidth,height=.40\textheight,keepaspectratio]{\detokenize{assets/module-05/original/tikz-m5-topic46-v9-from43-v8-08.png}}
\par\smallskip{\small 图示：转动金属杆各处线速度不同，动生电动势必须沿半径积分}
\end{center}
\end{minipage}\par

若 \(\boldsymbol\omega\) 与 \(\mathbf B\)
反向，端电压极性反转。法拉第圆盘具有相同的径向电压，但其外电路通常含滑动电刷：几何闭合曲线可以不动，而圆盘材料仍在运动。此时若机械套用
\(-d\Phi_B/dt\) 会错误地得到零，直接积分 \(\mathbf u\times\mathbf B\)
才能保留真正的动生项。

\subsubsection{2.
匀速转动线圈}\label{m5:ux5300ux901fux8f6cux52a8ux7ebfux5708}

面积为 \(A\)、匝数为 \(N\)
的线圈在匀强磁场中匀速转动。设线圈正法向与磁场夹角为

\[
\theta=\omega t+\theta_0.
\]

单匝磁通为 \(\Phi_B=BA\cos\theta\)，总磁链为

\[
\Psi=N\Phi_B=NBA\cos\theta.
\]

因此

\[
\boxed{
\mathcal E=-\frac{d\Psi}{dt}
=NBA\omega\sin\theta
}.
\]

\par\noindent\begin{minipage}{\linewidth}
\begin{center}
\includegraphics[width=\linewidth,height=.40\textheight,keepaspectratio]{\detokenize{assets/module-05/theory/tikz-m5-topic46-v9-from43-v8-09-margin12-v11.pdf}}
\par\smallskip{\small 图示：转动线圈的磁链与电动势}
\end{center}
\end{minipage}\par

磁链取极值时变化率为零，电动势过零；磁链过零时变化最快，电动势大小达到最大值
\(NBA\omega\)。二者相差四分之一周期。

\par\addvspace{.4\baselineskip}
\noindent\begin{minipage}{\linewidth}

\textbf{例 M5-075：机械功率与电功率}

匝数 \(N\)、单匝面积 \(A\) 的线圈在匀强外磁场 \(B\) 中以角速度
\(\omega=\dot\theta\) 转动。\(\theta\) 为正法向与外场夹角，正电流 \(i\)
使磁矩沿正法向。求有符号电动势及沿 \(\theta\)
增大方向的电磁力矩，检验机械功率与电功率关系。

\end{minipage}\par

\par\noindent\begin{minipage}{\linewidth}
\textbf{解答}

\begin{center}
\includegraphics[width=\linewidth,height=.40\textheight,keepaspectratio]{\detokenize{assets/module-05/examples/M5-075.pdf}}
\par\smallskip{\small 图 M5-075：机械功率与电功率}
\end{center}
\end{minipage}\par

图 M5-075 中 \(\theta\)
随规定转向增加，绿色法向决定正电流的磁矩。电磁力矩与感应电动势的符号配对后，机械输出功率与电路吸收功率相反。

取 \(\theta\)
增大的转动方向为正，并把正电流定义为使磁矩沿线圈正法向。电磁力矩沿转轴的有符号分量为

\[
\tau_{\mathrm{em}}=-NIAB\sin\theta.
\]

负号表示在发电状态下电磁力矩阻碍转动。忽略损耗并由外力矩维持匀速时，外力矩为
\(\tau_{\mathrm{\text{外}}}=-\tau_{\mathrm{em}}\)，故

\[
P_{\mathrm{\text{机械}}}=\tau_{\mathrm{\text{外}}}\omega
=I(NBA\omega\sin\theta)
=I\mathcal E.
\]

\subsection{八、使用条件与常见错误}\label{m5:ux516bux4f7fux7528ux6761ux4ef6ux4e0eux5e38ux89c1ux9519ux8bef}

\begin{enumerate}
\def\labelenumi{\arabic{enumi}.}
\tightlist
\item
  总电动势的基本式是
  \(\oint(\mathbf E+\mathbf u\times\mathbf B)\cdot d\mathbf l\)。感生项与动生项是它的两个来源，可以同时存在。
\item
  \(d\Phi_B/dt\)
  是移动曲面磁通的全导数；\(\partial\mathbf B/\partial t\)
  是固定空间点上的局部变化率，二者不能混用。
\item
  磁通量必须带符号。先定回路正向与正法向，再使用负号；不要先凭手势猜方向后又重复加负号。
\item
  ``切割磁感线''只能作辅助图像，不能代替洛伦兹力积分。整体线框在均匀磁场中平移时，每段都可能受磁力，但总电动势为零。
\item
  开路运动杆有端电压而无持续电流；闭合后才可写
  \(I=\mathcal E/R\)。若有自感，还必须写暂态微分方程。
\item
  磁力对单个电荷不做功，不等于电路中没有能量转化。动生问题的能量源通常是维持导体运动的外力。
\item
  有滑动接点或不同部分以不同速度运动时，应逐段使用各自的
  \(\mathbf u\)。若几何边界速度与导体速度不一致，不可机械套用磁通法。
\item
  感生电场不一定只存在于 \(\partial\mathbf B/\partial t\ne0\)
  的局部区域。环路包围的总变化磁通决定其环量。
\item
  \(\mathcal E=-d(N\Phi_B)/dt\)
  中应使用磁链而非单匝磁通；多匝线圈不能漏掉 \(N\)。
\item
  本专题默认准静态、细导线和非相对论速度。高速运动、显著辐射或分布参数效应需要更完整的电磁场理论。
\end{enumerate}

\subsection{练习与解答}\label{m5:ux7ec3ux4e60ux4e0eux89e3ux7b54-4}

\par\addvspace{.4\baselineskip}
\noindent\begin{minipage}{\linewidth}

\subsubsection{作业一：滑杆的电动势、电流与外力}\label{m5:ux4f5cux4e1aux4e00ux6ed1ux6746ux7684ux7535ux52a8ux52bfux7535ux6d41ux4e0eux5916ux529b}

\textbf{题号：M5-076}

\textbf{层级：} 基础　　\textbf{方法：} 法拉第定律

设 \(B\text{、}v\text{、}\ell\text{、}R>0\)；从杆的 \(-x\) 端到 \(+x\)
端定义电动势正向。

长度为 \(\ell\) 的金属杆沿两条平行导轨以速度 \(\mathbf v=v\mathbf j\)
向右运动，这里把 \(+y\) 方向定义为``向右''。导轨平面垂直匀强磁场
\(\mathbf B=B\mathbf k\)，杆沿 \(+x\) 方向放置。取
\(\mathbf i,\mathbf j,\mathbf k\) 为右手正交基，导轨与电阻 \(R\)
构成闭合回路。忽略回路自感及杆、导轨电阻，求杆两端的电动势、回路电流和维持匀速所需外力。

\end{minipage}\par

\par\noindent\begin{minipage}{\linewidth}
\textbf{解答}

\begin{center}
\includegraphics[width=\linewidth,height=.40\textheight,keepaspectratio]{\detokenize{assets/module-05/examples/M5-076.pdf}}
\par\smallskip{\small 图 M5-076：滑杆的电动势、电流与外力}
\end{center}
\end{minipage}\par

图 M5-076 在本题取 \(+x\) 向上、\(+y\) 向右、\(+z\)
向纸内，构成右手基。杆中正电流向上，安培力向左，维持匀速的外力向右。

杆内正电荷所受磁力方向由

\[
\mathbf v\times\mathbf B
=v\mathbf j\times B\mathbf k
=vB\mathbf i
\]

确定，所以 \(+x\) 端电势高。沿杆从 \(-x\) 端积分至 \(+x\) 端：

\[
\boxed{\mathcal E=B\ell v}.
\]

回路电流大小为

\[
\boxed{I=\frac{B\ell v}{R}}.
\]

杆上安培力为

\[
\mathbf F_B
=I\ell\mathbf i\times B\mathbf k
=-I\ell B\mathbf j
=-\frac{B^2\ell^2v}{R}\mathbf j.
\]

它阻碍杆沿 \(+y\) 方向运动，故维持匀速所需外力为

\[
\boxed{
\mathbf F_{\mathrm{\text{外}}}
=\frac{B^2\ell^2v}{R}\mathbf j
}.
\]

功率检验为

\[
\mathbf F_{\mathrm{\text{外}}}\cdot\mathbf v
=\frac{B^2\ell^2v^2}{R}
=I^2R,
\]

说明机械功率完整转化为焦耳热。

\par\addvspace{.4\baselineskip}
\noindent\begin{minipage}{\linewidth}

\subsubsection{作业二：有限变化磁场区域内外的感生电场}\label{m5:ux4f5cux4e1aux4e8cux6709ux9650ux53d8ux5316ux78c1ux573aux533aux57dfux5185ux5916ux7684ux611fux751fux7535ux573a}

\textbf{题号：M5-077}

\textbf{层级：} 基础　　\textbf{方法：} 法拉第定律

空间无外加静电场或自由电荷，采用无限长圆柱和严格轴对称模型，正环向为
\(+\mathbf e_{\phi}\)。

半径为 \(a\) 的无限长圆柱区域内存在均匀磁场

\[
\mathbf B(t)=\beta t\,\mathbf e_z,
\]

圆柱外磁场为零，其中 \(\beta>0\)。求任意半径 \(r\)
处的感生电场大小和方向。

\end{minipage}\par

\par\noindent\begin{minipage}{\linewidth}
\textbf{解答}

\begin{center}
\includegraphics[width=\linewidth,height=.40\textheight,keepaspectratio]{\detokenize{assets/module-05/examples/M5-077.pdf}}
\par\smallskip{\small 图 M5-077：有限变化磁场区域内外的感生电场}
\end{center}
\end{minipage}\par

图 M5-077 的圆柱磁区边界为 \(a\)；在内部安培式圆周包围面积
\(\pi r^2\)，外部则只包围 \(\pi a^2\)。两段环量给出连续的切向电场。

由轴对称性，\(\mathbf E=E_{\varphi}(r)\mathbf e_{\phi}\)。规定从 \(+z\)
方向观察逆时针为正向。

当 \(r<a\) 时，

\[
2\pi rE_{\varphi}
=-\frac d{dt}(\beta t\,\pi r^2),
\]

所以

\[
E_{\varphi}=-\frac{\beta r}{2}.
\]

当 \(r>a\) 时，圆周包围的磁通为 \(\beta t\,\pi a^2\)，故

\[
2\pi rE_{\varphi}
=-\frac d{dt}(\beta t\,\pi a^2),
\]

从而

\[
E_{\varphi}=-\frac{\beta a^2}{2r}.
\]

因此

\[
\boxed{
\mathbf E(r)=
\begin{cases}
-\dfrac{\beta r}{2}\mathbf e_{\phi}, & r<a,\\[6pt]
-\dfrac{\beta a^2}{2r}\mathbf e_{\phi}, & r>a.
\end{cases}
}
\]

负号表示电场沿顺时针方向，在 \(r=a\) 处两式连续。

\par\addvspace{.4\baselineskip}
\noindent\begin{minipage}{\linewidth}

\subsubsection{作业三：矩形线框进入匀强磁场区域}\label{m5:ux4f5cux4e1aux4e09ux77e9ux5f62ux7ebfux6846ux8fdbux5165ux5300ux5f3aux78c1ux573aux533aux57df}

\textbf{题号：M5-078}

\textbf{层级：} 基础　　\textbf{方法：} 法拉第定律

磁场区为沿运动方向宽于线框的有限矩形区域，边界平直并垂直速度；忽略边缘场。取逆时针为正电流，线框各段电阻包含在总电阻
\(R\) 内。

高为 \(\ell\)、沿运动方向宽为 \(a\) 的矩形线框，以速度 \(v\)
向右进入垂直纸面向里的匀强磁场区域。线框总电阻为
\(R\)，忽略自感。在线框部分进入、完全进入和部分离开三个阶段，求电动势和电流方向；在线框部分进入时求安培阻力大小。

\end{minipage}\par

\par\noindent\begin{minipage}{\linewidth}
\textbf{解答}

\begin{center}
\includegraphics[width=\linewidth,height=.40\textheight,keepaspectratio]{\detokenize{assets/module-05/examples/M5-078.pdf}}
\par\smallskip{\small 图 M5-078：矩形线框进入匀强磁场区域}
\end{center}
\end{minipage}\par

图 M5-078
表示进入阶段，蓝色叉号为纸内磁场；右侧框边中电流向上以产生纸外感应场，安培力向左。离开时磁通变化反向，电流随之反向。

在线框部分进入时，设进入磁区的宽度为 \(x\)，则向纸里的磁通量大小为

\[
\Phi=B\ell x,
\qquad
\frac{dx}{dt}=v.
\]

电动势大小为

\[
\boxed{\mathcal E=B\ell v}.
\]

\par\noindent\begin{minipage}{\linewidth}
\begin{center}
\includegraphics[width=\linewidth,height=.40\textheight,keepaspectratio]{\detokenize{assets/module-05/process/M5-078a.pdf}}
\par\smallskip{\small 图 M5-078a：进入、完全进入、离开时的磁通变化}
\end{center}
\end{minipage}\par

图 M5-078a
自上而下表示进入、完全进入与离开固定磁区；上图右边电流向上，中图无电流，下图左边电流向上，后者对应顺时针环流。

向纸里的磁通正在增加，感应电流应产生向纸外的磁场，因此从纸面正面观察，电流沿逆时针方向。电流大小为

\[
I=\frac{B\ell v}{R}.
\]

只有磁区中的前竖边受到不被另一竖边抵消的水平安培力，其大小为

\[
\boxed{
F_B=I\ell B
=\frac{B^2\ell^2v}{R}
},
\]

方向向左，阻碍进入。

线框完全进入均匀磁区后，面积和磁场均不变，\(d\Phi/dt=0\)，故
\(\mathcal E=I=0\)。

部分离开时，向纸里的磁通减小，感应电流应维持向纸里的磁场，因此电流沿顺时针方向；电动势大小仍为
\(B\ell v\)，安培力仍阻碍运动。

\par\addvspace{.4\baselineskip}
\noindent\begin{minipage}{\linewidth}

\subsubsection{作业四：转动金属杆的端电压}\label{m5:ux4f5cux4e1aux56dbux8f6cux52a8ux91d1ux5c5eux6746ux7684ux7aefux7535ux538b}

\textbf{题号：M5-079}

\textbf{层级：} 基础　　\textbf{方法：} 法拉第定律、机电耦合

\(B\text{、}\omega\text{、}a>0\)，杆开路，已达到无径向传导电流的电荷分离状态；忽略边缘场和接触电势。

长度为 \(a\) 的金属杆在匀强磁场 \(\mathbf B=B\mathbf e_z\) 中绕一端
\(O\) 以恒定角速度 \(\boldsymbol\omega=\omega\mathbf e_z\)
转动。求外端相对转轴端的电势，并讨论 \(\boldsymbol\omega\)
反向时的极性。

\end{minipage}\par

\par\noindent\begin{minipage}{\linewidth}
\textbf{解答}

\begin{center}
\includegraphics[width=\linewidth,height=.40\textheight,keepaspectratio]{\detokenize{assets/module-05/examples/M5-079.pdf}}
\par\smallskip{\small 图 M5-079：转动金属杆的端电压}
\end{center}
\end{minipage}\par

图 M5-079 中转轴端为 \(O\)、外端为 \(P\)；\(\mathbf v\times\mathbf B\)
沿杆向外，正电荷聚在外端。每一点的线速度不同，必须沿杆积分。

距转轴 \(r\) 处的速度为

\[
\mathbf u=\omega r\mathbf e_{\phi}.
\]

所以

\[
\mathbf u\times\mathbf B
=\omega rB\mathbf e_r,
\]

沿杆由转轴指向外端。端电压为

\[
V_{\mathrm{\text{边}}}-V_O
=\int_0^a(\mathbf u\times\mathbf B)\cdot d\mathbf l
=\int_0^aB\omega r\,dr.
\]

故

\[
\boxed{
V_{\mathrm{\text{边}}}-V_O=\frac12B\omega a^2
}.
\]

当 \(\boldsymbol\omega\) 反向时，\(\mathbf u\times\mathbf B\)
反向，外端变为低电势端，电压符号反转。

\par\addvspace{.4\baselineskip}
\noindent\begin{minipage}{\linewidth}

\subsubsection{作业五：匀速转动线圈的感应电动势}\label{m5:ux4f5cux4e1aux4e94ux5300ux901fux8f6cux52a8ux7ebfux5708ux7684ux611fux5e94ux7535ux52a8ux52bf}

\textbf{题号：M5-080}

\textbf{层级：} 基础　　\textbf{方法：} 法拉第定律、机电耦合

线圈刚性，外场恒定、均匀，不计线圈自感、导线截面和负载反作用。电动势正向由瞬时正法向按右手定则规定。

面积为 \(A\)、匝数为 \(N\) 的线圈在匀强磁场 \(B\)
中绕位于线圈平面内的轴匀速转动，角速度为 \(\omega\)。取 \(t=0\)
时线圈正法向与磁场同向。求磁链和感应电动势，并指出二者达到极值和过零的时刻。

\end{minipage}\par

\par\noindent\begin{minipage}{\linewidth}
\textbf{解答}

\begin{center}
\includegraphics[width=\linewidth,height=.40\textheight,keepaspectratio]{\detokenize{assets/module-05/examples/M5-080.pdf}}
\par\smallskip{\small 图 M5-080：匀速转动线圈的感应电动势}
\end{center}
\end{minipage}\par

图 M5-080 中绿色法向与外场的夹角就是
\(\omega t\)。磁链取余弦，电动势是其负时间导数，因此过零和极值相差四分之一周期。

正法向与磁场夹角为 \(\theta=\omega t\)，故磁链为

\[
\boxed{\Psi(t)=NBA\cos\omega t}.
\]

感应电动势为

\[
\boxed{
\mathcal E(t)
=-\frac{d\Psi}{dt}
=NBA\omega\sin\omega t
}.
\]

\par\noindent\begin{minipage}{\linewidth}
\begin{center}
\includegraphics[width=\linewidth,height=.40\textheight,keepaspectratio]{\detokenize{assets/module-05/process/M5-080a.pdf}}
\par\smallskip{\small 图 M5-080a：磁链和电动势相差四分之一周期}
\end{center}
\end{minipage}\par

图 M5-080a
的两条曲线只比较相位；显示时共同缩放并作纵向平移，蓝线达到峰值时红线通过自身中线，四分之一周期后交换。

当

\[
t=\frac{k\pi}{\omega}
\]

时，磁链取正、负极值，电动势为零；当

\[
t=\frac{(2k+1)\pi}{2\omega}
\]

时，磁链为零，电动势大小达到最大值 \(NBA\omega\)。两者相差四分之一周期。

\par\addvspace{.4\baselineskip}
\noindent\begin{minipage}{\linewidth}

\subsubsection{作业六：磁场变化与圆环膨胀同时发生}\label{m5:ux4f5cux4e1aux516dux78c1ux573aux53d8ux5316ux4e0eux5706ux73afux81a8ux80c0ux540cux65f6ux53d1ux751f}

\textbf{题号：M5-081}

\textbf{层级：} 基础　　\textbf{方法：} 法拉第定律

圆环始终保持平面圆形，取几何半径和场的代数轴向分量求全导数；\(a_0>0\)，不要求变化项同号。

一个柔性圆形导线环的半径为 \(a(t)\)，处在垂直环面的空间均匀磁场 \(B(t)\)
中。取固定的 \(+z\) 方向为正法向，\(B(t)\)
表示沿此方向的有符号分量，求总电动势，并分别指出感生项和动生项。若

\[
B(t)=B_0+\alpha t,
\qquad
a(t)=a_0+ut,
\]

写出 \(t=0\) 时的电动势。

\end{minipage}\par

\par\noindent\begin{minipage}{\linewidth}
\textbf{解答}

\begin{center}
\includegraphics[width=\linewidth,height=.40\textheight,keepaspectratio]{\detokenize{assets/module-05/examples/M5-081.pdf}}
\par\smallskip{\small 图 M5-081：磁场变化与圆环膨胀同时发生}
\end{center}
\end{minipage}\par

图 M5-081
中半径和磁场同时改变；第一项保留瞬时面积，第二项保留瞬时磁场。代入
\(t=0\) 后分别得到两个有限贡献。

磁通量为

\[
\Phi_B(t)=\pi a^2(t)B(t).
\]

因此

\[
\begin{aligned}
\mathcal E
&=-\frac{d\Phi_B}{dt}\\
&=-\pi a^2\frac{dB}{dt}
-2\pi aB\frac{da}{dt}.
\end{aligned}
\]

\par\noindent\begin{minipage}{\linewidth}
\begin{center}
\includegraphics[width=\linewidth,height=.40\textheight,keepaspectratio]{\detokenize{assets/module-05/process/M5-081a.pdf}}
\par\smallskip{\small 图 M5-081a：半径变化带来环形面积增量}
\end{center}
\end{minipage}\par

图 M5-081a 的径向运动扫过一圈窄环形面积，面积增量与周长 \(2\pi a\)
成正比，这正是动生项中的系数。

所以

\[
\boxed{
\mathcal E_{\mathrm{\text{感生}}}
=-\pi a^2\frac{dB}{dt},
\qquad
\mathcal E_{\mathrm{\text{动生}}}
=-2\pi aB\frac{da}{dt}
}.
\]

在 \(t=0\) 时，

\[
\boxed{
\mathcal E(0)
=-\pi a_0^2\alpha
-2\pi a_0B_0u
}.
\]

若 \(\alpha>0\)、\(u>0\) 且 \(B_0>0\)，在 \(t=0\) 两项都使沿 \(+z\)
的正向磁通增加，故总电动势沿所规定正向的反方向；一般情形应保留上式中各代数项的符号。

\par\addvspace{.4\baselineskip}
\noindent\begin{minipage}{\linewidth}

\subsubsection{作业七：非均匀磁场中的整体平移}\label{m5:ux4f5cux4e1aux4e03ux975eux5747ux5300ux78c1ux573aux4e2dux7684ux6574ux4f53ux5e73ux79fb}

\textbf{题号：M5-082}

\textbf{层级：} 基础　　\textbf{方法：} 法拉第定律

矩形位于 \(xy\) 平面，较长边平行 \(y\)；\(X(t)=X_0+vt\)，\(k\)
为有符号常量，正边界方向按 \(+z\) 配对。

矩形线框高为 \(\ell\)，沿 \(x\) 方向宽为 \(a\)，以速度 \(v\mathbf e_x\)
在磁场

\[
\mathbf B(x)=(B_0+kx)\mathbf e_z
\]

中平移。设线框左边的位置为 \(X(t)\)，取 \(+z\)
为正法向。求线框电动势，并说明为什么均匀磁场极限 \(k=0\) 时结果为零。

\end{minipage}\par

\par\noindent\begin{minipage}{\linewidth}
\textbf{解答}

\begin{center}
\includegraphics[width=\linewidth,height=.40\textheight,keepaspectratio]{\detokenize{assets/module-05/examples/M5-082.pdf}}
\par\smallskip{\small 图 M5-082：非均匀静磁场中的平移线框}
\end{center}
\end{minipage}\par

图 M5-082
用移动矩形框表示两个竖边在不同位置取样磁场；蓝色直线表示给定线性场
\(B_0+kx\)，两边的差为 \(ka\)，均匀部分抵消。

磁通量为

\[
\begin{aligned}
\Phi_B
&=\ell\int_X^{X+a}(B_0+kx)\,dx\\
&=\ell\left[B_0a+kXa+\frac12ka^2\right].
\end{aligned}
\]

由于 \(dX/dt=v\)，

\[
\frac{d\Phi_B}{dt}=k\ell av.
\]

因此

\[
\boxed{\mathcal E=-k\ell av}.
\]

也可逐段计算动生电动势：两条竖边分别位于 \(X\) 与 \(X+a\)，其贡献之差为

\[
\mathcal E
=\ell vB(X)-\ell vB(X+a)
=-k\ell av.
\]

当 \(k=0\)
时，两条竖边贡献大小相等、符号相反，严格抵消；这就是整个闭合线框在均匀静磁场中平移时电动势为零的原因。

\section{专题四十七：自感、互感与耦合电路（基础必讲与拔高选讲）}\label{m5:ux4e13ux9898ux56dbux5341ux4e03ux81eaux611fux4e92ux611fux4e0eux8026ux5408ux7535ux8defux57faux7840ux5fc5ux8bb2ux4e0eux62d4ux9ad8ux9009ux8bb2}

电磁感应把``磁场的变化''转化为回路中的电动势；电感则进一步回答：磁场由哪些电流建立，磁链如何随各回路电流变化，磁场储存了多少能量。处理这类问题时，必须始终分清磁通与磁链、感应电动势与元件端电压，以及互感的大小与符号。

除特别说明外，以下假定回路相对静止、几何形状固定，系统处于准静态条件，介质线性且满足互易性。若存在铁磁饱和、磁滞、运动部件或非互易介质，应回到磁链函数和场能定义，不能直接套用常量
\(L\)、\(M\) 的公式。

\subsection{一、磁通、磁链与电感的定义}\label{m5:ux4e00ux78c1ux901aux78c1ux94feux4e0eux7535ux611fux7684ux5b9aux4e49}

\subsubsection{1.
磁通不等于磁链}\label{m5:ux78c1ux901aux4e0dux7b49ux4e8eux78c1ux94fe}

穿过某一匝回路的有向磁通为

\[
\Phi=\int_S\mathbf B\cdot d\mathbf S.
\]

含 \(N\) 匝的线圈中，各匝磁通未必相同。总磁链定义为

\[
\boxed{\Psi=\sum_{k=1}^{N}\Phi_k}.
\]

只有各匝都链接同一磁通 \(\Phi\) 时，才可以写成
\(\Psi=N\Phi\)。磁通和磁链的单位都是韦伯，但它们代表的对象不同；电感定义使用的是磁链。

先规定每个回路的正电流方向，再用右手定则确定与之配对的正法向。磁通、磁链和互感的正负都依赖这套方向约定。

\subsubsection{2.
自感与互感的下标}\label{m5:ux81eaux611fux4e0eux4e92ux611fux7684ux4e0bux6807}

回路 1 的电流 \(i_1\) 产生并链接回路 1 的磁链记作

\[
\Psi_{11}=L_1i_1.
\]

回路 1 的电流产生磁场，而回路 2 接收到的磁链记作

\[
\boxed{\Psi_{21}=M_{21}i_1}.
\]

本讲义采用``第一个下标表示接收回路，第二个下标表示源回路''的约定。相应地，\(\Psi_{12}=M_{12}i_2\)。

\par\noindent\begin{minipage}{\linewidth}
\begin{center}
\includegraphics[width=\linewidth,height=.40\textheight,keepaspectratio]{\detokenize{assets/module-05/original/tikz-m5-topic47-v9-from44-v8-01.png}}
\par\smallskip{\small 图示：磁链、自感与互感的下标约定}
\end{center}
\end{minipage}\par

在线性系统中，两个回路的总磁链可写成

\[
\boxed{
\begin{aligned}
\Psi_1&=L_1i_1+M_{12}i_2,\\
\Psi_2&=M_{21}i_1+L_2i_2.
\end{aligned}}
\]

在静止、线性、互易的系统中，后文将证明
\(M_{12}=M_{21}\)。若只反转一个回路的正方向，\(M_{12}\) 与 \(M_{21}\)
会同时变号，自感 \(L_1,L_2\) 不变。

\par\addvspace{.4\baselineskip}
\noindent\begin{minipage}{\linewidth}

\textbf{例 M5-083：磁链的计算}

一个 \(100\) 匝的紧密线圈中，每匝都链接同方向磁通
\(2.0\times10^{-5}\ \mathrm{Wb}\)。求磁链，说明这个数值与单匝磁通的区别；各匝可视为链接同一近似均匀磁场。

\end{minipage}\par

\par\noindent\begin{minipage}{\linewidth}
\textbf{解答}

\begin{center}
\includegraphics[width=\linewidth,height=.40\textheight,keepaspectratio]{\detokenize{assets/module-05/examples/M5-083.pdf}}
\par\smallskip{\small 图 M5-083：磁链的计算}
\end{center}
\end{minipage}\par

图 M5-083
的每一匝都链接同一轴向磁通；求总磁链时要把各匝的磁通相加，不能把单匝磁通直接当磁链。

一个 \(100\) 匝线圈中，每匝都通过 \(2.0\times10^{-5}\ \mathrm{Wb}\)
的同向磁通，则

\[
\Psi=N\Phi=2.0\times10^{-3}\ \mathrm{Wb}.
\]

若只有 \(80\) 匝完全链接该磁通，其余 \(20\)
匝的磁通可忽略，则磁链只能写成求和，结果为
\(1.6\times10^{-3}\ \mathrm{Wb}\)，不能仍取 \(N\Phi\)。

\subsubsection{3.
感应电动势与端电压的符号}\label{m5:ux611fux5e94ux7535ux52a8ux52bfux4e0eux7aefux7535ux538bux7684ux7b26ux53f7}

沿回路规定正向的感应电动势为

\[
\boxed{\mathcal E_k=-\frac{d\Psi_k}{dt}}.
\]

对理想电感采用被动符号约定，即电流流入电压正端时，外电路加在元件上的端电压为

\[
\boxed{v_k=\frac{d\Psi_k}{dt}}.
\]

两式的符号不同并不矛盾：\(\mathcal E_k\)
描述感应场沿规定回路正向的做功，\(v_k\)
描述外电路为了维持电流变化而施加的端电压。

当几何和介质固定且 \(L,M\) 为常量时，

\[
\begin{aligned}
v_1&=L_1\frac{di_1}{dt}+M_{12}\frac{di_2}{dt},\\
v_2&=M_{21}\frac{di_1}{dt}+L_2\frac{di_2}{dt}.
\end{aligned}
\]

若自感随位置或时间改变，则必须保留全导数：

\[
v=\frac{d(Li)}{dt}=L\frac{di}{dt}+i\frac{dL}{dt}.
\]

\subsection{二、自感的计算与磁场能量}\label{m5:ux4e8cux81eaux611fux7684ux8ba1ux7b97ux4e0eux78c1ux573aux80fdux91cf}

\subsubsection{1.
线性电感、微分电感与能量}\label{m5:ux7ebfux6027ux7535ux611fux5faeux5206ux7535ux611fux4e0eux80fdux91cf}

在线性、无磁滞系统中，

\[
\boxed{L=\frac{\Psi}{i}},
\qquad
[L]=\mathrm{H}=\frac{\mathrm{Wb}}{\mathrm A}
=\frac{\mathrm{V\,s}}{\mathrm A}.
\]

建立电流时，外电源输入理想电感的功率为

\[
p=vi=i\frac{d\Psi}{dt}.
\]

因而磁场能量为

\[
W_m=\int_0^{\Psi} i(\Psi')\,d\Psi'.
\]

在线性情形 \(\Psi=Li\) 下，

\[
\boxed{W_m=\frac12Li^2=\frac12\Psi i}.
\]

线性介质中的场能还可写成

\[
\boxed{W_m=\frac12\int_V\mathbf B\cdot\mathbf H\,dV}.
\]

若磁化曲线非线性但单值，则微分电感与割线电感不同：

\[
L_{\mathrm d}=\frac{d\Psi}{di},
\qquad
L_{\mathrm s}=\frac{\Psi}{i}.
\]

这时磁共能为 \(W_m'=\int_0^i\Psi(i')\,di'\)；一般
\(W_m\ne W_m'\)，只有线性系统中二者都等于
\(Li^2/2\)。存在磁滞时，\(\Psi\) 不再是 \(i\)
的单值函数，还要计入每周磁滞损耗。

\subsubsection{2.
自感与互感的物理图像}\label{m5:ux81eaux611fux4e0eux4e92ux611fux7684ux7269ux7406ux56feux50cf}

根据楞次定律，自感电动势阻碍的是\textbf{电流的变化}，而不是电流本身。接通电源时，它抵抗电流增大；断开电源时，它又试图维持原有电流。把自感类比为力学中的惯性有助于形成初步直觉，但不能把二者等同：自感取决于完整回路的几何、匝数与磁介质，还可能随部件运动而改变。

在相同电流下，\(L\) 越大，系统储存的磁场能 \(W_m=Li^2/2\)
越多。互感则描述一个回路的电流对另一个回路磁链的影响，是变压器、感应耦合和无线能量传输的基础。互感并不直接等于``传输效率''；实际效率还受线圈电阻、漏磁、频率、负载和辐射损耗等因素影响。

\subsubsection{3.
长直螺线管的自感}\label{m5:ux957fux76f4ux87baux7ebfux7ba1ux7684ux81eaux611f}

设长直螺线管长度为 \(\ell\)、总匝数为 \(N\)、截面积为
\(S\)，内部填充磁导率为 \(\mu\) 的线性均匀介质，并满足 \(\ell\)
远大于横向尺度。忽略边缘漏磁，由安培环路定理，

\[
B\simeq\mu\frac{N}{\ell}i.
\]

单匝磁通为 \(\Phi_1=BS\)，总磁链为

\[
\Psi=N\Phi_1
=\frac{\mu N^2S}{\ell}i.
\]

所以

\[
\boxed{L=\frac{\mu N^2S}{\ell}}
=\mu n^2S\ell,
\qquad n=\frac N\ell.
\]

\par\noindent\begin{minipage}{\linewidth}
\begin{center}
\includegraphics[width=\linewidth,height=.40\textheight,keepaspectratio]{\detokenize{assets/module-05/original/tikz-m5-topic47-v9-from44-v8-02.png}}
\par\smallskip{\small 图示：长直螺线管由磁场逐步得到单匝磁通、总磁链与自感}
\end{center}
\end{minipage}\par

也可由场能检验：内部磁场近似均匀，

\[
W_m=\frac{B^2}{2\mu}S\ell
=\frac12\frac{\mu N^2S}{\ell}i^2
=\frac12Li^2.
\]

\par\addvspace{.4\baselineskip}
\noindent\begin{minipage}{\linewidth}

\textbf{例 M5-084：匝数变化}

理想长螺线管的长度 \(\ell\)、截面 \(S\) 和线性磁导率 \(\mu\)
不变，将总匝数 \(N\)
加倍。比较同一电流下磁场、磁链和自感的变化倍数，明确磁链中第二个匝数因子的来源。

\end{minipage}\par

\par\noindent\begin{minipage}{\linewidth}
\textbf{解答}

\begin{center}
\includegraphics[width=\linewidth,height=.40\textheight,keepaspectratio]{\detokenize{assets/module-05/examples/M5-084.pdf}}
\par\smallskip{\small 图 M5-084：匝数变化}
\end{center}
\end{minipage}\par

图 M5-084
的几何长度与截面保持不变；匝数增加既增强场，又增加链接磁通的匝数，所以自感含有两个
\(N\)。

保持 \(\ell,S,\mu\) 不变，把匝数加倍，则 \(B\)
在同一电流下加倍，磁链还要再乘一次匝数，所以 \(L\) 变为原来的 \(4\)
倍。只从 \(B\propto N\) 出发而漏掉磁链中的另一个 \(N\)，会误得二倍关系。

\subsubsection{4.
同轴电缆的外部自感}\label{m5:ux540cux8f74ux7535ux7f06ux7684ux5916ux90e8ux81eaux611f}

高频趋肤近似下，内导体的电流 \(I\) 流在 \(r=a\)
的表面，大小相反的回流流在外导体内表面 \(r=b\)。安培环路定理给出

\[
B(r)=
\begin{cases}
0,&r<a,\\[2mm]
\dfrac{\mu_0I}{2\pi r},&a<r<b,\\[2mm]
0,&r>b.
\end{cases}
\]

单位长度内的磁场能为

\[
\frac{W_m}{\ell}
=\int_a^b\frac{B^2}{2\mu_0}\,2\pi r\,dr
=\frac{\mu_0I^2}{4\pi}\ln\frac ba.
\]

与 \(W_m/\ell=L'I^2/2\) 比较，得到

\[
\boxed{L'=\frac{\mu_0}{2\pi}\ln\frac ba}.
\]

\par\noindent\begin{minipage}{\linewidth}
\begin{center}
\includegraphics[width=\linewidth,height=.40\textheight,keepaspectratio]{\detokenize{assets/module-05/original/tikz-m5-topic47-v9-from44-v8-03.png}}
\par\smallskip{\small 图示：同轴电缆横截面的表面电流、环向磁场与场能积分}
\end{center}
\end{minipage}\par

这里得到的是两导体之间的\textbf{外部自感}。低频时磁场可以进入导体，导体内部还会贡献内自感；其数值依赖电流在导体截面内的分布，不能继续使用上述表面电流模型。

\par\addvspace{.4\baselineskip}
\noindent\begin{minipage}{\linewidth}

\textbf{例 M5-085：极限检验}

理想同轴线两导体的相邻表面半径为 \(a<b\)，介质磁导率为
\(\mu\)，忽略导体内部磁场。已知单位长度外部自感
\(L^{\prime}=\mu\ln(b/a)/(2\pi)\)。判断 \(b\to a^+\)
时的极限，并从储能区域消失解释它。

\end{minipage}\par

\par\noindent\begin{minipage}{\linewidth}
\textbf{解答}

\begin{center}
\includegraphics[width=\linewidth,height=.40\textheight,keepaspectratio]{\detokenize{assets/module-05/examples/M5-085.pdf}}
\par\smallskip{\small 图 M5-085：极限检验}
\end{center}
\end{minipage}\par

图 M5-085 的场能储存在 \(a<r<b\)
的间隙中；将间隙压到零时积分区间消失，外部自感也应趋于零。

当 \(b\to a^+\) 时，磁场所占空间趋于零，\(\ln(b/a)\to0\)，故
\(L'\to0\)；若计算结果在此极限不趋零，通常是积分边界或回流位置处理错误。

\subsection{三、互感、方向与互易性}\label{m5:ux4e09ux4e92ux611fux65b9ux5411ux4e0eux4e92ux6613ux6027}

\subsubsection{1.
互感的大小不是拓扑量}\label{m5:ux4e92ux611fux7684ux5927ux5c0fux4e0dux662fux62d3ux6251ux91cf}

互感表示一个回路电流对另一回路磁链的影响。即使两个回路完全没有``互锁''，互感也可以非零；即使拓扑关系不变，平移、转动或改变间距也会连续改变互感。真空中的细导线回路，互感由几何、尺度和
\(\mu_0\)
决定；存在磁性介质时还依赖磁导率分布。互感的符号则由两个回路的正方向约定决定。

因此，较准确的表述是：

\[
\boxed{M\text{ 是依赖几何、介质和方向约定的有符号耦合参数。}}
\]

\subsubsection{2.
长直导线与共面矩形线框}\label{m5:ux957fux76f4ux5bfcux7ebfux4e0eux5171ux9762ux77e9ux5f62ux7ebfux6846}

设长直导线电流为 \(I\)，矩形线框平行于导线的边长为
\(\ell\)，垂直于导线的宽度为 \(a\)，最近边距导线为 \(d\)。距导线 \(x\)
处，

\[
B(x)=\frac{\mu_0I}{2\pi x}.
\]

取面积微元 \(dS=\ell\,dx\)，则

\[
\Phi=\int_d^{d+a}B(x)\ell\,dx
=\frac{\mu_0I\ell}{2\pi}\ln\left(1+\frac ad\right).
\]

若线框正法向与图示磁场同向，则

\[
\boxed{M=\frac{\mu_0\ell}{2\pi}\ln\left(1+\frac ad\right)}.
\]

反转线框正方向时，\(M\) 整体变号。

\par\noindent\begin{minipage}{\linewidth}
\begin{center}
\includegraphics[width=\linewidth,height=.40\textheight,keepaspectratio]{\detokenize{assets/module-05/original/tikz-m5-topic47-v9-from44-v8-04.png}}
\par\smallskip{\small 图示：长直导线的非均匀磁场按距离分层穿过矩形线框}
\end{center}
\end{minipage}\par

\par\addvspace{.4\baselineskip}
\noindent\begin{minipage}{\linewidth}

\textbf{例 M5-086：远距离近似}

真空中长直导线与共面矩形线框的互感大小为
\(M=\mu_0\ell\ln(1+a/d)/(2\pi)\)，其中 \(d\) 是最近边距离，\(a\)
是垂直导线的宽度，\(\ell\) 是平行边长。若
\(a\ll d\)，求主阶近似，并用线框中近似均匀的磁场检验。

\end{minipage}\par

\par\noindent\begin{minipage}{\linewidth}
\textbf{解答}

\begin{center}
\includegraphics[width=\linewidth,height=.40\textheight,keepaspectratio]{\detokenize{assets/module-05/examples/M5-086.pdf}}
\par\smallskip{\small 图 M5-086：远距离近似}
\end{center}
\end{minipage}\par

图 M5-086 的直导线到观测区最近距离在本题记为 \(d\)；当线框宽度
\(a\ll d\) 时，其范围内磁场近似不变，积分退化为面积近似。

当 \(a\ll d\) 时，\(\ln(1+a/d)\simeq a/d\)，于是

\[
M\simeq\frac{\mu_0a\ell}{2\pi d}.
\]

这正等于把线框内磁场近似取为最近尺度上的均匀值
\(B\simeq\mu_0I/(2\pi d)\) 后所得的 \(M\simeq B(a\ell)/I\)。

\subsubsection{3.
诺伊曼公式证明互感互易性}\label{m5:ux8bfaux4f0aux66fcux516cux5f0fux8bc1ux660eux4e92ux611fux4e92ux6613ux6027}

先考虑真空中两个相互分离的有向闭合细导线回路 \(C_1,C_2\)。回路 1 在回路
2 上产生的磁矢势为

\[
\mathbf A_1(\mathbf r_2)
=\frac{\mu_0I_1}{4\pi}
\oint_{C_1}\frac{d\mathbf l_1}{r_{12}},
\qquad
r_{12}=|\mathbf r_2-\mathbf r_1|.
\]

由斯托克斯定理，回路 2 接收的磁通为

\[
\Phi_{21}
=\int_{S_2}\mathbf B_1\cdot d\mathbf S_2
=\oint_{C_2}\mathbf A_1\cdot d\mathbf l_2.
\]

代入磁矢势，得到诺伊曼公式

\[
\boxed{
M_{21}=\frac{\mu_0}{4\pi}
\oint_{C_1}\oint_{C_2}
\frac{d\mathbf l_1\cdot d\mathbf l_2}{r_{12}}
}.
\]

交换 \(1,2\) 后，利用

\[
d\mathbf l_1\cdot d\mathbf l_2
=d\mathbf l_2\cdot d\mathbf l_1,
\qquad r_{12}=r_{21},
\]

便有

\[
\boxed{M_{12}=M_{21}}.
\]

多匝线圈只需对各匝成对求和，结论不变。

\par\noindent\begin{minipage}{\linewidth}
\begin{center}
\includegraphics[width=\linewidth,height=.40\textheight,keepaspectratio]{\detokenize{assets/module-05/theory/tikz-m5-topic47-v9-from44-v8-05.pdf}}
\par\smallskip{\small 图示：诺伊曼积分与互感互易性}
\end{center}
\end{minipage}\par

诺伊曼公式直接适用于真空中的定常闭合细导线。在线性、静止、互易介质中，更一般的电磁互易定理仍保证
\(M_{12}=M_{21}\)。若介质被外加偏置而呈非互易性，或系统含时变、主动元件，则不能未经检验就使用互感对称性。

\subsubsection{4.
能量路径证明互感互易性}\label{m5:ux80fdux91cfux8defux5f84ux8bc1ux660eux4e92ux611fux4e92ux6613ux6027}

还可从磁场能的状态函数性质得到同一结论。由
\(p=i_1\dot\Psi_1+i_2\dot\Psi_2\)，从 \((0,0)\) 建立到 \((I_1,I_2)\)。

先建立 \(i_1\)，再保持 \(i_1=I_1\) 建立 \(i_2\)，有

\[
W_A=\frac12L_1I_1^2+\frac12L_2I_2^2+M_{12}I_1I_2.
\]

交换建流顺序，有

\[
W_B=\frac12L_1I_1^2+\frac12L_2I_2^2+M_{21}I_1I_2.
\]

若最终磁场状态相同且能量与路径无关，则 \(W_A=W_B\) 对任意 \(I_1,I_2\)
成立，故

\[
\boxed{M_{12}=M_{21}}.
\]

\par\noindent\begin{minipage}{\linewidth}
\begin{center}
\includegraphics[width=\linewidth,height=.40\textheight,keepaspectratio]{\detokenize{assets/module-05/theory/tikz-m5-topic47-v9-from44-v8-06.pdf}}
\par\smallskip{\small 图示：互易线圈的磁能与建立路径}
\end{center}
\end{minipage}\par

这一路径证明也清楚说明了适用条件：如果系统存在磁滞或非互易过程，磁场建立过程可能与路径有关，上述证明的前提便不成立。

\subsection{四、点标法与耦合电路方程}\label{m5:ux56dbux70b9ux6807ux6cd5ux4e0eux8026ux5408ux7535ux8defux65b9ux7a0b}

\subsubsection{1.
点端的物理含义}\label{m5:ux70b9ux7aefux7684ux7269ux7406ux542bux4e49}

点标法只记录两个线圈绕向之间的相对极性，不表示点端在空间中一定彼此靠近。采用以下统一参考约定：

\begin{itemize}
\tightlist
\item
  两个参考电流都流入点端，两个端电压都由点端量到非点端，则互感项取
  \(+|M|\)；
\item
  一个参考电流流入点端，另一个流出点端，则互感项取 \(-|M|\)。
\end{itemize}

\par\noindent\begin{minipage}{\linewidth}
\begin{center}
\includegraphics[width=\linewidth,height=.40\textheight,keepaspectratio]{\detokenize{assets/module-05/theory/tikz-m5-topic47-v9-from44-v8-07.pdf}}
\par\smallskip{\small 图示：点标法与端口参考方向}
\end{center}
\end{minipage}\par

于是，在两电流都流入点端、两电压都以点端为正端时，

\[
\boxed{
\begin{aligned}
v_1&=L_1\dot i_1+|M|\dot i_2,\\
v_2&=|M|\dot i_1+L_2\dot i_2.
\end{aligned}}
\]

若只翻转一个电流参考方向或一个电压参考方向，必须重新按参考方向写方程，不能只凭图形感觉改一个符号。

\par\addvspace{.4\baselineskip}
\noindent\begin{minipage}{\linewidth}

\textbf{例 M5-087：开路次级的感应电压}

两线性互易线圈互感大小 \(M=0.30\ \mathrm H\)。原边电流按流入点端方向以
\(2.0\ \mathrm{A/s}\)
增加，次级开路。次级电压从点端量到非点端，求感应电压的大小和符号，忽略寄生电容。

\end{minipage}\par

\par\noindent\begin{minipage}{\linewidth}
\textbf{解答}

\begin{center}
\includegraphics[width=\linewidth,height=.40\textheight,keepaspectratio]{\detokenize{assets/module-05/examples/M5-087.pdf}}
\par\smallskip{\small 图 M5-087：开路次级的感应电压}
\end{center}
\end{minipage}\par

图 M5-087
的电流参考流入原边点端，次级电压从点端量到非点端。按这一共同点标约定，互感电压为正。

\(M=0.30\ \mathrm H\)，原边电流按流入点端方向以
\(di_1/dt=2.0\ \mathrm{A/s}\) 增加，副边开路且电压由点端量到非点端，则

\[
v_2=M\frac{di_1}{dt}=0.60\ \mathrm V.
\]

若把副边电压参考方向反转，数值不变而代数符号变为 \(-0.60\ \mathrm V\)。

\subsubsection{2. 电感矩阵}\label{m5:ux7535ux611fux77e9ux9635}

在互易系统中，把互感符号并入 \(M\)，可统一写成

\[
\boxed{
\begin{pmatrix}\Psi_1\\\Psi_2\end{pmatrix}
=
\underbrace{
\begin{pmatrix}L_1&M\\M&L_2\end{pmatrix}}_{\boldsymbol L}
\begin{pmatrix}i_1\\i_2\end{pmatrix}}
\]

和

\[
\boxed{
\begin{pmatrix}v_1\\v_2\end{pmatrix}
=\boldsymbol L
\begin{pmatrix}\dot i_1\\\dot i_2\end{pmatrix}}
\]

其中第二式要求 \(L_1,L_2,M\)
不随时间、电流和位置改变。若线圈电阻不可忽略，则端口方程应写成

\[
u_1=R_1i_1+v_1,
\qquad
u_2=R_2i_2+v_2.
\]

\par\noindent\begin{minipage}{\linewidth}
\begin{center}
\includegraphics[width=\linewidth,height=.40\textheight,keepaspectratio]{\detokenize{assets/module-05/original/tikz-m5-topic47-v9-from44-v8-08.png}}
\par\smallskip{\small 图示：磁链、端电压与电流变化率的电感矩阵写法}
\end{center}
\end{minipage}\par

当 \(L_1L_2-M^2>0\) 时，

\[
\begin{pmatrix}\dot i_1\\\dot i_2\end{pmatrix}
=\frac1{L_1L_2-M^2}
\begin{pmatrix}L_2&-M\\-M&L_1\end{pmatrix}
\begin{pmatrix}v_1\\v_2\end{pmatrix}.
\]

这个逆矩阵形式常用于列耦合电路的状态方程。

\subsection{五、磁场能、耦合系数与互感上限}\label{m5:ux4e94ux78c1ux573aux80fdux8026ux5408ux7cfbux6570ux4e0eux4e92ux611fux4e0aux9650}

\subsubsection{1.
两耦合线圈的磁场能}\label{m5:ux4e24ux8026ux5408ux7ebfux5708ux7684ux78c1ux573aux80fd}

互易、线性系统的磁场能为

\[
\boxed{
W_m=\frac12L_1i_1^2+Mi_1i_2+\frac12L_2i_2^2
=\frac12\mathbf i^{\mathsf T}\boldsymbol L\mathbf i}.
\]

互感项的正负取决于参考方向；不能在已经使用有符号 \(M\)
后再任意添加一个正负号。

把能量配方：

\[
W_m=
\frac12L_1\left(i_1+\frac M{L_1}i_2\right)^2
+\frac12\left(L_2-\frac{M^2}{L_1}\right)i_2^2.
\]

无源系统对任意 \(i_1,i_2\) 都必须有 \(W_m\ge0\)，因此

\[
\boxed{M^2\le L_1L_2}.
\]

定义非负耦合系数

\[
\boxed{k=\frac{|M|}{\sqrt{L_1L_2}}},
\qquad 0\le k\le1.
\]

\par\noindent\begin{minipage}{\linewidth}
\begin{center}
\includegraphics[width=\linewidth,height=.40\textheight,keepaspectratio]{\detokenize{assets/module-05/original/tikz-m5-topic47-v9-from44-v8-09.png}}
\par\smallskip{\small 图示：电感矩阵的能量非负性给出互感上限和耦合系数范围}
\end{center}
\end{minipage}\par

\(k\)
只表示耦合强弱，不携带绕向符号；互感的正负仍由方向或点标约定给出。\(k=1\)
表示理想全耦合，电感矩阵行列式为零，两个磁链受到额外约束，并非两个完全独立的储能自由度。实际装置因漏磁通常有
\(k<1\)。

\par\addvspace{.4\baselineskip}
\noindent\begin{minipage}{\linewidth}

\textbf{例 M5-088：检查互感是否可能}

两个静止线性无源互易线圈给定
\(L_1=4\ \mathrm H\text{、}L_2=9\ \mathrm H\text{、}M=3\ \mathrm H\)。检验这些参数能否对应非负磁场能，并求耦合系数，不要求两个自感相同。

\end{minipage}\par

\par\noindent\begin{minipage}{\linewidth}
\textbf{解答}

\begin{center}
\includegraphics[width=\linewidth,height=.40\textheight,keepaspectratio]{\detokenize{assets/module-05/examples/M5-088.pdf}}
\par\smallskip{\small 图 M5-088：检查互感是否可能}
\end{center}
\end{minipage}\par

图 M5-088
的两个绕组储能由电感矩阵描述；用行列式和耦合系数检验给定互感，不应只要求
\(M<L_1\) 或 \(M<L_2\)。

若 \(L_1=4\ \mathrm H\)、\(L_2=9\ \mathrm H\)、\(M=3\ \mathrm H\)，则

\[
k=\frac3{\sqrt{4\times9}}=\frac12.
\]

此系统允许的最大互感绝对值为 \(6\ \mathrm H\)；若题目给出
\(|M|>6\ \mathrm H\)，便违反无源系统的能量非负性。

\subsection{六、耦合电感的连接与典型边界条件}\label{m5:ux516dux8026ux5408ux7535ux611fux7684ux8fdeux63a5ux4e0eux5178ux578bux8fb9ux754cux6761ux4ef6}

\subsubsection{1.
串联助磁与串联反磁}\label{m5:ux4e32ux8054ux52a9ux78c1ux4e0eux4e32ux8054ux53cdux78c1}

串联时同一电流通过两个线圈。若该电流同时流入两个点端，互感磁链相助，

\[
\boxed{L_{\text{助}}=L_1+L_2+2|M|}.
\]

若同一电流流入一个点端、流出另一个点端，互感磁链相消，

\[
\boxed{L_{\text{反}}=L_1+L_2-2|M|}.
\]

\par\noindent\begin{minipage}{\linewidth}
\begin{center}
\includegraphics[width=\linewidth,height=.40\textheight,keepaspectratio]{\detokenize{assets/module-05/theory/tikz-m5-topic47-v9-from44-v8-10.pdf}}
\par\smallskip{\small 图示：串联线圈的助磁与反磁}
\end{center}
\end{minipage}\par

两式既可由端电压相加得到，也可把 \(i_1=\pm i_2=i\) 代入磁场能，再与
\(W_m=L_{\mathrm{eq}}i^2/2\) 比较得到。

\par\addvspace{.4\baselineskip}
\noindent\begin{minipage}{\linewidth}

\textbf{例 M5-089：串联助磁与反磁的数值检验}

两耦合线圈的
\(L_1=4\ \mathrm H\text{、}L_2=9\ \mathrm H\text{、}|M|=3\ \mathrm H\)。分别让同一个串联电流对两线圈助磁和反磁，求两个端口等效电感，说明互感项为何各出现两次。

\end{minipage}\par

\par\noindent\begin{minipage}{\linewidth}
\textbf{解答}

\begin{center}
\includegraphics[width=\linewidth,height=.40\textheight,keepaspectratio]{\detokenize{assets/module-05/examples/M5-089.pdf}}
\par\smallskip{\small 图 M5-089：串联助磁与反磁的数值检验}
\end{center}
\end{minipage}\par

图 M5-089
先给出两个线圈各自的点端；串联时比较同一串联电流是否同时流入点端，助磁取正互感项，反磁取负互感项。

把两个绕组的电压正向都取为沿串联电流的压降。每个绕组的磁链均有一个互感项，所以
\[
v_1=(L_1\pm M)\frac{di}{dt},\qquad
v_2=(L_2\pm M)\frac{di}{dt}.
\] 总压降为两式相加，得到 \(v=(L_1+L_2\pm2M)\,di/dt\)。\(2M\)
来自两个绕组各自的互感电压，并非重复计算同一绕组。

\(L_1=4\ \mathrm H\)、\(L_2=9\ \mathrm H\)、\(|M|=3\ \mathrm H\) 时，

\[
L_{\text{助}}=19\ \mathrm H,
\qquad
L_{\text{反}}=7\ \mathrm H.
\]

\subsubsection{2.
并联耦合电感}\label{m5:ux5e76ux8054ux8026ux5408ux7535ux611f}

若两线圈点端相接、非点端相接，支路电流都取流入点端，且公共端电压为
\(v\)，则

\[
v=L_1\dot i_1+M\dot i_2
=M\dot i_1+L_2\dot i_2,
\qquad i=i_1+i_2.
\]

当 \(k<1\) 时反解可得

\[
\boxed{
L_{\parallel,+}
=\frac{L_1L_2-M^2}{L_1+L_2-2M}}
\qquad(M=|M|).
\]

若翻转一个线圈再并联，相当于把 \(M\) 换成 \(-|M|\)，于是

\[
\boxed{
L_{\parallel,-}
=\frac{L_1L_2-M^2}{L_1+L_2+2|M|}}.
\]

\(k=1\)
时电感矩阵奇异，上式所用的逆矩阵推导失效；应回到磁链约束，并检验并联连接、外加端电压和初始条件是否相容，不能把上式的形式极限当作一般并联结论。

\subsubsection{3.
次级开路与理想短路}\label{m5:ux6b21ux7ea7ux5f00ux8defux4e0eux7406ux60f3ux77edux8def}

次级开路时 \(i_2=0\)，所以

\[
v_1=L_1\dot i_1,
\qquad
v_2=M\dot i_1.
\]

次级为理想短路且忽略电阻时，\(v_2=0\)。若初始电流为零，则

\[
L_2\dot i_2+M\dot i_1=0
\quad\Longrightarrow\quad
i_2=-\frac M{L_2}i_1.
\]

代回原边方程，

\[
\boxed{
v_1=\left(L_1-\frac{M^2}{L_2}\right)\dot i_1
=L_1(1-k^2)\dot i_1}.
\]

次级电流产生的磁场抵消原磁通变化，因此原边看到的等效电感下降。实际短路线圈含电阻，响应还会有衰减时间尺度，不能始终满足上面的理想关系。

\subsection{七、使用条件与常见错误}\label{m5:ux4e03ux4f7fux7528ux6761ux4ef6ux4e0eux5e38ux89c1ux9519ux8bef}

\begin{enumerate}
\def\labelenumi{\arabic{enumi}.}
\tightlist
\item
  \textbf{磁通与磁链混用。}
  多匝线圈应先判断各匝磁通是否相同，再决定能否使用 \(\Psi=N\Phi\)。
\item
  \textbf{把感应电动势直接当作同向端电压。}
  \(\mathcal E=-d\Psi/dt\)，而被动符号约定下 \(v=d\Psi/dt\)。
\item
  \textbf{忘记全导数。} 只有几何和介质固定、系统线性时，才有
  \(v=L\,di/dt\)；一般应写 \(v=d\Psi/dt\)。
\item
  \textbf{把互感当作非负量。} \(|M|\) 描述大小，\(M\)
  的符号由回路方向决定；\(k\) 才按定义取非负值。
\item
  \textbf{把互感说成拓扑量。}
  互感会随间距、姿态和磁介质连续改变，不由``是否互锁''单独决定。
\item
  \textbf{无条件使用 \(M_{12}=M_{21}\)。}
  竞赛常见的静止线性互易系统满足该式；磁滞、非互易介质或主动时变系统须另行分析。
\item
  \textbf{忽略内自感。}
  同轴电缆的高频表面电流模型只给外部自感；低频电流进入导体后还存在内自感。
\item
  \textbf{直接计算理想细丝的自感。}
  诺伊曼互感积分对相互分离的两回路有限，但把两个回路重合来求零半径细丝的自感会出现短程发散；真实导线必须保留有限横截面或给出内部场模型。
\item
  \textbf{串、并联时只看线圈画法。}
  必须先标点端和参考电流，再判断互感项正负。
\item
  \textbf{越过准静态范围。}
  当结构尺度不能忽略电磁传播时间时，电流与磁链不再是瞬时集中参数关系，应改用分布参数和电磁波方法。
\end{enumerate}

\subsection{八、练习与解答}\label{m5:ux516bux7ec3ux4e60ux4e0eux89e3ux7b54-1}

\par\addvspace{.4\baselineskip}
\noindent\begin{minipage}{\linewidth}

\subsubsection{作业一：同轴电缆的单位长度自感}\label{m5:ux4f5cux4e1aux4e00ux540cux8f74ux7535ux7f06ux7684ux5355ux4f4dux957fux5ea6ux81eaux611f}

\textbf{题号：M5-090}

\textbf{层级：} 复赛　　\textbf{方法：} 自感互感

导体间是真空；采用准TEM的表面电流近似，忽略两导体材料内的磁场及端部效应。

一根无限长同轴电缆的内导体半径为 \(a\)，外导体内半径为 \(b\)。总电流
\(I\)
沿内导体流去，并沿外导体返回。假设高频趋肤效应使电流只分布在内导体外表面和外导体内表面。求单位长度自感
\(L'\)。

\end{minipage}\par

\par\noindent\begin{minipage}{\linewidth}
\textbf{解答}

\begin{center}
\includegraphics[width=\linewidth,height=.40\textheight,keepaspectratio]{\detokenize{assets/module-05/examples/M5-090.pdf}}
\par\smallskip{\small 图 M5-090：同轴电缆的单位长度自感}
\end{center}
\end{minipage}\par

图 M5-090 中高频表面电流使磁场只留在两导体之间；单位长度场能对环形面积元
\(2\pi r\,dr\) 积分，再与 \(L^{\prime} I^2/2\) 比较。

对半径为 \(r\) 的同心圆使用安培环路定理。\(r<a\) 时未包围电流；\(a<r<b\)
时包围电流为 \(I\)；\(r>b\) 时去、回电流代数和为零。因此

\[
B(r)=
\begin{cases}
0,&r<a,\\
\dfrac{\mu_0I}{2\pi r},&a<r<b,\\
0,&r>b.
\end{cases}
\]

取长度为 \(\ell\) 的一段，两导体之间径向截面的磁通为

\[
\Phi=\int_a^bB(r)\ell\,dr
=\frac{\mu_0I\ell}{2\pi}\ln\frac ba.
\]

在该单回路模型中 \(\Psi=\Phi\)，故

\[
L=\frac\Psi I
=\frac{\mu_0\ell}{2\pi}\ln\frac ba,
\]

从而

\[
\boxed{L'=\frac L\ell
=\frac{\mu_0}{2\pi}\ln\frac ba}.
\]

由场能法计算也得到相同结果。该答案不包含导体内部自感。

\par\addvspace{.4\baselineskip}
\noindent\begin{minipage}{\linewidth}

\subsubsection{作业二：长直导线与矩形线框的互感}\label{m5:ux4f5cux4e1aux4e8cux957fux76f4ux5bfcux7ebfux4e0eux77e9ux5f62ux7ebfux6846ux7684ux4e92ux611f}

\textbf{题号：M5-091}

\textbf{层级：} 复赛　　\textbf{方法：} 毕奥---萨伐尔、自感互感

空间为真空，\(d\text{、}a\text{、}\ell>0\)，线框为单匝；只求互感大小，方向由另选回路正向确定。

一根很长直导线与一个共面矩形线框相邻。矩形线框垂直于导线的边长为
\(a\)，平行于导线的边长为 \(\ell\)，最近边距导线为
\(d\)。设直导线所属回路的回流支路足够远，其在矩形区域内的磁场可忽略。求互感的大小。

\end{minipage}\par

\par\noindent\begin{minipage}{\linewidth}
\textbf{解答}

\begin{center}
\includegraphics[width=\linewidth,height=.40\textheight,keepaspectratio]{\detokenize{assets/module-05/examples/M5-091.pdf}}
\par\smallskip{\small 图 M5-091：长直导线与矩形线框的互感}
\end{center}
\end{minipage}\par

图 M5-091 的长导线在本题与矩形线框共面；将场点距离由 \(d\) 扫到
\(d+a\)，每一窄条的面积为 \(\ell\,dr\)。远处回流支路的贡献按题设忽略。

令长直导线通电流 \(I\)。距离导线 \(r\) 处，

\[
B(r)=\frac{\mu_0I}{2\pi r}.
\]

在矩形线框内取长为 \(\ell\)、宽为 \(dr\) 的面积条，\(dS=\ell\,dr\)。于是

\[
\Phi
=\int_d^{d+a}\frac{\mu_0I}{2\pi r}\ell\,dr
=\frac{\mu_0I\ell}{2\pi}
\ln\frac{d+a}{d}.
\]

所以互感大小为

\[
\boxed{|M|=\frac{\mu_0\ell}{2\pi}
\ln\left(1+\frac ad\right)}.
\]

若题目规定线框正法向，则再根据长直导线磁场与正法向的关系决定 \(M\)
的正负。

\par\addvspace{.4\baselineskip}
\noindent\begin{minipage}{\linewidth}

\subsubsection{作业三：长直螺线管的自感与场能检验}\label{m5:ux4f5cux4e1aux4e09ux957fux76f4ux87baux7ebfux7ba1ux7684ux81eaux611fux4e0eux573aux80fdux68c0ux9a8c}

\textbf{题号：M5-092}

\textbf{层级：} 复赛　　\textbf{方法：} 自感互感

介质线性、各向同性、无磁滞，\(\mu\) 为常量，绕组紧密均匀且 \(\ell\)
远大于截面横向尺度。

长度为 \(\ell\)、总匝数为 \(N\)、截面积为 \(S\)
的长直螺线管内填充磁导率为 \(\mu\)
的线性介质。忽略边缘效应，求自感，并分别由电路公式与场能密度计算电流为
\(I\) 时的磁场能。

\end{minipage}\par

\par\noindent\begin{minipage}{\linewidth}
\textbf{解答}

\begin{center}
\includegraphics[width=\linewidth,height=.40\textheight,keepaspectratio]{\detokenize{assets/module-05/examples/M5-092.pdf}}
\par\smallskip{\small 图 M5-092：长直螺线管的自感与场能检验}
\end{center}
\end{minipage}\par

图 M5-092 的长螺线管内部取均匀场，储能体积为
\(S\ell\)；磁链法与场能法应得到完全相同的 \(L\)。

内部磁场为

\[
B=\mu\frac N\ell I.
\]

磁链为

\[
\Psi=NBS=\frac{\mu N^2S}{\ell}I,
\]

故

\[
\boxed{L=\frac{\mu N^2S}{\ell}}.
\]

电路公式给出

\[
W_m=\frac12LI^2
=\frac{\mu N^2S}{2\ell}I^2.
\]

另一方面，场能密度为 \(u_m=B^2/(2\mu)\)，有效体积为 \(S\ell\)，所以

\[
W_m=\frac{B^2}{2\mu}S\ell
=\frac{\mu N^2S}{2\ell}I^2.
\]

两种方法完全一致。

\par\addvspace{.4\baselineskip}
\noindent\begin{minipage}{\linewidth}

\subsubsection{作业四：点标法与串联等效电感}\label{m5:ux4f5cux4e1aux56dbux70b9ux6807ux6cd5ux4e0eux4e32ux8054ux7b49ux6548ux7535ux611f}

\textbf{题号：M5-093}

\textbf{层级：} 复赛　　\textbf{方法：} 等效变换、自感互感

各线圈电阻与寄生电容忽略；助磁是指同一个串联电流在两绕组中的磁通贡献相加，反磁则相减。

两个线性耦合线圈满足
\(L_1=4\ \mathrm H\)、\(L_2=9\ \mathrm H\)、\(|M|=3\ \mathrm H\)。求串联助磁和串联反磁时的等效电感，并求耦合系数。

\end{minipage}\par

\par\noindent\begin{minipage}{\linewidth}
\textbf{解答}

\begin{center}
\includegraphics[width=\linewidth,height=.40\textheight,keepaspectratio]{\detokenize{assets/module-05/examples/M5-093.pdf}}
\par\smallskip{\small 图 M5-093：点标法与串联等效电感}
\end{center}
\end{minipage}\par

图 M5-093
的橙点规定互感符号，实际串联导线及电流关系见后续过程图。两个互感压降在助磁时相加、反磁时相减。

串联助磁时，同一电流同时流入两点端：

\[
L_{\text{助}}=L_1+L_2+2|M|
=4+9+6=19\ \mathrm H.
\]

\par\noindent\begin{minipage}{\linewidth}
\begin{center}
\includegraphics[width=\linewidth,height=.40\textheight,keepaspectratio]{\detokenize{assets/module-05/process/M5-093a.pdf}}
\par\smallskip{\small 图 M5-093a：点端连接决定串联互感项的符号}
\end{center}
\end{minipage}\par

图 M5-093a
的橙点都位于各绕组的左端，同一向右电流同时流入两点端，所以两个互感项同号相加。

串联反磁时，同一电流一入点端、一出点端：

\[
L_{\text{反}}=L_1+L_2-2|M|
=4+9-6=7\ \mathrm H.
\]

耦合系数为

\[
\boxed{k=\frac{|M|}{\sqrt{L_1L_2}}
=\frac3{\sqrt{4\times9}}=\frac12}.
\]

\par\addvspace{.4\baselineskip}
\noindent\begin{minipage}{\linewidth}

\subsubsection{作业五：由能量证明互感上限}\label{m5:ux4f5cux4e1aux4e94ux7531ux80fdux91cfux8bc1ux660eux4e92ux611fux4e0aux9650}

\textbf{题号：M5-094}

\textbf{层级：} 复赛　　\textbf{方法：} 能量法、自感互感

\(L_1\text{、}L_2>0\)，电流变量可独立取任意实值，\(M\)
按共同参考方向为有符号互感。

两个静止线性互易线圈的磁场能为

\[
W_m=\frac12L_1i_1^2+Mi_1i_2+\frac12L_2i_2^2.
\]

只利用无源系统磁场能非负，证明
\(M^2\le L_1L_2\)。当等号成立时，求使理想模型磁场能为零的一组非零电流比。

\end{minipage}\par

\par\noindent\begin{minipage}{\linewidth}
\textbf{解答}

\begin{center}
\includegraphics[width=\linewidth,height=.40\textheight,keepaspectratio]{\detokenize{assets/module-05/examples/M5-094.pdf}}
\par\smallskip{\small 图 M5-094：由能量证明互感上限}
\end{center}
\end{minipage}\par

图 M5-094 用流入点端的两个电流作二次型变量；配方后沿 \(i_1=-Mi_2/L_1\)
的方向检查剩余能量，得到互感上限。

配方得

\[
W_m=
\frac12L_1\left(i_1+\frac M{L_1}i_2\right)^2
+\frac12\left(L_2-\frac{M^2}{L_1}\right)i_2^2.
\]

第一项非负。为了使任意 \(i_2\) 下总能量都非负，必须有

\[
L_2-\frac{M^2}{L_1}\ge0.
\]

因此

\[
\boxed{M^2\le L_1L_2}.
\]

等号成立时，第二项恒为零；再令第一项为零，得到

\[
\boxed{\frac{i_1}{i_2}=-\frac M{L_1}
=-\frac{L_2}{M}}.
\]

这说明理想全耦合模型中存在不建立净磁链的电流组合，电感矩阵因而奇异。真实线圈的漏磁会使该组合仍具有少量储能。

\par\addvspace{.4\baselineskip}
\noindent\begin{minipage}{\linewidth}

\subsubsection{作业六：次级开路与短路的反作用}\label{m5:ux4f5cux4e1aux516dux6b21ux7ea7ux5f00ux8defux4e0eux77edux8defux7684ux53cdux4f5cux7528}

\textbf{题号：M5-095}

\textbf{层级：} 复赛　　\textbf{方法：} 自感互感

所有电压从点端量到非点端，初始时刻为 \(t=0\)；\(\alpha\)
为给定常量，副边短路连接不含电阻与外加电源。

两个耦合线圈满足
\(L_1=4\ \mathrm H\)、\(L_2=9\ \mathrm H\)、\(M=3\ \mathrm H\)，参考电流都流入点端。原边电流以恒定变化率
\(di_1/dt=\alpha\) 增加。

\begin{enumerate}
\def\labelenumi{\arabic{enumi}.}
\tightlist
\item
  次级开路时求 \(v_2\) 和 \(v_1\)。
\item
  次级理想短路、初始电流均为零且忽略电阻时，求 \(i_2\) 与 \(i_1\)
  的关系，并求原边端电压。
\end{enumerate}

\end{minipage}\par

\par\noindent\begin{minipage}{\linewidth}
\textbf{解答}

\begin{center}
\includegraphics[width=\linewidth,height=.40\textheight,keepaspectratio]{\detokenize{assets/module-05/examples/M5-095.pdf}}
\par\smallskip{\small 图 M5-095：次级开路与短路的反作用}
\end{center}
\end{minipage}\par

图 M5-095 画出次级短路状态，仍把 \(i_2\)
参考方向取为流入点端。由初始零磁链得 \(i_2\)
为负，它的磁场反抗原边增长。

次级开路时 \(i_2=0\)，所以

\[
\boxed{v_2=M\alpha=3\alpha},
\qquad
\boxed{v_1=L_1\alpha=4\alpha}.
\]

次级理想短路时 \(v_2=0\)：

\[
M\dot i_1+L_2\dot i_2=0.
\]

结合零初始条件积分，

\[
\boxed{i_2=-\frac M{L_2}i_1=-\frac13i_1}.
\]

\par\noindent\begin{minipage}{\linewidth}
\begin{center}
\includegraphics[width=\linewidth,height=.40\textheight,keepaspectratio]{\detokenize{assets/module-05/process/M5-095a.pdf}}
\par\smallskip{\small 图 M5-095a：短路次级电流抵抗原边磁链增长}
\end{center}
\end{minipage}\par

图 M5-095a
的负斜率表示参考为流入点端的次级电流实际为负；这与原边电流增加时的楞次反作用一致。

原边电压为

\[
v_1=L_1\dot i_1+M\dot i_2
=\left(L_1-\frac{M^2}{L_2}\right)\alpha,
\]

故

\[
\boxed{v_1=3\alpha}.
\]

开路时原边等效电感为 \(4\ \mathrm H\)，理想短路时降为 \(3\ \mathrm H\)。

\par\addvspace{.4\baselineskip}
\noindent\begin{minipage}{\linewidth}

\subsubsection{作业七：耦合电感的同极性并联}\label{m5:ux4f5cux4e1aux4e03ux8026ux5408ux7535ux611fux7684ux540cux6781ux6027ux5e76ux8054}

\textbf{题号：M5-096}

\textbf{层级：} 复赛　　\textbf{方法：} 自感互感

采用线性无源模型并忽略电阻，各绕组初始电流为零；只求端口增量等效，不另叠加持久环流。

两个耦合线圈的点端相接、非点端相接，\(L_1,L_2\) 与 \(M>0\) 均为常量，且
\(M^2<L_1L_2\)。两支路电流均取流入点端，求端口等效电感。

\end{minipage}\par

\par\noindent\begin{minipage}{\linewidth}
\textbf{解答}

\begin{center}
\includegraphics[width=\linewidth,height=.40\textheight,keepaspectratio]{\detokenize{assets/module-05/examples/M5-096.pdf}}
\par\smallskip{\small 图 M5-096：耦合电感的同极性并联}
\end{center}
\end{minipage}\par

图 M5-096
中上轨连接两个点端、下轨连接非点端，所以两绕组电压相同。先联立各自的互感方程，再把支路电流相加求端口电感。

两支路端电压相同：

\[
\begin{pmatrix}v\\v\end{pmatrix}
=
\begin{pmatrix}L_1&M\\M&L_2\end{pmatrix}
\begin{pmatrix}\dot i_1\\\dot i_2\end{pmatrix}.
\]

反解得

\[
\begin{pmatrix}\dot i_1\\\dot i_2\end{pmatrix}
=\frac v{L_1L_2-M^2}
\begin{pmatrix}L_2-M\\L_1-M\end{pmatrix}.
\]

总电流 \(i=i_1+i_2\)，故

\[
\dot i
=v\frac{L_1+L_2-2M}{L_1L_2-M^2}.
\]

由 \(v=L_{\mathrm{eq}}\dot i\)，得到

\[
\boxed{L_{\mathrm{eq}}
=\frac{L_1L_2-M^2}{L_1+L_2-2M}}.
\]

若翻转一个线圈再并联，只需把式中的 \(M\) 换为 \(-M\)。

\par\addvspace{.4\baselineskip}
\noindent\begin{minipage}{\linewidth}

\subsubsection{作业八：可动铁芯与电磁力（拔高）}\label{m5:ux4f5cux4e1aux516bux53efux52a8ux94c1ux82afux4e0eux7535ux78c1ux529bux62d4ux9ad8}

\textbf{题号：M5-097}

\textbf{层级：} 复赛　　\textbf{方法：} 自感互感

取 \(x\) 为给定机械广义坐标，磁路线性、无磁滞；力的正方向与 \(x\)
增大方向一致。

一个线性电磁装置的自感随广义坐标 \(x\) 变化为
\(L(x)\)，理想电流源把线圈电流恒定在 \(I\)。忽略机械耗散，求沿 \(x\)
方向的电磁力。

\end{minipage}\par

\par\noindent\begin{minipage}{\linewidth}
\textbf{解答}

\begin{center}
\includegraphics[width=\linewidth,height=.40\textheight,keepaspectratio]{\detokenize{assets/module-05/examples/M5-097.pdf}}
\par\smallskip{\small 图 M5-097：可动铁芯与电磁力}
\end{center}
\end{minipage}\par

图 M5-097 的 \(x\) 表示铁芯插入深度；若插入使电感增大，力就沿
\(+x\)。固定电流时应对磁共能求位置偏导。

电流受控时应使用磁共能

\[
W_m'(I,x)=\int_0^I\Psi(i,x)\,di.
\]

线性系统中 \(\Psi=L(x)i\)，所以

\[
W_m'=\frac12L(x)I^2.
\]

保持电流不变时，广义力为

\[
\boxed{
F_x=\left.\frac{\partial W_m'}{\partial x}\right|_I
=\frac12I^2\frac{dL}{dx}}.
\]

因此装置倾向于向自感增大的方向运动。若控制量是磁链而不是电流，则应在保持磁链不变的条件下使用磁场能求导，不能直接照搬上式。

\section{专题四十八：感生电场的空间分布、边值约束与叠加法（拔高选讲）}\label{m5:ux4e13ux9898ux56dbux5341ux516bux611fux751fux7535ux573aux7684ux7a7aux95f4ux5206ux5e03ux8fb9ux503cux7ea6ux675fux4e0eux53e0ux52a0ux6cd5ux62d4ux9ad8ux9009ux8bb2}

法拉第定律首先给出的是感生电场的\textbf{环量}或\textbf{旋度}。只有在对称性足够强时，才可以从一个环量方程直接读出每一点的电场。面对非圆形磁区、多个变化磁区或导体边界，必须把旋度、散度、边界条件与全局条件同时使用。

\subsection{一、环量不能单独确定电场}\label{m5:ux4e00ux73afux91cfux4e0dux80fdux5355ux72ecux786eux5b9aux7535ux573a}

固定回路满足

\[
\oint_C\mathbf E\cdot d\mathbf l
=-\frac{d}{dt}\int_S\mathbf B\cdot d\mathbf S,
\]

局域形式为

\[
\boxed{\nabla\times\mathbf E=-\frac{\partial\mathbf B}{\partial t}}.
\]

但旋度相同的两个场可以相差任意梯度场。若 \(\mathbf E_0\) 是一个特解，则

\[
\mathbf E=\mathbf E_0-\nabla\varphi
\]

仍有相同的旋度。因此，完整确定 \(\mathbf E\) 还要结合

\[
\boxed{\nabla\cdot\mathbf E=\frac{\rho}{\varepsilon_0}}
\]

以及边界条件、无穷远条件和导体上的电荷重排。上式是\textbf{真空中的微观形式}；在线性介质的宏观问题中，应改用

\[
\nabla\cdot\mathbf D=\rho_f
\]

并与本构关系共同求解，不能把自由电荷密度直接代入
\(\nabla\cdot\mathbf E=\rho/\varepsilon_0\)。用势表示时，这一结构写成

\[
\boxed{\mathbf E=-\frac{\partial\mathbf A}{\partial t}-\nabla\varphi},
\qquad
\mathbf B=\nabla\times\mathbf A.
\]

其中 \(-\partial\mathbf A/\partial t\)
提供由变化磁场决定的旋度，\(-\nabla\varphi\) 则由电荷与边界调节。

\par\noindent\begin{minipage}{\linewidth}
\begin{center}
\includegraphics[width=\linewidth,height=.40\textheight,keepaspectratio]{\detokenize{assets/module-05/theory/tikz-m5-topic48-v9-from45-v8-01.pdf}}
\par\smallskip{\small 图示：旋度约束与感生电场的不唯一性}
\end{center}
\end{minipage}\par

\par\addvspace{.4\baselineskip}
\noindent\begin{minipage}{\linewidth}

\textbf{例 M5-098：同一旋度并不对应唯一电场}

整个平面内 \(\partial_t\mathbf B=\dot B\mathbf e_z\)
为给定常向量。已知一个电场特解
\(\mathbf E_0=\dot B(y\mathbf e_x-x\mathbf e_y)/2\)。在其上叠加任意恒定场
\(\mathbf E_c\)，验证旋度不变，并据此说明仅给法拉第旋度约束能否唯一确定各点电场。

\end{minipage}\par

\par\noindent\begin{minipage}{\linewidth}
\textbf{解答}

\begin{center}
\includegraphics[width=\linewidth,height=.40\textheight,keepaspectratio]{\detokenize{assets/module-05/examples/M5-098.pdf}}
\par\smallskip{\small 图 M5-098：法拉第旋度约束与电场的不唯一性}
\end{center}
\end{minipage}\par

图 M5-098
的蓝色特解加上一个向右的恒定无旋场后成为绿色电场；两者旋度相同而各点场值不同，说明还需要边界与电荷信息。

对题设恒定附加电场，直接取
\(\delta\varphi=-\mathbf E_c\cdot\mathbf r\)，可得
\(-\nabla\delta\varphi=\mathbf E_c\)、\(\nabla\times\mathbf E_c=0\)、\(\nabla\cdot\mathbf E_c=0\)。它不改变法拉第定律中的旋度，也不在区域内增加体电荷；是否允许实际叠加，还须满足导体等边界条件。下文调和函数说法是对此例的推广，不是本题计算额外必需的工具。

设整个平面内
\(\partial\mathbf B/\partial t=\dot B\,\mathbf e_z\)。一个特解为

\[
\mathbf E_0=-\frac{\dot B}{2}\,\mathbf e_z\times\mathbf r.
\]

对任意调和函数 \(\varphi\)，场 \(\mathbf E_0-\nabla\varphi\)
仍满足相同的法拉第方程。只有给定电荷分布和边界条件后，\(\varphi\)
才能确定。

\subsection{二、边界上的约束}\label{m5:ux4e8cux8fb9ux754cux4e0aux7684ux7ea6ux675f}

取一个跨过静止分界面的无穷小矩形回路。若
\(\partial\mathbf B/\partial t\)
在界面附近有限，则右端磁通变化随矩形面积趋于零，故切向电场连续：

\[
\boxed{\mathbf n\times(\mathbf E_2-\mathbf E_1)=\mathbf0}.
\]

法向分量由高斯定律决定：

\[
\boxed{\mathbf n\cdot(\varepsilon_2\mathbf E_2-\varepsilon_1\mathbf E_1)=\sigma_f}.
\]

所以``感生电场沿磁区边界切向''并不是普遍定律。只有对称性、导体约束或题设进一步保证法向分量为零时，才能作此判断。理想化的阶跃磁场还可能隐含奇异表面电流；使用其结果时应明确这是薄边界近似。

\par\addvspace{.4\baselineskip}
\noindent\begin{minipage}{\linewidth}

\textbf{例 M5-099：导体表面对电场的修正}

一静止理想导体置于变化磁场附近，假定题设允许其内部保持总电场为零，且界面附近
\(\partial_t\mathbf B\)
有限。问导体表面外侧的切向总电场应满足什么条件，以及表面电荷所产生的无旋场怎样修正自由空间涡旋场。

\end{minipage}\par

\par\noindent\begin{minipage}{\linewidth}
\textbf{解答}

\begin{center}
\includegraphics[width=\linewidth,height=.40\textheight,keepaspectratio]{\detokenize{assets/module-05/examples/M5-099.pdf}}
\par\smallskip{\small 图 M5-099：导体表面对电场的修正}
\end{center}
\end{minipage}\par

图 M5-099
的细长矩形跨过静止导体边界；让其高度趋于零，便得到总切向场连续。表面电荷须修正自由空间涡旋场以满足导体边界。

一块静止理想导体置于变化磁场附近。导体内若保持
\(\mathbf E=\mathbf0\)，则外侧切向电场必须在表面趋于零。为满足这一条件，表面电荷会产生附加的
\(-\nabla\varphi\)，所以自由空间中由对称性得到的涡旋场不能未经修改直接延伸到导体内部。

\subsection{三、椭圆磁区内的感生电场}\label{m5:ux4e09ux692dux5706ux78c1ux533aux5185ux7684ux611fux751fux7535ux573a}

设无限长椭圆柱区域

\[
\frac{x^2}{a^2}+\frac{y^2}{b^2}<1
\]

内存在均匀轴向磁场 \(\mathbf B=B(t)\mathbf e_z\)。仅给出这一区域内的
\(\mathbf B\) 仍不足以唯一确定
\(\mathbf E\)。下面研究一个可解析的\textbf{内区边值模型}：内区无自由电荷，电场在中心正则，并由边界上的电荷或外部装置保证椭圆边界满足

\[
\mathbf E\cdot\mathbf n=0.
\]

这是附加边界条件，不是单由反射对称性推出的结论。在这些条件下，寻找关于
\(x,y\) 的一次场

\[
E_x=Ay,\qquad E_y=Cx.
\]

法拉第方程要求

\[
\frac{\partial E_y}{\partial x}-\frac{\partial E_x}{\partial y}
=C-A=-\dot B.
\]

椭圆边界的法向与

\[
\left(\frac{x}{a^2},\frac{y}{b^2}\right)
\]

平行。代入已给定的切向边界条件，得到

\[
\frac{A}{a^2}+\frac{C}{b^2}=0.
\]

联立得

\[
\boxed{
\mathbf E_{\text{内}}
=\frac{\dot B}{a^2+b^2}
\left(a^2y\,\mathbf e_x-b^2x\,\mathbf e_y\right)}.
\]

直接计算可检验 \(\nabla\cdot\mathbf E=0\)，且

\[
\nabla\times\mathbf E=-\dot B\,\mathbf e_z.
\]

因此该式是上述内区边值问题的解。若题目只是给出自由空间中的阶跃椭圆磁区，却没有给出边界电荷或边界条件，还必须继续求外区场，并用切向连续条件完成匹配；此时不能把上式当作普遍的``椭圆磁区公式''。

当 \(a=b=R\) 时，退化为圆域结果

\[
\mathbf E=-\frac{\dot B}{2}r\,\mathbf e_{\phi}.
\]

\par\noindent\begin{minipage}{\linewidth}
\begin{center}
\includegraphics[width=\linewidth,height=.40\textheight,keepaspectratio]{\detokenize{assets/module-05/original/tikz-m5-topic48-v9-from45-v8-04.png}}
\par\smallskip{\small 图示：椭圆磁区内的感生电场沿同心椭圆切向分布}
\end{center}
\end{minipage}\par

\par\addvspace{.4\baselineskip}
\noindent\begin{minipage}{\linewidth}

\textbf{例 M5-100：长轴端和短轴端的场强}

椭圆区域 \(x^2/a^2+y^2/b^2<1\) 内有均匀 \(\dot B_z\)，给定内区边值解
\(\mathbf E=\dot B_z(a^2y\mathbf e_x-b^2x\mathbf e_y)/(a^2+b^2)\)。其中
\(a>b>0\)，边界另满足 \(\mathbf E\cdot\mathbf n=0\)。求长轴端
\((a,0)\)、短轴端 \((0,b)\) 的场强并比较。

\end{minipage}\par

\par\noindent\begin{minipage}{\linewidth}
\textbf{解答}

\begin{center}
\includegraphics[width=\linewidth,height=.40\textheight,keepaspectratio]{\detokenize{assets/module-05/examples/M5-100.pdf}}
\par\smallskip{\small 图 M5-100：长轴端和短轴端的场强}
\end{center}
\end{minipage}\par

图 M5-100
的右端与上端分别是长轴端和短轴端；两个红色切向箭头按给定边值解计算，短轴端场强更大，不能用到中心的距离直接套圆域公式。

在 \((a,0)\) 处

\[
|\mathbf E|=\frac{|\dot B|ab^2}{a^2+b^2},
\]

在 \((0,b)\) 处

\[
|\mathbf E|=\frac{|\dot B|a^2b}{a^2+b^2}.
\]

若
\(a>b\)，短轴端的场更强；这与同一闭合椭圆上的线元长度和曲率分布有关，不能把圆域公式中的
\(r\) 机械换成到中心的距离。

\subsection{四、叠加法与远场判断}\label{m5:ux56dbux53e0ux52a0ux6cd5ux4e0eux8fdcux573aux5224ux65ad}

麦克斯韦方程在线性介质中对 \(\mathbf E\)、\(\mathbf B\)
是线性的。若变化磁场可写成

\[
\frac{\partial\mathbf B}{\partial t}
=\sum_j\frac{\partial\mathbf B_j}{\partial t},
\]

则在同一组边界条件下

\[
\boxed{\mathbf E=\sum_j\mathbf E_j}.
\]

对半径为 \(a_j\)、圆心为 \(\mathbf R_j\) 的小圆磁区，在其外部令

\[
\boldsymbol\rho_j=\mathbf r-\mathbf R_j,
\]

则

\[
\mathbf E_j
=-\frac{a_j^2\dot B_j}{2|\boldsymbol\rho_j|^2}
\,\mathbf e_z\times\boldsymbol\rho_j.
\]

若多个磁区的总变化磁通率

\[
\sum_j\pi a_j^2\dot B_j
\]

不为零，远处的主项按 \(1/r\) 衰减；若总变化磁通率为零，\(1/r\)
主项抵消，远场衰减更快。这是判断叠加结果是否合理的有力检验。

\par\noindent\begin{minipage}{\linewidth}
\begin{center}
\includegraphics[width=\linewidth,height=.40\textheight,keepaspectratio]{\detokenize{assets/module-05/theory/tikz-m5-topic48-v9-from45-v8-06.pdf}}
\par\smallskip{\small 图示：两反向变化磁区的电场叠加}
\end{center}
\end{minipage}\par

\subsection{五、非轴对称问题的标准步骤}\label{m5:ux4e94ux975eux8f74ux5bf9ux79f0ux95eeux9898ux7684ux6807ux51c6ux6b65ux9aa4}

\begin{enumerate}
\def\labelenumi{\arabic{enumi}.}
\tightlist
\item
  先画出真正具有的平移、转动和镜面对称性，不凭图形``近似像圆''套公式。
\item
  写出 \(\nabla\times\mathbf E=-\partial\mathbf B/\partial t\) 与
  \(\nabla\cdot\mathbf E=\rho/\varepsilon_0\)。
\item
  在各光滑区域内选取满足对称性的最低阶试探场。
\item
  用切向连续、法向跃变、导体等势或无穷远衰减条件定系数。
\item
  用任意闭合回路的环量检验法拉第定律，并检查量纲、极限和连续性。
\end{enumerate}

\subsection{六、使用条件与常见错误}\label{m5:ux516dux4f7fux7528ux6761ux4ef6ux4e0eux5e38ux89c1ux9519ux8bef}

\begin{enumerate}
\def\labelenumi{\arabic{enumi}.}
\tightlist
\item
  法拉第定律给出环量，不直接给出任意点的场强。
\item
  \(\partial\mathbf B/\partial t=0\) 的区域内仍可能有
  \(\mathbf E\ne0\)；此时只是局域旋度为零。
\item
  旋度为零不等于全局环量一定为零；区域若有孔或排除了奇点，仍须单独检查闭合路径。
\item
  叠加前必须保证各子问题采用同一介质和边界条件。
\item
  边界切向连续并不意味着法向连续；法向跃变由面电荷决定。
\item
  非圆形磁区一般不能只用一条积分回路求出整个场分布。
\end{enumerate}

\subsection{练习与解答}\label{m5:ux7ec3ux4e60ux4e0eux89e3ux7b54-5}

\par\addvspace{.4\baselineskip}
\noindent\begin{minipage}{\linewidth}

\subsubsection{作业一：椭圆磁区内场的检验}\label{m5:ux4f5cux4e1aux4e00ux692dux5706ux78c1ux533aux5185ux573aux7684ux68c0ux9a8c}

\textbf{题号：M5-101}

\textbf{层级：} 复赛　　\textbf{方法：} 法拉第定律、边界条件

设 \(a\text{、}b>0\)，内区无自由电荷，中心正则，并另给椭圆边界条件
\(\mathbf E\cdot\mathbf n=0\)；这里检验该内区边值解，不把它当作任意自由空间椭圆磁区的通用公式。

对本专题得到的

\[
\mathbf E=\frac{\dot B}{a^2+b^2}
(a^2y\mathbf e_x-b^2x\mathbf e_y)
\]

分别计算散度、旋度，并证明电场与任意同心椭圆相切。

\end{minipage}\par

\par\noindent\begin{minipage}{\linewidth}
\textbf{解答}

\begin{center}
\includegraphics[width=\linewidth,height=.40\textheight,keepaspectratio]{\detokenize{assets/module-05/examples/M5-101.pdf}}
\par\smallskip{\small 图 M5-101：椭圆磁区内场的检验}
\end{center}
\end{minipage}\par

图 M5-101 中椭圆族的法向与 \((x/a^2,y/b^2)\)
平行，代入场分量后法向投影相消，说明红色场方向确与同心椭圆相切。

散度为

\[
\nabla\cdot\mathbf E
=\frac{\partial E_x}{\partial x}+\frac{\partial E_y}{\partial y}=0.
\]

旋度的 \(z\) 分量为

\[
(\nabla\times\mathbf E)_z
=\frac{\partial E_y}{\partial x}-\frac{\partial E_x}{\partial y}
=-\frac{b^2\dot B}{a^2+b^2}-\frac{a^2\dot B}{a^2+b^2}
=-\dot B.
\]

\par\noindent\begin{minipage}{\linewidth}
\begin{center}
\includegraphics[width=\linewidth,height=.40\textheight,keepaspectratio]{\detokenize{assets/module-05/process/M5-101a.pdf}}
\par\smallskip{\small 图 M5-101a：椭圆边界上切向场的两个端点}
\end{center}
\end{minipage}\par

图 M5-101a
的右端电场向下、上端电场向右，两个方向都与各自边界法向垂直；一般点的零法向投影由右侧内积检验。

同心椭圆 \(x^2/a^2+y^2/b^2=c\) 的法向与

\[
\mathbf n_0=\frac{x}{a^2}\mathbf e_x+\frac{y}{b^2}\mathbf e_y
\]

平行。于是

\[
\mathbf E\cdot\mathbf n_0
=\frac{\dot Bxy}{a^2+b^2}-\frac{\dot Bxy}{a^2+b^2}=0.
\]

故电场处处与同心椭圆相切。

\par\addvspace{.4\baselineskip}
\noindent\begin{minipage}{\linewidth}

\subsubsection{作业二：两个相反变化的圆磁区}\label{m5:ux4f5cux4e1aux4e8cux4e24ux4e2aux76f8ux53cdux53d8ux5316ux7684ux5706ux78c1ux533a}

\textbf{题号：M5-102}

\textbf{层级：} 复赛　　\textbf{方法：} 法拉第定律、叠加法

两圆柱磁区平行 \(z\)
轴、截面不重叠，场均匀且外部为零；采用无外加无旋场的自由空间解。考察距离
\(r\gg d,a\) 处。

两个半径均为 \(a\) 的圆磁区相距 \(d\)，一个满足
\(\dot B=\beta\)，另一个满足 \(\dot B=-\beta\)。说明远处电场的 \(1/r\)
主项为何消失，并估计远场的衰减阶数。

\end{minipage}\par

\par\noindent\begin{minipage}{\linewidth}
\textbf{解答}

\begin{center}
\includegraphics[width=\linewidth,height=.40\textheight,keepaspectratio]{\detokenize{assets/module-05/examples/M5-102.pdf}}
\par\smallskip{\small 图 M5-102：两个相反变化的圆磁区}
\end{center}
\end{minipage}\par

图 M5-102 的两个磁区面积相等、变化率相反，远处分别产生的环向 \(1/r\)
项抵消，剩下与中心间距成正比的下一阶项。

每个圆磁区在远处等效贡献一个与变化磁通率

\[
\dot\Phi_j=\pi a^2\dot B_j
\]

成正比的 \(1/r\) 环向场。两区满足

\[
\dot\Phi_1+\dot\Phi_2=\pi a^2\beta-\pi a^2\beta=0,
\]

所以总场的 \(1/r\) 单极型主项抵消。把两个场在 \(d/r\ll1\)
下展开，下一项比单区远场多一个因子 \(d/r\)，故

\[
\boxed{|\mathbf E_{\text{远}}|\sim\frac{a^2|\beta|d}{r^2}}.
\]

方向由两个子场的矢量和决定，不再是绕任一圆心的纯环向场。

\par\addvspace{.4\baselineskip}
\noindent\begin{minipage}{\linewidth}

\subsubsection{作业三：为什么一条矩形环路不够}\label{m5:ux4f5cux4e1aux4e09ux4e3aux4ec0ux4e48ux4e00ux6761ux77e9ux5f62ux73afux8defux4e0dux591f}

\textbf{题号：M5-103}

\textbf{层级：} 复赛　　\textbf{方法：} 法拉第定律、对称性检验

只给出矩形磁区内外的磁场变化率，没有另给保证每边局部场恒定的边界条件。

某矩形区域内
\(\partial B_z/\partial t=\beta\)，区域外为零。有人取恰好包围该矩形的回路，求得总环量
\(-\beta S\)，再把它平均分给四条边，断言每条边上切向电场恒定且大小相同。指出错误。

\end{minipage}\par

\par\noindent\begin{minipage}{\linewidth}
\textbf{解答}

\begin{center}
\includegraphics[width=\linewidth,height=.40\textheight,keepaspectratio]{\detokenize{assets/module-05/examples/M5-103.pdf}}
\par\smallskip{\small 图 M5-103：为什么一条矩形环路不够}
\end{center}
\end{minipage}\par

图 M5-103
的四段积分相加只给出一个总量；即使总环量已知，也不能把它平均分到不等价的边或同一条边的各点。

法拉第定律只给出

\[
\oint_C\mathbf E\cdot d\mathbf l=-\beta S,
\]

这是四条边积分的总和。矩形只有离散镜面对称性，长边与短边并不等价；即使同一条边上，场强也未必处处相同。要确定分布，还需散度方程、各区域的边界条件和无穷远条件。将总环量平均分配给边长不同、局部环境不同的四段，没有任何对称性依据。

\par\addvspace{.4\baselineskip}
\noindent\begin{minipage}{\linewidth}

\subsubsection{作业四：导体边界上的切向条件}\label{m5:ux4f5cux4e1aux56dbux5bfcux4f53ux8fb9ux754cux4e0aux7684ux5207ux5411ux6761ux4ef6}

\textbf{题号：M5-104}

\textbf{层级：} 复赛　　\textbf{方法：} 法拉第定律、边界条件

分界面静止、场在两侧分片光滑，磁场时间导数无奇异面项；采用跨界面细长矩形的极限。

证明在 \(\partial\mathbf B/\partial t\)
有限时，静止分界面两侧电场切向分量连续。

\end{minipage}\par

\par\noindent\begin{minipage}{\linewidth}
\textbf{解答}

\begin{center}
\includegraphics[width=\linewidth,height=.40\textheight,keepaspectratio]{\detokenize{assets/module-05/examples/M5-104.pdf}}
\par\smallskip{\small 图 M5-104：导体边界上的切向条件}
\end{center}
\end{minipage}\par

图 M5-104
的两条长边紧贴界面两侧；有限的磁通变化率使右端随面积趋零，短边贡献也趋零，留下切向连续条件。

取跨过界面的窄矩形回路，长边 \(\Delta l\) 平行界面，短边高为
\(h\)。由法拉第定律

\[
\oint\mathbf E\cdot d\mathbf l
=-\frac d{dt}\int\mathbf B\cdot d\mathbf S.
\]

令 \(h\to0\)。右端与面积 \(h\Delta l\)
同阶而趋零，两条短边积分也趋零，只剩

\[
(E_{1t}-E_{2t})\Delta l=0.
\]

所以

\[
\boxed{E_{1t}=E_{2t}},
\]

矢量形式即 \(\mathbf n\times(\mathbf E_2-\mathbf E_1)=\mathbf0\)。

\par\addvspace{.4\baselineskip}
\noindent\begin{minipage}{\linewidth}

\subsubsection{作业五：无限长变化磁场条带}\label{m5:ux4f5cux4e1aux4e94ux65e0ux9650ux957fux53d8ux5316ux78c1ux573aux6761ux5e26}

\textbf{题号：M5-105}

\textbf{层级：} 复赛　　\textbf{方法：} 法拉第定律、边值约束

取标准右手直角坐标，不叠加额外恒定电场；结论只要求条带内部，条带外还受全局边界选择影响。

无限长条带 \(|y|<H/2\) 内有均匀磁场
\(B_z(t)\)，条带外磁场为零。条带内无自由电荷，且不叠加破坏系统对称性的外加静电场。已知

\[
-\frac{dB_z}{dt}=\alpha
\]

为正常数。条带沿 \(x\)
方向无限延伸。求条带内部可由对称性唯一确定的感生电场，并说明适用条件。

\end{minipage}\par

\par\noindent\begin{minipage}{\linewidth}
\textbf{解答}

\begin{center}
\includegraphics[width=\linewidth,height=.40\textheight,keepaspectratio]{\detokenize{assets/module-05/examples/M5-105.pdf}}
\par\smallskip{\small 图 M5-105：无限变化磁场条带的感生场}
\end{center}
\end{minipage}\par

图 M5-105
的两条水平线是磁条带边界；平移与反射对称性使电场只依赖y，并关于中线反号，红箭头明确标出两个半区的方向。

系统沿 \(x\) 平移不变，故 \(\mathbf E\) 只依赖 \(y\)。条带内
\(\nabla\cdot\mathbf E=0\) 给出 \(dE_y/dy=0\)；再结合绕 \(z\) 轴转过
\(\pi\) 后系统不变，有 \(\mathbf E(-y)=-\mathbf E(y)\)，因此常量 \(E_y\)
只能为零，且 \(E_x(-y)=-E_x(y)\)。于是

\[
\mathbf E=E_x(y)\mathbf e_x,
\qquad E_x(0)=0.
\]

法拉第方程的 \(z\) 分量为

\[
(\nabla\times\mathbf E)_z
=-\frac{dE_x}{dy}
=\alpha.
\]

积分得

\[
\boxed{\mathbf E=-\alpha y\,\mathbf e_x,
\qquad |y|<\frac H2}.
\]

上下半条带内电场方向相反。该结果依赖条带沿 \(x\)
无限长；若长度有限，端部破坏平移对称性，场还会出现 \(y\) 分量。

\par\addvspace{.4\baselineskip}
\noindent\begin{minipage}{\linewidth}

\subsubsection{作业六：正方形与等边三角形磁区中心附近的场}\label{m5:ux4f5cux4e1aux516dux6b63ux65b9ux5f62ux4e0eux7b49ux8fb9ux4e09ux89d2ux5f62ux78c1ux533aux4e2dux5fc3ux9644ux8fd1ux7684ux573a}

\textbf{题号：M5-106}

\textbf{层级：} 复赛　　\textbf{方法：} 法拉第定律、局部展开

题中 \(\alpha\)
为常量，内区无自由电荷，中心场正则；假定电荷与边界装置均保持各磁区完整旋转对称性，不叠加破坏对称性的静电场。

正方形或等边三角形磁区内有均匀
\(-dB_z/dt=\alpha\)。设几何中心为原点，讨论中心附近感生电场的一阶项。

\end{minipage}\par

\par\noindent\begin{minipage}{\linewidth}
\textbf{解答}

\begin{center}
\includegraphics[width=\linewidth,height=.40\textheight,keepaspectratio]{\detokenize{assets/module-05/examples/M5-106.pdf}}
\par\smallskip{\small 图 M5-106：正方形与等边三角形磁区中心附近的场}
\end{center}
\end{minipage}\par

图 M5-106
的正方形中心与等边三角形中心都具有至少三重旋转对称性。一次场中的各向异性部分被对称性排除，旋度再确定环向系数。

中心附近把电场展开为

\[
\mathbf E=\boldsymbol M\mathbf r+O(r^2).
\]

正方形的四重旋转对称性和等边三角形的三重旋转对称性都排除了具有优选轴向的对称线性部分；再结合中心无电荷时
\(\nabla\cdot\mathbf E=0\)，一阶项只能是刚性旋转型：

\[
\mathbf E=c(-y\mathbf e_x+x\mathbf e_y)+O(r^2).
\]

其旋度为 \(2c\mathbf e_z\)，与 \(\alpha\mathbf e_z\) 比较得
\(c=\alpha/2\)。因此

\[
\boxed{
\mathbf E=\frac\alpha2(-y\mathbf e_x+x\mathbf e_y)+O(r^2)}.
\]

这一结论只保证中心附近的一阶行为。离边界较近时，正方形与三角形的高阶修正不同，不能把上式当作整个区域内的精确解。

\par\addvspace{.4\baselineskip}
\noindent\begin{minipage}{\linewidth}

\subsubsection{作业七：两个相邻圆磁区中心连线上的场}\label{m5:ux4f5cux4e1aux4e03ux4e24ux4e2aux76f8ux90bbux5706ux78c1ux533aux4e2dux5fc3ux8fdeux7ebfux4e0aux7684ux573a}

\textbf{题号：M5-107}

\textbf{层级：} 复赛　　\textbf{方法：} 法拉第定律、叠加法

两磁区是无限长圆柱的横截面，取无附加无旋场的轴对称单圆解后线性叠加，\(\beta>0\)。

两个半径均为 \(a\) 的圆磁区，圆心位于 \(x=\pm c\)，且 \(c>a\)。左区满足
\(\dot B_z=\beta\)，右区满足 \(\dot B_z=-\beta\)。求原点处感生电场。

\end{minipage}\par

\par\noindent\begin{minipage}{\linewidth}
\textbf{解答}

\begin{center}
\includegraphics[width=\linewidth,height=.40\textheight,keepaspectratio]{\detokenize{assets/module-05/examples/M5-107.pdf}}
\par\smallskip{\small 图 M5-107：两个相邻圆磁区中心连线上的场}
\end{center}
\end{minipage}\par

图 M5-107
的原点位于两个磁区之外且到两圆心等距。左区增强、右区减弱，右手定则给出两个电场在
\(O\) 都向下，因此相加。

原点位于两个圆区外。单个圆磁区在外部产生

\[
\mathbf E_j
=-\frac{a^2\dot B_j}{2\rho_j^2}
\mathbf e_z\times\boldsymbol\rho_j.
\]

对左圆，\(\boldsymbol\rho_1=c\mathbf e_x\)，所以

\[
\mathbf E_1=-\frac{a^2\beta}{2c}\mathbf e_y.
\]

对右圆，\(\boldsymbol\rho_2=-c\mathbf e_x\) 且 \(\dot B_2=-\beta\)，故

\[
\mathbf E_2=-\frac{a^2\beta}{2c}\mathbf e_y.
\]

两场同向叠加：

\[
\boxed{\mathbf E(0)=-\frac{a^2\beta}{c}\mathbf e_y}.
\]

这里``变化磁通率相反''并不意味着两圆之间处处相消；在中点处，两者的绕行方向恰好给出同向电场。

\par\addvspace{.4\baselineskip}
\noindent\begin{minipage}{\linewidth}

\subsubsection{作业八：转动的非均匀磁场图样}\label{m5:ux4f5cux4e1aux516bux8f6cux52a8ux7684ux975eux5747ux5300ux78c1ux573aux56feux6837}

\textbf{题号：M5-108}

\textbf{层级：} 复赛　　\textbf{方法：} 法拉第定律、坐标变换

\(B_0\) 的单位为 \(\mathrm{T/m}\)，\(\omega\)
为常量；仅求一个内区特解，未指定圆域边界电势，因此不要求唯一性或特定边界法向条件。

半径为 \(R\) 的圆域内

\[
B_z(r,\theta,t)=B_0r\cos(\theta-\omega t).
\]

在圆心附近无自由电荷，并要求场在 \(r=0\) 正则。求一个满足法拉第方程和
\(\nabla\cdot\mathbf E=0\) 的内区电场。

\end{minipage}\par

\par\noindent\begin{minipage}{\linewidth}
\textbf{解答}

\begin{center}
\includegraphics[width=\linewidth,height=.40\textheight,keepaspectratio]{\detokenize{assets/module-05/examples/M5-108.pdf}}
\par\smallskip{\small 图 M5-108：转动磁场图样的一个正则内区电场}
\end{center}
\end{minipage}\par

图 M5-108
示出t=0时左侧向纸内、右侧向纸外的线性磁场分布；转动会使各点磁场随时间变化。下方检验一个正则电场的散度和旋度，边界未给时不能断言唯一。

令 \(\psi=\theta-\omega t\)，按角向结构设

\[
E_r=Ar^2\cos\psi,
\qquad
E_{\theta}=Cr^2\sin\psi.
\]

散度条件给出

\[
\nabla\cdot\mathbf E
=(3A+C)r\cos\psi=0,
\]

故 \(C=-3A\)。旋度为

\[
(\nabla\times\mathbf E)_z
=(3C+A)r\sin\psi.
\]

另一方面

\[
-\frac{\partial B_z}{\partial t}
=-B_0\omega r\sin\psi.
\]

\par\noindent\begin{minipage}{\linewidth}
\begin{center}
\includegraphics[width=\linewidth,height=.40\textheight,keepaspectratio]{\detokenize{assets/module-05/process/M5-108a.pdf}}
\par\smallskip{\small 图 M5-108a：转动图样在t=0时的一个二次场}
\end{center}
\end{minipage}\par

图 M5-108a
的箭头由右侧二次场计算，横向分量处处非负，纵向分量在相邻象限变号；对这些分量求偏导可检验散度、旋度及中心正则性。

因此 \(3C+A=-B_0\omega\)。联立得

\[
A=\frac{B_0\omega}{8},
\qquad
C=-\frac{3B_0\omega}{8}.
\]

所以

\[
\boxed{
\mathbf E
=\frac{B_0\omega r^2}{8}
\left[\cos(\theta-\omega t)\mathbf e_r
-3\sin(\theta-\omega t)\mathbf e_{\theta}\right]}.
\]

若还要求与圆域外部精确匹配，须继续求外区无旋、无散场并使用边界连续条件；内区局部方程本身不能给出外区。

\section{专题四十九：线性无源元件的场与电路描述（基础必讲与拔高选讲）}\label{m5:ux4e13ux9898ux56dbux5341ux4e5dux7ebfux6027ux65e0ux6e90ux5143ux4ef6ux7684ux573aux4e0eux7535ux8defux63cfux8ff0ux57faux7840ux5fc5ux8bb2ux4e0eux62d4ux9ad8ux9009ux8bb2}

电阻、电容和电感都是集总二端元件，但能量关系不同：电阻耗散能量，电容用电荷描述状态并储存电场能，电感用磁链描述状态并储存磁场能。以下采用被动符号约定，列出三类元件的本构关系、功率和能量，并注明线性模型的适用条件。

除特别说明外，均采用被动符号约定：参考电流 \(i\) 流入参考电压 \(u\)
的正端。因此元件吸收的瞬时功率统一写成

\[
\boxed{p=ui}.
\]

\(p>0\) 表示元件从外电路吸收能量，\(p<0\)
表示元件向外电路释放能量。参考方向改变时，元件本身的物理状态不变，但方程中的代数符号必须随之调整。

\subsection{一、电阻、电容和电感的共同表示}\label{m5:ux4e00ux7535ux963bux7535ux5bb9ux548cux7535ux611fux7684ux5171ux540cux8868ux793a}

\subsubsection{1.
本构关系、状态量与能量}\label{m5:ux672cux6784ux5173ux7cfbux72b6ux6001ux91cfux4e0eux80fdux91cf}

理想线性电阻、电容和电感的本构关系分别为

\[
\boxed{u=Ri},
\]

\[
\boxed{q=Cu,\qquad i=\frac{dq}{dt}=C\frac{du}{dt}},
\]

\[
\boxed{\Psi=Li,\qquad u=\frac{d\Psi}{dt}=L\frac{di}{dt}}.
\]

后两式中写成最后一个等号，要求 \(C\)、\(L\)
为常量。若几何位置、介质状态或元件参数随时间改变，应保留全导数。

\par\noindent\begin{minipage}{\linewidth}
\begin{center}
\includegraphics[width=\linewidth,height=.40\textheight,keepaspectratio]{\detokenize{assets/module-05/original/tikz-m5-topic49-v9-from46-v8-01.png}}
\par\smallskip{\small 图示：同一参考方向下电阻、电容和电感的本构关系与能量}
\end{center}
\end{minipage}\par

三类元件的差别列于下表：

{\def\LTcaptype{none} % do not increment counter
\begin{longtable}[]{@{}lllll@{}}
\toprule\noalign{}
元件 & 本构关系 & 状态量 & 可恢复能量 & 理想作用 \\
\midrule\noalign{}
\endhead
\bottomrule\noalign{}
\endlastfoot
电阻 & \(u=Ri\) & 无独立储能状态 & 无 & 耗散 \\
电容 & \(q=Cu\) & \(q\) 或 \(u\) & \(W_e=Cu^2/2\) & 储存电场能 \\
电感 & \(\Psi=Li\) & \(\Psi\) 或 \(i\) & \(W_m=Li^2/2\) & 储存磁场能 \\
\end{longtable}
}

相应单位为

\[
[R]=\Omega=\frac{\mathrm V}{\mathrm A},
\]

\[
[C]=\mathrm F=\frac{\mathrm C}{\mathrm V}
=\frac{\mathrm{A\,s}}{\mathrm V},
\]

\[
[L]=\mathrm H=\frac{\mathrm{Wb}}{\mathrm A}
=\frac{\mathrm{V\,s}}{\mathrm A}.
\]

\par\addvspace{.4\baselineskip}
\noindent\begin{minipage}{\linewidth}

\textbf{例 M5-109：量纲检查}

线性理想电容、电感采用被动参考方向，本构关系分别为 \(i=C\,du/dt\) 与
\(u=L\,di/dt\)。由这两个关系给出 \(C\text{、}L\)
的SI基本电学单位组合，并检验它们分别与法拉、亨利一致。

\end{minipage}\par

\par\noindent\begin{minipage}{\linewidth}
\textbf{解答}

\begin{center}
\includegraphics[width=\linewidth,height=.40\textheight,keepaspectratio]{\detokenize{assets/module-05/examples/M5-109.pdf}}
\par\smallskip{\small 图 M5-109：量纲检查}
\end{center}
\end{minipage}\par

图 M5-109
的电容使用电流流入正端的参考方向，从端口关系检查电容单位；对电感的
\(u=L\,di/dt\) 作同样检验。

两类单位要分别检验： \[
[C]=\mathrm C/\mathrm V
=\mathrm{A\cdot s/V}=\mathrm F,\qquad
[L]=\mathrm{V/(A/s)}=\mathrm{V\cdot s/A}=\mathrm H.
\]
电容电压的时间导数乘法拉得到安培，电感电流的时间导数乘亨利得到伏特，与图中两个本构式逐项一致。

由 \(i=C\,du/dt\) 可得

\[
[C]\frac{[u]}{[t]}
=\frac{\mathrm{A\,s}}{\mathrm V}\frac{\mathrm V}{\mathrm s}
=\mathrm A.
\]

若把电容关系误写成 \(u=C\,di/dt\)，量纲立即不成立。

\subsubsection{2.
``线性''与``无源''不能混用}\label{m5:ux7ebfux6027ux4e0eux65e0ux6e90ux4e0dux80fdux6df7ux7528}

线性描述输入与响应是否满足叠加原理。若元件关系以算子 \(\mathcal F\)
表示，则

\[
\boxed{\mathcal F(ax_1+bx_2)=a\mathcal F(x_1)+b\mathcal F(x_2)}.
\]

无源描述能量：元件不能在没有先前储能的情况下向外提供净能量。若 \(W(t)\)
是元件在时刻 \(t\) 的储能，则

\[
W(t)-W(t_0)
\le\int_{t_0}^{t}u(\tau)i(\tau)\,d\tau.
\]

等号与不等号之间的差额代表耗散。

\par\noindent\begin{minipage}{\linewidth}
\begin{center}
\includegraphics[width=\linewidth,height=.40\textheight,keepaspectratio]{\detokenize{assets/module-05/theory/tikz-m5-topic49-v9-from46-v8-02.pdf}}
\par\smallskip{\small 图示：线性与无源是不同的约束}
\end{center}
\end{minipage}\par

例如，\(u=ai^3\) 且 \(a>0\)
时，\(p=ai^4\ge0\)，它可以是无源电阻，但不是线性元件。反之，\(u=Ri\) 且
\(R<0\) 虽然是线性关系，却会在 \(i\ne0\) 时持续向外送出功率
\(-Ri^2\)，不是孤立的无源电阻模型。

``线性''也不自动等于``参数不随时间变化''。关系 \(q=C(t)u\) 对 \(q,u\)
仍可线性，却是时变元件。后续相量和阻抗方法通常还需要系统参数不随时间变化。

\par\addvspace{.4\baselineskip}
\noindent\begin{minipage}{\linewidth}

\textbf{例 M5-110：判断属性}

一个无记忆二端元件按被动参考方向满足
\(u=2i^3\)，其中数值系数带使单位相容的单位。用瞬时功率判断它是否能凭空供能，再用比例性判断它是否线性；不要将无源与线性当作同一条件。

\end{minipage}\par

\par\noindent\begin{minipage}{\linewidth}
\textbf{解答}

\begin{center}
\includegraphics[width=\linewidth,height=.40\textheight,keepaspectratio]{\detokenize{assets/module-05/examples/M5-110.pdf}}
\par\smallskip{\small 图 M5-110：无源条件与线性条件的区别}
\end{center}
\end{minipage}\par

图 M5-110 的弯曲关系直接表示电压 \(u\) 随电流 \(i\)
的非线性本构。功率非负保证耗能性，却不保证电压与电流成正比。

元件满足 \(u=2i^3\)。它有
\(ui=2i^4\ge0\)，因此不会凭空供能；但若把电流加倍，电压变为八倍而不是两倍，故它无源但非线性。

\subsection{二、电阻：由局域输运到端口耗散}\label{m5:ux4e8cux7535ux963bux7531ux5c40ux57dfux8f93ux8fd0ux5230ux7aefux53e3ux8017ux6563}

\subsubsection{1.
局域欧姆定律}\label{m5:ux5c40ux57dfux6b27ux59c6ux5b9aux5f8b}

各向同性、均匀、线性的欧姆介质满足

\[
\boxed{\mathbf J=\sigma\mathbf E},
\qquad
\boxed{\mathbf E=\rho\mathbf J},
\qquad
\rho=\frac1\sigma.
\]

其中 \(\sigma\) 是电导率，\(\rho\) 是电阻率。对长度为
\(\ell\)、等截面积为 \(S\) 的均匀直导体，若端部效应可以忽略，导体内
\(\mathbf E\) 与 \(\mathbf J\) 近似均匀，则

\[
U=E\ell,
\qquad
I=JS.
\]

因此

\[
\boxed{R=\rho\frac{\ell}{S}}.
\]

该几何公式不是任意形状导体的普遍公式。截面变化、接触电阻、非均匀材料或电流拥挤存在时，应先解
\(\nabla\cdot\mathbf J=0\)
与材料本构关系，再由端口电压和总电流定义等效电阻。

\subsubsection{2.
经典弛豫模型}\label{m5:ux7ecfux5178ux5f1bux8c6bux6a21ux578b}

设金属中自由电子数密度为 \(n\)、电子质量为 \(m\)、电荷为
\(-e\)，用平均弛豫时间 \(\tau\) 描述散射。平均漂移速度满足

\[
m\frac{d\mathbf v_d}{dt}
=-e\mathbf E-\frac m\tau\mathbf v_d.
\]

若 \(\mathbf E\) 为恒定电场，解为

\[
\mathbf v_d(t)
=\mathbf v_d(0)e^{-t/\tau}
-\frac{e\tau}{m}\mathbf E\left(1-e^{-t/\tau}\right).
\]

经过若干个弛豫时间后，

\[
\mathbf v_d\longrightarrow-\frac{e\tau}{m}\mathbf E.
\]

电流密度方向与传统电流方向一致：

\[
\mathbf J=-ne\mathbf v_d
=\frac{ne^2\tau}{m}\mathbf E.
\]

因此经典模型给出

\[
\boxed{\sigma=\frac{ne^2\tau}{m}}.
\]

\par\noindent\begin{minipage}{\linewidth}
\begin{center}
\includegraphics[width=\linewidth,height=.40\textheight,keepaspectratio]{\detokenize{assets/module-05/original/tikz-m5-topic49-v9-from46-v8-03.png}}
\par\smallskip{\small 图示：经典弛豫模型从电子平均运动逐步得到欧姆定律}
\end{center}
\end{minipage}\par

这个推导给出直观机制，但不是现代固体导电理论的全部。它忽略能带结构、费米统计和复杂散射机制；在高频下还必须保留
\(d\mathbf v_d/dt\)，电流与电场不再简单同相。竞赛中的直流或低频普通金属题，才常把
\(\sigma\) 视为实常量。

\par\addvspace{.4\baselineskip}
\noindent\begin{minipage}{\linewidth}

\textbf{例 M5-111：弛豫建立}

恒定电场在 \(t=0\) 突然接通，初始平均漂移为零，经典弛豫模型的电流为
\(J(t)=\sigma_0E_0(1-e^{-t/\tau})\)。求 \(t=\tau\)
时达到稳态的比例，说明弛豫时间是否等于完全建立所需时间。

\end{minipage}\par

\par\noindent\begin{minipage}{\linewidth}
\textbf{解答}

\begin{center}
\includegraphics[width=\linewidth,height=.40\textheight,keepaspectratio]{\detokenize{assets/module-05/examples/M5-111.pdf}}
\par\smallskip{\small 图 M5-111：弛豫建立}
\end{center}
\end{minipage}\par

图 M5-111 中 \(t/\tau=1\) 的虚线落在建立曲线上，其高度约为稳态虚线的
\(0.632\)；它不是稳态已经完全建立的时刻。

若初始平均漂移速度为零，则在 \(t=\tau\) 时

\[
\mathbf v_d(\tau)
=-\frac{e\tau}{m}\mathbf E(1-e^{-1}),
\]

只达到稳态漂移速度的约 \(63.2\%\)，而不是在一次碰撞后突然达到稳态。

\subsubsection{3.
焦耳热与功率守恒}\label{m5:ux7126ux8033ux70edux4e0eux529fux7387ux5b88ux6052}

单位体积、单位时间内电场对载流子做功为

\[
\boxed{p_V=\mathbf J\cdot\mathbf E}.
\]

对各向同性欧姆导体，

\[
p_V=\sigma E^2=\rho J^2\ge0.
\]

对均匀直导体积分，

\[
P=\int_V\mathbf J\cdot\mathbf E\,dV
=JE(S\ell)
=UI.
\]

再使用 \(U=IR\)，得到

\[
\boxed{P=UI=I^2R=\frac{U^2}{R}}.
\]

\par\noindent\begin{minipage}{\linewidth}
\begin{center}
\includegraphics[width=\linewidth,height=.40\textheight,keepaspectratio]{\detokenize{assets/module-05/original/tikz-m5-topic49-v9-from46-v8-04.png}}
\par\smallskip{\small 图示：局域焦耳热密度积分为电阻的端口功率}
\end{center}
\end{minipage}\par

各向异性材料应写成
\(\mathbf J=\boldsymbol\sigma\mathbf E\)。无源性要求耗散部分满足
\(\mathbf E\cdot\boldsymbol\sigma\mathbf E\ge0\)；不能把标量公式逐分量机械套用。

\par\addvspace{.4\baselineskip}
\noindent\begin{minipage}{\linewidth}

\textbf{例 M5-112：功率选择公式}

理想欧姆电阻 \(R=10\ \Omega\) 两端电压为
\(20\ \mathrm V\)。采用被动参考方向，求电流和耗散功率，并以
\(UI\text{、}I^2R\text{、}U^2/R\) 三种关系核对数值。

\end{minipage}\par

\par\noindent\begin{minipage}{\linewidth}
\textbf{解答}

\begin{center}
\includegraphics[width=\linewidth,height=.40\textheight,keepaspectratio]{\detokenize{assets/module-05/examples/M5-112.pdf}}
\par\smallskip{\small 图 M5-112：功率选择公式}
\end{center}
\end{minipage}\par

图 M5-112
的电流、电压都取被动参考方向，三种功率式必须算出同一个正值；其中
\(U^2/R\) 可直接用已知量求功率。

\(R=10\ \Omega\) 的电阻两端电压为 \(20\ \mathrm V\)，则
\(I=2\ \mathrm A\)，

\[
P=UI=40\ \mathrm W
=I^2R
=\frac{U^2}{R}.
\]

三种写法必须给出同一结果。

\subsubsection{4.
温度与自热边界}\label{m5:ux6e29ux5ea6ux4e0eux81eaux70edux8fb9ux754c}

在不太宽的温区内，金属电阻常近似为

\[
R(T)=R_0[1+\alpha(T-T_0)].
\]

若温度由外界固定，则对某一确定温度仍可使用线性关系
\(u=R(T)i\)。若焦耳热显著改变温度，\(R\)
会通过温度间接依赖电流和时间，电学方程必须与热平衡方程联立，此时二端口不再是常参数线性电阻。

\subsection{三、电容：电荷、端电压与电场能}\label{m5:ux4e09ux7535ux5bb9ux7535ux8377ux7aefux7535ux538bux4e0eux7535ux573aux80fd}

\subsubsection{1. 定义与方向}\label{m5:ux5b9aux4e49ux4e0eux65b9ux5411}

取参考电压正端极板上的自由电荷为 \(q\)。线性电容定义为

\[
\boxed{q=Cu},
\qquad C>0.
\]

由于参考电流流入正端，电荷守恒给出

\[
i=\frac{dq}{dt}.
\]

若 \(C\) 为常量，

\[
\boxed{i=C\frac{du}{dt}}.
\]

若极板运动或介电常数随时间改变，则

\[
\boxed{i=\frac{d(Cu)}{dt}
=C\frac{du}{dt}+u\frac{dC}{dt}}.
\]

后一项对应机械运动或介质变化引起的电荷重新分配，不能漏掉。

\par\addvspace{.4\baselineskip}
\noindent\begin{minipage}{\linewidth}

\textbf{例 M5-113：恒流充电}

理想电容 \(C=2.0\ \mu\mathrm F\) 的电流参考流入正端；从 \(t=0\)
起输入恒流 \(i=4.0\ \mu\mathrm A\)，初始电压记为 \(u_0\)。求电压变化率和
\(u_C(t)\)，说明初值在积分中的作用。

\end{minipage}\par

\par\noindent\begin{minipage}{\linewidth}
\textbf{解答}

\begin{center}
\includegraphics[width=\linewidth,height=.40\textheight,keepaspectratio]{\detokenize{assets/module-05/examples/M5-113.pdf}}
\par\smallskip{\small 图 M5-113：恒流充电}
\end{center}
\end{minipage}\par

图 M5-113 的恒流流入左侧正端，极板电荷线性增加；电压斜率为
\(i/C\)，初始电压则以积分常数保留。

全程沿用题设的任意初值 \(u_C(0)=u_0\)。由
\(du_C/dt=I/C\)，代入参数得到斜率 \(2.0\,\mathrm{V/s}\)。从0到 \(t\)
积分， \[
u_C(t)-u_0=\int_0^t2.0\,\mathrm{V/s}\,dt',
\qquad
\boxed{u_C(t)=u_0+(2.0\,\mathrm{V/s})t}.
\] 代入 \(t=0\)
回到题设初值；后续电压变化量为斜率乘经历时间，不能在未说明时另设初始电压。

\subsubsection{2.
平行板电容器}\label{m5:ux5e73ux884cux677fux7535ux5bb9ux5668}

设两平行导体板的正对面积为 \(S\)、间距为 \(d\)，板间填充介电常数为
\(\varepsilon\) 的线性均匀介质。忽略边缘效应，极板自由面电荷密度为
\(\sigma_f=Q/S\)。

由 \(\nabla\cdot\mathbf D=\rho_f\) 的高斯定律，板间

\[
D=\sigma_f,
\qquad
E=\frac{D}{\varepsilon}
=\frac{Q}{\varepsilon S}.
\]

端电压为

\[
U=Ed=\frac{Qd}{\varepsilon S}.
\]

所以

\[
\boxed{C=\frac QU=\frac{\varepsilon S}{d}}.
\]

\par\noindent\begin{minipage}{\linewidth}
\begin{center}
\includegraphics[width=\linewidth,height=.40\textheight,keepaspectratio]{\detokenize{assets/module-05/original/tikz-m5-topic49-v9-from46-v8-05.png}}
\par\smallskip{\small 图示：平行板电容由自由电荷、板间电场和端电压逐层推出}
\end{center}
\end{minipage}\par

该式要求 \(d\) 远小于极板横向尺度。若边缘场不可忽略，电容仍可定义为
\(C=Q/U\)，但不能再直接用 \(\varepsilon S/d\)。

\subsubsection{3.
电容储能与场能}\label{m5:ux7535ux5bb9ux50a8ux80fdux4e0eux573aux80fd}

从零电荷缓慢充电到 \(Q\)，外界所做的功为

\[
W_e=\int_0^Q u(q)\,dq
=\int_0^Q\frac qC\,dq.
\]

故

\[
\boxed{
W_e=\frac{Q^2}{2C}
=\frac12QU
=\frac12CU^2}.
\]

线性介质中的电场能也可写成

\[
\boxed{W_e=\frac12\int_V\mathbf E\cdot\mathbf D\,dV}.
\]

对理想平行板电容，代入 \(E=U/d\) 和体积 \(Sd\)，恰好得到
\(W_e=CU^2/2\)。

瞬时功率为

\[
p=ui=uC\frac{du}{dt}
=\frac d{dt}\left(\frac12Cu^2\right).
\]

因此理想电容的 \(p\)
可以为正也可以为负：充电时吸收能量，放电时归还能量。它是无源元件，却不是耗散元件。

\par\addvspace{.4\baselineskip}
\noindent\begin{minipage}{\linewidth}

\textbf{例 M5-114：符号判断}

理想电容采用电流流入电压正端的被动方向，某时段内 \(u_C>0\) 且
\(du_C/dt<0\)。判断电流和端口功率的符号，说明实际电流的流向及储能是在增加还是向外释放。

\end{minipage}\par

\par\noindent\begin{minipage}{\linewidth}
\textbf{解答}

\begin{center}
\includegraphics[width=\linewidth,height=.40\textheight,keepaspectratio]{\detokenize{assets/module-05/examples/M5-114.pdf}}
\par\smallskip{\small 图 M5-114：符号判断}
\end{center}
\end{minipage}\par

图 M5-114 的箭头只是正电流参考，计算得到 \(i<0\)
意味实际电流由正端流出，电容把储能输送给外电路。

若 \(u>0\) 但 \(du/dt<0\)，则 \(i=C\,du/dt<0\)，所以
\(p=ui<0\)。此时电流从参考电压正端流出，电容正在向外电路释放能量。

\subsubsection{4.
电容电压连续性的条件}\label{m5:ux7535ux5bb9ux7535ux538bux8fdeux7eedux6027ux7684ux6761ux4ef6}

在开关时刻 \(t_0\) 两侧积分 \(i=C\,du_C/dt\)：

\[
C[u_C(t_0^+)-u_C(t_0^-)]
=\int_{t_0^-}^{t_0^+}i(t)\,dt.
\]

若电流在 \(t_0\)
附近有界，或至少不含有限面积的狄拉克冲激，则当积分区间缩至零时右端为零，从而

\[
\boxed{u_C(t_0^+)=u_C(t_0^-)}.
\]

\par\noindent\begin{minipage}{\linewidth}
\begin{center}
\includegraphics[width=\linewidth,height=.40\textheight,keepaspectratio]{\detokenize{assets/module-05/theory/M5-R20-continuity-C.pdf}}
\par\smallskip{\small 图示：有限瞬时电流保证电容电压连续，电压跳变则要求冲激电流}
\end{center}
\end{minipage}\par

所以``电容电压不能突变''不是无条件公理。若理想电压源直接强迫电容电压跳变
\(\Delta U\)，理想模型要求

\[
i_{\text{冲}}=C\Delta U\,\delta(t-t_0).
\]

真实电路中的串联电阻、寄生电感、有限带宽和辐射会把冲激展宽并限制峰值。

\subsection{四、电感：磁链、端电压与磁场能}\label{m5:ux56dbux7535ux611fux78c1ux94feux7aefux7535ux538bux4e0eux78c1ux573aux80fd}

\subsubsection{1. 本构关系}\label{m5:ux672cux6784ux5173ux7cfb}

在线性、固定几何、无磁滞的单线圈模型中，

\[
\boxed{\Psi=Li}.
\]

在被动符号约定下，线圈端电压等于磁链变化率：

\[
\boxed{u=\frac{d\Psi}{dt}}.
\]

当 \(L\) 为常量，

\[
\boxed{u=L\frac{di}{dt}}.
\]

若线圈位置、磁路气隙或介质状态改变，则

\[
u=\frac{d(Li)}{dt}
=L\frac{di}{dt}+i\frac{dL}{dt}.
\]

注意，沿规定回路正向的自感电动势是
\(\mathcal E=-d\Psi/dt\)，而端电压在被动符号约定下为
\(u=d\Psi/dt\)。两者不应混写。

\subsubsection{2.
长直螺线管模型回顾}\label{m5:ux957fux76f4ux87baux7ebfux7ba1ux6a21ux578bux56deux987e}

长度为 \(\ell\)、总匝数为 \(N\)、截面积为 \(S\)
的长直螺线管内填充线性均匀介质 \(\mu\)。忽略边缘漏磁，

\[
B\simeq\mu\frac N\ell i.
\]

单匝磁通为 \(\Phi=BS\)，总磁链为 \(\Psi=N\Phi\)，故

\[
\boxed{L=\frac\Psi i=\frac{\mu N^2S}{\ell}}.
\]

\par\noindent\begin{minipage}{\linewidth}
\begin{center}
\includegraphics[width=\linewidth,height=.40\textheight,keepaspectratio]{\detokenize{assets/module-05/original/tikz-m5-topic49-v9-from46-v8-07.png}}
\par\smallskip{\small 图示：长直螺线管由磁场、单匝磁通和总磁链得到自感}
\end{center}
\end{minipage}\par

这一推导在专题四十七已有完整讨论，此处保留它是为了显示电路参数 \(L\)
如何从场与几何产生。真实线圈还含绕组电阻、漏磁、寄生电容和磁芯损耗，只有在相应影响可忽略的频率和幅值范围内，才可近似为单一理想电感。

\par\addvspace{.4\baselineskip}
\noindent\begin{minipage}{\linewidth}

\textbf{例 M5-115：恒压激励}

理想电感 \(L=0.50\ \mathrm H\) 采用被动参考方向，从 \(t=0\) 起施加恒压
\(u=2.0\ \mathrm V\)，初始电流为 \(i_0\)。求电流变化率和
\(i(t)\)，忽略电阻与磁饱和。

\end{minipage}\par

\par\noindent\begin{minipage}{\linewidth}
\textbf{解答}

\begin{center}
\includegraphics[width=\linewidth,height=.40\textheight,keepaspectratio]{\detokenize{assets/module-05/examples/M5-115.pdf}}
\par\smallskip{\small 图 M5-115：恒压激励}
\end{center}
\end{minipage}\par

图 M5-115
的恒定正端电压使流入正端的电流线性增加；初始电流不可在求导过程中丢掉。

全程沿用题设的任意初值 \(i_L(0)=i_0\)。由
\(di_L/dt=U/L\)，代入参数得到斜率 \(4.0\,\mathrm{A/s}\)。从0到 \(t\)
积分， \[
i_L(t)-i_0=\int_0^t4.0\,\mathrm{A/s}\,dt',
\qquad
\boxed{i_L(t)=i_0+(4.0\,\mathrm{A/s})t}.
\] 代入 \(t=0\)
回到题设初值；电流线性建立是理想电感在恒定端电压下的响应。

\subsubsection{3. 电感储能}\label{m5:ux7535ux611fux50a8ux80fd}

瞬时功率为

\[
p=ui=Li\frac{di}{dt}.
\]

从零电流建立到 \(I\)，

\[
W_m=\int_0^I Li\,di,
\]

所以

\[
\boxed{W_m=\frac12LI^2}.
\]

在线性磁介质中，

\[
\boxed{W_m=\frac12\int_V\mathbf B\cdot\mathbf H\,dV}.
\]

与电容相同，理想电感的瞬时功率可以为正或负。\(i\) 与 \(di/dt\)
同号时吸收能量，异号时释放能量；理想电感本身不把能量不可逆地变成热。

\subsubsection{4.
电感电流连续性的条件}\label{m5:ux7535ux611fux7535ux6d41ux8fdeux7eedux6027ux7684ux6761ux4ef6}

在开关时刻两侧积分 \(u=L\,di_L/dt\)：

\[
L[i_L(t_0^+)-i_L(t_0^-)]
=\int_{t_0^-}^{t_0^+}u(t)\,dt.
\]

若电压不含有限面积的冲激，则

\[
\boxed{i_L(t_0^+)=i_L(t_0^-)}.
\]

\par\noindent\begin{minipage}{\linewidth}
\begin{center}
\includegraphics[width=\linewidth,height=.40\textheight,keepaspectratio]{\detokenize{assets/module-05/theory/M5-R20-continuity-L.pdf}}
\par\smallskip{\small 图示：有限瞬时电压保证电感电流连续，电流跳变则要求冲激电压}
\end{center}
\end{minipage}\par

若理想电流源强迫电感电流跳变 \(\Delta I\)，理想模型要求

\[
u_{\text{冲}}=L\Delta I\,\delta(t-t_0).
\]

真实线圈中，寄生电容、绕组绝缘和电磁辐射会限制电压峰值；贸然切断有电流的电感，可能产生高压电弧。

\subsection{五、开关瞬间、直流稳态与网络能量}\label{m5:ux4e94ux5f00ux5173ux77acux95f4ux76f4ux6d41ux7a33ux6001ux4e0eux7f51ux7edcux80fdux91cf}

\subsubsection{1.
先继承状态，再求变化率}\label{m5:ux5148ux7ee7ux627fux72b6ux6001ux518dux6c42ux53d8ux5316ux7387}

在没有冲激激励的普通开关题中，

\[
\boxed{u_C(t_0^+)=u_C(t_0^-)},
\qquad
\boxed{i_L(t_0^+)=i_L(t_0^-)}.
\]

这些是开关瞬间的初始约束。随后应根据切换后的拓扑列基尔霍夫方程，求
\(du_C/dt\)、\(di_L/dt\) 及其他未知量。

\subsubsection{2.
直流稳态等效}\label{m5:ux76f4ux6d41ux7a33ux6001ux7b49ux6548}

只有经过足够长时间并达到直流稳态后，才有

\[
\frac{du_C}{dt}=0,
\qquad
\frac{di_L}{dt}=0.
\]

因此

\[
\boxed{i_C=0},
\qquad
\boxed{u_L=0}.
\]

理想电容在直流稳态相当于开路，理想电感在直流稳态相当于短路。

\par\noindent\begin{minipage}{\linewidth}
\begin{center}
\includegraphics[width=\linewidth,height=.40\textheight,keepaspectratio]{\detokenize{assets/module-05/original/tikz-m5-topic49-v9-from46-v8-09.png}}
\par\smallskip{\small 图示：开关瞬间继承电容电压和电感电流，直流稳态再作开短路等效}
\end{center}
\end{minipage}\par

不能在 \(t_0^+\)
直接把所有电容替换成开路、所有电感替换成短路。例如，一个初始未充电的电容在接通电源瞬间可以有非零充电电流；一个带初始电流的电感在切换后仍会维持该电流。

\par\addvspace{.4\baselineskip}
\noindent\begin{minipage}{\linewidth}

\textbf{例 M5-116：初始状态优先}

理想电容在 \(t=0^-\) 的电压为
\(5\ \mathrm V\)，切换开关时没有冲激电流，且电容值不变。求
\(u_C(0^+)\)，说明能否用直流稳态的``电容开路''模型把这个初始电压直接置零。

\end{minipage}\par

\par\noindent\begin{minipage}{\linewidth}
\textbf{解答}

\begin{center}
\includegraphics[width=\linewidth,height=.40\textheight,keepaspectratio]{\detokenize{assets/module-05/examples/M5-116.pdf}}
\par\smallskip{\small 图 M5-116：初始状态优先}
\end{center}
\end{minipage}\par

图 M5-116
两极板的电荷不能由有限电流在零时间内突然改变，所以开关刚切换时先继承原电压，之后才由新回路决定变化率。

电容在 \(t=0^-\) 的电压为
\(5\ \mathrm V\)，且没有冲激电流，则无论开关如何切换，都先有
\(u_C(0^+)=5\ \mathrm V\)，不能因为``电容相当于开路''而把初始电压写成零。

\subsubsection{3.
网络功率与能量平衡}\label{m5:ux7f51ux7edcux529fux7387ux4e0eux80fdux91cfux5e73ux8861}

对由理想 \(R,C,L\) 组成的网络，把所有无源元件都按被动符号约定记功率，则

\[
\sum_{R,C,L} p
=\sum_R R_ki_k^2
+\frac d{dt}\left(
\sum_C\frac12C_ku_k^2
+\sum_L\frac12L_ki_k^2
\right).
\]

若再把电源按同一套端口参考方向计入，基尔霍夫定律保证整个网络所有支路瞬时功率的代数和为零。于是电源提供的功率，一部分成为电阻热，一部分改变电场能和磁场能。这是检查微分方程符号最可靠的方法之一。

\subsection{六、串并联等效及其条件}\label{m5:ux516dux4e32ux5e76ux8054ux7b49ux6548ux53caux5176ux6761ux4ef6}

\subsubsection{1. 电阻}\label{m5:ux7535ux963b}

串联电阻具有相同支路电流，端电压相加：

\[
\boxed{R_s=\sum_jR_j}.
\]

并联电阻具有相同端电压，支路电流相加：

\[
\boxed{\frac1{R_p}=\sum_j\frac1{R_j}}.
\]

\subsubsection{2. 电容}\label{m5:ux7535ux5bb9}

并联电容端电压相同，总电荷为各支路正端电荷之和：

\[
\boxed{C_p=\sum_jC_j}.
\]

串联电容中，若中间节点没有外接支路，且其初始净自由电荷为零，则各电容电荷大小相等，故

\[
\boxed{\frac1{C_s}=\sum_j\frac1{C_j}}.
\]

若中间节点带有非零初始净电荷，端口的增量电容仍满足上式，但各电容的绝对电压还含由初始电荷决定的偏置，不能只靠等效电容恢复内部状态。

\subsubsection{3. 无互感电感}\label{m5:ux65e0ux4e92ux611fux7535ux611f}

若各线圈之间互感可以忽略，串联电感有

\[
\boxed{L_s=\sum_jL_j}.
\]

并联电感在同端电压下满足

\[
\boxed{\frac1{L_p}=\sum_j\frac1{L_j}}.
\]

并联支路的初始电流分配还必须满足基尔霍夫电流定律。等效电感只决定端口总电流的变化，不唯一决定各支路中原有的恒定环流。若线圈之间有互感，上述两式均应改用含互感项的电感矩阵，不能直接套用。

\par\noindent\begin{minipage}{\linewidth}
\begin{center}
\includegraphics[width=\linewidth,height=.40\textheight,keepaspectratio]{\detokenize{assets/module-05/theory/tikz-m5-topic49-v9-from46-v8-10.pdf}}
\par\smallskip{\small 图示：串并联规则来自端口约束}
\end{center}
\end{minipage}\par

\par\addvspace{.4\baselineskip}
\noindent\begin{minipage}{\linewidth}

\textbf{例 M5-117：两个相同元件}

分别把两个相同电阻 \(R\)、两个相同理想电容
\(C\)、两个相同且互不耦合的理想电感 \(L\)
作串联与并联。求六个等效参数；电容串联中间节点初始净电荷为零，电感连接采用相容的零初始状态。

\end{minipage}\par

\par\noindent\begin{minipage}{\linewidth}
\textbf{解答}

\begin{center}
\includegraphics[width=\linewidth,height=.40\textheight,keepaspectratio]{\detokenize{assets/module-05/examples/M5-117.pdf}}
\par\smallskip{\small 图 M5-117：两个相同元件}
\end{center}
\end{minipage}\par

图 M5-117
画出两个相同电容串联的中间节点；零初始节点净电荷使两个电容带等量电荷，端电压相加，因此电容串联规则与电阻不同。

把六种连接都画在图 M5-117a
中。左列只有一条支路，电流相同，两个元件的电压相加；右列两个元件跨接同一对端点
\(A,B\)，电压相同，各支路电流相加。

\begin{center}
\includegraphics[width=\linewidth,height=.40\textheight,keepaspectratio]{\detokenize{assets/module-05/process/M5-117a.pdf}}\par\smallskip{\small 图 M5-117a：R、C、L 的六种连接与等效参数}
\end{center}

上排电阻用 \(u=Ri\) 得串联 \(R_{\rm eq}=2R\)、并联
\(R_{\rm eq}=R/2\)。中排电容串联时，图 M5-117
的中间节点净电荷为零，所以两个电容电荷等量，\(u=q/C+q/C\) 给出
\(C_{\rm eq}=C/2\)；并联则两支路电荷相加，\(q=Cu+Cu\) 给出
\(C_{\rm eq}=2C\)。下排互不耦合的电感由 \(u=L\,di/dt\) 得串联
\(L_{\rm eq}=2L\)、并联 \(L_{\rm eq}=L/2\)。电容规则来自 \(q=Cu\)
和连接条件。

\subsection{七、理想模型与集总参数近似的边界}\label{m5:ux4e03ux7406ux60f3ux6a21ux578bux4e0eux96c6ux603bux53c2ux6570ux8fd1ux4f3cux7684ux8fb9ux754c}

\subsubsection{1.
集总参数近似}\label{m5:ux96c6ux603bux53c2ux6570ux8fd1ux4f3c}

把一个实际器件表示成单一 \(R\)、\(C\) 或 \(L\)，隐含以下条件：

\begin{enumerate}
\def\labelenumi{\arabic{enumi}.}
\tightlist
\item
  器件和连接导线的特征尺度远小于信号在一个显著变化时间内传播的距离，节点可近似视为等势体。
\item
  电磁场能或耗散主要集中在预期区域，寄生参数相对较小。
\item
  材料处于给定本构关系的线性范围，温度、磁饱和和介质损耗等影响可以忽略或另行建模。
\item
  辐射损耗可以忽略，电路变量可用有限个状态量描述。
\end{enumerate}

对角频率为 \(\omega\) 的信号，常用的量级判据是器件尺度 \(a\) 满足

\[
a\ll\lambda=\frac{2\pi v}{\omega},
\]

其中 \(v\)
是电磁扰动在相应介质或传输结构中的传播速度。这里只是必要的量级判断；高品质因数、强耦合或陡峭脉冲还可能提出更严格要求。

\subsubsection{2.
实际元件的寄生参数}\label{m5:ux5b9eux9645ux5143ux4ef6ux7684ux5bc4ux751fux53c2ux6570}

\begin{itemize}
\tightlist
\item
  实际电阻还含引线电感和分布电容。
\item
  实际电容存在漏电电阻、等效串联电阻和引线电感。
\item
  实际电感存在绕组电阻、匝间电容、漏磁和磁芯损耗。
\end{itemize}

随着频率升高，一个器件可能跨过自谐振频率而表现出与低频名称相反的电抗特征。因此``这是电容''或``这是电感''只在指定频率范围内有意义。

\subsubsection{3.
理想源与理想储能元件的冲突}\label{m5:ux7406ux60f3ux6e90ux4e0eux7406ux60f3ux50a8ux80fdux5143ux4ef6ux7684ux51b2ux7a81}

不同初始电压的两个理想电容若被零电阻导线瞬间并联，理想模型会要求冲激电流；只用初、末电场能比较还会发现一部分能量``消失''。这部分能量在真实电路中转化为电阻热、电磁辐射或寄生电感中的暂态能量。若把所有这些机制都删掉，模型本身不足以描述切换过程。

同理，有电流的理想电感若被理想开关瞬间断开，会要求无限高的冲激电压。遇到这类矛盾，应主动补回最重要的寄生参数，而不是把无穷大当作真实可测量值。

\subsection{八、练习与解答}\label{m5:ux516bux7ec3ux4e60ux4e0eux89e3ux7b54-2}

\par\addvspace{.4\baselineskip}
\noindent\begin{minipage}{\linewidth}

\subsubsection{\texorpdfstring{作业一：理想 \(LC\)
电路的自由振荡}{作业一：理想 LC 电路的自由振荡}}\label{m5:ux4f5cux4e1aux4e00ux7406ux60f3-lc-ux7535ux8defux7684ux81eaux7531ux632fux8361}

\textbf{题号：M5-118}

\textbf{层级：} 基础　　\textbf{方法：} 电路暂态、电磁振荡

\(L\text{、}C>0\)，初始电感电流为零；正电流沿电容初始正极板的放电方向定义，故
\(i=-\dot q\)。忽略电阻、辐射和互感。

一个初始电流为零的理想电感 \(L\) 与一个初始电压为 \(U_0\) 的理想电容
\(C\) 串联成闭合回路。闭合开关后，求电容正极板电荷 \(q(t)\)
和放电方向电流 \(i(t)\)，并证明总能量守恒。

\end{minipage}\par

\par\noindent\begin{minipage}{\linewidth}
\textbf{解答}

\begin{center}
\includegraphics[width=\linewidth,height=.40\textheight,keepaspectratio]{\detokenize{assets/module-05/examples/M5-118.pdf}}
\par\smallskip{\small 图 M5-118：理想 \(LC\) 电路的自由振荡}
\end{center}
\end{minipage}\par

图 M5-118 的 \(q\) 是电容上极板电荷，正电流沿上边向右放电，所以
\(i=-\dot q\)。两个储能元件交替交换能量。

规定 \(q(t)\) 为初始带正电极板上的电荷，\(i\) 沿该极板放电方向，因此

\[
i=-\frac{dq}{dt}.
\]

沿电流方向列回路方程，电容电压驱动电感电流变化：

\[
L\frac{di}{dt}=\frac qC.
\]

代入 \(i=-dq/dt\)：

\[
\frac{d^2q}{dt^2}+\frac1{LC}q=0.
\]

令

\[
\omega_0=\frac1{\sqrt{LC}}.
\]

\par\noindent\begin{minipage}{\linewidth}
\begin{center}
\includegraphics[width=\linewidth,height=.40\textheight,keepaspectratio]{\detokenize{assets/module-05/process/M5-118a.pdf}}
\par\smallskip{\small 图 M5-118a：电场能和磁场能交替交换}
\end{center}
\end{minipage}\par

图 M5-118a
中一条储能曲线过零时另一条达到最大，两条非负曲线的和始终落在虚线高度。

初始条件为

\[
q(0)=CU_0,
\qquad
i(0)=0
\quad\Longrightarrow\quad
\dot q(0)=0.
\]

所以

\[
\boxed{q(t)=CU_0\cos\omega_0t},
\]

\[
\boxed{i(t)=CU_0\omega_0\sin\omega_0t
=U_0\sqrt{\frac CL}\sin\omega_0t}.
\]

总能量为

\[
W=\frac{q^2}{2C}+\frac12Li^2.
\]

代入解：

\[
W=\frac12CU_0^2\cos^2\omega_0t
+\frac12CU_0^2\sin^2\omega_0t.
\]

故

\[
\boxed{W=\frac12CU_0^2=\text{常量}}.
\]

振荡周期为

\[
\boxed{T=2\pi\sqrt{LC}}.
\]

\par\addvspace{.4\baselineskip}
\noindent\begin{minipage}{\linewidth}

\subsubsection{作业二：经典弛豫模型的暂态电流}\label{m5:ux4f5cux4e1aux4e8cux7ecfux5178ux5f1bux8c6bux6a21ux578bux7684ux6682ux6001ux7535ux6d41}

\textbf{题号：M5-119}

\textbf{层级：} 基础　　\textbf{方法：} 电路暂态、弛豫方程

\(e>0\) 为电子电荷绝对值，\(m\) 为电子质量，\(n\)
为恒定载流子数密度，\(\tau>0\) 为弛豫时间。

金属中的平均漂移速度满足

\[
m\dot{\mathbf v}_d=-e\mathbf E-\frac m\tau\mathbf v_d.
\]

\(t=0\) 前外场为零且平均漂移速度为零；\(t=0\) 后突然施加恒定均匀电场
\(\mathbf E_0\)。求
\(\mathbf v_d(t)\)、\(\mathbf J(t)\)，并给出稳态电导率。

\end{minipage}\par

\par\noindent\begin{minipage}{\linewidth}
\textbf{解答}

\begin{center}
\includegraphics[width=\linewidth,height=.40\textheight,keepaspectratio]{\detokenize{assets/module-05/examples/M5-119.pdf}}
\par\smallskip{\small 图 M5-119：经典弛豫模型的暂态电流}
\end{center}
\end{minipage}\par

图 M5-119 的电流从零单调趋向稳态虚线，时间尺度是
\(\tau\)；电子漂移方向与电场相反，乘负电荷后电流方向与电场相同。

方程改写为

\[
\dot{\mathbf v}_d+\frac1\tau\mathbf v_d
=-\frac e m\mathbf E_0.
\]

齐次解为 \(\mathbf A e^{-t/\tau}\)，常数特解为
\(-e\tau\mathbf E_0/m\)。利用 \(\mathbf v_d(0)=0\)，得

\[
\boxed{
\mathbf v_d(t)
=-\frac{e\tau}{m}\mathbf E_0(1-e^{-t/\tau})}.
\]

因此

\[
\boxed{
\mathbf J(t)
=\frac{ne^2\tau}{m}\mathbf E_0(1-e^{-t/\tau})}.
\]

\(t\gg\tau\) 时，

\[
\mathbf J\to\sigma\mathbf E_0,
\qquad
\boxed{\sigma=\frac{ne^2\tau}{m}}.
\]

电流以时间常数 \(\tau\) 指数建立。

\par\addvspace{.4\baselineskip}
\noindent\begin{minipage}{\linewidth}

\subsubsection{作业三：电容电压为何通常连续}\label{m5:ux4f5cux4e1aux4e09ux7535ux5bb9ux7535ux538bux4e3aux4f55ux901aux5e38ux8fdeux7eed}

\textbf{题号：M5-120}

\textbf{层级：} 基础　　\textbf{方法：} 电路暂态、状态连续性

电容值恒定，\(t=0\) 没有冲激电流；电压和电流采用被动参考方向。

一个 \(C=4.0\ \mu\mathrm F\) 的电容在 \(t=0^-\) 时电压为
\(3.0\ \mathrm V\)。\(t=0\) 后流入正端的电流为

\[
i(t)=2.0\ \mathrm{mA}\,e^{-t/(2.0\ \mathrm{ms})},
\qquad t>0.
\]

求 \(u_C(0^+)\) 和 \(u_C(\infty)\)。

\end{minipage}\par

\par\noindent\begin{minipage}{\linewidth}
\textbf{解答}

\begin{center}
\includegraphics[width=\linewidth,height=.40\textheight,keepaspectratio]{\detokenize{assets/module-05/examples/M5-120.pdf}}
\par\smallskip{\small 图 M5-120：电容电压为何通常连续}
\end{center}
\end{minipage}\par

图 M5-120 的电流流入正极板；初始瞬间电压连续，随后总输入电荷 \(i_0\tau\)
使电压增加 \(1.0\ \mathrm V\)。

给定电流在 \(t=0\) 有限且不含冲激，所以

\[
\boxed{u_C(0^+)=u_C(0^-)=3.0\ \mathrm V}.
\]

电容获得的总电荷为

\[
\Delta q=\int_0^{\infty} i(t)\,dt
=(2.0\times10^{-3})(2.0\times10^{-3})
=4.0\times10^{-6}\ \mathrm C.
\]

\par\noindent\begin{minipage}{\linewidth}
\begin{center}
\includegraphics[width=\linewidth,height=.40\textheight,keepaspectratio]{\detokenize{assets/module-05/process/M5-120a.pdf}}
\par\smallskip{\small 图 M5-120a：有限总电荷使电容电压从3V升至4V}
\end{center}
\end{minipage}\par

图 M5-120a 的纵坐标扣去了初始
\(3\ \mathrm V\)；输入电流积分有限，故曲线只再增加
\(1\ \mathrm V\)，而不是回到零。

故电压增量为

\[
\Delta u=\frac{\Delta q}{C}
=\frac{4.0\ \mu\mathrm C}{4.0\ \mu\mathrm F}
=1.0\ \mathrm V.
\]

所以

\[
\boxed{u_C(\infty)=4.0\ \mathrm V}.
\]

\par\addvspace{.4\baselineskip}
\noindent\begin{minipage}{\linewidth}

\subsubsection{作业四：断开电感所需的电压冲量}\label{m5:ux4f5cux4e1aux56dbux65adux5f00ux7535ux611fux6240ux9700ux7684ux7535ux538bux51b2ux91cf}

\textbf{题号：M5-121}

\textbf{层级：} 基础　　\textbf{方法：} 电路暂态、电感电流连续性

端电压以原正电流流入端为正端，电感值恒定，下降过程采用线性近似；给出的负电压表示其实际极性反向。

一个 \(L=0.20\ \mathrm H\) 的理想电感原有电流
\(5.0\ \mathrm A\)。若某理想控制装置试图在 \(10\ \mu\mathrm s\)
内使电流近似匀速降为零，求所需平均端电压；并求电压的时间积分。

\end{minipage}\par

\par\noindent\begin{minipage}{\linewidth}
\textbf{解答}

\begin{center}
\includegraphics[width=\linewidth,height=.40\textheight,keepaspectratio]{\detokenize{assets/module-05/examples/M5-121.pdf}}
\par\smallskip{\small 图 M5-121：断开电感所需的电压冲量}
\end{center}
\end{minipage}\par

图 M5-121
的正电压与正电流按被动方向规定；电流下降意味着所需端电压为负。时间越短，平均反向电压越大。

近似匀速变化时

\[
\frac{di}{dt}
=\frac{0-5.0}{10\times10^{-6}}
=-5.0\times10^5\ \mathrm{A/s}.
\]

因此

\[
\boxed{u=L\frac{di}{dt}
=-1.0\times10^5\ \mathrm V}.
\]

\par\noindent\begin{minipage}{\linewidth}
\begin{center}
\includegraphics[width=\linewidth,height=.40\textheight,keepaspectratio]{\detokenize{assets/module-05/process/M5-121a.pdf}}
\par\smallskip{\small 图 M5-121a：电流下降所需的负电压冲量}
\end{center}
\end{minipage}\par

图 M5-121a
的阴影在时间轴下方，表示实际电压反于图中正端参考；矩形面积等于电感电流变化要求的负磁链变化。

负号表示实际端电压方向与原被动参考方向相反。无论具体波形如何，只要电流从
\(5.0\ \mathrm A\) 变到零，电压冲量都满足

\[
\int u\,dt=L\Delta i
=(0.20)(-5.0).
\]

故

\[
\boxed{\int u\,dt=-1.0\ \mathrm{V\,s}}.
\]

缩短切断时间会提高电压峰值，但不会改变所需的总电压冲量。

\par\addvspace{.4\baselineskip}
\noindent\begin{minipage}{\linewidth}

\subsubsection{作业五：两电容并联后的能量去向}\label{m5:ux4f5cux4e1aux4e94ux4e24ux7535ux5bb9ux5e76ux8054ux540eux7684ux80fdux91cfux53bbux5411}

\textbf{题号：M5-122}

\textbf{层级：} 基础　　\textbf{方法：} 能量法、电路暂态

忽略外接导线对最终电荷守恒的影响，待暂态完全结束后比较储能；有限电阻模型下能量差最终计入焦耳热。

两个相同电容均为 \(C\)。电容 1 初始电压为 \(U_0\)，电容 2
初始不带电。现用含很小但非零电阻的导线将两者同极性并联，求最终共同电压、初始与最终电场能，并解释能量差。

\end{minipage}\par

\par\noindent\begin{minipage}{\linewidth}
\textbf{解答}

\begin{center}
\includegraphics[width=\linewidth,height=.40\textheight,keepaspectratio]{\detokenize{assets/module-05/examples/M5-122.pdf}}
\par\smallskip{\small 图 M5-122：两电容并联后的能量去向}
\end{center}
\end{minipage}\par

图 M5-122 中两电容同极性并联，终态电压相等；用总电荷守恒先求
\(U_f\)，再分别计算两电容的能量，差值属于暂态耗散。

总自由电荷守恒：

\[
CU_0+0=(C+C)U_f.
\]

因此

\[
\boxed{U_f=\frac{U_0}{2}}.
\]

初始能量为

\[
W_i=\frac12CU_0^2.
\]

\par\noindent\begin{minipage}{\linewidth}
\begin{center}
\includegraphics[width=\linewidth,height=.40\textheight,keepaspectratio]{\detokenize{assets/module-05/process/M5-122a.pdf}}
\par\smallskip{\small 图 M5-122a：两电容的能量账目}
\end{center}
\end{minipage}\par

图 M5-122a
的蓝柱等于绿柱与橙柱之和；连接过程的有限电阻把电压差驱动的暂态电流能量耗散掉。

最终能量为

\[
W_f=2\cdot\frac12C\left(\frac{U_0}{2}\right)^2
=\frac14CU_0^2.
\]

所以

\[
\boxed{W_i-W_f=\frac14CU_0^2}.
\]

这部分能量在非零连接电阻中转化为焦耳热；若把电阻取得很小，电流峰值和充放电时间常数会改变，但在只含电阻作为耗散机制的模型中，总焦耳热仍为上述能量差。若把连接电阻、寄生电感和辐射全部删掉，理想模型无法完整描述切换暂态，不能据此断言能量凭空消失。

\par\addvspace{.4\baselineskip}
\noindent\begin{minipage}{\linewidth}

\subsubsection{作业六：串联电容的初始电荷条件}\label{m5:ux4f5cux4e1aux516dux4e32ux8054ux7535ux5bb9ux7684ux521dux59cbux7535ux8377ux6761ux4ef6}

\textbf{题号：M5-123}

\textbf{层级：} 基础　　\textbf{方法：} 电路暂态、初始条件

定义左电容左极板电荷为 \(q_1\)、右电容左极板电荷为
\(q_2\)；中间节点包含左电容右极板与右电容左极板。

电容 \(C_1,C_2\)
串联，二者之间的中间节点不与其他支路相连。证明中间节点净自由电荷守恒；若其初始净电荷为零，求端口等效电容。若初始净电荷不为零，等效电容公式是否失效？

\end{minipage}\par

\par\noindent\begin{minipage}{\linewidth}
\textbf{解答}

\begin{center}
\includegraphics[width=\linewidth,height=.40\textheight,keepaspectratio]{\detokenize{assets/module-05/examples/M5-123.pdf}}
\par\smallskip{\small 图 M5-123：串联电容的初始电荷条件}
\end{center}
\end{minipage}\par

图 M5-123 的 \(M\)
是孤立中间节点，电荷只能在其相连的两极板间重新分布；非零初始 \(Q_M\)
带来端口偏置，不改变微分电容。

中间节点没有导电通路与外界交换自由电荷，所以其净自由电荷是常量。若初始净电荷为零，则两只电容面向中间节点的极板电荷代数和始终为零，故两电容的串联电荷大小相等，记为
\(q\)。

总端电压为

\[
u=\frac q{C_1}+\frac q{C_2}.
\]

由 \(q=C_su\) 得

\[
\boxed{
\frac1{C_s}=\frac1{C_1}+\frac1{C_2}},
\]

即

\[
\boxed{C_s=\frac{C_1C_2}{C_1+C_2}}.
\]

若中间节点初始净电荷不为零，各电容绝对电荷不再大小相等，内部电压会带有由初始电荷决定的偏置。但在初始电荷保持不变时，对端口小变化仍有

\[
dq=C_s\,du,
\]

所以同一 \(C_s\)
仍是端口的增量电容。失效的是``只凭等效电容即可确定各电容绝对电压''的说法，而不是端口增量关系本身。

\par\addvspace{.4\baselineskip}
\noindent\begin{minipage}{\linewidth}

\subsubsection{作业七：由能量验证长直螺线管的自感}\label{m5:ux4f5cux4e1aux4e03ux7531ux80fdux91cfux9a8cux8bc1ux957fux76f4ux87baux7ebfux7ba1ux7684ux81eaux611f}

\textbf{题号：M5-124}

\textbf{层级：} 基础　　\textbf{方法：} 能量法、自感

磁导率 \(\mu\) 为线性常量，无磁滞；场能积分体积取管内 \(S\ell\)。

长度为 \(\ell\)、总匝数为 \(N\)、截面积为 \(S\)
的长直螺线管内填充磁导率为 \(\mu\)
的线性介质。忽略边缘效应，分别由磁链和场能方法求其自感。

\end{minipage}\par

\par\noindent\begin{minipage}{\linewidth}
\textbf{解答}

\begin{center}
\includegraphics[width=\linewidth,height=.40\textheight,keepaspectratio]{\detokenize{assets/module-05/examples/M5-124.pdf}}
\par\smallskip{\small 图 M5-124：由能量验证长直螺线管的自感}
\end{center}
\end{minipage}\par

图 M5-124 的场能集中在管内体积
\(S\ell\)，与磁链法中的同一均匀场对应。两条计算路线相互检验匝数平方和长度因子。

内部磁场为

\[
B=\mu\frac N\ell I.
\]

由磁链法，

\[
\Psi=NBS
=\frac{\mu N^2S}{\ell}I,
\]

所以

\[
\boxed{L=\frac{\mu N^2S}{\ell}}.
\]

由场能法，磁场能密度为 \(B^2/(2\mu)\)，有效体积为 \(S\ell\)，故

\[
W_m=\frac{B^2}{2\mu}S\ell
=\frac{\mu N^2S}{2\ell}I^2.
\]

与 \(W_m=LI^2/2\) 比较，再次得到

\[
\boxed{L=\frac{\mu N^2S}{\ell}}.
\]

两种方法一致，说明集总参数中的储能与场能积分完全对应。

\section{专题五十：电感计算、可变电感、磁能与磁共能（基础必讲与拔高选讲）}\label{m5:ux4e13ux9898ux4e94ux5341ux7535ux611fux8ba1ux7b97ux53efux53d8ux7535ux611fux78c1ux80fdux4e0eux78c1ux5171ux80fdux57faux7840ux5fc5ux8bb2ux4e0eux62d4ux9ad8ux9009ux8bb2}

\begin{quote}
阅读准备：广义坐标只是一种统一记号 用 \(x\)
表示气隙、位移或转角时，\(x\) 就是广义坐标；与它对应的广义力满足机械功
\(\delta W_m=F_x\,\delta x\)。对磁能或磁共能求偏导时，要在下标中明确保持磁链还是保持电流。
\end{quote}

\par\addvspace{.4\baselineskip}
\noindent\begin{minipage}{\linewidth}

\textbf{例 M5-125：可变电感的固定电流力}

线性电磁装置的电感为 \(L(x)>0\)，理想电流源把电流固定为 \(I\)，且
\(L^{\prime}(x)>0\)。在准静态、无磁滞和无机械耗散条件下，用磁共能求沿
\(x\) 的电磁力及其方向。

\end{minipage}\par

\par\noindent\begin{minipage}{\linewidth}
\textbf{解答}

\begin{center}
\includegraphics[width=\linewidth,height=.40\textheight,keepaspectratio]{\detokenize{assets/module-05/examples/M5-125.pdf}}
\par\smallskip{\small 图 M5-125：可变电感的固定电流力}
\end{center}
\end{minipage}\par

图 M5-125 的绿色 \(+x\)
是插入方向；在电流固定时，磁共能随插入增加，电磁力沿插入方向做功。

若线性电感 \(L(x)\) 随位置增加，在固定电流下磁共能 \(W'_m=L(x)i^2/2\)
增大，故 \(F_x=(\partial W'_m/\partial x)_i=i^2L'(x)/2\)。

当线圈位置、气隙、磁芯插入深度或回路形状发生变化时，电感不再是常数。此时感应电压、机械力和能量交换必须由磁链的全微分统一处理。只写
\(v=L\,di/dt\) 会漏掉几何变化产生的电压。

\subsection{一、可变电感的端口方程}\label{m5:ux4e00ux53efux53d8ux7535ux611fux7684ux7aefux53e3ux65b9ux7a0b}

设机械广义坐标为 \(x\)，在线性磁路中

\[
\Psi=L(x)i.
\]

按被动符号约定，电感端电压为

\[
v=\frac{d\Psi}{dt},
\]

故

\[
\boxed{
v=L(x)\frac{di}{dt}
+i\frac{dL}{dx}\frac{dx}{dt}}.
\]

第一项由电流变化产生，第二项由几何变化产生。若采用感应电动势记号，则
\(\mathcal E_{\rm ind}=-d\Psi/dt\)，符号与被动端电压相反。

\par\addvspace{.4\baselineskip}
\noindent\begin{minipage}{\linewidth}

\textbf{例 M5-126：电流恒定但仍有端电压}

理想电流源维持线圈电流 \(I_0\) 不变，机械运动使电感从 \(L(x)\)
随位置改变。采用被动端口方向，已知运动速度
\(\dot x\)，求端电压，解释为何 \(\dot i=0\) 时它仍可以非零。

\end{minipage}\par

\par\noindent\begin{minipage}{\linewidth}
\textbf{解答}

\begin{center}
\includegraphics[width=\linewidth,height=.40\textheight,keepaspectratio]{\detokenize{assets/module-05/examples/M5-126.pdf}}
\par\smallskip{\small 图 M5-126：电流恒定但仍有端电压}
\end{center}
\end{minipage}\par

图 M5-126
中铁芯运动改变电感和磁链，即使绕组电流恒定，电源仍须提供端电压，以补偿磁链的时间变化。

若理想电流源维持 \(i=I_0\)，而铁芯运动使 \(L\) 增大，则

\[
v=I_0\frac{dL}{dx}\dot x\ne0.
\]

电源仍需交换能量，不能因 \(di/dt=0\) 就断言电感电压为零。

\subsection{二、磁能、磁共能与广义力}\label{m5:ux4e8cux78c1ux80fdux78c1ux5171ux80fdux4e0eux5e7fux4e49ux529b}

对一般单端口磁系统，磁场能定义为

\[
W_m(\Psi,x)=\int_0^{\Psi} i(\Psi',x)\,d\Psi',
\]

磁共能定义为

\[
W_m'(i,x)=\int_0^i\Psi(i',x)\,di'.
\]

二者满足

\[
\boxed{W_m+W_m'=i\Psi}.
\]

在准静态、无耗散条件下，广义力可写为

\[
\boxed{
F_x=-\left.\frac{\partial W_m}{\partial x}\right|_{\Psi}
=\left.\frac{\partial W_m'}{\partial x}\right|_i}.
\]

线性系统中 \(\Psi=L(x)i\)，因此

\[
W_m=W_m'=\frac12L(x)i^2
\]

只是数值相等，所用独立变量不同。固定电流时

\[
\boxed{F_x=\frac12i^2\frac{dL}{dx}}.
\]

所以系统趋向电感增大的位置。

\par\noindent\begin{minipage}{\linewidth}
\begin{center}
\includegraphics[width=\linewidth,height=.40\textheight,keepaspectratio]{\detokenize{assets/module-05/original/tikz-m5-topic50-v9-from47-v8-02.png}}
\par\smallskip{\small 图示：磁能与磁共能在磁链电流图上的几何意义}
\end{center}
\end{minipage}\par

\subsection{三、功率守恒给出的力公式}\label{m5:ux4e09ux529fux7387ux5b88ux6052ux7ed9ux51faux7684ux529bux516cux5f0f}

端口输入功率为

\[
p_e=vi=i\frac{d(Li)}{dt}
=Li\dot i+i^2\frac{dL}{dx}\dot x.
\]

磁场能变化率为

\[
\frac{dW_m}{dt}
=\frac{d}{dt}\left(\frac12Li^2\right)
=Li\dot i+\frac12i^2\frac{dL}{dx}\dot x.
\]

二者之差是机械功率：

\[
p_{\rm mech}=p_e-\frac{dW_m}{dt}
=\frac12i^2\frac{dL}{dx}\dot x.
\]

因此再次得到

\[
F_x\dot x=p_{\rm mech},
\qquad
F_x=\frac12i^2\frac{dL}{dx}.
\]

\par\noindent\begin{minipage}{\linewidth}
\begin{center}
\includegraphics[width=\linewidth,height=.40\textheight,keepaspectratio]{\detokenize{assets/module-05/theory/tikz-m5-topic50-v9-from47-v8-03.pdf}}
\par\smallskip{\small 图示：可变电感的电端口功率分配}
\end{center}
\end{minipage}\par

\subsection{四、磁路、磁阻与气隙力}\label{m5:ux56dbux78c1ux8defux78c1ux963bux4e0eux6c14ux9699ux529b}

若磁通主要沿截面积近似恒定的磁路闭合，可写

\[
\mathcal R_m=\oint\frac{dl}{\mu S},
\qquad
\Phi=\frac{Ni}{\mathcal R_m},
\qquad
\boxed{L=\frac{N^2}{\mathcal R_m}}.
\]

磁阻串联时直接相加。若高磁导率铁芯中有长度为 \(g\)、面积为 \(S\)
的均匀气隙，且忽略边缘漏磁，则气隙通常占主要磁阻：

\[
\mathcal R_m\simeq\frac{g}{\mu_0S},
\qquad
L(g)\simeq\frac{\mu_0N^2S}{g}.
\]

取 \(g\) 增大的方向为正，固定电流时

\[
\boxed{
F_g=\frac12i^2\frac{dL}{dg}
=-\frac{\mu_0N^2Si^2}{2g^2}}.
\]

负号表示电磁力总是拉小气隙。用气隙磁场 \(B\simeq\mu_0Ni/g\)
表示，力的大小为

\[
\boxed{|F|=\frac{B^2}{2\mu_0}S}.
\]

其中 \(B^2/(2\mu_0)\) 可理解为磁压力。

\par\noindent\begin{minipage}{\linewidth}
\begin{center}
\includegraphics[width=\linewidth,height=.40\textheight,keepaspectratio]{\detokenize{assets/module-05/theory/tikz-m5-topic50-v9-from47-v8-04.pdf}}
\par\smallskip{\small 图示：单气隙磁路与吸力}
\end{center}
\end{minipage}\par

\par\addvspace{.4\baselineskip}
\noindent\begin{minipage}{\linewidth}

\textbf{例 M5-127：气隙减半}

高磁导率电磁铁只有一个主导气隙
\(g\)，忽略漏磁、边缘效应及铁芯磁阻，电流由理想电流源固定。将气隙减为原来一半，比较电感和吸力的倍数，并说明若改由一般固定电压源供电为何不能直接沿用该比例。

\end{minipage}\par

\par\noindent\begin{minipage}{\linewidth}
\textbf{解答}

\begin{center}
\includegraphics[width=\linewidth,height=.40\textheight,keepaspectratio]{\detokenize{assets/module-05/examples/M5-127.pdf}}
\par\smallskip{\small 图 M5-127：气隙减半}
\end{center}
\end{minipage}\par

图 M5-127 的橙色 \(g\)
是单个气隙长度；在忽略铁芯磁阻且电流固定时，自感与吸力具有不同的反比幂次。

在电流由理想电流源保持不变、漏磁可忽略时，\(L\propto1/g\)，吸力大小
\(|F|\propto1/g^2\)。气隙减半，自感加倍，吸力增为四倍。若线圈由固定电压源经电阻供电，电流可能随运动变化，不能直接使用这一比例。

\subsection{五、固定电流与固定磁链为何不矛盾}\label{m5:ux4e94ux56faux5b9aux7535ux6d41ux4e0eux56faux5b9aux78c1ux94feux4e3aux4f55ux4e0dux77dbux76fe}

在线性系统中，把磁能写成磁链变量：

\[
W_m(\Psi,x)=\frac{\Psi^2}{2L(x)}.
\]

固定磁链时

\[
F_x=-\left.\frac{\partial W_m}{\partial x}\right|_{\Psi}
=\frac{\Psi^2}{2L^2}\frac{dL}{dx}
=\frac12i^2\frac{dL}{dx}.
\]

与固定电流下的共能法结果相同。常见错误是把 \(W_m=Li^2/2\) 对 \(x\)
求导时既声称固定磁链，又把 \(i\) 当常量；这混用了两套独立变量。

\subsection{六、使用条件与常见错误}\label{m5:ux516dux4f7fux7528ux6761ux4ef6ux4e0eux5e38ux89c1ux9519ux8bef-1}

\begin{enumerate}
\def\labelenumi{\arabic{enumi}.}
\tightlist
\item
  可变几何必须从 \(v=d\Psi/dt\) 出发，不能只写 \(L\,di/dt\)。
\item
  固定电流用共能对位置求偏导；固定磁链用磁能对位置求负偏导。
\item
  线性系统中能量与共能数值相等，不代表求导时可以混用控制变量。
\item
  磁路法要求漏磁和边缘效应可忽略，截面上的磁通近似均匀。
\item
  铁磁材料进入饱和或出现磁滞后，\(L\) 不再是只由位置决定的常数。
\item
  由 \(B^2/(2\mu_0)\)
  得到的磁压力适用于近似真空或空气气隙；介质受力还需说明热力学约束。
\end{enumerate}

\subsection{练习与解答}\label{m5:ux7ec3ux4e60ux4e0eux89e3ux7b54-6}

\par\addvspace{.4\baselineskip}
\noindent\begin{minipage}{\linewidth}

\subsubsection{作业一：可变电感的端电压}\label{m5:ux4f5cux4e1aux4e00ux53efux53d8ux7535ux611fux7684ux7aefux7535ux538b}

\textbf{题号：M5-128}

\textbf{层级：} 基础　　\textbf{方法：} 自感互感、机电耦合

机械速度为常量 \(v_x\)，端电压按被动方向记为
\(u_L\)；\(x(t)=x_0+v_xt\)，且所研究时间内 \(L(x)>0\)。

线性电感满足 \(L(x)=L_0+\alpha x\)，线圈电流为
\(i(t)=I_0\cos\omega t\)，机械坐标为 \(x(t)=x_0+v_xt\)。求端电压。

\end{minipage}\par

\par\noindent\begin{minipage}{\linewidth}
\textbf{解答}

\begin{center}
\includegraphics[width=\linewidth,height=.40\textheight,keepaspectratio]{\detokenize{assets/module-05/examples/M5-128.pdf}}
\par\smallskip{\small 图 M5-128：可变电感的端电压}
\end{center}
\end{minipage}\par

图 M5-128 中铁芯位置和电流同时改变，端电压必须保留两个乘积求导项；图中的
\(+x\) 方向决定速度参数 \(v_x\) 的正负。

采用被动方向，线圈磁链为 \(\Psi=L(x)i\)。对乘积求全导数：

\[
u_L=\frac{d\Psi}{dt}=L(x)\dot i+i\frac{dL}{dx}\dot x.
\]

由 \(\dot i=-\omega I_0\sin\omega t\)、\(\dot x=v_x\) 及
\(L^{\prime}(x)=\alpha\)，得

\[
\boxed{u_L(t)=-\omega I_0[L_0+\alpha(x_0+v_xt)]\sin\omega t
+\alpha v_xI_0\cos\omega t}.
\]

第一项来自电流变化，第二项来自铁芯运动；\(v_x=0\)
时才退化为普通固定电感关系。

\par\addvspace{.4\baselineskip}
\noindent\begin{minipage}{\linewidth}

\subsubsection{作业二：气隙电磁铁的吸力}\label{m5:ux4f5cux4e1aux4e8cux6c14ux9699ux7535ux78c1ux94c1ux7684ux5438ux529b}

\textbf{题号：M5-129}

\textbf{层级：} 基础　　\textbf{方法：} 自感互感、机电耦合

以气隙长度 \(g\)
为机械坐标，增大气隙为正方向；要求给出有符号力及吸力大小。

匝数为 \(N\) 的线圈绕在高磁导率铁芯上，磁路中有单个气隙，长度为
\(g\)、面积为 \(S\)。忽略铁芯磁阻、漏磁和边缘效应，电流恒为
\(I\)。求自感和吸力。

\end{minipage}\par

\par\noindent\begin{minipage}{\linewidth}
\textbf{解答}

\begin{center}
\includegraphics[width=\linewidth,height=.40\textheight,keepaspectratio]{\detokenize{assets/module-05/examples/M5-129.pdf}}
\par\smallskip{\small 图 M5-129：气隙电磁铁的吸力}
\end{center}
\end{minipage}\par

图 M5-129 的气隙长度增大时电感减小，所以以 \(g\)
为坐标的力为负，实际作用是闭合气隙。

气隙磁阻为

\[
\mathcal R_m=\frac{g}{\mu_0S}.
\]

故

\[
\boxed{L=\frac{N^2}{\mathcal R_m}=\frac{\mu_0N^2S}{g}}.
\]

\par\noindent\begin{minipage}{\linewidth}
\begin{center}
\includegraphics[width=\linewidth,height=.40\textheight,keepaspectratio]{\detokenize{assets/module-05/process/M5-129a.pdf}}
\par\smallskip{\small 图 M5-129a：固定电流下缩短气隙会增强吸力}
\end{center}
\end{minipage}\par

图 M5-129a
的正值曲线只表示吸力大小；有符号广义力沿气隙坐标的负方向，不能把大小曲线误当作力的正负。

取 \(g\) 增大为正方向，

\[
F_g=\frac12I^2\frac{dL}{dg}
=-\frac{\mu_0N^2SI^2}{2g^2}.
\]

所以吸力大小为

\[
\boxed{|F|=\frac{\mu_0N^2SI^2}{2g^2}=\frac{B^2S}{2\mu_0}}.
\]

\par\addvspace{.4\baselineskip}
\noindent\begin{minipage}{\linewidth}

\subsubsection{作业三：固定磁链下的力}\label{m5:ux4f5cux4e1aux4e09ux56faux5b9aux78c1ux94feux4e0bux7684ux529b}

\textbf{题号：M5-130}

\textbf{层级：} 基础　　\textbf{方法：} 自感互感、机电耦合

线圈与外界无变化外磁通耦合，无外加电源，\(L(x)>0\)；固定磁链是总磁链条件，并非固定电流。

某无电阻闭合线圈的自感为 \(L(x)\)，初始磁链为
\(\Psi_0\)。在准静态运动中磁链保持不变，求广义力。

\end{minipage}\par

\par\noindent\begin{minipage}{\linewidth}
\textbf{解答}

\begin{center}
\includegraphics[width=\linewidth,height=.40\textheight,keepaspectratio]{\detokenize{assets/module-05/examples/M5-130.pdf}}
\par\smallskip{\small 图 M5-130：固定磁链下的力}
\end{center}
\end{minipage}\par

图 M5-130
中无电阻回路固定的是总磁链，电流随电感改变；对磁能求位置导数后，力仍趋向电感增大的方向。

固定磁链时应使用

\[
W_m(\Psi_0,x)=\frac{\Psi_0^2}{2L(x)}.
\]

因此

\[
\boxed{
F_x=-\frac{\partial W_m}{\partial x}
=\frac{\Psi_0^2}{2L^2}\frac{dL}{dx}}.
\]

又因 \(i=\Psi_0/L\)，也可写成

\[
\boxed{F_x=\frac12i^2\frac{dL}{dx}}.
\]

\par\addvspace{.4\baselineskip}
\noindent\begin{minipage}{\linewidth}

\subsubsection{作业四：非线性磁性系统的共能}\label{m5:ux4f5cux4e1aux56dbux975eux7ebfux6027ux78c1ux6027ux7cfbux7edfux7684ux5171ux80fd}

\textbf{题号：M5-131}

\textbf{层级：} 基础　　\textbf{方法：} 自感互感、机电耦合

材料关系单值、无磁滞，磁共能以下限 \(i=0\) 为零；\(b\)
为常量，\(a(x)>0\)。

某磁性装置在固定位置 \(x\) 的磁链关系为 \(\Psi=ai+bi^3\)，其中
\(a,b>0\)。求磁共能；若 \(a=a(x)\)、\(b\)
与位置无关，求固定电流下的广义力。

\end{minipage}\par

\par\noindent\begin{minipage}{\linewidth}
\textbf{解答}

\begin{center}
\includegraphics[width=\linewidth,height=.40\textheight,keepaspectratio]{\detokenize{assets/module-05/examples/M5-131.pdf}}
\par\smallskip{\small 图 M5-131：非线性磁性系统的共能}
\end{center}
\end{minipage}\par

图 M5-131 的曲线下方面积是
\(\int_0^i\Psi\,di\)，对应磁共能。固定电流求导时，只有题中随位置变化的
\(a(x)\) 产生力项。

磁共能为

\[
W_m'(i,x)=\int_0^i(ai'+bi'^3)\,di'
=\frac12ai^2+\frac14bi^4.
\]

固定电流时

\[
F_x=\left.\frac{\partial W_m'}{\partial x}\right|_i
=\frac12i^2\frac{da}{dx}.
\]

故

\[
\boxed{F_x=\frac12i^2a'(x)}.
\]

不能把非线性系统的能量写成 \(i\Psi/2\)；只有线性关系下才有该形式。

\par\addvspace{.4\baselineskip}
\noindent\begin{minipage}{\linewidth}

\subsubsection{作业五：孤立长直导线的自感为何发散}\label{m5:ux4f5cux4e1aux4e94ux5b64ux7acbux957fux76f4ux5bfcux7ebfux7684ux81eaux611fux4e3aux4f55ux53d1ux6563}

\textbf{题号：M5-132}

\textbf{层级：} 基础　　\textbf{方法：}
毕奥---萨伐尔、自感互感、机电耦合

空间为真空，\(R>a>0\)；只计算 \(a<r<R\)
的单位长度磁场能，不预设真实回流支路。

半径为 \(a\) 的长直导线通电流 \(I\)，暂以同轴圆柱面 \(r=R\)
作为磁场能积分的外边界。忽略导线内部磁场，求单位长度外部自感，并讨论
\(R\to\infty\)。

\end{minipage}\par

\par\noindent\begin{minipage}{\linewidth}
\textbf{解答}

\begin{center}
\includegraphics[width=\linewidth,height=.40\textheight,keepaspectratio]{\detokenize{assets/module-05/examples/M5-132.pdf}}
\par\smallskip{\small 图 M5-132：孤立长直导线的自感为何发散}
\end{center}
\end{minipage}\par

图 M5-132 的同轴外环在本题是人为设定的积分截断半径
\(R\)；若不指定回流，外场能的对数积分随截断半径增加而发散。

导线外

\[
B(r)=\frac{\mu_0 I}{2\pi r}.
\]

单位长度磁场能为

\[
W_m'=\int_a^R\frac{B^2}{2\mu_0}2\pi r\,dr
=\frac{\mu_0 I^2}{4\pi}\ln\frac Ra.
\]

由 \(W_m'=L'I^2/2\)，

\[
\boxed{L'_{\text{外}}=\frac{\mu_0}{2\pi}\ln\frac Ra}.
\]

当 \(R\to\infty\)
时对数发散。这说明``单根孤立导线的自感''没有独立有限值；真实电流必须有回流路径，自感属于完整闭合回路。

\par\addvspace{.4\baselineskip}
\noindent\begin{minipage}{\linewidth}

\subsubsection{作业六：平行双导线传输线的单位长度自感}\label{m5:ux4f5cux4e1aux516dux5e73ux884cux53ccux5bfcux7ebfux4f20ux8f93ux7ebfux7684ux5355ux4f4dux957fux5ea6ux81eaux611f}

\textbf{题号：M5-133}

\textbf{层级：} 复赛　　\textbf{方法：}
毕奥---萨伐尔、自感互感、机电耦合

两根半径均为 \(a\)、圆心距为 \(D>2a\) 的平行圆柱导线，在磁导率为 \(\mu\)
的均匀介质中通等大反向电流 \(I\)，忽略导体内部磁场能。

\begin{enumerate}
\def\labelenumi{(\arabic{enumi})}
\item
  先讨论
  \(D\gg a\)，用各导线附近的近似圆周磁场求回路单位长度外部自感的主阶，说明为何两根线的能量都应计入。
\item
  \textbf{精确结果选学：}
  理想高频表面电流、准TEM横截面模型的双极坐标边界解给出 \[
  W_m'=\frac{\mu I^2}{2\pi}\operatorname{arcosh}\frac{D}{2a},
  \qquad \operatorname{arcosh}z=\ln(z+\sqrt{z^2-1}).
  \]
  把这一单位长度磁能作为已知，求精确外部自感，并检验与第(1)问一致。实际表面电流不均匀；两个孤立导线的近场近似不能当作近距离精确解。
\end{enumerate}

\end{minipage}\par

\par\noindent\begin{minipage}{\linewidth}
\textbf{解答}

\begin{center}
\includegraphics[width=\linewidth,height=.40\textheight,keepaspectratio]{\detokenize{assets/module-05/examples/M5-133.pdf}}
\par\smallskip{\small 图 M5-133：平行双导线传输线的单位长度自感}
\end{center}
\end{minipage}\par

图 M5-133 的两导线电流等大反向，圆心距 \(D\) 与导线半径 \(a\)
不可混淆。精确外部解使用双极坐标；大间距时再取对数近似。

先看主阶近似。每根导线附近 \(B\simeq\mu I/(2\pi r)\)，从线径尺度 \(a\)
积分到间距尺度 \(D\)，一根线的单位长度磁能约为 \[
\int_a^D\frac{B^2}{2\mu}2\pi r\,dr
=\frac{\mu I^2}{4\pi}\ln\frac Da.
\] 两根线贡献相同，故 \(W_m'\simeq\mu I^2\ln(D/a)/(2\pi)\)，再用
\(W_m'=L'I^2/2\) 得
\(L'\simeq\mu\ln(D/a)/\pi\)。远场重叠区只影响该大对数近似的低阶部分。以下精确式使用题干已给的横截面磁能结论，不要求先掌握双极坐标。

令二极坐标为

\[
x=\frac{c\sinh\tau}{\cosh\tau-\cos\sigma},
\qquad
y=\frac{c\sin\sigma}{\cosh\tau-\cos\sigma}.
\]

\(\tau=\pm\tau_0\)
是两族等半径圆。把它们取作两导体表面，可由圆的半径和圆心位置得到

\[
\boxed{\cosh\tau_0=\frac{D}{2a}}.
\]

横截面的静电势只依赖 \(\tau\)，故在两导体分别带单位长度电荷
\(\pm\lambda\) 时，电势差为 \(V=\lambda\tau_0/(\pi\varepsilon)\)，从而

\[
C'=\frac{\lambda}{V}=\frac{\pi\varepsilon}{\tau_0}
=\frac{\pi\varepsilon}{\operatorname{arcosh}(D/2a)}.
\]

对无耗均匀介质中的横电磁模式，电场与磁场具有相同的二维等势结构；单位长度磁场能为

\[
W_m'=\frac{\mu I^2\tau_0}{2\pi}.
\]

由 \(W_m'=L'I^2/2\) 得

\[
\boxed{
L'=\frac{\mu}{\pi}\operatorname{arcosh}\frac{D}{2a}}.
\]

当 \(D\gg a\) 时，

\[
\operatorname{arcosh}\frac{D}{2a}
=\ln\left(\frac{D+\sqrt{D^2-4a^2}}{2a}\right)
\simeq\ln\frac Da,
\]

所以

\[
\boxed{L'\simeq\frac\mu\pi\ln\frac Da}.
\]

并且

\[
L'C'=\mu\varepsilon=\frac1{v_p^2},
\]

与横电磁波沿线传播速度 \(v_p=1/\sqrt{\mu\varepsilon}\)
一致。这里假定导体为等势面、介质均匀，并只计算导体外部磁场能；导体内部自感和高频色散均未计入。

\par\addvspace{.4\baselineskip}
\noindent\begin{minipage}{\linewidth}

\subsubsection{作业七：同轴线的电感与电容乘积}\label{m5:ux4f5cux4e1aux4e03ux540cux8f74ux7ebfux7684ux7535ux611fux4e0eux7535ux5bb9ux4e58ux79ef}

\textbf{题号：M5-134}

\textbf{层级：} 基础　　\textbf{方法：} 自感互感、机电耦合

半径 \(a\) 是内导体外半径，\(b>a\)
是外导体内半径；介质线性、均匀，忽略导体内部自感与端部效应，讨论准TEM传播。

同轴线内外导体半径分别为 \(a,b\)，介质为
\(\varepsilon,\mu\)。求单位长度电感、电容，并检验传播速度关系。

\end{minipage}\par

\par\noindent\begin{minipage}{\linewidth}
\textbf{解答}

\begin{center}
\includegraphics[width=\linewidth,height=.40\textheight,keepaspectratio]{\detokenize{assets/module-05/examples/M5-134.pdf}}
\par\smallskip{\small 图 M5-134：同轴线的电感与电容乘积}
\end{center}
\end{minipage}\par

图 M5-134
的两导体间隙同时储存电场能和磁场能；电感、电容中的对数因子互为倒数，因此乘积只取决于介质。

由磁场能或磁通法，

\[
L'=\frac\mu{2\pi}\ln\frac ba.
\]

由静电场解，

\[
C'=\frac{2\pi\varepsilon}{\ln(b/a)}.
\]

因此

\[
\boxed{L'C'=\mu\varepsilon},
\]

从而

\[
\boxed{v_p=\frac1{\sqrt{L'C'}}=\frac1{\sqrt{\mu\varepsilon}}}.
\]

这个检验只涉及外部电感；低频下若电流进入导体内部，还需加内自感。

\par\addvspace{.4\baselineskip}
\noindent\begin{minipage}{\linewidth}

\subsubsection{作业八：嵌套长螺线管的互感}\label{m5:ux4f5cux4e1aux516bux5d4cux5957ux957fux87baux7ebfux7ba1ux7684ux4e92ux611f}

\textbf{题号：M5-135}

\textbf{层级：} 基础　　\textbf{方法：} 自感互感、机电耦合

两绕组参考方向使其轴向场同向，内管的截面完全包含于外管，介质线性且两管均足够长。

内外两个同轴长螺线管完全重叠，单位长度匝数分别为
\(n_1,n_2\)，共同重叠长度为 \(\ell\)，内管截面积为 \(S\)，介质磁导率为
\(\mu\)。忽略漏磁，求互感。

\end{minipage}\par

\par\noindent\begin{minipage}{\linewidth}
\textbf{解答}

\begin{center}
\includegraphics[width=\linewidth,height=.40\textheight,keepaspectratio]{\detokenize{assets/module-05/examples/M5-135.pdf}}
\par\smallskip{\small 图 M5-135：嵌套长螺线管的互感}
\end{center}
\end{minipage}\par

图 M5-135 的内外螺线管共享长度 \(\ell\)，互链面积由较小的内管截面 \(S\)
决定，而不是取外管面积。

线圈 1 通电流 \(i_1\) 时，重叠区域内

\[
B_1=\mu n_1i_1.
\]

线圈 2 每匝磁通为 \(B_1S\)，总匝数为 \(n_2\ell\)，故

\[
\Psi_{21}=n_2\ell\,B_1S
=\mu n_1n_2S\ell\,i_1.
\]

所以

\[
\boxed{M=\mu n_1n_2S\ell}.
\]

若两线圈绕向相反，则有符号互感为负；互感大小不变。

\par\addvspace{.4\baselineskip}
\noindent\begin{minipage}{\linewidth}

\subsubsection{作业九：两套同轴螺线管的轴向吸力}\label{m5:ux4f5cux4e1aux4e5dux4e24ux5957ux540cux8f74ux87baux7ebfux7ba1ux7684ux8f74ux5411ux5438ux529b}

\textbf{题号：M5-136}

\textbf{层级：} 基础　　\textbf{方法：} 自感互感、机电耦合

线性模型中各线圈自感随套叠位置的变化忽略，仅互感随共同重叠长度 \(x\)
改变；参考电流正向取产生同向轴向场。

两套长螺线管部分套叠，重叠长度为 \(x\)，公共截面积为 \(S\)，匝密度为
\(n_1,n_2\)，电流由理想电流源固定为
\(I_1,I_2\)。忽略端部效应，求轴向力。

\end{minipage}\par

\par\noindent\begin{minipage}{\linewidth}
\textbf{解答}

\begin{center}
\includegraphics[width=\linewidth,height=.40\textheight,keepaspectratio]{\detokenize{assets/module-05/examples/M5-136.pdf}}
\par\smallskip{\small 图 M5-136：两套同轴螺线管的轴向吸力}
\end{center}
\end{minipage}\par

图 M5-136 的共同区间长度为
\(x\)，箭头表示内绕组向增加重叠的方向运动；互感随共同重叠长度线性增加。

由上一题，随重叠长度变化的互感为

\[
M(x)=\mu n_1n_2Sx.
\]

\par\noindent\begin{minipage}{\linewidth}
\begin{center}
\includegraphics[width=\linewidth,height=.40\textheight,keepaspectratio]{\detokenize{assets/module-05/process/M5-136a.pdf}}
\par\smallskip{\small 图 M5-136a：互感随重叠长度增加}
\end{center}
\end{minipage}\par

图 M5-136a
的橙色内绕组移动时，新增重叠片段链接额外磁通；在端部效应忽略的区间，单位新增长度对应同样的互感增量。

两线圈的共能中与位置有关的部分为

\[
W_m'=M(x)I_1I_2.
\]

所以

\[
\boxed{F_x=\frac{\partial W_m'}{\partial x}
=\mu n_1n_2SI_1I_2}.
\]

若电流绕向使 \(I_1I_2>0\)，力使重叠长度增加；若绕向相反，力的方向反转。

\par\addvspace{.4\baselineskip}
\noindent\begin{minipage}{\linewidth}

\subsubsection{作业十：分支磁路的磁通分配}\label{m5:ux4f5cux4e1aux5341ux5206ux652fux78c1ux8defux7684ux78c1ux901aux5206ux914d}

\textbf{题号：M5-137}

\textbf{层级：} 基础　　\textbf{方法：} 安培环路、法拉第定律、自感互感

材料为磁导率 \(\mu\) 的线性各向同性介质，截面均为
\(S\)；两支路连接在相同磁势节点之间，不另计共同主干磁阻。

同一磁动势驱动两条并联磁支路。两支路材料和截面积相同，长度分别为
\(\ell\) 与 \(3\ell\)，忽略漏磁。求两支路磁通之比和等效磁阻。

\end{minipage}\par

\par\noindent\begin{minipage}{\linewidth}
\textbf{解答}

\begin{center}
\includegraphics[width=\linewidth,height=.40\textheight,keepaspectratio]{\detokenize{assets/module-05/examples/M5-137.pdf}}
\par\smallskip{\small 图 M5-137：两并联磁支路的磁通分配}
\end{center}
\end{minipage}\par

图 M5-137
的两个并联磁阻承受同一磁势差，蓝色与绿色箭头分别是两支路磁通。磁通按磁阻倒数分流，较短支路磁通较大。

支路磁阻为

\[
\mathcal R_1=\frac\ell{\mu S},
\qquad
\mathcal R_2=\frac{3\ell}{\mu S}=3\mathcal R_1.
\]

并联支路两端磁势降相同，所以

\[
\frac{\Phi_1}{\Phi_2}
=\frac{\mathcal R_2}{\mathcal R_1}=3.
\]

因此

\[
\boxed{\Phi_1:\Phi_2=3:1}.
\]

等效磁阻为

\[
\boxed{\mathcal R_{\rm eq}
=\frac{\mathcal R_1\mathcal R_2}{\mathcal R_1+\mathcal R_2}
=\frac34\mathcal R_1}.
\]

\par\addvspace{.4\baselineskip}
\noindent\begin{minipage}{\linewidth}

\subsubsection{作业十一：正四面体面回路的互感}\label{m5:ux4f5cux4e1aux5341ux4e00ux6b63ux56dbux9762ux4f53ux9762ux56deux8defux7684ux4e92ux611f}

\textbf{题号：M5-138}

\textbf{层级：} 基础　　\textbf{方法：} 自感互感、机电耦合

正四面体的四个三角形面均按外法向取正。沿其中一个面的边界通电流
\(I\)，其磁场穿过源面及其余三个面。给定源面磁通
\(\Phi_{\rm self}=LI\)，其中 \(L\) 为有限值。

用``穿过闭合面的总磁通为零''和四面体的三重对称性，求它与任一相邻面的带符号互感
\(M\) 及大小。

四个面的磁通采用同一组具有小而有限线径的共棱导线约定，公共棱附近完全按相同规则处理。本题只使用给定的有限
\(L\)，不要求计算导线附近的自感细节。

\end{minipage}\par

\par\noindent\begin{minipage}{\linewidth}
\textbf{解答}

\begin{center}
\includegraphics[width=\linewidth,height=.40\textheight,keepaspectratio]{\detokenize{assets/module-05/examples/M5-138.pdf}}
\par\smallskip{\small 图 M5-138：正四面体面回路的互感}
\end{center}
\end{minipage}\par

图 M5-138 的源面为
\(ABC\)，各面正向由各自外法向规定。闭合四面体的总磁通为零，另三面由对称性各分担源面磁通的负三分之一。

取四面体各面的外法向为正。源回路所在面的自磁通大小为 \(LI\)。由于

\[
\oiint_{\partial V}\mathbf B\cdot d\mathbf S=0,
\]

其余三个全等面上的磁通之和必须为 \(-LI\)。由三重对称性，每个相邻面得到

\[
\Phi_{\text{邻}}=-\frac13LI.
\]

故按外法向定向时

\[
M=-\frac L3,
\]

互感大小为

\[
\boxed{|M|=\frac L3}.
\]

若把相邻面的回路正向反转，\(M\) 的符号随之反转。

严格的零线径细丝模型在导线自身附近有对数奇异，自感和跨公共棱的互感都不能分别视为有限量；上面的
\(|M|/L=1/3\)
依赖保持四面体对称性的共同正则化，也不适用于任意实际接线方式。

\par\addvspace{.4\baselineskip}
\noindent\begin{minipage}{\linewidth}

\subsubsection{作业十二：小线圈转动引起的互感电动势}\label{m5:ux4f5cux4e1aux5341ux4e8cux5c0fux7ebfux5708ux8f6cux52a8ux5f15ux8d77ux7684ux4e92ux611fux7535ux52a8ux52bf}

\textbf{题号：M5-139}

\textbf{层级：} 基础　　\textbf{方法：}
毕奥---萨伐尔、法拉第定律、自感互感

\(I_1\ne0\)
保持恒定，初始法向与外场同向；开路电动势正向由小线圈瞬时正法向配对。

一个面积为 \(S\)、匝数为 \(N_2\)
的小线圈位于大线圈近中心区域。大线圈通恒定电流
\(I_1\)，在小线圈位置产生近似均匀磁场 \(B_1\)。小线圈以角速度 \(\omega\)
转动，其法向与 \(\mathbf B_1\) 的夹角为
\(\theta=\omega t\)。求互感与开路感应电动势。

\end{minipage}\par

\par\noindent\begin{minipage}{\linewidth}
\textbf{解答}

\begin{center}
\includegraphics[width=\linewidth,height=.40\textheight,keepaspectratio]{\detokenize{assets/module-05/examples/M5-139.pdf}}
\par\smallskip{\small 图 M5-139：小线圈转动引起的互感电动势}
\end{center}
\end{minipage}\par

图 M5-139
中绿色小线圈法向转动，外部蓝场由大线圈的恒定电流产生；互感只因投影面积的余弦改变。

小线圈磁链为

\[
\Psi_{21}=N_2B_1S\cos\theta=M(\theta)I_1.
\]

所以

\[
\boxed{M(\theta)=\frac{N_2B_1S}{I_1}\cos\theta=M_0\cos\theta}.
\]

因 \(I_1\) 恒定，开路感应电动势为

\[
\boxed{\mathcal E_2
=-\frac{d\Psi_{21}}{dt}
=N_2B_1S\omega\sin\omega t}.
\]

它来自互感随姿态变化，而不是源电流变化。

\par\addvspace{.4\baselineskip}
\noindent\begin{minipage}{\linewidth}

\subsubsection{作业十三：小圆线圈与大圆线圈之间的轴向力}\label{m5:ux4f5cux4e1aux5341ux4e09ux5c0fux5706ux7ebfux5708ux4e0eux5927ux5706ux7ebfux5708ux4e4bux95f4ux7684ux8f74ux5411ux529b}

\textbf{题号：M5-140}

\textbf{层级：} 基础　　\textbf{方法：}
毕奥---萨伐尔、自感互感、机电耦合

以沿 \(+z\) 产生磁矩为小线圈 \(I_2\)
的正方向，大线圈正电流同样取产生正轴向场的方向；\(I_1\text{、}I_2\)
均由理想源保持。

半径为 \(R\)、匝数为 \(N\) 的大圆线圈通恒定电流 \(I_1\)。面积为 \(S\)
的单匝小线圈与其同轴，位于轴线上距大线圈中心 \(z\) 处，通恒定电流
\(I_2\)。小线圈足够小，可认为其截面内磁场均匀。求轴向力。

\end{minipage}\par

\par\noindent\begin{minipage}{\linewidth}
\textbf{解答}

\begin{center}
\includegraphics[width=\linewidth,height=.40\textheight,keepaspectratio]{\detokenize{assets/module-05/examples/M5-140.pdf}}
\par\smallskip{\small 图 M5-140：小圆线圈与大圆线圈之间的轴向力}
\end{center}
\end{minipage}\par

图 M5-140 的轴上 \(P\) 放置面积为 \(S\) 的小线圈，正电流取使磁矩沿
\(+z\)。当 \(z>0\) 且两磁矩同向时，导数为负，力指向大线圈。

大线圈轴线磁场为

\[
B_1(z)=\frac{\mu_0NI_1R^2}{2(R^2+z^2)^{3/2}}.
\]

若两线圈正法向同向，互感为

\[
M(z)=\frac{SB_1(z)}{I_1}
=\frac{\mu_0NSR^2}{2(R^2+z^2)^{3/2}}.
\]

固定电流时用磁共能求力：

\[
F_z=I_1I_2\frac{dM}{dz}.
\]

因此

\[
\boxed{
F_z=-\frac{3\mu_0NSR^2z}{2(R^2+z^2)^{5/2}}I_1I_2}.
\]

当 \(I_1I_2>0\) 时，\(z>0\)
处的力指向大线圈，表现为吸引；反向电流使力反向。

\par\addvspace{.4\baselineskip}
\noindent\begin{minipage}{\linewidth}

\subsubsection{作业十四：长螺线管绕组承受的磁压力}\label{m5:ux4f5cux4e1aux5341ux56dbux957fux87baux7ebfux7ba1ux7ed5ux7ec4ux627fux53d7ux7684ux78c1ux538bux529b}

\textbf{题号：M5-141}

\textbf{层级：} 基础　　\textbf{方法：} 自感互感、机电耦合

绕组用理想薄电流面近似；计算作用场须取电流面两侧磁场的平均值，避免把自场全部重复计入。

真空长螺线管单位长度匝数为 \(n\)，电流为
\(I\)。忽略端部效应，求绕组承受的径向磁压力，并解释方向。

\end{minipage}\par

\par\noindent\begin{minipage}{\linewidth}
\textbf{解答}

\begin{center}
\includegraphics[width=\linewidth,height=.40\textheight,keepaspectratio]{\detokenize{assets/module-05/examples/M5-141.pdf}}
\par\smallskip{\small 图 M5-141：长螺线管绕组承受的磁压力}
\end{center}
\end{minipage}\par

图 M5-141
中管内磁场能随绕组半径增加而变化；把绕组局部看作电流面并取两侧平均场，得到向外的磁压力。

管内磁场为

\[
B=\mu_0nI,
\]

管外近似为零。真空磁场的法向应力差给出压力

\[
p_m=\frac{B_{\text{内}}^2-B_{\text{外}}^2}{2\mu_0}.
\]

所以

\[
\boxed{p_m=\frac12\mu_0n^2I^2}.
\]

该压力把绕组向外撑开，形成环向拉应力。它与把线圈简单想象成``相邻平行电流互相吸引''并不矛盾：局部导线间作用与整个连续绕组的宏观应力必须在同一几何中合成，磁应力张量给出的整体结果是径向外压。

\par\addvspace{.4\baselineskip}
\noindent\begin{minipage}{\linewidth}

\subsubsection{作业十五：半无限长螺线管端面中心的磁场}\label{m5:ux4f5cux4e1aux5341ux4e94ux534aux65e0ux9650ux957fux87baux7ebfux7ba1ux7aefux9762ux4e2dux5fc3ux7684ux78c1ux573a}

\textbf{题号：M5-142}

\textbf{层级：} 复赛　　\textbf{方法：} 自感互感、机电耦合

空间为真空，绕组均匀且紧密，取电流方向使轴线场沿
\(+z\)，仅计算端面中心，不是把全空间磁场都作半值处理。

半径为 \(a\)、单位长度匝数为 \(n\) 的螺线管占据 \(z>0\)
的半无限空间，电流为 \(I\)。求端面中心 \(O\)
处的轴向磁场，并与无限长螺线管内部磁场比较。

\end{minipage}\par

\par\noindent\begin{minipage}{\linewidth}
\textbf{解答}

\begin{center}
\includegraphics[width=\linewidth,height=.40\textheight,keepaspectratio]{\detokenize{assets/module-05/examples/M5-142.pdf}}
\par\smallskip{\small 图 M5-142：半无限长螺线管端面中心的磁场}
\end{center}
\end{minipage}\par

图 M5-142 的 \(O\) 是半无限长螺线管端面中心；从 \(z=0\) 到 \(+\infty\)
积分圆环轴线场，只得到无限长管内场的一半。

把位于 \(z\) 至 \(z+dz\) 的薄层看成 \(n\,dz\) 匝同轴圆电流。单匝半径为
\(a\) 的圆电流在轴线上距圆心 \(z\) 处产生

\[
dB_z
=\frac{\mu_0I a^2}{2(a^2+z^2)^{3/2}}\,n\,dz.
\]

各薄层在端面中心产生的磁场方向相同，因此

\[
B_O
=\frac{\mu_0nI}{2}
\int_0^{\infty}
\frac{a^2\,dz}{(a^2+z^2)^{3/2}}.
\]

注意到

\[
\frac{d}{dz}\left(\frac{z}{\sqrt{a^2+z^2}}\right)
=\frac{a^2}{(a^2+z^2)^{3/2}},
\]

所以积分值为 \(1\)，从而

\[
\boxed{B_O=\frac12\mu_0nI}.
\]

无限长螺线管深处的磁场为
\(B_{\infty}=\mu_0nI\)，故端面中心磁场恰为其一半。这个结果只描述端面中心的局域磁场，不能据此把整根有限螺线管的自感简单减半；自感必须由全体磁链或总磁能计算。

\section{专题五十一：电磁---机械耦合系统的建模（基础必讲与拔高选讲）}\label{m5:ux4e13ux9898ux4e94ux5341ux4e00ux7535ux78c1ux673aux68b0ux8026ux5408ux7cfbux7edfux7684ux5efaux6a21ux57faux7840ux5fc5ux8bb2ux4e0eux62d4ux9ad8ux9009ux8bb2}

滑杆、导轨和储能元件组成的系统要同时满足电路方程、机械方程和能量守恒。选定参考方向与机械坐标后，可按具体元件和约束化简。

\subsection{一、杆---轨模型}\label{m5:ux4e00ux6746ux8f68ux6a21ux578b}

设长度为 \(\ell\) 的导体杆沿导轨运动，速度
\(v=\dot x\)。匀强磁场垂直导轨平面，记

\[
G=B\ell.
\]

选取正方向，使杆向正向运动时产生正电动势
\(Gv\)，正电流产生反向安培力。串联支路含 \(R,L,C\)，电容电荷为
\(q\)，并有 \(\dot q=i\)。方程组为

\[
\boxed{
\begin{aligned}
L\dot i+Ri+\frac qC&=Gv,\\
m\dot v&=F_{\rm ext}-Gi,\\
\dot q&=i.
\end{aligned}}
\]

这里 \(Gv\) 是机械端向电端的耦合项，\(Gi\)
是电端对机械端的反作用；二者系数相同，保证能量交换一致。

\par\noindent\begin{minipage}{\linewidth}
\begin{center}
\includegraphics[width=\linewidth,height=.40\textheight,keepaspectratio]{\detokenize{assets/module-05/theory/tikz-m5-topic51-v9-from48-v8-01.pdf}}
\par\smallskip{\small 图示：RLC滑杆的统一方向与方程}
\end{center}
\end{minipage}\par

\subsection{二、能量守恒是首要检验}\label{m5:ux4e8cux80fdux91cfux5b88ux6052ux662fux9996ux8981ux68c0ux9a8c}

电路方程乘以 \(i\)，机械方程乘以 \(v\)，相加得到

\[
\frac d{dt}
\left(
\frac12mv^2+\frac12Li^2+\frac{q^2}{2C}
\right)
=F_{\rm ext}v-Ri^2.
\]

即

\[
\boxed{
\frac{dE_{\text{总}}}{dt}
=P_{\text{外力}}-P_{\text{焦耳热}}}.
\]

耦合功率 \(Gvi\)
在电路和机械方程中一正一负，内部抵消。若列出的方程不能产生这一抵消，通常说明电流方向或安培力方向写错。

\subsection{三、只含电容：等效惯性}\label{m5:ux4e09ux53eaux542bux7535ux5bb9ux7b49ux6548ux60efux6027}

令 \(R=0\)、\(L=0\)，则

\[
\frac qC=Gv,
\qquad
i=\dot q=CG\dot v.
\]

代入机械方程：

\[
m\dot v=F_{\rm ext}-CG^2\dot v.
\]

所以

\[
\boxed{(m+CG^2)\dot v=F_{\rm ext}}.
\]

电容并未真正增加物体质量，但电容能

\[
\frac{q^2}{2C}=\frac12CG^2v^2
\]

与速度平方成正比，在动力学中表现为附加惯性 \(CG^2\)。

\par\addvspace{.4\baselineskip}
\noindent\begin{minipage}{\linewidth}

\textbf{例 M5-143：恒力加速}

质量为 \(m\)、长度 \(\ell\) 的滑杆在垂直导轨的匀强磁场 \(B\)
中向右运动，回路只有理想电容
\(C\)，杆和导轨电阻、自感均忽略。初始杆静止、电容不带电，从 \(t=0\)
起施恒力 \(F_0\)。求加速度、速度和电荷，定义 \(G=B\ell\)。

\end{minipage}\par

\par\noindent\begin{minipage}{\linewidth}
\textbf{解答}

\begin{center}
\includegraphics[width=\linewidth,height=.40\textheight,keepaspectratio]{\detokenize{assets/module-05/examples/M5-143.pdf}}
\par\smallskip{\small 图 M5-143：恒力加速}
\end{center}
\end{minipage}\par

图 M5-143
中电容电荷与滑杆速度通过动生电压相联系；电流随加速度变化，安培力因此表现为附加惯性
\(CG^2\)。

取图 M5-143 滑杆向右为正，\(G=B\ell\)，电荷 \(q\) 按图中使电容电压满足
\(q/C=Gv\) 的极板取正。恒力从静止开始作用，故

\[
a=\frac{F_0}{m+CG^2},\qquad
v(t)=\frac{F_0t}{m+CG^2},\qquad
q(t)=CGv(t)=\frac{CGF_0t}{m+CG^2}.
\]

与无电容时相比，加速度减小；无电阻时外力做功全部储存在机械动能与电场能中。

\subsection{四、只含电感：电磁---机械振荡}\label{m5:ux56dbux53eaux542bux7535ux611fux7535ux78c1ux673aux68b0ux632fux8361}

令 \(R=0\)，去掉电容支路，并令 \(F_{\rm ext}=0\)：

\[
L\dot i=Gv,
\qquad
m\dot v=-Gi.
\]

对第二式求导并代入第一式：

\[
\ddot v+\frac{G^2}{mL}v=0.
\]

因此

\[
\boxed{\omega_0=\frac{G}{\sqrt{mL}}=\frac{B\ell}{\sqrt{mL}}}.
\]

能量在

\[
\frac12mv^2
\quad\text{与}\quad
\frac12Li^2
\]

之间周期性交换。

\subsection{五、只含电阻：磁阻尼与终端速度}\label{m5:ux4e94ux53eaux542bux7535ux963bux78c1ux963bux5c3cux4e0eux7ec8ux7aefux901fux5ea6}

令 \(L=0\) 且无电容，则

\[
i=\frac{Gv}{R}.
\]

机械方程成为

\[
m\dot v=F_{\rm ext}-\frac{G^2}{R}v.
\]

电磁阻尼系数为

\[
\boxed{\gamma_{\rm em}=\frac{G^2}{R}}.
\]

若 \(F_{\rm ext}=F_0\) 为常量，则

\[
\boxed{
v(t)=v_{\infty}+(v_0-v_{\infty})e^{-t/\tau}},
\]

其中

\[
\boxed{v_{\infty}=\frac{F_0R}{G^2}},
\qquad
\boxed{\tau=\frac{mR}{G^2}}.
\]

若 \(F_{\rm ext}=0\)，速度指数衰减，但总位移有限：

\[
\boxed{x_{\infty}-x_0=v_0\tau=\frac{mRv_0}{B^2\ell^2}}.
\]

\par\noindent\begin{minipage}{\linewidth}
\begin{center}
\includegraphics[width=\linewidth,height=.40\textheight,keepaspectratio]{\detokenize{assets/module-05/theory/tikz-m5-topic51-v9-from48-v8-03.pdf}}
\par\smallskip{\small 图示：电容、电感、电阻对应三种动力学}
\end{center}
\end{minipage}\par

\subsection{六、电阻与电感并存：二阶暂态}\label{m5:ux516dux7535ux963bux4e0eux7535ux611fux5e76ux5b58ux4e8cux9636ux6682ux6001}

无电容时

\[
L\dot i+Ri=Gv.
\]

对该式求导，并用 \(m\dot v=F_{\rm ext}-Gi\)，得到

\[
\boxed{
L\ddot i+R\dot i+\frac{G^2}{m}i
=\frac{G}{m}F_{\rm ext}}.
\]

无外力时，这是阻尼振子方程。其固有角频率和阻尼参数为

\[
\omega_0=\frac{G}{\sqrt{mL}},
\qquad
2\beta=\frac RL.
\]

欠阻尼条件为

\[
\boxed{R<2G\sqrt{\frac Lm}}.
\]

这说明机械惯性与电感共同提供两个储能自由度，电阻则耗散总能量。

\subsection{七、行波磁场与同步速度}\label{m5:ux4e03ux884cux6ce2ux78c1ux573aux4e0eux540cux6b65ux901fux5ea6}

设外加磁场具有行波形式

\[
B(x,t)=B_0\cos(kx-\omega t).
\]

沿 \(+x\) 方向运动的回路感受到的相位变化率为

\[
\frac d{dt}(kx-\omega t)=kv-\omega.
\]

所以感应电动势与滑差频率

\[
\boxed{\omega_{\rm slip}=\omega-kv}
\]

成正比。同步速度为

\[
\boxed{v_s=\frac\omega k}.
\]

在 \(v=v_s\)
时，回路与磁场图样同速，磁通在回路参考系中不再变化，理想模型中的感应电流和平均电磁推力消失。实际感应电动机要产生推力，必须保留非零滑差。

\par\noindent\begin{minipage}{\linewidth}
\begin{center}
\includegraphics[width=\linewidth,height=.40\textheight,keepaspectratio]{\detokenize{assets/module-05/theory/tikz-m5-topic51-v9-from48-v8-04.pdf}}
\par\smallskip{\small 图示：行波磁场与同步速度}
\end{center}
\end{minipage}\par

\subsection{八、使用条件与常见错误}\label{m5:ux516bux4f7fux7528ux6761ux4ef6ux4e0eux5e38ux89c1ux9519ux8bef-1}

\begin{enumerate}
\def\labelenumi{\arabic{enumi}.}
\tightlist
\item
  电路方程和机械方程必须使用同一套方向约定，使 \(Gvi\)
  在总能量方程中抵消。
\item
  电感存在时，电流不能瞬时等于 \(Gv/R\)；必须保留暂态。
\item
  电容产生``等效惯性''只是在给定约束下的动力学等效，不是静质量改变。
\item
  无耗散的电感耦合系统一般振荡而不自动停下。
\item
  电阻磁阻尼力与速度成正比只适用于 \(B,\ell,R\) 恒定且准静态的滑杆模型。
\item
  行波磁场的同步点是零滑差点，不是最大推力点。
\end{enumerate}

\subsection{练习与解答}\label{m5:ux7ec3ux4e60ux4e0eux89e3ux7b54-7}

\par\addvspace{.4\baselineskip}
\noindent\begin{minipage}{\linewidth}

\subsubsection{作业一：电容滑杆的加速度}\label{m5:ux4f5cux4e1aux4e00ux7535ux5bb9ux6ed1ux6746ux7684ux52a0ux901fux5ea6}

\textbf{题号：M5-144}

\textbf{层级：} 复赛　　\textbf{方法：} 自感互感、机电耦合

导轨水平、固定且光滑，\(B\text{、}m\text{、}\ell\text{、}C\text{、}F_0>0\)，忽略电阻、自感、重力及端部效应；初始电容不带电，向右为正。

质量为 \(m\)、长度为 \(\ell\) 的滑杆处于垂直导轨平面的匀强磁场 \(B\)
中，回路只含电容 \(C\)。杆从静止开始受恒力
\(F_0\)，求加速度、速度和电容电荷。

\end{minipage}\par

\par\noindent\begin{minipage}{\linewidth}
\textbf{解答}

\begin{center}
\includegraphics[width=\linewidth,height=.40\textheight,keepaspectratio]{\detokenize{assets/module-05/examples/M5-144.pdf}}
\par\smallskip{\small 图 M5-144：电容滑杆的加速度}
\end{center}
\end{minipage}\par

图 M5-144
的杆向右加速，电流给电容充电并产生向左安培力；电容初始不带电，故速度和电荷都从零建立。

由 \(q/C=B\ell v\)，

\[
i=\dot q=CB\ell\dot v.
\]

机械方程为

\[
m\dot v=F_0-B\ell i
=F_0-CB^2\ell^2\dot v.
\]

\par\noindent\begin{minipage}{\linewidth}
\begin{center}
\includegraphics[width=\linewidth,height=.40\textheight,keepaspectratio]{\detokenize{assets/module-05/process/M5-144a.pdf}}
\par\smallskip{\small 图 M5-144a：匀加速同时建立机械能与电容能}
\end{center}
\end{minipage}\par

图 M5-144a
的线性速度意味着电荷也线性建立；输入功不仅增加杆的动能，也增加与速度平方成正比的电容场能。

所以

\[
\boxed{a=\dot v=\frac{F_0}{m+CB^2\ell^2}}.
\]

由初始静止，

\[
\boxed{v(t)=\frac{F_0t}{m+CB^2\ell^2}},
\]

\[
\boxed{q(t)=CB\ell v(t)
=\frac{CB\ell F_0t}{m+CB^2\ell^2}}.
\]

\par\addvspace{.4\baselineskip}
\noindent\begin{minipage}{\linewidth}

\subsubsection{作业二：无电阻电感滑杆的振荡}\label{m5:ux4f5cux4e1aux4e8cux65e0ux7535ux963bux7535ux611fux6ed1ux6746ux7684ux632fux8361}

\textbf{题号：M5-145}

\textbf{层级：} 复赛　　\textbf{方法：} 网络方程、自感互感、机电耦合

质量为 \(m\)、长度 \(\ell\)
的滑杆沿光滑水平导轨运动，匀强磁场垂直导轨，大小为 \(B\)。定义
\(G=B\ell>0\)，固定向右为 \(+x\)，正电流沿杆向上，对应安培力的有符号分量
\(F_B=-Gi\)；忽略电阻、重力和外力。

回路只含自感 \(L\)。初始时杆速度为 \(v_0\)、电流为零，随后无外力。求
\(v(t)\) 与 \(i(t)\)。

\end{minipage}\par

\par\noindent\begin{minipage}{\linewidth}
\textbf{解答}

\begin{center}
\includegraphics[width=\linewidth,height=.40\textheight,keepaspectratio]{\detokenize{assets/module-05/examples/M5-145.pdf}}
\par\smallskip{\small 图 M5-145：无电阻电感滑杆的振荡}
\end{center}
\end{minipage}\par

图 M5-145
的电流正方向使安培力阻碍正速度，机械能先转为电感磁能，再反向释放，形成无阻尼振荡。

速度、电流的正向不随实际运动反向而改变。无电阻时，\(m\dot v=-Gi\)、\(L\dot i=Gv\)，因此
\[
\frac{d}{dt}\left(\frac12mv^2+\frac12Li^2\right)=0.
\] 记 \(\omega_0=G/\sqrt{mL}\)。这是一种可逆的能量交换。例如
\(\pi/2<\omega_0t<\pi\) 时，\(v<0\)、\(i>0\)，磁力仍沿
\(-x\)、与瞬时速度同向，磁场储能正在返还机械能。不能把本题的磁力全程称为阻力。

方程为

\[
L\dot i=Gv,
\qquad
m\dot v=-Gi.
\]

令 \(\omega_0=G/\sqrt{mL}\)，由初始条件得到

\[
\boxed{v(t)=v_0\cos\omega_0t}.
\]

再由机械方程

\[
i=-\frac mG\dot v,
\]

\par\noindent\begin{minipage}{\linewidth}
\begin{center}
\includegraphics[width=\linewidth,height=.40\textheight,keepaspectratio]{\detokenize{assets/module-05/process/M5-145a.pdf}}
\par\smallskip{\small 图 M5-145a：速度和电流的相位对应}
\end{center}
\end{minipage}\par

图 M5-145a
按实际正负号画出归一化速度和电流，横轴就是零值线；两者相差四分之一周期，可直接判断能量交换方向。

故

\[
\boxed{i(t)=v_0\sqrt{\frac mL}\sin\omega_0t}.
\]

能量检验：

\[
\frac12mv^2+\frac12Li^2=\frac12mv_0^2.
\]

\par\addvspace{.4\baselineskip}
\noindent\begin{minipage}{\linewidth}

\subsubsection{作业三：电阻磁阻尼的停止距离}\label{m5:ux4f5cux4e1aux4e09ux7535ux963bux78c1ux963bux5c3cux7684ux505cux6b62ux8dddux79bb}

\textbf{题号：M5-146}

\textbf{层级：} 复赛　　\textbf{方法：} 网络方程、自感互感、机电耦合

滑杆质量 \(m\)、有效长度 \(\ell\)，磁场大小 \(B\) 垂直光滑水平导轨，定义
\(G=B\ell\)；忽略回路自感、杆轨电阻及重力，初始位置为零。

滑杆回路只有电阻 \(R\)。撤去外力时杆速度为
\(v_0\)。求速度随时间和最终位移。

\end{minipage}\par

\par\noindent\begin{minipage}{\linewidth}
\textbf{解答}

\begin{center}
\includegraphics[width=\linewidth,height=.40\textheight,keepaspectratio]{\detokenize{assets/module-05/examples/M5-146.pdf}}
\par\smallskip{\small 图 M5-146：电阻磁阻尼的停止距离}
\end{center}
\end{minipage}\par

图 M5-146
的电阻把杆的动能持续变成焦耳热；安培力与速度成正比且反向，所以速度指数衰减，停止距离有限。

\[
m\dot v=-\frac{B^2\ell^2}{R}v.
\]

定义

\[
\tau=\frac{mR}{B^2\ell^2},
\]

\par\noindent\begin{minipage}{\linewidth}
\begin{center}
\includegraphics[width=\linewidth,height=.40\textheight,keepaspectratio]{\detokenize{assets/module-05/process/M5-146a.pdf}}
\par\smallskip{\small 图 M5-146a：指数减速与有限停止距离}
\end{center}
\end{minipage}\par

图 M5-146a
的蓝色速度趋零而绿色位移趋有限水平线；速度曲线从零到无穷的面积就是最终位移。

则

\[
\boxed{v(t)=v_0e^{-t/\tau}}.
\]

总位移为

\[
x_{\infty}-x_0=\int_0^{\infty} v_0e^{-t/\tau}dt
=v_0\tau.
\]

所以

\[
\boxed{x_{\infty}-x_0=\frac{mRv_0}{B^2\ell^2}}.
\]

\par\addvspace{.4\baselineskip}
\noindent\begin{minipage}{\linewidth}

\subsubsection{作业四：含电阻与电感时的临界阻尼}\label{m5:ux4f5cux4e1aux56dbux542bux7535ux963bux4e0eux7535ux611fux65f6ux7684ux4e34ux754cux963bux5c3c}

\textbf{题号：M5-147}

\textbf{层级：} 复赛　　\textbf{方法：} 网络方程、自感互感、机电耦合

滑杆质量为 \(m\)、长度 \(\ell\)，磁场大小 \(B\)
垂直光滑水平导轨；回路串联理想电感 \(L>0\) 与总电阻 \(R\)，定义
\(G=B\ell\)，初始速度非零且电流为零。

无外力、无电容时，求使电流暂态达到临界阻尼的电阻。

\end{minipage}\par

\par\noindent\begin{minipage}{\linewidth}
\textbf{解答}

\begin{center}
\includegraphics[width=\linewidth,height=.40\textheight,keepaspectratio]{\detokenize{assets/module-05/examples/M5-147.pdf}}
\par\smallskip{\small 图 M5-147：含电阻与电感时的临界阻尼}
\end{center}
\end{minipage}\par

图 M5-147 中回路串联 \(R\) 与
\(L\)，滑杆仍构成电磁反馈。消去速度后由二阶方程的判别式给出临界阻尼。

电流满足

\[
L\ddot i+R\dot i+\frac{G^2}{m}i=0.
\]

特征方程为

\[
Lr^2+Rr+\frac{G^2}{m}=0.
\]

\par\noindent\begin{minipage}{\linewidth}
\begin{center}
\includegraphics[width=\linewidth,height=.40\textheight,keepaspectratio]{\detokenize{assets/module-05/process/M5-147a.pdf}}
\par\smallskip{\small 图 M5-147a：三种阻尼由判别式划分}
\end{center}
\end{minipage}\par

图 M5-147a
用无量纲阻尼比划分三个区域，临界点是特征根合并的位置；初始速度只影响振幅，不改变临界电阻。

临界阻尼要求判别式为零：

\[
R^2-4L\frac{G^2}{m}=0.
\]

故

\[
\boxed{R_c=2G\sqrt{\frac Lm}=2B\ell\sqrt{\frac Lm}}.
\]

\par\addvspace{.4\baselineskip}
\noindent\begin{minipage}{\linewidth}

\subsubsection{作业五：行波磁场的滑差}\label{m5:ux4f5cux4e1aux4e94ux884cux6ce2ux78c1ux573aux7684ux6ed1ux5dee}

\textbf{题号：M5-148}

\textbf{层级：} 复赛　　\textbf{方法：} 自感互感、机电耦合

\(k\text{、}\omega>0\)，定义波速 \(v_s=\omega/k\)，\(v\)
是实验室中回路的有符号速度。问观察频率时取绝对值，问滑差和相位推进时保留符号。

某周期磁场沿 \(+x\) 传播，\(B=B_0\cos(kx-\omega t)\)。回路以恒速 \(v\)
沿 \(+x\) 运动。求回路看到的角频率，并讨论
\(v<0\)、\(0<v<v_s\)、\(v=v_s\) 和 \(v>v_s\) 四种情况。

\end{minipage}\par

\par\noindent\begin{minipage}{\linewidth}
\textbf{解答}

\begin{center}
\includegraphics[width=\linewidth,height=.40\textheight,keepaspectratio]{\detokenize{assets/module-05/examples/M5-148.pdf}}
\par\smallskip{\small 图 M5-148：行波磁场的滑差}
\end{center}
\end{minipage}\par

图 M5-148 的红色回路跟随位置
\(x=vt+x_0\)，把它代入蓝色行波相位后，时间因子变为
\((\omega-kv)t\)；同步时相位不随时间变化。

沿回路轨迹 \(x=vt+x_0\)，相位为

\[
\theta(t)=(kv-\omega)t+kx_0.
\]

故回路看到的角频率大小为

\[
\boxed{|\omega-kv|}.
\]

当 \(v<0\) 时，回路逆向运动，频率大于 \(\omega\)；当 \(0<v<v_s\)
时，同向追赶使频率降低；当 \(v=v_s=\omega/k\) 时，磁场图样相对静止；当
\(v>v_s\) 时，相位变化反向，感应电流的相位推进方向也随之反转。

\par\addvspace{.4\baselineskip}
\noindent\begin{minipage}{\linewidth}

\subsubsection{作业六：含电阻与电容的滑杆：最终速度与电荷}\label{m5:ux4f5cux4e1aux516dux542bux7535ux963bux4e0eux7535ux5bb9ux7684ux6ed1ux6746ux6700ux7ec8ux901fux5ea6ux4e0eux7535ux8377}

\textbf{题号：M5-149}

\textbf{层级：} 复赛　　\textbf{方法：} 网络方程、自感互感、机电耦合

滑杆有效长度为 \(\ell\)，匀强磁场大小 \(B\) 垂直光滑水平导轨，定义
\(G=B\ell\)；所有电阻计入 \(R>0\)，向右为正，不计重力与辐射。

质量为 \(m\) 的滑杆初速度为 \(v_0\)，回路串联电阻 \(R\) 与电容
\(C\)，无电感、无外力，初始电容不带电。求最终速度、最终电荷和弛豫时间。

\end{minipage}\par

\par\noindent\begin{minipage}{\linewidth}
\textbf{解答}

\begin{center}
\includegraphics[width=\linewidth,height=.40\textheight,keepaspectratio]{\detokenize{assets/module-05/examples/M5-149.pdf}}
\par\smallskip{\small 图 M5-149：含电阻与电容的滑杆：最终速度与电荷}
\end{center}
\end{minipage}\par

图 M5-149 的电流给串联电容充电，同时减小杆的动量；将机械方程积分可得
\(mv+Gq\) 守恒，终态电流为零但速度一般非零。

方程为

\[
R\dot q+\frac qC=Gv,
\qquad
m\dot v=-G\dot q.
\]

第二式积分并用初始条件得

\[
mv+Gq=mv_0.
\]

终态 \(\dot q=0\)，故 \(q_{\infty}/C=Gv_{\infty}\)。联立得到

\[
\boxed{v_{\infty}=\frac{m}{m+CG^2}v_0},
\]

\[
\boxed{q_{\infty}=\frac{mCG}{m+CG^2}v_0}.
\]

\par\noindent\begin{minipage}{\linewidth}
\begin{center}
\includegraphics[width=\linewidth,height=.40\textheight,keepaspectratio]{\detokenize{assets/module-05/process/M5-149a.pdf}}
\par\smallskip{\small 图 M5-149a：RC滑杆趋向非零终端速度}
\end{center}
\end{minipage}\par

图 M5-149a
的蓝线扣除了终端速度；速度的减少量与绿色电荷增长量由同一个动量积分联系，不能把纵坐标零误读为杆最后静止。

把 \(v=v_0-Gq/m\) 代回电路方程：

\[
R\dot q+\left(\frac1C+\frac{G^2}{m}\right)q=Gv_0.
\]

所以

\[
\boxed{\tau=\frac{R}{1/C+G^2/m}
=\frac{mRC}{m+CG^2}}.
\]

\(R\to0^+\) 是奇异极限：对每个
\(R>0\)，上述终态都成立，且弛豫时间趋于零；但若直接令
\(R=0\)，电路方程退化为代数约束 \(q=CGv\)。当初态 \(q(0)=0\) 而
\(v(0)=v_0\ne0\)
时，该约束与初态不相容，理想模型要求瞬时电荷跃变。真实系统必须恢复寄生电感、辐射或有限开关时间，因而不能把
\(R=0\) 直接代入暂态解。

\par\addvspace{.4\baselineskip}
\noindent\begin{minipage}{\linewidth}

\subsubsection{\texorpdfstring{作业七：矩形回路进入磁场时的 \(RL\)
暂态}{作业七：矩形回路进入磁场时的 RL 暂态}}\label{m5:ux4f5cux4e1aux4e03ux77e9ux5f62ux56deux8defux8fdbux5165ux78c1ux573aux65f6ux7684-rl-ux6682ux6001}

\textbf{题号：M5-150}

\textbf{层级：} 复赛　　\textbf{方法：} 自感互感、机电耦合

矩形回路沿速度方向的宽度为 \(a\)，场区边界锐利且固定，只求 \(0<t<a/v\)
的进入阶段。磁场向纸内，取逆时针电流为正，定义 \(G=B\ell\)。

高为 \(\ell\) 的矩形回路以恒速 \(v\) 进入匀强磁场，回路电阻为
\(R\)、自感为 \(L\)。在进入阶段磁通以速率 \(B\ell v\) 增加，且
\(i(0)=0\)。求电流、磁阻力和外力功率，并写出能量平衡。

\end{minipage}\par

\par\noindent\begin{minipage}{\linewidth}
\textbf{解答}

\begin{center}
\includegraphics[width=\linewidth,height=.40\textheight,keepaspectratio]{\detokenize{assets/module-05/examples/M5-150.pdf}}
\par\smallskip{\small 图 M5-150：矩形回路进入磁场时的 \(RL\) 暂态}
\end{center}
\end{minipage}\par

图 M5-150
在进入阶段由外力保持恒速，电流从零建立。外力功率一部分变焦耳热，一部分建立磁场能，不能一开始就全当热。

取感应电流正向使其产生阻力，则

\[
L\dot i+Ri=B\ell v=Gv.
\]

解得

\[
\boxed{i(t)=\frac{Gv}{R}\left(1-e^{-Rt/L}\right)}.
\]

磁阻力为

\[
\boxed{F_B(t)=-Gi(t)
=-\frac{G^2v}{R}\left(1-e^{-Rt/L}\right)}.
\]

\par\noindent\begin{minipage}{\linewidth}
\begin{center}
\includegraphics[width=\linewidth,height=.40\textheight,keepaspectratio]{\detokenize{assets/module-05/process/M5-150a.pdf}}
\par\smallskip{\small 图 M5-150a：外功在磁能与焦耳热间分配}
\end{center}
\end{minipage}\par

图 M5-150a
中蓝线始终等于红、绿两线之和；最初电流小，外功中用于建立磁能的比例较大，随后焦耳热占主导。

维持匀速所需外力功率为

\[
P_{\rm ext}=-F_Bv=Gvi.
\]

由电路方程乘以 \(i\)，

\[
\boxed{P_{\rm ext}=Ri^2+\frac d{dt}\left(\frac12Li^2\right)}.
\]

进入初期外功主要建立磁场能；达到稳态后储能不再变化，外功全部转为焦耳热。

设回路沿运动方向的宽度为 \(a\)，以上恒定磁通变化率只适用于进入阶段
\(0<t<a/v\)。回路完全进入匀强磁场后，\(d\Phi/dt=0\)，电流改为

\[
i(t)=i(a/v)\exp\left[-\frac RL\left(t-\frac av\right)\right].
\]

\par\addvspace{.4\baselineskip}
\noindent\begin{minipage}{\linewidth}

\subsubsection{\texorpdfstring{作业八：重力、恒定电源与 \(RL\)
回路共同作用的滑杆}{作业八：重力、恒定电源与 RL 回路共同作用的滑杆}}\label{m5:ux4f5cux4e1aux516bux91cdux529bux6052ux5b9aux7535ux6e90ux4e0e-rl-ux56deux8defux5171ux540cux4f5cux7528ux7684ux6ed1ux6746}

\textbf{题号：M5-151}

\textbf{层级：} 复赛　　\textbf{方法：} 自感互感、机电耦合

有效杆长为 \(\ell\)，磁场大小为 \(B\)，定义
\(G=B\ell>0\)；\(m\text{、}R_t\text{、}L\text{、}g>0\)，初值为
\(v(0)=v_0\text{、}i(0)=i_0\)。只按给定参考方向求解，允许终端速度为负。

竖直导轨上的滑杆质量为 \(m\)，向下为正。回路总电阻为 \(R_t\)、自感为
\(L\)，另有恒定电源电动势 \(\mathcal E_0\)；方向约定使电路方程为

\[
\mathcal E_0+Gv=L\dot i+R_ti,
\]

机械方程为
\(m\dot v=mg-Gi\)。求终端电流和终端速度；给出任意初始条件下求最大速度的完整方法。

\end{minipage}\par

\par\noindent\begin{minipage}{\linewidth}
\textbf{解答}

\begin{center}
\includegraphics[width=\linewidth,height=.40\textheight,keepaspectratio]{\detokenize{assets/module-05/examples/M5-151.pdf}}
\par\smallskip{\small 图 M5-151：重力与恒定电源共同驱动的滑杆}
\end{center}
\end{minipage}\par

图 M5-151
的滑杆结构在本题整体取向下为机械正方向，橙色磁力反向；按题给方程求终态，再由
\(\dot v=0\) 的所有有限时刻与终态比较最大速度。

终态满足 \(\dot i=\dot v=0\)，所以

\[
\boxed{i_{\infty}=\frac{mg}{G}},
\]

\[
\boxed{v_{\infty}=\frac{R_tmg/G-\mathcal E_0}{G}}.
\]

对电路方程求导并消去 \(\dot v\)：

\[
\boxed{
L\ddot i+R_t\dot i+\frac{G^2}{m}i=Gg}.
\]

设特征根

\[
r_{1,2}=\frac{-R_t\pm\sqrt{R_t^2-4LG^2/m}}{2L},
\]

当判别式为正、\(r_1\ne r_2\) 时，

\[
i(t)=i_{\infty}+Ae^{r_1t}+Be^{r_2t}.
\]

\(A,B\) 由 \(i(0)\) 和

\[
\dot i(0)=\frac{\mathcal E_0+Gv(0)-R_ti(0)}L
\]

确定。速度由

\[
\boxed{v(t)=\frac{L\dot i+R_ti-\mathcal E_0}{G}}
\]

恢复。速度极值满足 \(\dot v=0\)，即

\[
\boxed{i(t_*)=\frac{mg}{G}}.
\]

若 \(r_1,r_2\) 为相异实根，便解

\[
Ae^{r_1t_*}+Be^{r_2t_*}=0
\]

并代回 \(v(t)\)。若判别式为零，令 \(r=-R_t/(2L)\)，通解应写成

\[
i(t)=i_{\infty}+(A+Bt)e^{rt}.
\]

此时速度极值由 \((A+Bt_*)e^{rt_*}=0\) 决定；若 \(B\ne0\)，候选时刻为
\(t_*=-A/B\)，还必须检查 \(t_*>0\) 且 \(\dot v\)
确实由正变负。若特征根为共轭复根，则把 \(i(t)\)
写成指数衰减的正弦形式，再取第一个满足 \(i=i_{\infty}\) 且 \(\dot v\)
由正变负的时刻。这样过阻尼、临界阻尼和可能出现超调的欠阻尼三种情形都没有遗漏。

为求全局最大值，还须比较初始速度
\(v_0\)、全部满足极大条件的有限时刻速度以及终端极限
\(v_{\infty}\)。若速度单调趋近更大的终端值但在有限时刻不达到，应明确给出上确界，而不虚构有限的最大时刻。

\par\addvspace{.4\baselineskip}
\noindent\begin{minipage}{\linewidth}

\subsubsection{作业九：恒力驱动柔性圆环扩张}\label{m5:ux4f5cux4e1aux4e5dux6052ux529bux9a71ux52a8ux67d4ux6027ux5706ux73afux6269ux5f20}

\textbf{题号：M5-152}

\textbf{层级：} 复赛　　\textbf{方法：} 自感互感、机电耦合

\(B\text{、}F\text{、}R\text{、}r_0\text{、}r_1>0\) 且
\(r_1>r_0\)；磁场恒定，正向取半径增大，\(F\)
是与周长对应的广义力，不是对半径作功的总力。

柔性圆形电阻回路处在垂直环面的匀强磁场中，电阻恒为
\(R\)。把沿圆周均匀施力的机械装置等效为广义力 \(F\)，使机械功为

\[
dW_{\rm mech}=F\,d(2\pi r).
\]

忽略惯性和自感，求半径从 \(r_0\) 增至 \(r_1\) 所需时间。

\end{minipage}\par

\par\noindent\begin{minipage}{\linewidth}
\textbf{解答}

\begin{center}
\includegraphics[width=\linewidth,height=.40\textheight,keepaspectratio]{\detokenize{assets/module-05/examples/M5-152.pdf}}
\par\smallskip{\small 图 M5-152：恒力驱动柔性圆环扩张}
\end{center}
\end{minipage}\par

图 M5-152 的绿色箭头表示扩张，磁阻力向内；广义力的功率是
\(F\,d(2\pi r)/dt\)，应与整个回路焦耳热相等。

磁通为 \(\Phi=B\pi r^2\)，电动势大小

\[
|\mathcal E|=2\pi Br\dot r.
\]

焦耳功率为

\[
P_J=\frac{\mathcal E^2}{R}
=\frac{4\pi^2B^2r^2\dot r^2}{R}.
\]

准静态时机械功率全部转为热：

\[
2\pi F\dot r=P_J.
\]

对 \(\dot r\ne0\)，

\[
\dot r=\frac{FR}{2\pi B^2r^2}.
\]

因此

\[
\boxed{
T=\frac{2\pi B^2}{3FR}(r_1^3-r_0^3)}.
\]

半径越大，同样的扩张速度引起的磁通变化越快，电磁阻力也越强。

\par\addvspace{.4\baselineskip}
\noindent\begin{minipage}{\linewidth}

\subsubsection{作业十：转动导电圆环的非线性磁阻尼}\label{m5:ux4f5cux4e1aux5341ux8f6cux52a8ux5bfcux7535ux5706ux73afux7684ux975eux7ebfux6027ux78c1ux963bux5c3c}

\textbf{题号：M5-153}

\textbf{层级：} 复赛　　\textbf{方法：} 自感互感、机电耦合

令
\(\theta(0)=0\)，即初始环面法向与外场同向，\(\omega(0)=\omega_0>0\)；转轴固定，圆环转动惯量为
\(J=mr^2/2\)。角速度关于角度的公式只在尚未停止转动的区间使用。

质量为 \(m\)、半径为 \(r\)、电阻为 \(R\)
的细圆环绕一条直径转动。匀强磁场 \(B\)
与转轴垂直。设环面法向与磁场夹角为 \(\theta\)，初始角速度为
\(\omega_0\)。忽略自感，求运动方程和角速度关于 \(\theta\) 的关系。

\end{minipage}\par

\par\noindent\begin{minipage}{\linewidth}
\textbf{解答}

\begin{center}
\includegraphics[width=\linewidth,height=.40\textheight,keepaspectratio]{\detokenize{assets/module-05/examples/M5-153.pdf}}
\par\smallskip{\small 图 M5-153：转动导电圆环的非线性磁阻尼}
\end{center}
\end{minipage}\par

图 M5-153 以 \(\theta=0\)
时法向平行外场为初始姿态，磁阻尼力矩反向于转动。用
\(\dot\omega=\omega\,d\omega/d\theta\)
在持续转动区间先对角度积分，再取渐近停转极限。

应区分角速度关于角度的零点与有限时刻停下。令
\(K=B^2\pi^2r^4/R\)，时间方程为
\(\dot\omega=-(K/J)\sin^2\theta\,\omega\)，故 \[
\omega(t)=\omega_0\exp\left[-\frac KJ\int_0^t\sin^2\theta(t')\,dt'\right]>0
\] 对任意有限 \(t\) 都成立。因此实际运动是
\(\theta\to\theta_*\)、\(\omega\to0\)
的渐近过程，不能越过角度公式的极限零点延拓为反转支路。

圆环磁通为

\[
\Phi=B\pi r^2\cos\theta.
\]

故

\[
i=\frac{B\pi r^2}{R}\dot\theta\sin\theta.
\]

磁矩为 \(\mu=i\pi r^2\)，阻尼力矩为

\[
\tau_B=-\mu B\sin\theta
=-\frac{B^2\pi^2r^4}{R}\dot\theta\sin^2\theta.
\]

圆环绕直径的转动惯量 \(J=mr^2/2\)，因此

\[
\boxed{
J\ddot\theta
=-\frac{B^2\pi^2r^4}{R}\dot\theta\sin^2\theta}.
\]

令 \(K=B^2\pi^2r^4/R\)，并用
\(\ddot\theta=\dot\theta\,d\dot\theta/d\theta\)，得到

\[
J\frac{d\dot\theta}{d\theta}=-K\sin^2\theta.
\]

从 \((\theta_0,\omega_0)\) 积分：

\[
\boxed{
\dot\theta(\theta)=\omega_0
-\frac{K}{2J}
\left[(\theta-\theta_0)-\frac12(\sin2\theta-\sin2\theta_0)\right]}.
\]

同时

\[
\frac d{dt}\left(\frac12J\dot\theta^2\right)=-i^2R,
\]

说明转动动能严格转化为焦耳热。

由 \(\ddot\theta=\dot\theta\,d\dot\theta/d\theta\) 得到的关系只在
\(\dot\theta\ne0\)、且仍在转动的区间有效。若 \(\omega(\theta)\)
的表达式在 \(\theta_*\)
有首个零点，它给出渐近停转的极限角度，不能理解为有限时刻突然停止后反转。

本题初始角度取 \(\theta_0=0\)，将其代入上式即得到主图中的简化表达。

\par\addvspace{.4\baselineskip}
\noindent\begin{minipage}{\linewidth}

\subsubsection{作业十一：矩形回路在行波磁场中的平均推力}\label{m5:ux4f5cux4e1aux5341ux4e00ux77e9ux5f62ux56deux8defux5728ux884cux6ce2ux78c1ux573aux4e2dux7684ux5e73ux5747ux63a8ux529b}

\textbf{题号：M5-154}

\textbf{层级：} 复赛　　\textbf{方法：} 自感互感、机电耦合

\(B_0\text{、}h\text{、}\ell\text{、}k\text{、}\omega\text{、}R>0\)，线框处于
\(xy\) 平面，竖边平行
\(y\)；无附加沿回路的静电电动势。需同时计入外场的感生项和线框运动项，不能只用
\(\mathbf v\times\mathbf B\)。

高为 \(h\)、沿 \(x\) 方向宽为 \(\ell\) 的矩形电阻回路以速度 \(v\)
运动，电阻为 \(R\)，自感忽略。外场为

\[
B_z(x,t)=B_0\cos(\omega t-kx).
\]

求平均电磁力。

\end{minipage}\par

\par\noindent\begin{minipage}{\linewidth}
\textbf{解答}

\begin{center}
\includegraphics[width=\linewidth,height=.40\textheight,keepaspectratio]{\detokenize{assets/module-05/examples/M5-154.pdf}}
\par\smallskip{\small 图 M5-154：矩形回路在行波磁场中的平均推力}
\end{center}
\end{minipage}\par

图 M5-154 的两个竖边分别采样行波的相位，两者相差
\(k\ell\)；动生项与感生项共同给出滑差 \(\Omega\)，平均力随滑差变号。

令左边位于 \(x(t)\)，相位 \(\phi=\omega t-kx\)。磁通为

\[
\Phi=\frac{hB_0}{k}
\left[\sin\phi-\sin(\phi-k\ell)\right].
\]

因为 \(\dot\phi=\omega-kv\)，

\[
\mathcal E
=-\frac{h(\omega-kv)}{k}
\left[B(x)-B(x+\ell)\right].
\]

两个竖边的合力可写成

\[
F_x=\frac{h^2(\omega-kv)}{kR}
\left[B(x+\ell)-B(x)\right]^2.
\]

\par\noindent\begin{minipage}{\linewidth}
\begin{center}
\includegraphics[width=\linewidth,height=.40\textheight,keepaspectratio]{\detokenize{assets/module-05/process/M5-154a.pdf}}
\par\smallskip{\small 图 M5-154a：行波平均推力在同步速度处变号}
\end{center}
\end{minipage}\par

图 M5-154a 的横轴交点是同步速度，图中取非零几何耦合因子；若
\(\sin(k\ell/2)=0\)，整个回路的平均推力恒为零。

一个周期内

\[
\left\langle[B(x+\ell)-B(x)]^2\right\rangle
=2B_0^2\sin^2\frac{k\ell}{2}.
\]

故

\[
\boxed{
\langle F_x\rangle
=\frac{2h^2B_0^2}{kR}
(\omega-kv)\sin^2\frac{k\ell}{2}}.
\]

当 \(v<v_s\) 时力沿波传播方向；当 \(v>v_s\) 时力反向；\(v=v_s=\omega/k\)
时平均力为零。

\section{专题五十二：涡流、磁扩散、趋肤效应与超导体（拔高选讲与决赛拓展）}\label{m5:ux4e13ux9898ux4e94ux5341ux4e8cux6da1ux6d41ux78c1ux6269ux6563ux8d8bux80a4ux6548ux5e94ux4e0eux8d85ux5bfcux4f53ux62d4ux9ad8ux9009ux8bb2ux4e0eux51b3ux8d5bux62d3ux5c55}

\begin{quote}
阅读准备：扩散与传播的时间尺度不同 扩散方程
\(\partial B/\partial t=D_m\nabla^2B\)
表示局部曲率驱动场逐渐摊平，特征穿透距离按 \(\sqrt{D_mt}\)
增长；它不像无耗波动那样以固定速度搬运波形。正弦稳态中用复指数只是在同时记录振幅衰减和相位滞后。
\end{quote}

\par\addvspace{.4\baselineskip}
\noindent\begin{minipage}{\linewidth}

\textbf{例 M5-155：磁扩散深度的时间标度}

在磁导率 \(\mu\)、电导率 \(\sigma\)
为常量的普通导体中，忽略位移电流并采用磁扩散近似。穿透深度尺度
\(d(t)\sim\sqrt{D_mt}\)、\(D_m=1/(\mu\sigma)\)。若经历时间增为四倍，深度变为几倍？解释良导体的短时穿透为何较浅。

\end{minipage}\par

\par\noindent\begin{minipage}{\linewidth}
\textbf{解答}

\begin{center}
\includegraphics[width=\linewidth,height=.40\textheight,keepaspectratio]{\detokenize{assets/module-05/examples/M5-155.pdf}}
\par\smallskip{\small 图 M5-155：磁扩散深度的时间标度}
\end{center}
\end{minipage}\par

图 M5-155
的两条衰减曲线示意不同穿透尺度；扩散过程的精确剖面随边界条件变化，但时间尺度关系始终是长度正比于时间平方根。

若经历时间扩大为原来的四倍，扩散深度约扩大为两倍；良导体中
\(D_m=1/(\mu\sigma)\) 较小，所以突然施加的磁场短时间内只进入表层。

变化磁场进入普通导体后会驱动涡流，涡流又反过来改变磁场的空间分布。普通导体中的``穿透''是带耗散的扩散过程；超导体中的迈斯纳效应则来自不同的物理方程。把二者区分清楚，是理解磁屏蔽、趋肤效应、电磁制动和磁悬浮的关键。

\subsection{一、普通导体中的磁扩散方程}\label{m5:ux4e00ux666eux901aux5bfcux4f53ux4e2dux7684ux78c1ux6269ux6563ux65b9ux7a0b}

在均匀线性导体中取

\[
\mathbf J=\sigma\mathbf E,
\qquad
\mathbf B=\mu\mathbf H.
\]

若满足磁准静态条件，可忽略位移电流：

\[
\nabla\times\mathbf B=\mu\mathbf J.
\]

对上式取旋度，并使用

\[
\nabla\times\nabla\times\mathbf B
=\nabla(\nabla\cdot\mathbf B)-\nabla^2\mathbf B
=-\nabla^2\mathbf B
\]

以及

\[
\nabla\times\mathbf E=-\frac{\partial\mathbf B}{\partial t},
\]

得到

\[
-\nabla^2\mathbf B
=\mu\sigma\nabla\times\mathbf E
=-\mu\sigma\frac{\partial\mathbf B}{\partial t}.
\]

因此

\[
\boxed{
\frac{\partial\mathbf B}{\partial t}
=D_m\nabla^2\mathbf B},
\qquad
\boxed{D_m=\frac1{\mu\sigma}}.
\]

这是扩散方程而不是波动方程。尺度为 \(a\)
的导体中，磁场重新分布的特征时间约为

\[
\boxed{\tau_m\sim\mu\sigma a^2}.
\]

\par\noindent\begin{minipage}{\linewidth}
\begin{center}
\includegraphics[width=\linewidth,height=.40\textheight,keepaspectratio]{\detokenize{assets/module-05/theory/M5-R20-diffusion.pdf}}
\par\smallskip{\small 图示：外加磁场随时间扩散进入普通导体}
\end{center}
\end{minipage}\par

\par\addvspace{.4\baselineskip}
\noindent\begin{minipage}{\linewidth}

\textbf{例 M5-156：导电率越大，短时间内穿透越慢}

比较几何、磁导率及施场历史相同但电导率不同的两块普通导体。在线性磁扩散近似下，判断较大电导率对应的扩散系数及同一短时刻的穿透深度，说明这种暂态排斥是否意味着长期完全排磁。

\end{minipage}\par

\par\noindent\begin{minipage}{\linewidth}
\textbf{解答}

\begin{center}
\includegraphics[width=\linewidth,height=.40\textheight,keepaspectratio]{\detokenize{assets/module-05/examples/M5-156.pdf}}
\par\smallskip{\small 图 M5-156：导电率越大，短时间内穿透越慢}
\end{center}
\end{minipage}\par

图 M5-156
的蓝色短穿透尺度对应较大的电导率，绿色较长尺度对应较小电导率；普通导体短时间的屏蔽不能等同于静态迈斯纳效应。

\(D_m=1/(\mu\sigma)\)，所以 \(\sigma\)
越大，磁扩散系数越小。高导体在突然加场后的短时间内表现出较强排斥，但只要电阻非零，涡流会衰减，磁场最终仍可进入。

\subsection{二、正弦场与趋肤深度}\label{m5:ux4e8cux6b63ux5f26ux573aux4e0eux8d8bux80a4ux6df1ux5ea6}

考虑导体占据 \(x>0\) 的半空间，外部施加沿 \(z\) 方向的正弦磁场。令复表示

\[
\mathbf B(x,t)=\Re\{\widetilde B(x)e^{i\omega t}\}\mathbf e_z.
\]

扩散方程给出

\[
\frac{d^2\widetilde B}{dx^2}=i\mu\sigma\omega\widetilde B.
\]

取随 \(x\) 衰减的根，得到

\[
\boxed{
\widetilde B(x)=B_0e^{-(1+i)x/\delta}},
\qquad
\boxed{\delta=\sqrt{\frac{2}{\mu\sigma\omega}}}.
\]

实场为

\[
\boxed{
B(x,t)=B_0e^{-x/\delta}
\cos\left(\omega t-\frac{x}{\delta}\right)}.
\]

\(\delta\) 是趋肤深度：振幅在深度 \(\delta\) 处降为表面的
\(e^{-1}\)，相位同时滞后 \(1\)
弧度。频率越高、导电率越大、磁导率越大，电流越集中在表面。

\par\noindent\begin{minipage}{\linewidth}
\begin{center}
\includegraphics[width=\linewidth,height=.40\textheight,keepaspectratio]{\detokenize{assets/module-05/theory/tikz-m5-topic52-v9-from49-v8-02.pdf}}
\par\smallskip{\small 图示：趋肤场的振幅衰减与相位滞后}
\end{center}
\end{minipage}\par

\subsection{三、涡流耗散与电磁制动}\label{m5:ux4e09ux6da1ux6d41ux8017ux6563ux4e0eux7535ux78c1ux5236ux52a8}

涡流的体功率密度为

\[
\boxed{p_J=\mathbf J\cdot\mathbf E=\frac{J^2}{\sigma}}.
\]

楞次定律决定涡流产生的磁场抵抗原磁通变化，因此运动导体通常受到与相对运动相反的平均力。机械功率最终转化为焦耳热：

\[
\boxed{-\langle\mathbf F_{\rm em}\rangle\cdot\mathbf v
=\int_V\left\langle\frac{J^2}{\sigma}\right\rangle dV}
\]

在忽略暂态储能变化时成立。

为了减小低频涡流损耗，铁芯常沿涡流可能闭合的方向切成彼此绝缘的薄片。薄片缩短了涡流回路的有效面积，并增大其电阻。高频下还需同时考虑趋肤效应、邻近效应和磁滞损耗。

\par\addvspace{.4\baselineskip}
\noindent\begin{minipage}{\linewidth}

\textbf{例 M5-157：为什么完整金属盘制动明显}

完整金属圆盘在非均匀磁场中转动会受到磁制动力矩。若沿径向开若干不贯通轮轴的绝缘窄槽，定性比较制动力矩变化，并说明电流路径、焦耳热与机械能之间的联系。

\end{minipage}\par

\par\noindent\begin{minipage}{\linewidth}
\textbf{解答}

\begin{center}
\includegraphics[width=\linewidth,height=.40\textheight,keepaspectratio]{\detokenize{assets/module-05/examples/M5-157.pdf}}
\par\smallskip{\small 图 M5-157：为什么完整金属盘制动明显}
\end{center}
\end{minipage}\par

图 M5-157
左侧完整圆盘允许较大的闭合涡流；右侧径向绝缘窄槽从盘缘切入而保留轮轴，阻碍跨槽电流，缩小或破坏原有闭合路径。

金属盘进入非均匀磁场时，不同位置磁通随运动变化，盘内形成闭合涡流。涡流与外磁场作用产生反向力矩；若把盘沿径向开槽，许多闭合电流路径被切断，制动力矩显著减小。

\subsection{四、理想导体与超导体不是同一概念}\label{m5:ux56dbux7406ux60f3ux5bfcux4f53ux4e0eux8d85ux5bfcux4f53ux4e0dux662fux540cux4e00ux6982ux5ff5}

若把普通导体形式地取到 \(\sigma\to\infty\)，由欧姆定律可有
\(\mathbf E\to0\)。法拉第定律给出

\[
\frac{\partial\mathbf B}{\partial t}=0.
\]

这表示磁场被``冻结''在初始状态：初始有磁场时，它不必自动排出。理想导体强调\textbf{磁通保持}。

超导体在进入超导态时会把体内磁场排斥到伦敦穿透深度范围内，即迈斯纳效应。它强调的是热力学平衡态中的\textbf{磁通排斥}，不能仅由``电阻为零''推出。

\subsection{五、伦敦方程与迈斯纳效应}\label{m5:ux4e94ux4f26ux6566ux65b9ux7a0bux4e0eux8fc8ux65afux7eb3ux6548ux5e94}

伦敦模型写成

\[
\frac{\partial\mathbf J_s}{\partial t}
=\frac{n_sq^2}{m}\mathbf E,
\]

\[
\boxed{
\nabla\times\mathbf J_s
=-\frac{n_sq^2}{m}\mathbf B
=-\frac{1}{\mu_0\lambda_L^2}\mathbf B},
\]

其中

\[
\boxed{\lambda_L=\sqrt{\frac{m}{\mu_0n_sq^2}}}
\]

是伦敦穿透深度。由 \(\nabla\times\mathbf B=\mu_0\mathbf J_s\)
取旋度，得到

\[
-\nabla^2\mathbf B
=\mu_0\nabla\times\mathbf J_s
=-\frac{\mathbf B}{\lambda_L^2},
\]

即

\[
\boxed{\nabla^2\mathbf B=\frac{\mathbf B}{\lambda_L^2}}.
\]

对占据 \(x>0\) 的超导半空间，若表面磁场为 \(B_0\mathbf e_z\)，则

\[
\boxed{B_z(x)=B_0e^{-x/\lambda_L}}.
\]

与普通导体的趋肤场不同，这里没有由频率决定的相位传播；\(\lambda_L\)
是超导凝聚态的材料长度尺度。

\par\noindent\begin{minipage}{\linewidth}
\begin{center}
\includegraphics[width=\linewidth,height=.40\textheight,keepaspectratio]{\detokenize{assets/module-05/original/tikz-m5-topic52-v9-from49-v8-04.png}}
\par\smallskip{\small 图示：普通导体、理想导体与超导体对磁场的不同响应}
\end{center}
\end{minipage}\par

\subsection{六、超导表面电流与磁压力}\label{m5:ux516dux8d85ux5bfcux8868ux9762ux7535ux6d41ux4e0eux78c1ux538bux529b}

由安培定律，若 \(\mathbf B=B_z(x)\mathbf e_z\)，则

\[
J_{s,y}=-\frac1{\mu_0}\frac{dB_z}{dx}
=\frac{B_0}{\mu_0\lambda_L}e^{-x/\lambda_L}.
\]

表面附近的超电流在厚度约 \(\lambda_L\) 的层内流动，其面电流为

\[
K_y=\int_0^{\infty} J_{s,y}\,dx
=\frac{B_0}{\mu_0}.
\]

该电流产生的磁场抵消体内外场。理想完全抗磁状态下，表面承受的磁压力数量级仍为

\[
\boxed{p_m=\frac{B_0^2}{2\mu_0}}.
\]

真实超导体还受临界温度、临界磁场和临界电流限制；第二类超导体在混合态允许量子化磁通涡旋进入，不能始终按完全排斥处理。

\subsection{七、使用条件与常见错误}\label{m5:ux4e03ux4f7fux7528ux6761ux4ef6ux4e0eux5e38ux89c1ux9519ux8bef-1}

\begin{enumerate}
\def\labelenumi{\arabic{enumi}.}
\tightlist
\item
  磁扩散方程建立在磁准静态近似下；频率高到位移电流或辐射不可忽略时，应回到完整麦克斯韦方程。
\item
  趋肤深度描述正弦稳态的幅度与相位，不能直接替代阶跃场的扩散长度
  \(\sim\sqrt{D_mt}\)。
\item
  高导电率意味着磁场短时间内进入较慢，不等于静态磁场永远不能进入。
\item
  理想导体保持初始磁通，超导体表现迈斯纳排斥；二者不能混为一谈。
\item
  伦敦方程是宏观唯象方程，不是由经典欧姆定律取 \(R=0\) 得到。
\item
  涡流制动力的方向由能量耗散和楞次定律共同约束，但具体大小仍需解电流分布。
\end{enumerate}

\subsection{练习与解答}\label{m5:ux7ec3ux4e60ux4e0eux89e3ux7b54-8}

\par\addvspace{.4\baselineskip}
\noindent\begin{minipage}{\linewidth}

\subsubsection{作业一：磁扩散时间尺度}\label{m5:ux4f5cux4e1aux4e00ux78c1ux6269ux6563ux65f6ux95f4ux5c3aux5ea6}

\textbf{题号：M5-158}

\textbf{层级：} 决赛　　\textbf{方法：} 磁场与感应

采用均匀线性普通导体的准静态磁扩散模型，忽略位移电流，要求量级而非依赖具体边界条件的数值系数。

一块厚度尺度为 \(a\) 的均匀导体，磁导率为 \(\mu\)、电导率为
\(\sigma\)。用量纲分析估计外磁场改变后内部磁场重新分布的时间尺度。

\end{minipage}\par

\par\noindent\begin{minipage}{\linewidth}
\textbf{解答}

\begin{center}
\includegraphics[width=\linewidth,height=.40\textheight,keepaspectratio]{\detokenize{assets/module-05/examples/M5-158.pdf}}
\par\smallskip{\small 图 M5-158：磁扩散时间尺度}
\end{center}
\end{minipage}\par

图 M5-158 的横向尺度在本题取为样品厚度
\(a\)，磁扩散方程中的二阶空间导数带来
\(a^{-2}\)，由此得到重新分布的时间尺度。

磁扩散方程为

\[
\frac{\partial B}{\partial t}
=\frac1{\mu\sigma}\nabla^2B.
\]

估计

\[
\frac{B}{\tau_m}\sim\frac1{\mu\sigma}\frac{B}{a^2},
\]

故

\[
\boxed{\tau_m\sim\mu\sigma a^2}.
\]

尺寸增大一倍，扩散时间约增为四倍。

\par\addvspace{.4\baselineskip}
\noindent\begin{minipage}{\linewidth}

\subsubsection{作业二：趋肤场的检验}\label{m5:ux4f5cux4e1aux4e8cux8d8bux80a4ux573aux7684ux68c0ux9a8c}

\textbf{题号：M5-159}

\textbf{层级：} 决赛　　\textbf{方法：} 磁场与感应

普通导体占据 \(x>0\)，\(\mu\text{、}\sigma\text{、}\omega>0\)
为常量，忽略位移电流；给定式为 \(z\) 向磁场，界面驱动为
\(B_0\cos\omega t\)。

验证

\[
B(x,t)=B_0e^{-x/\delta}\cos(\omega t-x/\delta)
\]

满足一维磁扩散方程，并求 \(\delta\)。

\end{minipage}\par

\par\noindent\begin{minipage}{\linewidth}
\textbf{解答}

\begin{center}
\includegraphics[width=\linewidth,height=.40\textheight,keepaspectratio]{\detokenize{assets/module-05/examples/M5-159.pdf}}
\par\smallskip{\small 图 M5-159：趋肤场的检验}
\end{center}
\end{minipage}\par

图 M5-159 的蓝线是 \(t=0\) 瞬时场，绿线是幅值包络；\(x=\delta\)
时相位仅滞后一弧度，瞬时值仍为正，随后才过零反号。

使用复场

\[
\widetilde B=B_0e^{-(1+i)x/\delta}e^{i\omega t}.
\]

有

\[
\frac{\partial\widetilde B}{\partial t}=i\omega\widetilde B,
\]

\par\noindent\begin{minipage}{\linewidth}
\begin{center}
\includegraphics[width=\linewidth,height=.40\textheight,keepaspectratio]{\detokenize{assets/module-05/process/M5-159a.pdf}}
\par\smallskip{\small 图 M5-159a：趋肤场同时衰减并随深度滞相}
\end{center}
\end{minipage}\par

图 M5-159a
的蓝色瞬时场会经过零点并反号，绿色包络仅表明振幅衰减；不能把包络误画成整个时变场。

\[
\frac{\partial^2\widetilde B}{\partial x^2}
=\frac{(1+i)^2}{\delta^2}\widetilde B
=\frac{2i}{\delta^2}\widetilde B.
\]

代入

\[
\frac{\partial\widetilde B}{\partial t}
=\frac1{\mu\sigma}\frac{\partial^2\widetilde B}{\partial x^2}
\]

得

\[
i\omega=\frac{2i}{\mu\sigma\delta^2}.
\]

所以

\[
\boxed{\delta=\sqrt{\frac2{\mu\sigma\omega}}}.
\]

\par\addvspace{.4\baselineskip}
\noindent\begin{minipage}{\linewidth}

\subsubsection{作业三：超导半空间的电流分布}\label{m5:ux4f5cux4e1aux4e09ux8d85ux5bfcux534aux7a7aux95f4ux7684ux7535ux6d41ux5206ux5e03}

\textbf{题号：M5-160}

\textbf{层级：} 决赛　　\textbf{方法：} 分布与奇性

取真空磁导率 \(\mu_0\)，\(x\) 为入内法向、\(z\)
为切向，超导体足够厚。总面电流定义为 \(K_y=\int_0^{\infty} J_y(x)dx\)。

超导体占据 \(x>0\)，内部磁场为
\(B_z(x)=B_0e^{-x/\lambda_L}\)。求超电流密度和总面电流。

\end{minipage}\par

\par\noindent\begin{minipage}{\linewidth}
\textbf{解答}

\begin{center}
\includegraphics[width=\linewidth,height=.40\textheight,keepaspectratio]{\detokenize{assets/module-05/examples/M5-160.pdf}}
\par\smallskip{\small 图 M5-160：超导半空间的电流分布}
\end{center}
\end{minipage}\par

图 M5-160 的 \(x\) 沿超导体内法向，\(B_z\)
随深度衰减；旋度方程中的负导数给出正 \(y\)
方向超电流，沿深度积分得总面电流。

由

\[
\nabla\times\mathbf B=\mu_0\mathbf J_s,
\]

且 \(\mathbf B=B_z(x)\mathbf e_z\)，有

\[
J_{s,y}=-\frac1{\mu_0}\frac{dB_z}{dx}
=\frac{B_0}{\mu_0\lambda_L}e^{-x/\lambda_L}.
\]

因此

\[
\boxed{\mathbf J_s(x)
=\frac{B_0}{\mu_0\lambda_L}e^{-x/\lambda_L}\mathbf e_y}.
\]

积分得面电流

\[
\boxed{
\mathbf K_s=\int_0^{\infty}\mathbf J_s\,dx
=\frac{B_0}{\mu_0}\mathbf e_y}.
\]

\par\addvspace{.4\baselineskip}
\noindent\begin{minipage}{\linewidth}

\subsubsection{作业四：理想导体与超导体的历史差别}\label{m5:ux4f5cux4e1aux56dbux7406ux60f3ux5bfcux4f53ux4e0eux8d85ux5bfcux4f53ux7684ux5386ux53f2ux5deeux522b}

\textbf{题号：M5-161}

\textbf{层级：} 决赛　　\textbf{方法：} 磁场与感应

比较相同有限样品的两种理想化历史；迈斯纳态要求外场低于适用临界场、无陷获磁通，样品尺度远大于穿透深度，所说内部场指远离表面的体内场。

样品最初处于均匀磁场 \(B_0\)
中。分别讨论把普通导体电导率理想化为无穷大、以及把样品冷却进入迈斯纳态后，体内磁场的结果。

\end{minipage}\par

\par\noindent\begin{minipage}{\linewidth}
\textbf{解答}

\begin{center}
\includegraphics[width=\linewidth,height=.40\textheight,keepaspectratio]{\detokenize{assets/module-05/examples/M5-161.pdf}}
\par\smallskip{\small 图 M5-161：理想导体与超导体的历史差别}
\end{center}
\end{minipage}\par

图 M5-161
左样品保留原有穿透场，右样品体内为零场；区别在于材料状态与历史，而非把电阻都写成零后就认为二者等价。

理想导体中 \(\mathbf E=0\)，法拉第定律给出

\[
\frac{\partial\mathbf B}{\partial t}=0.
\]

因此初始磁场被保持，体内仍可有
\(B_0\)。超导体进入迈斯纳平衡态后，体内磁场在距表面若干个 \(\lambda_L\)
之外指数趋零。故

\[
\boxed{
\text{理想导体：保持初始磁通；}
\qquad
\text{超导体：迈斯纳排斥。}}
\]

\par\addvspace{.4\baselineskip}
\noindent\begin{minipage}{\linewidth}

\subsubsection{作业五：薄片化为何减小涡流}\label{m5:ux4f5cux4e1aux4e94ux8584ux7247ux5316ux4e3aux4f55ux51cfux5c0fux6da1ux6d41}

\textbf{题号：M5-162}

\textbf{层级：} 决赛　　\textbf{方法：} 磁场与感应

薄片取向须使主磁通可沿片内通过而跨片涡流受到绝缘层阻断；讨论低频工程近似，不要求所有损耗都降为零。

定性说明把铁芯沿适当方向分割成彼此绝缘的薄片，为何能降低涡流损耗；同时说明这一措施不能消除哪两类损耗。

\end{minipage}\par

\par\noindent\begin{minipage}{\linewidth}
\textbf{解答}

\begin{center}
\includegraphics[width=\linewidth,height=.40\textheight,keepaspectratio]{\detokenize{assets/module-05/examples/M5-162.pdf}}
\par\smallskip{\small 图 M5-162：薄片化为何减小涡流}
\end{center}
\end{minipage}\par

图 M5-162
的绝缘间隙阻断跨片大电流回路；薄片本身仍可能有局部涡流，磁滞与绕组电阻也不会随叠片自动消失。

绝缘层切断了大尺度闭合涡流回路，使每个回路包围的变化磁通减小，并使回路电阻增大，因此涡流幅度和
\(J^2/\sigma\)
损耗都下降。薄片化不能消除铁磁材料本身的磁滞损耗，也不能完全消除高频下由趋肤效应和邻近效应引起的非均匀电流损耗。

\par\addvspace{.4\baselineskip}
\noindent\begin{minipage}{\linewidth}

\subsubsection{作业六：有微小电阻的闭合线圈如何释放磁能}\label{m5:ux4f5cux4e1aux516dux6709ux5faeux5c0fux7535ux963bux7684ux95edux5408ux7ebfux5708ux5982ux4f55ux91caux653eux78c1ux80fd}

\textbf{题号：M5-163}

\textbf{层级：} 决赛　　\textbf{方法：} 电路暂态、磁能

\(L\text{、}R>0\)
为常量，初始储能只有自感磁能；无外加电源、变化外磁通与机械运动，忽略辐射和寄生电容。

自感为 \(L\)、电阻为 \(R\) 的闭合线圈初始电流为
\(I_0\)，随后与外界无电磁耦合。求电流和累计焦耳热。

\end{minipage}\par

\par\noindent\begin{minipage}{\linewidth}
\textbf{解答}

\begin{center}
\includegraphics[width=\linewidth,height=.40\textheight,keepaspectratio]{\detokenize{assets/module-05/examples/M5-163.pdf}}
\par\smallskip{\small 图 M5-163：有微小电阻的闭合线圈如何释放磁能}
\end{center}
\end{minipage}\par

图 M5-163
的电感在本题通过有限电阻闭合放电；电流正向不变而端电压反向，减少的全部初始磁能最终变为焦耳热。

回路方程为

\[
L\dot i+Ri=0.
\]

所以

\[
\boxed{i(t)=I_0e^{-tR/L}}.
\]

从 \(0\) 到 \(t\) 的焦耳热为

\[
Q(t)=\int_0^tRi^2dt
=RI_0^2\int_0^te^{-2Rt'/L}dt'.
\]

故

\[
\boxed{Q(t)=\frac12LI_0^2\left(1-e^{-2Rt/L}\right)}.
\]

当 \(t\to\infty\)，

\[
Q(\infty)=\frac12LI_0^2,
\]

恰等于初始磁场能。只有严格 \(R=0\) 时电流才可永久保持。

\par\addvspace{.4\baselineskip}
\noindent\begin{minipage}{\linewidth}

\subsubsection{作业七：定电流近似下超导平面上方大电流环的悬浮与小振动}\label{m5:ux4f5cux4e1aux4e03ux5b9aux7535ux6d41ux8fd1ux4f3cux4e0bux8d85ux5bfcux5e73ux9762ux4e0aux65b9ux5927ux7535ux6d41ux73afux7684ux60acux6d6eux4e0eux5c0fux632fux52a8}

\textbf{题号：M5-164}

\textbf{层级：} 决赛　　\textbf{方法：} 磁场与感应

平面取 \(z=0\)，圆环在 \(z=h>0\)，近似镜像环位于 \(z=-h\)
且电流反向；只计磁斥力与重力，忽略恒流装置的机械影响。

半径为 \(r\)、质量为 \(M\) 的水平圆环位于完全抗磁平面上方，高度为
\(h\)，且
\(r\gg h\)。在所研究的小振幅范围内，环中电流由理想恒流装置维持为
\(I\)，或其相对变化可以忽略。用镜像电流近似求平衡高度和竖直小振动周期。

\end{minipage}\par

\par\noindent\begin{minipage}{\linewidth}
\textbf{解答}

\begin{center}
\includegraphics[width=\linewidth,height=.40\textheight,keepaspectratio]{\detokenize{assets/module-05/examples/M5-164.pdf}}
\par\smallskip{\small 图 M5-164：恒流大圆环在抗磁平面上方的悬浮}
\end{center}
\end{minipage}\par

图 M5-164
中镜像电流与实环反向，近邻的长直电流段彼此排斥；先求整体向上力，再在恒流条件下对高度求导得到小振动刚度。

完全抗磁平面可用位于平面下方、方向相反的镜像电流环代替。两环间距为
\(2h\)。在 \(h\ll r\)
时，可把相对的短弧近似为反向平行电流，单位长度斥力约为

\[
f=\frac{\mu_0I^2}{2\pi(2h)}.
\]

乘以圆周 \(2\pi r\)，上推力为

\[
F_m(h)=\frac{\mu_0I^2r}{2h}.
\]

平衡条件 \(F_m(h_0)=Mg\) 给出

\[
\boxed{h_0=\frac{\mu_0I^2r}{2Mg}}.
\]

\par\noindent\begin{minipage}{\linewidth}
\begin{center}
\includegraphics[width=\linewidth,height=.40\textheight,keepaspectratio]{\detokenize{assets/module-05/process/M5-164a.pdf}}
\par\smallskip{\small 图 M5-164a：恒流镜像斥力与重力的交点}
\end{center}
\end{minipage}\par

图 M5-164a
的下降磁斥力曲线与水平重力线相交；在交点上移时斥力减小，产生向下恢复力，因此恒流近似下竖直平衡稳定。

令 \(h=h_0+\xi\)，则

\[
F_m(h)\simeq Mg\left(1-\frac\xi{h_0}\right).
\]

所以

\[
M\ddot\xi=-\frac{Mg}{h_0}\xi,
\]

从而

\[
\boxed{\omega_0=\sqrt{\frac g{h_0}}},
\qquad
\boxed{T=2\pi\sqrt{\frac{h_0}{g}}}.
\]

该结果同时依赖 \(r\gg h\)
的局部平行线近似和定电流近似；一般高度应使用完整互感 \(M(h)\) 求力。

若圆环是与外界隔离的超导闭合环，守恒量是总磁通或更一般的磁通子，而不是电流。应写成

\[
\Psi_0=L_{\rm eff}(h)I(h),
\qquad
U(h)=\frac{\Psi_0^2}{2L_{\rm eff}(h)}+Mgh,
\]

再由 \(F_h=-dU/dh\) 求平衡，并由 \(U''(h_0)/M\)
求小振动角频率；不能直接套用上面的定电流结果。

\par\addvspace{.4\baselineskip}
\noindent\begin{minipage}{\linewidth}

\subsubsection{作业八：磁偶极子在完全抗磁平面上方的悬浮}\label{m5:ux4f5cux4e1aux516bux78c1ux5076ux6781ux5b50ux5728ux5b8cux5168ux6297ux78c1ux5e73ux9762ux4e0aux65b9ux7684ux60acux6d6e}

\textbf{题号：M5-165}

\textbf{层级：} 决赛　　\textbf{方法：} 多极展开

完全抗磁半空间为 \(z<0\)，小磁体可视为固定点偶极子，尺寸远小于
\(h\)，磁矩不随高度变化；忽略其他磁源，重力加速度为 \(g\)。

质量为 \(M\)、磁矩为 \(\boldsymbol\mu=\mu\mathbf e_z\)
的小磁体位于完全抗磁平面上方高度
\(h\)，磁矩垂直平面。用镜像偶极子法求平衡高度。

\end{minipage}\par

\par\noindent\begin{minipage}{\linewidth}
\textbf{解答}

\begin{center}
\includegraphics[width=\linewidth,height=.40\textheight,keepaspectratio]{\detokenize{assets/module-05/examples/M5-165.pdf}}
\par\smallskip{\small 图 M5-165：磁偶极子在完全抗磁平面上方的悬浮}
\end{center}
\end{minipage}\par

图 M5-165 的镜像偶极矩反向，镜像距离为
\(2h\)；求实偶极子在镜像场中的力时不能再多乘或少乘一个二分之一。

完全抗磁平面的镜像偶极子位于平面下方 \(h\) 处，磁矩为
\(-\mu\mathbf e_z\)。两个反平行同轴偶极子间距为 \(2h\)，斥力大小为

\[
F_m=\frac{3\mu_0\mu^2}{2\pi(2h)^4}
=\frac{3\mu_0\mu^2}{32\pi h^4}.
\]

平衡条件 \(F_m=Mg\)，故

\[
\boxed{
h_0=\left(\frac{3\mu_0\mu^2}{32\pi Mg}\right)^{1/4}}.
\]

若磁矩平行于平面，镜像方向和相互作用系数不同，不能沿用该式。

\section{专题五十三：转动和滚动导体中的磁通约束（决赛拓展）}\label{m5:ux4e13ux9898ux4e94ux5341ux4e09ux8f6cux52a8ux548cux6edaux52a8ux5bfcux4f53ux4e2dux7684ux78c1ux901aux7ea6ux675fux51b3ux8d5bux62d3ux5c55}

转动或滚动导体同时含有几何约束、动生电动势、磁通守恒和机械能变化。可选一个机械广义坐标，把磁通、电流和力矩都写成该坐标的函数。

\subsection{一、零电阻闭合回路中的磁通冻结}\label{m5:ux4e00ux96f6ux7535ux963bux95edux5408ux56deux8defux4e2dux7684ux78c1ux901aux51bbux7ed3}

若闭合回路电阻为零、自感为 \(L(q)\)，外磁通为
\(\Phi_{\rm ext}(q,t)\)，回路方程为

\[
\frac d{dt}\left[Li+\Phi_{\rm ext}\right]=0.
\]

因此

\[
\boxed{Li+\Phi_{\rm ext}=\text{常量}}.
\]

这不是说外磁通不变，而是自感磁链与外磁通之和保持。电流会自动调节，以抵抗总磁通的变化。

\par\noindent\begin{minipage}{\linewidth}
\begin{center}
\includegraphics[width=\linewidth,height=.40\textheight,keepaspectratio]{\detokenize{assets/module-05/original/tikz-m5-topic53-v9-from50-v8-01.png}}
\par\smallskip{\small 图示：零电阻回路中自感磁链补偿外磁通变化}
\end{center}
\end{minipage}\par

\subsection{二、滚动圆柱上的闭合线框}\label{m5:ux4e8cux6edaux52a8ux5706ux67f1ux4e0aux7684ux95edux5408ux7ebfux6846}

半径为 \(r\)、质量为 \(m\) 的实心圆柱沿倾角为 \(\alpha\)
的斜面无滑动滚下。圆柱上固定一匝面积为 \(2rh\) 的矩形闭合线框，自感为
\(L\)、电阻为零。匀强磁场 \(B\)
竖直，初始时线框法向与磁场同向、电流为零。令圆柱转角为 \(\psi=x/r\)。

外磁通为

\[
\Phi_{\rm ext}=2Brh\cos\psi.
\]

由总磁通守恒，

\[
Li+2Brh\cos\psi=2Brh,
\]

所以

\[
\boxed{i(\psi)=\frac{2Brh}{L}(1-\cos\psi)}.
\]

磁力矩阻碍转动，大小为

\[
\tau_m=i(2rh)B\sin\psi
=\frac{4B^2r^2h^2}{L}(1-\cos\psi)\sin\psi.
\]

实心圆柱关于接触点的等效转动惯量为

\[
J_{\rm eff}=mr^2+\frac12mr^2=\frac32mr^2.
\]

于是运动方程为

\[
\boxed{
\frac32mr^2\ddot\psi
=mgr\sin\alpha
-\frac{4B^2r^2h^2}{L}(1-\cos\psi)\sin\psi}.
\]

对应有效势可取

\[
\boxed{
U_{\rm eff}(\psi)
=\frac{2B^2r^2h^2}{L}(1-\cos\psi)^2
-mgr\sin\alpha\,\psi}.
\]

\par\noindent\begin{minipage}{\linewidth}
\begin{center}
\includegraphics[width=\linewidth,height=.40\textheight,keepaspectratio]{\detokenize{assets/module-05/theory/tikz-m5-topic53-v9-from50-v8-02.pdf}}
\par\smallskip{\small 图示：滚动线框的磁通与力矩}
\end{center}
\end{minipage}\par

平衡点满足 \(U_{\rm eff}'(\psi_0)=0\)；稳定性由
\(U_{\rm eff}''(\psi_0)>0\) 判断。若有稳定点，小振动角频率为

\[
\boxed{
\omega_s^2=\frac{U_{\rm eff}''(\psi_0)}{J_{\rm eff}}}.
\]

\subsection{三、两个零电阻耦合回路}\label{m5:ux4e09ux4e24ux4e2aux96f6ux7535ux963bux8026ux5408ux56deux8def}

两个静止线性回路的磁链为

\[
\begin{pmatrix}\Psi_1\\\Psi_2\end{pmatrix}
=\begin{pmatrix}L_1&M\\M&L_2\end{pmatrix}
\begin{pmatrix}i_1\\i_2\end{pmatrix}
+\begin{pmatrix}\Phi_{1e}\\\Phi_{2e}\end{pmatrix}.
\]

若两回路电阻均为零，则 \(\Psi_1,\Psi_2\) 分别守恒。只要
\(L_1L_2-M^2>0\)，电流可由矩阵反演得到：

\[
\boxed{
\begin{pmatrix}i_1\\i_2\end{pmatrix}
=-\frac1{L_1L_2-M^2}
\begin{pmatrix}L_2&-M\\-M&L_1\end{pmatrix}
\begin{pmatrix}\Delta\Phi_{1e}\\\Delta\Phi_{2e}\end{pmatrix}}.
\]

该式适用于线性互易条件下的超导耦合环、移动互感和磁通锁定模型。

\subsection{四、练习与解答}\label{m5:ux56dbux7ec3ux4e60ux4e0eux89e3ux7b54}

\par\addvspace{.4\baselineskip}
\noindent\begin{minipage}{\linewidth}

\subsubsection{作业一：滚动线框的平衡与稳定性}\label{m5:ux4f5cux4e1aux4e00ux6edaux52a8ux7ebfux6846ux7684ux5e73ux8861ux4e0eux7a33ux5b9aux6027}

\textbf{题号：M5-166}

\textbf{层级：} 决赛　　\textbf{方法：} 能量法、法拉第定律、机电耦合

半径 \(r\)、质量 \(m\) 的实心圆柱在倾角 \(\alpha\)
的斜面上无滑动滚动，固定一匝面积 \(2rh\) 的矩形零电阻线框，自感为常量
\(L\)。匀强磁场 \(B\) 竖直，初始法向与磁场同向且电流为零，取下滚转角
\(\psi=x/r\)。

对本专题滚动圆柱模型，写出平衡方程和小振动稳定条件。

\end{minipage}\par

\par\noindent\begin{minipage}{\linewidth}
\textbf{解答}

\begin{center}
\includegraphics[width=\linewidth,height=.40\textheight,keepaspectratio]{\detokenize{assets/module-05/examples/M5-166.pdf}}
\par\smallskip{\small 图 M5-166：滚动线框的平衡与稳定性}
\end{center}
\end{minipage}\par

图 M5-166 的滚动约束使平移和转动共用坐标
\(\psi=x/r\)；先由磁通冻结消去电流，再对同一个有效势判定平衡与稳定性。

先从三个物理量建立有效势。初始磁通为
\(2Brh\)、电流为零，零电阻回路的总磁链守恒，故 \[
LI+2Brh\cos\psi=2Brh,\qquad
I(\psi)=\frac{2Brh}{L}(1-\cos\psi).
\] 下滚距离 \(x=r\psi\) 对应重力势能变化
\(-mgr\sin\alpha\,\psi\)；将电流代入磁能 \(LI^2/2\)，得到 \[
U_{\rm eff}(\psi)=\frac{2B^2r^2h^2}{L}(1-\cos\psi)^2-mgr\sin\alpha\,\psi.
\] 无滑动时平动也贡献动能。实心圆柱自身惯量为 \(mr^2/2\)，所以
\(T=\tfrac12J_{\rm eff}\dot\psi^2\)、\(J_{\rm eff}=mr^2+mr^2/2=3mr^2/2\)。以下平衡式用
\(U_{\rm eff}'=0\)，稳定条件用 \(U_{\rm eff}''>0\)，小振动用
\(\omega_s^2=U_{\rm eff}''/J_{\rm eff}\)。

由 \(U_{\rm eff}'(\psi_0)=0\)，

\[
\boxed{
mgr\sin\alpha
=\frac{4B^2r^2h^2}{L}
(1-\cos\psi_0)\sin\psi_0}.
\]

二阶导数为

\[
U_{\rm eff}''(\psi)
=\frac{4B^2r^2h^2}{L}
\left[\sin^2\psi+(1-\cos\psi)\cos\psi\right].
\]

\par\noindent\begin{minipage}{\linewidth}
\begin{center}
\includegraphics[width=\linewidth,height=.40\textheight,keepaspectratio]{\detokenize{assets/module-05/process/M5-166a.pdf}}
\par\smallskip{\small 图 M5-166a：平衡和小振动使用同一有效势}
\end{center}
\end{minipage}\par

图 M5-166a
将磁链、电流、有效势和滚动惯量对应起来；左图极小值为稳定平衡，极大值为不稳定平衡。曲线使用的示意参数已另行标明。

故稳定条件为

\[
\boxed{
\sin^2\psi_0+(1-\cos\psi_0)\cos\psi_0>0}.
\]

若成立，

\[
\boxed{
\omega_s^2
=\frac{8B^2h^2}{3mL}
\left[\sin^2\psi_0+(1-\cos\psi_0)\cos\psi_0\right]}.
\]

\par\addvspace{.4\baselineskip}
\noindent\begin{minipage}{\linewidth}

\subsubsection{作业二：两个超导耦合环对外磁通变化的响应}\label{m5:ux4f5cux4e1aux4e8cux4e24ux4e2aux8d85ux5bfcux8026ux5408ux73afux5bf9ux5916ux78c1ux901aux53d8ux5316ux7684ux54cdux5e94}

\textbf{题号：M5-167}

\textbf{层级：} 决赛　　\textbf{方法：} 法拉第定律、自感互感、机电耦合

回路静止、线性、无电阻，互感采用共同点标参考且
\(|M|<L\)，外磁通改变前各回路总磁链均为零。

两个相同超导环满足自感 \(L\)、互感 \(M\)，初始电流为零。外界随后只给回路
1 增加外磁通 \(\Delta\Phi\)。求两环电流。

\end{minipage}\par

\par\noindent\begin{minipage}{\linewidth}
\textbf{解答}

\begin{center}
\includegraphics[width=\linewidth,height=.40\textheight,keepaspectratio]{\detokenize{assets/module-05/examples/M5-167.pdf}}
\par\smallskip{\small 图 M5-167：两个超导耦合环对外磁通变化的响应}
\end{center}
\end{minipage}\par

图 M5-167
的参考方向均流入点端；两个零电阻闭合环各自保持总磁链，外磁通只加在第一环，第二环仍通过互感产生反向响应。

磁链守恒给出

\[
Li_1+Mi_2+\Delta\Phi=0,
\]

\[
Mi_1+Li_2=0.
\]

解得

\[
\boxed{i_1=-\frac{L}{L^2-M^2}\Delta\Phi},
\]

\[
\boxed{i_2=\frac{M}{L^2-M^2}\Delta\Phi}.
\]

第二环虽未直接穿过新增外磁通，仍会因互感产生电流。

\par\addvspace{.4\baselineskip}
\noindent\begin{minipage}{\linewidth}

\subsubsection{作业三：椭圆轨道上的感生电场加速与磁场偏转}\label{m5:ux4f5cux4e1aux4e09ux692dux5706ux8f68ux9053ux4e0aux7684ux611fux751fux7535ux573aux52a0ux901fux4e0eux78c1ux573aux504fux8f6c}

\textbf{题号：M5-168}

\textbf{层级：} 决赛　　\textbf{方法：} 洛伦兹力、法拉第定律、机电耦合

粒子受无摩擦绝缘轨道约束，采用非相对论模型，忽略辐射；原磁区在 \(+z\)
方向增强，轨道处可另施恒定法向场 \(B_2\)，它不改变感生电场。

圆心处半径 \(R\) 的圆域内有均匀磁场，满足
\(dB_1/dt=k>0\)。其外有椭圆轨道

\[
\frac{x^2}{a^2}+\frac{y^2}{b^2}=1,
\qquad a>b>R.
\]

带正电粒子质量为 \(m\)、电荷为 \(q\)，从左端点 \(A=(-a,0)\)
静止释放，沿使电场做正功的方向运动。它转过 \(n\) 个整周后再到达上端点
\(C=(0,b)\)。若希望粒子在 \(C\) 点恰由垂直轨道面的磁场 \(B_2\)
提供全部法向力，使轨道支持力为零，求 \(B_2\)。

\end{minipage}\par

\par\noindent\begin{minipage}{\linewidth}
\textbf{解答}

\begin{center}
\includegraphics[width=\linewidth,height=.40\textheight,keepaspectratio]{\detokenize{assets/module-05/examples/M5-168.pdf}}
\par\smallskip{\small 图 M5-168：椭圆轨道上的感生电场加速与磁场偏转}
\end{center}
\end{minipage}\par

图 M5-168 从左端 \(A\) 顺时针到上端
\(C\)，每个整周再加四分之一周都使电场做正功。\(C\)
处法向向下，正电荷向右运动时应取 \(B_2\) 向纸外。

圆磁区内的 \(B_1\) 沿纸面向外且增强，因此在 \(r>R\)
的轨道处感生电场沿顺时针方向，并满足

\[
|E_{\theta}|=\frac{kR^2}{2r}.
\]

若以逆时针为方位角正向，则 \(E_{\theta}=-kR^2/(2r)\)。粒子从 \(A\)
顺时针运动到 \(C\)，经过 \(n\) 个整周后再走四分之一周，因此

\[
\Delta\theta=-\left(2\pi n+\frac{\pi}{2}\right).
\]

椭圆上电场一般不与位移平行；其切向投影沿题定顺时针方向做正功。极坐标中
\(d\mathbf l=\mathbf e_r\,dr+\mathbf e_{\phi} r\,d\phi\)，所以
\(\mathbf E\cdot d\mathbf l=E_{\phi} r\,d\phi\)，故做功为

\[
W=q\int E_{\theta} r\,d\theta
=qkR^2\pi\left(n+\frac14\right).
\]

\par\noindent\begin{minipage}{\linewidth}
\begin{center}
\includegraphics[width=\linewidth,height=.40\textheight,keepaspectratio]{\detokenize{assets/module-05/process/M5-168a.pdf}}
\par\smallskip{\small 图 M5-168a：到达上端点时的速度、法向与曲率}
\end{center}
\end{minipage}\par

图 M5-168a 的 \(C\)
点速度向右，椭圆内法向向下；用局部曲率半径而非短半轴作为向心半径，再确定纸外
\(B_2\)。

所以

\[
v_C^2=\frac{2W}{m}
=\frac{qkR^2\pi}{2m}(4n+1).
\]

椭圆在上端点的曲率半径为

\[
\rho_C=\frac{a^2}{b}.
\]

支持力为零时

\[
qv_CB_2=\frac{mv_C^2}{\rho_C}.
\]

因此

\[
\boxed{
B_2=\frac{bR}{a^2}
\sqrt{\frac{mk\pi(4n+1)}{2q}}}.
\]

在 \(C\) 点粒子速度沿 \(+x\) 方向，曲率中心位于下方。对正电荷，\(B_2\)
应沿纸面向外，使 \(q\mathbf v\times\mathbf B_2\) 沿 \(-y\)
指向曲率中心。

\section{专题五十四：正则动量、规范与对称性（决赛拓展）}\label{m5:ux4e13ux9898ux4e94ux5341ux56dbux6b63ux5219ux52a8ux91cfux89c4ux8303ux4e0eux5bf9ux79f0ux6027ux51b3ux8d5bux62d3ux5c55}

\begin{quote}
阅读准备：从牛顿方程到拉格朗日方程 对广义坐标
\(q_i\)，欧拉---拉格朗日方程为
\(d(\partial L/\partial\dot q_i)/dt-\partial L/\partial q_i=0\)。若
\(L\) 不显含某个坐标，该坐标是循环坐标，相应的正则动量
\(\partial L/\partial\dot q_i\) 守恒；它不必等于机械动量。
\end{quote}

\par\addvspace{.4\baselineskip}
\noindent\begin{minipage}{\linewidth}

\textbf{例 M5-169：循环坐标与正则动量}

一维自由粒子的拉格朗日量为
\(L=m\dot x^2/2\)，无外力且质量不变。由欧拉---拉格朗日方程证明
\(p_x=m\dot x\) 守恒；再说明含电磁势时，共轭动量为何一般不同于机械动量。

\end{minipage}\par

\par\noindent\begin{minipage}{\linewidth}
\textbf{解答}

\begin{center}
\includegraphics[width=\linewidth,height=.40\textheight,keepaspectratio]{\detokenize{assets/module-05/examples/M5-169.pdf}}
\par\smallskip{\small 图 M5-169：循环坐标与自由粒子的动量}
\end{center}
\end{minipage}\par

图 M5-169
中粒子在两个时刻具有相同的水平速度，位置变化不会改变动量。坐标不显含于拉格朗日量时，其共轭动量守恒。

自由粒子的 \(L=m\dot x^2/2\) 不显含 \(x\)，故
\(\partial L/\partial\dot x=m\dot x\)
守恒；加入电磁势后，守恒量会多出势相关项。

带电粒子处于磁场中时，平移对称性对应的守恒量一般是含磁矢势的正则动量，而非单独的机械动量
\(m\mathbf v\)。磁矢势与规范变换见专题四十五，洛伦兹力与粒子运动见专题四十一。

\subsection{一、最小耦合拉格朗日量与正则动量}\label{m5:ux4e00ux6700ux5c0fux8026ux5408ux62c9ux683cux6717ux65e5ux91cfux4e0eux6b63ux5219ux52a8ux91cf}

质量为 \(m\)、电荷为 \(q\) 的粒子在给定电磁势 \((\varphi,\mathbf A)\)
中的拉格朗日量为

\[
\boxed{
\mathcal L(\mathbf r,\mathbf v,t)
=\frac12m\mathbf v^2+q\mathbf A(\mathbf r,t)\cdot\mathbf v
-q\varphi(\mathbf r,t)}.
\]

对速度求偏导，得到正则动量

\[
\boxed{
\mathbf p_{\rm can}
=\frac{\partial\mathcal L}{\partial\mathbf v}
=m\mathbf v+q\mathbf A}.
\]

其中 \(m\mathbf v\) 是机械动量，\(q\mathbf A\)
是由最小耦合产生的势项。正则动量依赖规范，不能把 \(q\mathbf A\)
单独解释成可直接测量的机械动量。

\subsection{二、由欧拉---拉格朗日方程恢复洛伦兹力}\label{m5:ux4e8cux7531ux6b27ux62c9ux62c9ux683cux6717ux65e5ux65b9ux7a0bux6062ux590dux6d1bux4f26ux5179ux529b}

以下采用爱因斯坦求和约定：重复出现的直角坐标指标自动求和，\(i\)
是固定的自由指标。

对第 \(i\) 个直角坐标分量，欧拉---拉格朗日方程为

\[
\frac d{dt}\frac{\partial\mathcal L}{\partial v_i}
-\frac{\partial\mathcal L}{\partial x_i}=0.
\]

逐项展开：

\[
m\dot v_i
+q\frac{\partial A_i}{\partial t}
+qv_j\frac{\partial A_i}{\partial x_j}
=qv_j\frac{\partial A_j}{\partial x_i}
-q\frac{\partial\varphi}{\partial x_i}.
\]

整理得

\[
m\dot v_i
=q\left(-\frac{\partial\varphi}{\partial x_i}
-\frac{\partial A_i}{\partial t}\right)
+qv_j\left(
\frac{\partial A_j}{\partial x_i}
-\frac{\partial A_i}{\partial x_j}
\right).
\]

利用

\[
\mathbf E=-\nabla\varphi-\frac{\partial\mathbf A}{\partial t},
\qquad
\mathbf B=\nabla\times\mathbf A,
\]

便恢复

\[
\boxed{m\dot{\mathbf v}=q(\mathbf E+\mathbf v\times\mathbf B)}.
\]

因而``势的拉格朗日描述''和``场的洛伦兹力描述''给出同一运动方程。

\subsection{三、循环坐标、守恒量与规范变换}\label{m5:ux4e09ux5faaux73afux5750ux6807ux5b88ux6052ux91cfux4e0eux89c4ux8303ux53d8ux6362}

若拉格朗日量不显含坐标 \(x\)，即 \(\partial\mathcal L/\partial x=0\)，则
\(x\) 是循环坐标，欧拉---拉格朗日方程立即给出

\[
\boxed{p_x^{\rm can}=mv_x+qA_x=\text{常量}}.
\]

\par\addvspace{.4\baselineskip}
\noindent\begin{minipage}{\linewidth}

\textbf{例 M5-170：匀强磁场中的平移对称规范}

非相对论带电粒子在匀强磁场 \(B\mathbf e_z\) 中运动，取
\(\mathbf A=-By\mathbf e_x\text{、}\varphi=0\)。写出最小耦合拉格朗日量，说明
\(x\) 为何是循环坐标，并求相应守恒的正则动量及其与机械动量的区别。

\end{minipage}\par

\par\noindent\begin{minipage}{\linewidth}
\textbf{解答}

\begin{center}
\includegraphics[width=\linewidth,height=.40\textheight,keepaspectratio]{\detokenize{assets/module-05/examples/M5-170.pdf}}
\par\smallskip{\small 图 M5-170：匀强磁场中的平移对称规范}
\end{center}
\end{minipage}\par

图 M5-170 的磁场向纸外，轨道弯曲时机械分量 \(mv_x\) 会变化；矢势项
\(-qBy\) 恰好补偿，使规范下的正则动量 \(p_x\) 保持不变。

本题所需的最小耦合拉格朗日量先写为 \[
\mathcal L=\frac m2(v_x^2+v_y^2+v_z^2)
+q\mathbf A\cdot\mathbf v-q\varphi
=\frac m2(v_x^2+v_y^2+v_z^2)-qByv_x.
\] 它不显含 \(x\)。欧拉---拉格朗日方程给
\(d(\partial\mathcal L/\partial v_x)/dt=\partial\mathcal L/\partial x=0\)，于是得到下文的
\(mv_x-qBy\) 守恒；这里守恒的是含势项的正则动量。

对

\[
\mathbf B=B\mathbf e_z
\]

取

\[
\mathbf A=-By\mathbf e_x,
\]

则拉格朗日量不显含 \(x\)，所以

\[
\boxed{mv_x-qBy=\text{常量}}.
\]

对时间求导得到

\[
m\dot v_x=qBv_y,
\]

与洛伦兹力的 \(x\) 分量完全一致。

规范变换

\[
\mathbf A'=\mathbf A+\nabla\chi,
\qquad
\varphi'=\varphi-\frac{\partial\chi}{\partial t}
\]

使拉格朗日量改变一个全时间导数：

\[
\mathcal L'=\mathcal L+q\frac{d\chi}{dt}.
\]

因而运动方程不变，但

\[
\mathbf p_{\rm can}'=\mathbf p_{\rm can}+q\nabla\chi.
\]

若新规范使 \(\mathcal L'\) 显含
\(x\)，不能据此误判物理平移对称性消失；诺特守恒量还要包含全导数项带来的修正。真正可观测的是轨道、机械动量和场，而不是某一规范下正则动量的单独数值。

\subsection{四、哈密顿量与能量守恒}\label{m5:ux56dbux54c8ux5bc6ux987fux91cfux4e0eux80fdux91cfux5b88ux6052}

用 \(\mathbf p_{\rm can}\) 消去速度：

\[
\mathbf v=\frac{\mathbf p_{\rm can}-q\mathbf A}{m}.
\]

哈密顿量为

\[
\boxed{
\mathcal H
=\mathbf p_{\rm can}\cdot\mathbf v-\mathcal L
=\frac{(\mathbf p_{\rm can}-q\mathbf A)^2}{2m}+q\varphi}.
\]

若外加势不显含时间，则 \(\mathcal H\) 守恒。代回机械速度可见

\[
\mathcal H=\frac12m v^2+q\varphi,
\]

磁场不做功这一事实与正则形式相容。

\subsection{五、使用条件与常见错误}\label{m5:ux4e94ux4f7fux7528ux6761ux4ef6ux4e0eux5e38ux89c1ux9519ux8bef}

\begin{enumerate}
\def\labelenumi{\arabic{enumi}.}
\tightlist
\item
  机械动量 \(m\mathbf v\) 与正则动量 \(m\mathbf v+q\mathbf A\)
  必须区分。
\item
  ``坐标不显含''是对所选拉格朗日量的判断；规范变换可能改变显式形式，但不改变物理对称性。
\item
  单独的 \(\mathbf A\) 和 \(\mathbf p_{\rm can}\)
  依赖规范，洛伦兹力、轨道和场强不依赖规范。
\item
  本专题采用经典点粒子模型；相对论情形应把自由粒子项改为
  \(-mc^2\sqrt{1-v^2/c^2}\)。
\item
  若用显含时间的规范描述静态场，正则哈密顿量的表达会随规范改变；物理能量是否守恒应由场的时间平移对称性判断。
\end{enumerate}

\subsection{六、练习与解答}\label{m5:ux516dux7ec3ux4e60ux4e0eux89e3ux7b54}

\par\addvspace{.4\baselineskip}
\noindent\begin{minipage}{\linewidth}

\subsubsection{作业一：由正则动量恢复洛伦兹力分量}\label{m5:ux4f5cux4e1aux4e00ux7531ux6b63ux5219ux52a8ux91cfux6062ux590dux6d1bux4f26ux5179ux529bux5206ux91cf}

\textbf{题号：M5-171}

\textbf{层级：} 决赛　　\textbf{方法：} 势与规范

粒子质量 \(m\)、电荷 \(q\) 均为常量，磁场 \(B\mathbf e_z\)
匀强且不随时间变化；采用最小耦合的非相对论拉格朗日量。

在 \(\mathbf A=-By\mathbf e_x\)、\(\varphi=0\) 的规范中，证明
\(mv_x-qBy\) 守恒，并恢复 \(x\) 方向运动方程。

\end{minipage}\par

\par\noindent\begin{minipage}{\linewidth}
\textbf{解答}

\begin{center}
\includegraphics[width=\linewidth,height=.40\textheight,keepaspectratio]{\detokenize{assets/module-05/examples/M5-171.pdf}}
\par\smallskip{\small 图 M5-171：由正则动量恢复洛伦兹力分量}
\end{center}
\end{minipage}\par

图 M5-171 中磁力的 \(x\) 分量为 \(qBv_y\)，与守恒式求导后的
\(m\dot v_x=qB\dot y\) 相符；不能把守恒量误写成单独的 \(mv_x\)。

拉格朗日量不显含 \(x\)，故

\[
p_x^{\rm can}=\frac{\partial\mathcal L}{\partial v_x}
=mv_x+qA_x=mv_x-qBy
\]

守恒。求导得

\[
m\dot v_x-qBv_y=0,
\]

即

\[
\boxed{m\dot v_x=qBv_y}.
\]

这与 \(q\mathbf v\times\mathbf B\) 的 \(x\) 分量相同。

\section{模块五小结}\label{m5:ux6a21ux5757ux4e94ux5c0fux7ed3}

{\def\LTcaptype{none} % do not increment counter
\begin{longtable}[]{@{}
  >{\raggedright\arraybackslash}p{(\linewidth - 4\tabcolsep) * \real{0.3333}}
  >{\raggedright\arraybackslash}p{(\linewidth - 4\tabcolsep) * \real{0.3333}}
  >{\raggedright\arraybackslash}p{(\linewidth - 4\tabcolsep) * \real{0.3333}}@{}}
\toprule\noalign{}
\begin{minipage}[b]{\linewidth}\raggedright
检验问题
\end{minipage} & \begin{minipage}[b]{\linewidth}\raggedright
结论或判断
\end{minipage} & \begin{minipage}[b]{\linewidth}\raggedright
回看位置
\end{minipage} \\
\midrule\noalign{}
\endhead
\bottomrule\noalign{}
\endlastfoot
磁力是否改变单个带电粒子的速率？ & 纯磁场不做功 & 专题四十一 \\
毕奥---萨伐尔与安培环路定理如何选择？ & 前者通用积分，后者依赖高对称 &
专题四十二、四十三 \\
磁标势在哪些区域失效？ & 含自由电流或拓扑环流的区域 & 专题四十五 \\
动生电动势的方向由什么决定？ & \(q\mathbf v\times\mathbf B\) 与闭路取向
& 专题四十六 \\
互感符号为何不能只看线圈位置？ & 还依赖两个参考电流和磁通取向 &
专题四十七 \\
磁共能法求力时固定什么？ & 按选定独立电流或磁链变量求偏导 &
专题五十、五十一 \\
\end{longtable}
}

\section{模块五综合题：变化磁场中的自相似线框}\label{m5:ux6a21ux5757ux4e94ux7efcux5408ux9898ux53d8ux5316ux78c1ux573aux4e2dux7684ux81eaux76f8ux4f3cux7ebfux6846}

\par\addvspace{.4\baselineskip}
\noindent\begin{minipage}{\linewidth}

\subsubsection{作业一：综合应用}\label{m5:ux4f5cux4e1aux4e00ux7efcux5408ux5e94ux7528}

\textbf{题号：M5-172}

\textbf{层级：} 复赛　　\textbf{方法：} 法拉第定律、自相似、机电耦合

规定 \(F_0\) 为边长 \(l\) 的普通正方形线框；\(F_{n+1}\) 由四个边长
\(l/3\) 的 \(F_n\)
副本放在四角，并只沿大正方形的四边中间三分之一相连。\(A\text{、}B\text{、}C\text{、}D\)
为按逆时针排列的外顶点，单位长度电阻
\(r>0\)、\(k>0\)。只计电阻与给定均匀磁场的准静态感应，忽略自感、互感与电磁辐射。

一个方形线框按比例 \(1/3\)
递归构造：每一级大正方形的四个角均由边长缩小为原来 \(1/3\)
的同形线框替代，相邻角部线框之间由长度为原边长 \(1/3\)
的直导线连接。导线单位长度的电阻为 \(r\)，最外层正方形边长为 \(l\)。

严格地说，无限分形线框定义为有限级线框 \(F_n\) 在 \(n\to\infty\)
时的极限。以下先对有限级电阻网络使用基尔霍夫定律，再取自相似不动点；不能把``无限细节''直接当成一条普通闭合导线。

\begin{enumerate}
\def\labelenumi{(\arabic{enumi})}
\item
  求相邻外顶点 \(A\text{、}B\) 之间的等效电阻 \(R_{AB}\)，以及对角外顶点
  \(A\text{、}C\) 之间的等效电阻 \(R_{AC}\)。
\item
  线框位于垂直其平面的匀强磁场中，磁感应强度满足
  \(B=kt\)。求最外层每条边中间那段长度为 \(l/3\)
  的直导线中的感应电流大小 \(I_c\)。
\item
  求整个无限自相似线框的焦耳热功率 \(P\)。
\end{enumerate}

\end{minipage}\par

\par\noindent\begin{minipage}{\linewidth}
\textbf{解答}

\begin{center}
\includegraphics[width=\linewidth,height=.40\textheight,keepaspectratio]{\detokenize{assets/module-05/examples/M5-172.pdf}}
\par\smallskip{\small 图 M5-172：自相似线框的电阻、感应电流与功率}
\end{center}
\end{minipage}\par

图 M5-172
实际画出有限二级近似，橙线是最外层下边中段；四角子网通过四条中段连接，端点按
\(A\to B\to C\to D\) 编号。有限级网络求解后再取极限。

\par\noindent\begin{minipage}{\linewidth}
\begin{center}
\includegraphics[width=\linewidth,height=.40\textheight,keepaspectratio]{\detokenize{assets/module-05/process/M5-172a.pdf}}
\par\smallskip{\small 图 M5-172a：三个外端口的Y形自相似等效}
\end{center}
\end{minipage}\par

图 M5-172a 的中间外角 \(B\) 对应 \(x\) 臂，两个对角顶点 \(A\text{、}C\)
对应两条 \(y\) 臂。整体缩小为三分之一时，三个臂电阻都同倍缩小。

\paragraph{第(1)问：等效电阻}\label{m5:ux7b2c1ux95eeux7b49ux6548ux7535ux963b}

选取三个外顶点作为端口，第四个外顶点不接外电路。由对角线反射对称性，这个三端电阻网络可等效为支路电阻为
\(x,y,y\) 的Y形网络。线框整体缩小为原来的 \(1/3\)
后，每段导线电阻均缩小为原来的
\(1/3\)，所以角部小线框的三个Y形支路分别为 \(x/3,y/3,y/3\)。

记相邻两个角部小线框之间一段宏观连接的等效电阻为

\[
R_1=\frac{y}{3}+\frac{lr}{3}+\frac{y}{3}
=\frac{2y}{3}+\frac{lr}{3}.
\]

当电流从相邻顶点 \(A\text{、}B\)
流入、流出时，中部环形网络的一条路径电阻为 \(R_1\)，另一条路径电阻为
\(3R_1\)，二者并联。因此

\[
x+y=\frac{2x}{3}+R_1\parallel 3R_1
=\frac{2x}{3}+\frac34\left(\frac{2y}{3}+\frac{lr}{3}\right).
\]

当电流从对角顶点 \(A\text{、}C\)
流入、流出时，中部环形网络分成两条电阻均为 \(2R_1\) 的对称路径，故

\[
2y=\frac{2x}{3}+2R_1\parallel 2R_1
=\frac{2x}{3}+\frac{2y}{3}+\frac{lr}{3}.
\]

联立两式得到

\[
x=\frac{3}{14}lr,\qquad y=\frac{5}{14}lr.
\]

所以

\[
\boxed{R_{AB}=x+y=\frac47lr},
\qquad
\boxed{R_{AC}=2y=\frac57lr}.
\]

\paragraph{第(2)问：中间直导线中的感应电流}\label{m5:ux7b2c2ux95eeux4e2dux95f4ux76f4ux5bfcux7ebfux4e2dux7684ux611fux5e94ux7535ux6d41}

记最外层正方形所围磁通量变化率的大小为

\[
\mathcal E(l)=\left|\frac{d\Phi_l}{dt}\right|=kl^2.
\]

边长为 \(s=l/3\) 的角部小线框满足

\[
\mathcal E(s)=ks^2=\frac19\mathcal E(l),
\qquad
R_{AC}(s)=\frac13R_{AC}(l)=\frac{5}{21}lr.
\]

这里不能把非保守感生电场产生的电动势任意均分给导线。将左下角边长 \(s\)
的子框放大，令
\(A=(0,0)\)、\(P=(s,0)\)、\(Q=(0,s)\)、\(T=(s,s)\)；\(P,Q\)
是与相邻中间直段相接的两个端点。

\par\noindent\begin{minipage}{\linewidth}
\begin{center}
\includegraphics[width=\linewidth,height=.40\textheight,keepaspectratio]{\detokenize{assets/module-05/process/M5-172e.pdf}}
\par\smallskip{\small 图 M5-172e：角部子框的两条完整路径与半因子}
\end{center}
\end{minipage}\par

图 M5-172e 取两条实际边界链
\(\Gamma_+:P\to A\to Q\)、\(\Gamma_-:P\to T\to Q\)，积分方向均为
\(P\to Q\)。先单独考虑没有外部端口注入的纯感应辅助态，定义
\(e_{\pm}=\sum_{\Gamma_{\pm}}R_ei_e^{(B)}\)。

绕子框中心 \(O_s\) 旋转 \(180^{\circ}\)，\(P,Q\) 交换、\(A,T\)
交换，\(\Gamma_+\) 因而映成反向的
\(\Gamma_-\)。均匀磁场与有限子网在该正旋转下不变，故
\(e_+=-e_-\)。两条路径相减构成顺时针闭合边界
\(P\to A\to Q\to T\to P\)，其法向纸内；外场向纸外增强，法拉第定律给 \[
e_+-e_-=-\frac{d(-Bs^2)}{dt}=\mathcal E(s).
\] 两式联立才得到 \(e_+=\mathcal E(s)/2\)。这个半因子不表示每一条边获得
\(\mathcal E(s)/4\)。

再单独叠加关掉感应激励、只从 \(P\) 注入并从 \(Q\) 抽出 \(I_c\)
的端口电流态。两条路径压降均为 \(R_{AC}(s)I_c\)。叠加后，沿外侧路径
\(\Gamma_+\) 的电阻压降为
\(R_{AC}(s)I_c+\mathcal E(s)/2\)。两种辅助态分别满足基尔霍夫关系，有限递归级数及其极限都可这样处理。

由四重旋转对称性，四段长度为 \(l/3\)
的中间直导线中电流大小相同，并沿同一环流方向。一个四分之一约化单元中的串联电阻为

\[
R_q=R_{AC}(s)+\frac{lr}{3}
=\frac{5}{21}lr+\frac13lr
=\frac47lr.
\]

四个全同约化方程相加必须还原最外层回路的法拉第方程，所以每个方程的宏观磁通项为
\(\mathcal E(l)/4\)。取电流正方向与楞次定律给出的环流方向一致，有

\[
\frac14\mathcal E(l)-\frac12\mathcal E(s)=I_cR_q.
\]

代入 \(\mathcal E(s)=\mathcal E(l)/9\) 和 \(R_q=4lr/7\)，得到

\[
\frac{7}{36}\mathcal E(l)=I_c\frac47lr,
\]

从而

\[
\boxed{
I_c=\frac{49}{144}\frac{\mathcal E(l)}{lr}
=\frac{49}{144}\frac{kl}{r}}.
\]

\par\noindent\begin{minipage}{\linewidth}
\begin{center}
\includegraphics[width=\linewidth,height=.40\textheight,keepaspectratio]{\detokenize{assets/module-05/process/M5-172b.pdf}}
\par\smallskip{\small 图 M5-172b：完整路径配对后的四分之一约化}
\end{center}
\end{minipage}\par

图 M5-172b 是完整路径配对后的约化电路，\(e_q\)
是约化激励而非某条导线分得的感应电压。两段串联电阻与同一个端口电流
\(I_c\) 对应。

量纲检验：\(kl\) 的量纲为 \(\mathrm{V/m}\)，而 \(r\) 的量纲为
\(\Omega/\mathrm m\)，所以 \(kl/r\) 的量纲确为电流。

\paragraph{第(3)问：整个线框的焦耳热功率}\label{m5:ux7b2c3ux95eeux6574ux4e2aux7ebfux6846ux7684ux7126ux8033ux70edux529fux7387}

先证明感应激励与端口电流在功率中没有交叉项。对任意一个有限级子线框，把支路电流写成

\[
i_e=i_e^{(B)}+I\,i_e^{(I)},
\]

其中 \(i_e^{(B)}\)
是只有变化磁场、没有外部端口注入时的支路电流；\(i_e^{(I)}\)
是关闭感应电动势后，由单位端口电流产生的支路电流。后者是纯电阻网络中的势流，因此对支路
\(e=(m,n)\) 有

\[
R_ei_e^{(I)}=V_m-V_n.
\]

于是焦耳热展开式中的交叉项系数满足

\[
\begin{aligned}
\sum_eR_ei_e^{(B)}i_e^{(I)}
&=\sum_{e=(m,n)}(V_m-V_n)i_{mn}^{(B)}\\
&=\sum_mV_m\sum_n i_{mn}^{(B)}=0.
\end{aligned}
\]

最后一步使用了基尔霍夫电流定律：纯感应电流在每个节点都没有外部注入，故
\(\sum_n i_{mn}^{(B)}=0\)。因此交叉项消失来自``环流与势流的功率正交性''，并不是因为只把磁场反向后总功率必然不变。

因此，对边长为 \(l\) 的自相似线框，可写成

\[
P(l;\mathcal E,I)=D_1\mathcal E^2+D_2I^2.
\]

当 \(I\) 是从一对对角端点注入的电流时，

\[
D_2=R_{AC}(l)=\frac57lr.
\]

现在令外部端口电流为零，并把大线框分成四个边长为 \(l/3\)
的小线框和四段长度为 \(l/3\) 的直导线。线性尺度缩小为 \(1/3\)
时，所有电阻缩小为 \(1/3\)；对给定感应电动势，功率系数 \(D_1\)
因而放大为 \(3D_1\)，而端口电流对应的电阻系数缩小为
\(D_2/3\)。同时，小线框的磁通感应电动势为 \(\mathcal E(l)/9\)。

利用刚刚证明的交叉项为零，每个角部小线框的功率为

\[
3D_1\left(\frac{\mathcal E(l)}9\right)^2
+\frac13D_2I_c^2,
\]

\par\noindent\begin{minipage}{\linewidth}
\begin{center}
\includegraphics[width=\linewidth,height=.40\textheight,keepaspectratio]{\detokenize{assets/module-05/process/M5-172c.pdf}}
\par\smallskip{\small 图 M5-172c：功率中的环流与端口势流正交}
\end{center}
\end{minipage}\par

图 M5-172c
的有限级图先满足每个节点电流守恒；无外部注入的感应环流与端口电势流的交叉项可按节点汇总为零，之后才取无限级极限。

每段中间直导线的功率为 \(I_c^2lr/3\)。所以

\[
D_1\mathcal E(l)^2
=4\left[
3D_1\left(\frac{\mathcal E(l)}9\right)^2
+\frac13D_2I_c^2
+\frac13lrI_c^2
\right].
\]

整理得

\[
\frac{23}{27}D_1\mathcal E(l)^2
=\frac43(D_2+lr)I_c^2
=\frac{16}{7}lrI_c^2.
\]

代入

\[
I_c=\frac{49}{144}\frac{\mathcal E(l)}{lr},
\]

\par\noindent\begin{minipage}{\linewidth}
\begin{center}
\includegraphics[width=\linewidth,height=.40\textheight,keepaspectratio]{\detokenize{assets/module-05/process/M5-172d.pdf}}
\par\smallskip{\small 图 M5-172d：四个角部子网的功率递推}
\end{center}
\end{minipage}\par

图 M5-172d
的四个角部子网各自承受缩小九倍的磁通激励，外加同一中段电流通过相邻对角端口；四条中段的焦耳热也必须计入。

可得

\[
\frac{23}{27}D_1\mathcal E(l)^2
=\frac{343}{1296}\frac{\mathcal E(l)^2}{lr},
\]

即

\[
D_1=\frac{343}{1104}\frac1{lr}.
\]

因此整个无限自相似线框的焦耳热功率为

\[
\boxed{
P=D_1\mathcal E(l)^2
=\frac{343}{1104}\frac{k^2l^3}{r}}.
\]

递推式中感应功率的自相似系数为
\(4/27<1\)，所以有限级线框的功率序列收敛；这也说明上述无限分形极限是良定义的。

最后一步须保留 \(D_1\) 中的 \(1/l\)：\(D_1\propto1/(lr)\)，而
\(\mathcal E(l)^2=k^2l^4\)，故
\(P\propto k^2l^3/r\)。这既与有限级网络复算收敛值一致，也满足
\([P]=\mathrm W\)；写成 \(k^2l^4/r\) 会多出一个长度量纲。
\clearpage
\endgroup
